\documentclass[12pt]{report}
\usepackage[utf8x]{inputenc}

%Warning: tipa declares many non-standard macros used by utf8x to
%interpret utf8 characters but extra packages might have to be added
%such as "textgreek" for Greek letters not already in tipa
%or "stmaryrd" for mathematical symbols.
%Utf8 codes missing a LaTeX interpretation can be defined by using
%\DeclareUnicodeCharacter{code}{interpretation}.
%Use coqdoc's option -p to add new packages or declarations.
\usepackage{tipa}

\usepackage[T1]{fontenc}
\usepackage{fullpage}
\usepackage{coqdoc}
\usepackage{amsmath,amssymb}
\usepackage{url}
\begin{document}
%%%%%%%%%%%%%%%%%%%%%%%%%%%%%%%%%%%%%%%%%%%%%%%%%%%%%%%%%%%%%%%%%
%% This file has been automatically generated with the command
%% /nix/store/5nhkcrnn3hzlx1rsgy6zxlg2lrla4f5y-coq-with-packages-9.1.1/bin-orig/coqdoc -latex -toc -interpolate -utf8 -Q . SECF Maps.v Imp.v Equiv.v Hoare.v Hoare2.v Preface.v Noninterference.v StaticIFC.v SpecCT.v Postscript.v Bib.v MapsTest.v ImpTest.v EquivTest.v HoareTest.v Hoare2Test.v PrefaceTest.v NoninterferenceTest.v StaticIFCTest.v SpecCTTest.v PostscriptTest.v BibTest.v 
%%%%%%%%%%%%%%%%%%%%%%%%%%%%%%%%%%%%%%%%%%%%%%%%%%%%%%%%%%%%%%%%%
\coqlibrary{SECF.Noninterference}{Library }{SECF.Noninterference}

\begin{coqdoccode}
\end{coqdoccode}
\section{Noninterference: Defining Secrecy and Secure Multi-Execution}


\begin{coqdoccode}
\coqdocemptyline
\coqdocnoindent
\coqdockw{Set} \coqdocvar{Warnings} "-notation-overridden,-parsing,-deprecated-hint-without-locality".\coqdoceol
\coqdocnoindent
\coqdockw{From} \coqdocvar{Stdlib} \coqdockw{Require} \coqdockw{Import} \coqdoclibrary{Bool.Bool}.\coqdoceol
\coqdocnoindent
\coqdockw{From} \coqdocvar{Stdlib} \coqdockw{Require} \coqdockw{Import} \coqdoclibrary{Init.Nat}.\coqdoceol
\coqdocnoindent
\coqdockw{From} \coqdocvar{Stdlib} \coqdockw{Require} \coqdockw{Import} \coqdoclibrary{Arith.Arith}.\coqdoceol
\coqdocnoindent
\coqdockw{From} \coqdocvar{Stdlib} \coqdockw{Require} \coqdockw{Import} \coqdoclibrary{Arith.EqNat}. \coqdockw{Import} \coqdocmodule{Nat}.\coqdoceol
\coqdocnoindent
\coqdockw{From} \coqdocvar{SECF} \coqdockw{Require} \coqdockw{Import} \coqref{SECF.Maps}{\coqdoclibrary{Maps}}.\coqdoceol
\coqdocnoindent
\coqdockw{From} \coqdocvar{SECF} \coqdockw{Require} \coqdockw{Import} \coqref{SECF.Imp}{\coqdoclibrary{Imp}}.\coqdoceol
\coqdocnoindent
\coqdockw{Set} \coqdocvar{Default} \coqdockw{Goal} \coqdocvar{Selector} "!".\coqdoceol
\coqdocemptyline
\coqdocnoindent
\coqdockw{From} \coqdocvar{Stdlib} \coqdockw{Require} \coqdockw{Import} \coqdoclibrary{Lia}.\coqdoceol
\coqdocemptyline
\end{coqdoccode}
Programmers have to be very careful about how information flows in
    the software they develop to prevent leaking secret data. For
    instance, in course management systems students shouldn't be able
    to obtain information about other student's grades. In crypto
    protocols the keys should be kept secret and not sent over the
    network in the clear. 

 Information-flow control tries to prevent leaking secret
    information.  But how does one formalize that a program doesn't
    leak any information about the secret inputs to public outputs? 

 We first investigate this question in the very simple setting of Rocq
    functions taking two arguments, one we call the public input and the other
    one we call the secret input. Our functions return a pair where the first
    element is the public output and the second one the secret output. 

 Say we have the following function working on natural numbers: 
\begin{coqdoccode}
\coqdocemptyline
\coqdocnoindent
\coqdockw{Definition} \coqdef{SECF.Noninterference.secure f}{secure\_f}{\coqdocdefinition{secure\_f}} (\coqdef{SECF.Noninterference.pi:1}{pi}{\coqdocbinder{pi}} \coqdef{SECF.Noninterference.si:2}{si}{\coqdocbinder{si}} : \coqexternalref{nat}{http://rocq-prover.org/doc/V9.1.1/stdlib//Corelib.Init.Datatypes}{\coqdocinductive{nat}}) : \coqexternalref{nat}{http://rocq-prover.org/doc/V9.1.1/stdlib//Corelib.Init.Datatypes}{\coqdocinductive{nat}}\coqexternalref{::type scope:x '*' x}{http://rocq-prover.org/doc/V9.1.1/stdlib//Corelib.Init.Datatypes}{\coqdocnotation{\ensuremath{\times}}}\coqexternalref{nat}{http://rocq-prover.org/doc/V9.1.1/stdlib//Corelib.Init.Datatypes}{\coqdocinductive{nat}} := \coqexternalref{::core scope:'(' x ',' x ',' '..' ',' x ')'}{http://rocq-prover.org/doc/V9.1.1/stdlib//Corelib.Init.Datatypes}{\coqdocnotation{(}}\coqref{SECF.Noninterference.pi:1}{\coqdocvariable{pi}}\coqexternalref{::nat scope:x '+' x}{http://rocq-prover.org/doc/V9.1.1/stdlib//Corelib.Init.Nat}{\coqdocnotation{+}}1\coqexternalref{::core scope:'(' x ',' x ',' '..' ',' x ')'}{http://rocq-prover.org/doc/V9.1.1/stdlib//Corelib.Init.Datatypes}{\coqdocnotation{,}} \coqref{SECF.Noninterference.pi:1}{\coqdocvariable{pi}}\coqexternalref{::nat scope:x '+' x}{http://rocq-prover.org/doc/V9.1.1/stdlib//Corelib.Init.Nat}{\coqdocnotation{+}}\coqref{SECF.Noninterference.si:2}{\coqdocvariable{si}}\coqexternalref{::nat scope:x '*' x}{http://rocq-prover.org/doc/V9.1.1/stdlib//Corelib.Init.Nat}{\coqdocnotation{\ensuremath{\times}}}2\coqexternalref{::core scope:'(' x ',' x ',' '..' ',' x ')'}{http://rocq-prover.org/doc/V9.1.1/stdlib//Corelib.Init.Datatypes}{\coqdocnotation{)}}.\coqdoceol
\coqdocemptyline
\end{coqdoccode}
This function seems intuitively secure, since the first output \coqdocvar{pi}+1, which
    we assume to be public, only depends on the public input \coqdocvar{pi}, but not on
    the secret input \coqdocvar{si}. The second output \coqdocvar{pi}+\coqdocvar{si}*2 depends on both the
    public input and the secret input, but that's okay, since we assume this
    second output to be secret. 

 Still, how can we mathematically define that this function is
    secure? Let's try it on a couple of inputs: 
\begin{coqdoccode}
\coqdocemptyline
\coqdocnoindent
\coqdockw{Example} \coqdef{SECF.Noninterference.example1 secure f}{example1\_secure\_f}{\coqdocdefinition{example1\_secure\_f}} : \coqref{SECF.Noninterference.secure f}{\coqdocdefinition{secure\_f}} 0 0 \coqexternalref{::type scope:x '=' x}{http://rocq-prover.org/doc/V9.1.1/stdlib//Corelib.Init.Logic}{\coqdocnotation{=}} \coqexternalref{::core scope:'(' x ',' x ',' '..' ',' x ')'}{http://rocq-prover.org/doc/V9.1.1/stdlib//Corelib.Init.Datatypes}{\coqdocnotation{(}}1\coqexternalref{::core scope:'(' x ',' x ',' '..' ',' x ')'}{http://rocq-prover.org/doc/V9.1.1/stdlib//Corelib.Init.Datatypes}{\coqdocnotation{,}}0\coqexternalref{::core scope:'(' x ',' x ',' '..' ',' x ')'}{http://rocq-prover.org/doc/V9.1.1/stdlib//Corelib.Init.Datatypes}{\coqdocnotation{)}}.\coqdoceol
\coqdocnoindent
\coqdockw{Proof}. \coqdoctac{reflexivity}. \coqdockw{Qed}.\coqdoceol
\coqdocemptyline
\coqdocnoindent
\coqdockw{Example} \coqdef{SECF.Noninterference.example2 secure f}{example2\_secure\_f}{\coqdocdefinition{example2\_secure\_f}} : \coqref{SECF.Noninterference.secure f}{\coqdocdefinition{secure\_f}} 0 1 \coqexternalref{::type scope:x '=' x}{http://rocq-prover.org/doc/V9.1.1/stdlib//Corelib.Init.Logic}{\coqdocnotation{=}} \coqexternalref{::core scope:'(' x ',' x ',' '..' ',' x ')'}{http://rocq-prover.org/doc/V9.1.1/stdlib//Corelib.Init.Datatypes}{\coqdocnotation{(}}1\coqexternalref{::core scope:'(' x ',' x ',' '..' ',' x ')'}{http://rocq-prover.org/doc/V9.1.1/stdlib//Corelib.Init.Datatypes}{\coqdocnotation{,}}2\coqexternalref{::core scope:'(' x ',' x ',' '..' ',' x ')'}{http://rocq-prover.org/doc/V9.1.1/stdlib//Corelib.Init.Datatypes}{\coqdocnotation{)}}.\coqdoceol
\coqdocnoindent
\coqdockw{Proof}. \coqdoctac{reflexivity}. \coqdockw{Qed}.\coqdoceol
\coqdocemptyline
\coqdocnoindent
\coqdockw{Example} \coqdef{SECF.Noninterference.example3 secure f}{example3\_secure\_f}{\coqdocdefinition{example3\_secure\_f}} : \coqref{SECF.Noninterference.secure f}{\coqdocdefinition{secure\_f}} 1 2 \coqexternalref{::type scope:x '=' x}{http://rocq-prover.org/doc/V9.1.1/stdlib//Corelib.Init.Logic}{\coqdocnotation{=}} \coqexternalref{::core scope:'(' x ',' x ',' '..' ',' x ')'}{http://rocq-prover.org/doc/V9.1.1/stdlib//Corelib.Init.Datatypes}{\coqdocnotation{(}}2\coqexternalref{::core scope:'(' x ',' x ',' '..' ',' x ')'}{http://rocq-prover.org/doc/V9.1.1/stdlib//Corelib.Init.Datatypes}{\coqdocnotation{,}}5\coqexternalref{::core scope:'(' x ',' x ',' '..' ',' x ')'}{http://rocq-prover.org/doc/V9.1.1/stdlib//Corelib.Init.Datatypes}{\coqdocnotation{)}}.\coqdoceol
\coqdocnoindent
\coqdockw{Proof}. \coqdoctac{reflexivity}. \coqdockw{Qed}.\coqdoceol
\coqdocemptyline
\end{coqdoccode}
In the last two cases the value of the public output is equal to the value
    of secret input. But that's just a coincidence, and has nothing to do with
    the public output leaking the secret input, which wasn't used at all in
    computing the public output. 
\begin{coqdoccode}
\coqdocemptyline
\end{coqdoccode}
\section{Naive attempt at defining secrecy}



 So a naive security definition, which we'll only use as a strawman, is one
    that simply requires that public outputs are different from secret inputs: 
\begin{coqdoccode}
\coqdocemptyline
\coqdocnoindent
\coqdockw{Definition} \coqdef{SECF.Noninterference.broken sec def}{broken\_sec\_def}{\coqdocdefinition{broken\_sec\_def}} (\coqdef{SECF.Noninterference.f:3}{f}{\coqdocbinder{f}} : \coqexternalref{nat}{http://rocq-prover.org/doc/V9.1.1/stdlib//Corelib.Init.Datatypes}{\coqdocinductive{nat}} \coqexternalref{::type scope:x '->' x}{http://rocq-prover.org/doc/V9.1.1/stdlib//Corelib.Init.Logic}{\coqdocnotation{\ensuremath{\rightarrow}}} \coqexternalref{nat}{http://rocq-prover.org/doc/V9.1.1/stdlib//Corelib.Init.Datatypes}{\coqdocinductive{nat}} \coqexternalref{::type scope:x '->' x}{http://rocq-prover.org/doc/V9.1.1/stdlib//Corelib.Init.Logic}{\coqdocnotation{\ensuremath{\rightarrow}}} \coqexternalref{nat}{http://rocq-prover.org/doc/V9.1.1/stdlib//Corelib.Init.Datatypes}{\coqdocinductive{nat}}\coqexternalref{::type scope:x '*' x}{http://rocq-prover.org/doc/V9.1.1/stdlib//Corelib.Init.Datatypes}{\coqdocnotation{\ensuremath{\times}}}\coqexternalref{nat}{http://rocq-prover.org/doc/V9.1.1/stdlib//Corelib.Init.Datatypes}{\coqdocinductive{nat}}) :=\coqdoceol
\coqdocindent{1.00em}
\coqdockw{\ensuremath{\forall}} \coqdef{SECF.Noninterference.pi:4}{pi}{\coqdocbinder{pi}} \coqdef{SECF.Noninterference.si:5}{si}{\coqdocbinder{si}}, \coqexternalref{fst}{http://rocq-prover.org/doc/V9.1.1/stdlib//Corelib.Init.Datatypes}{\coqdocdefinition{fst}} (\coqref{SECF.Noninterference.f:3}{\coqdocvariable{f}} \coqref{SECF.Noninterference.pi:4}{\coqdocvariable{pi}} \coqref{SECF.Noninterference.si:5}{\coqdocvariable{si}}) \coqexternalref{::type scope:x '<>' x}{http://rocq-prover.org/doc/V9.1.1/stdlib//Corelib.Init.Logic}{\coqdocnotation{\ensuremath{\not=}}} \coqref{SECF.Noninterference.si:5}{\coqdocvariable{si}}.\coqdoceol
\coqdocemptyline
\end{coqdoccode}
As discussed above, this definition would reject our secure
    function above as insecure: 
\begin{coqdoccode}
\coqdocemptyline
\coqdocnoindent
\coqdockw{Lemma} \coqdef{SECF.Noninterference.broken sec def rejects secure f}{broken\_sec\_def\_rejects\_secure\_f}{\coqdoclemma{broken\_sec\_def\_rejects\_secure\_f}} : \coqexternalref{::type scope:'x7E' x}{http://rocq-prover.org/doc/V9.1.1/stdlib//Corelib.Init.Logic}{\coqdocnotation{\ensuremath{\lnot}}}\coqref{SECF.Noninterference.broken sec def}{\coqdocdefinition{broken\_sec\_def}} \coqref{SECF.Noninterference.secure f}{\coqdocdefinition{secure\_f}}.\coqdoceol
\coqdocnoindent
\coqdockw{Proof}. \coqdoctac{intros} \coqdocvar{Hc}. \coqdoctac{apply} (\coqdocvar{Hc} 0 1). \coqdoctac{reflexivity}. \coqdockw{Qed}.\coqdoceol
\coqdocemptyline
\end{coqdoccode}
Even worse, this broken definition of security would allow insecure
    functions, such as the following one whose public output is \coqdocvar{si}+1: 
\begin{coqdoccode}
\coqdocemptyline
\coqdocnoindent
\coqdockw{Definition} \coqdef{SECF.Noninterference.insecure f}{insecure\_f}{\coqdocdefinition{insecure\_f}} (\coqdef{SECF.Noninterference.pi:6}{pi}{\coqdocbinder{pi}} \coqdef{SECF.Noninterference.si:7}{si}{\coqdocbinder{si}} : \coqexternalref{nat}{http://rocq-prover.org/doc/V9.1.1/stdlib//Corelib.Init.Datatypes}{\coqdocinductive{nat}}) : \coqexternalref{nat}{http://rocq-prover.org/doc/V9.1.1/stdlib//Corelib.Init.Datatypes}{\coqdocinductive{nat}}\coqexternalref{::type scope:x '*' x}{http://rocq-prover.org/doc/V9.1.1/stdlib//Corelib.Init.Datatypes}{\coqdocnotation{\ensuremath{\times}}}\coqexternalref{nat}{http://rocq-prover.org/doc/V9.1.1/stdlib//Corelib.Init.Datatypes}{\coqdocinductive{nat}} := \coqexternalref{::core scope:'(' x ',' x ',' '..' ',' x ')'}{http://rocq-prover.org/doc/V9.1.1/stdlib//Corelib.Init.Datatypes}{\coqdocnotation{(}}\coqref{SECF.Noninterference.si:7}{\coqdocvariable{si}}\coqexternalref{::nat scope:x '+' x}{http://rocq-prover.org/doc/V9.1.1/stdlib//Corelib.Init.Nat}{\coqdocnotation{+}}1\coqexternalref{::core scope:'(' x ',' x ',' '..' ',' x ')'}{http://rocq-prover.org/doc/V9.1.1/stdlib//Corelib.Init.Datatypes}{\coqdocnotation{,}} \coqref{SECF.Noninterference.pi:6}{\coqdocvariable{pi}}\coqexternalref{::nat scope:x '+' x}{http://rocq-prover.org/doc/V9.1.1/stdlib//Corelib.Init.Nat}{\coqdocnotation{+}}\coqref{SECF.Noninterference.si:7}{\coqdocvariable{si}}\coqexternalref{::nat scope:x '*' x}{http://rocq-prover.org/doc/V9.1.1/stdlib//Corelib.Init.Nat}{\coqdocnotation{\ensuremath{\times}}}2\coqexternalref{::core scope:'(' x ',' x ',' '..' ',' x ')'}{http://rocq-prover.org/doc/V9.1.1/stdlib//Corelib.Init.Datatypes}{\coqdocnotation{)}}.\coqdoceol
\coqdocemptyline
\end{coqdoccode}
This function's public output is never equal to its secret input, yet an
    attacker can easily compute one from the other by just subtracting 1. So
    the secret is entirely leaked, yet our broken definition accepts this: 
\begin{coqdoccode}
\coqdocemptyline
\coqdocnoindent
\coqdockw{Lemma} \coqdef{SECF.Noninterference.broken sec def accepts insecure f}{broken\_sec\_def\_accepts\_insecure\_f}{\coqdoclemma{broken\_sec\_def\_accepts\_insecure\_f}} : \coqref{SECF.Noninterference.broken sec def}{\coqdocdefinition{broken\_sec\_def}} \coqref{SECF.Noninterference.insecure f}{\coqdocdefinition{insecure\_f}}.\coqdoceol
\coqdocnoindent
\coqdockw{Proof}.\coqdoceol
\coqdocindent{1.00em}
\coqdoctac{unfold} \coqref{SECF.Noninterference.broken sec def}{\coqdocdefinition{broken\_sec\_def}}. \coqdoctac{intros} \coqdocvar{pi} \coqdocvar{si}. \coqdoctac{induction} \coqdocvar{si} \coqdockw{as} [| \coqdocvar{si'} \coqdocvar{IH}].\coqdoceol
\coqdocindent{1.00em}
- \coqdoctac{simpl}. \coqdoctac{intros} \coqdocvar{contra}. \coqdoctac{discriminate} \coqdocvar{contra}.\coqdoceol
\coqdocindent{1.00em}
- \coqdoctac{simpl} \coqdoctac{in} *. \coqdoctac{intro} \coqdocvar{Hc}. \coqdoctac{injection} \coqdocvar{Hc} \coqdockw{as} \coqdocvar{Hc}. \coqdoctac{apply} \coqdocvar{IH}. \coqdoctac{apply} \coqdocvar{Hc}.\coqdoceol
\coqdocnoindent
\coqdockw{Qed}.\coqdoceol
\coqdocemptyline
\end{coqdoccode}
This attempt at defining secure information flow by looking at how
    inputs and outputs are related for a single execution of the
    program was a complete failure. In fact, it is well known in the
    formal security research community that secure information flow
    \textit{cannot} be defined by looking at just one single program execution. 
\begin{coqdoccode}
\coqdocemptyline
\end{coqdoccode}
\section{Noninterference for pure functions}



 The simplest correct way to define secure information flow is a
    property called \textit{noninterference} \coqdocvar{Sabelfeld} \coqexternalref{and}{http://rocq-prover.org/doc/V9.1.1/stdlib//Corelib.Init.Logic}{\coqdocinductive{and}} \coqdocvar{Myers} 2003 (in Bib.v),
    which in its most standard form looks at \textit{two} program executions:
    for two different secret inputs the public outputs should not change: 
\begin{coqdoccode}
\coqdocemptyline
\coqdocnoindent
\coqdockw{Definition} \coqdef{SECF.Noninterference.noninterferent}{noninterferent}{\coqdocdefinition{noninterferent}} \{\coqdef{SECF.Noninterference.PI:8}{PI}{\coqdocbinder{PI}} \coqdef{SECF.Noninterference.SI:9}{SI}{\coqdocbinder{SI}} \coqdef{SECF.Noninterference.PO:10}{PO}{\coqdocbinder{PO}} \coqdef{SECF.Noninterference.SO:11}{SO}{\coqdocbinder{SO}} : \coqdockw{Type}\} (\coqdef{SECF.Noninterference.f:12}{f}{\coqdocbinder{f}}:\coqref{SECF.Noninterference.PI:8}{\coqdocvariable{PI}}\coqexternalref{::type scope:x '->' x}{http://rocq-prover.org/doc/V9.1.1/stdlib//Corelib.Init.Logic}{\coqdocnotation{\ensuremath{\rightarrow}}}\coqref{SECF.Noninterference.SI:9}{\coqdocvariable{SI}}\coqexternalref{::type scope:x '->' x}{http://rocq-prover.org/doc/V9.1.1/stdlib//Corelib.Init.Logic}{\coqdocnotation{\ensuremath{\rightarrow}}}\coqref{SECF.Noninterference.PO:10}{\coqdocvariable{PO}}\coqexternalref{::type scope:x '*' x}{http://rocq-prover.org/doc/V9.1.1/stdlib//Corelib.Init.Datatypes}{\coqdocnotation{\ensuremath{\times}}}\coqref{SECF.Noninterference.SO:11}{\coqdocvariable{SO}}) :=\coqdoceol
\coqdocindent{1.00em}
\coqdockw{\ensuremath{\forall}} (\coqdef{SECF.Noninterference.pi:13}{pi}{\coqdocbinder{pi}}:\coqref{SECF.Noninterference.PI:8}{\coqdocvariable{PI}}) (\coqdef{SECF.Noninterference.si1:14}{si1}{\coqdocbinder{si1}} \coqdef{SECF.Noninterference.si2:15}{si2}{\coqdocbinder{si2}}:\coqref{SECF.Noninterference.SI:9}{\coqdocvariable{SI}}), \coqexternalref{fst}{http://rocq-prover.org/doc/V9.1.1/stdlib//Corelib.Init.Datatypes}{\coqdocdefinition{fst}} (\coqref{SECF.Noninterference.f:12}{\coqdocvariable{f}} \coqref{SECF.Noninterference.pi:13}{\coqdocvariable{pi}} \coqref{SECF.Noninterference.si1:14}{\coqdocvariable{si1}}) \coqexternalref{::type scope:x '=' x}{http://rocq-prover.org/doc/V9.1.1/stdlib//Corelib.Init.Logic}{\coqdocnotation{=}} \coqexternalref{fst}{http://rocq-prover.org/doc/V9.1.1/stdlib//Corelib.Init.Datatypes}{\coqdocdefinition{fst}} (\coqref{SECF.Noninterference.f:12}{\coqdocvariable{f}} \coqref{SECF.Noninterference.pi:13}{\coqdocvariable{pi}} \coqref{SECF.Noninterference.si2:15}{\coqdocvariable{si2}}).\coqdoceol
\coqdocemptyline
\end{coqdoccode}
This definition prevents secret inputs from interfering with public
    outputs in any way. At the same time it allows secret inputs to
    influence secret outputs and also public inputs to influence both
    public and secret outputs:


                                ┌───╮
                                │ f │
                           pi ─>┼───┼─> po
                                │╲  │
                                │ ╲ │
                                │  ╲│
                           si ─>┼───┼─> so
                                └───╯


 The definition above defines noninterference for arbitrary types
    of inputs and outputs, so we can instantiate them to \coqexternalref{nat}{http://rocq-prover.org/doc/V9.1.1/stdlib//Corelib.Init.Datatypes}{\coqdocinductive{nat}} when
    looking at our example functions above: 
\begin{coqdoccode}
\coqdocemptyline
\coqdocnoindent
\coqdockw{Lemma} \coqdef{SECF.Noninterference.noninterferent secure f}{noninterferent\_secure\_f}{\coqdoclemma{noninterferent\_secure\_f}} : \coqref{SECF.Noninterference.noninterferent}{\coqdocdefinition{noninterferent}} \coqref{SECF.Noninterference.secure f}{\coqdocdefinition{secure\_f}}.\coqdoceol
\coqdocnoindent
\coqdockw{Proof}. \coqdoctac{unfold} \coqref{SECF.Noninterference.noninterferent}{\coqdocdefinition{noninterferent}}, \coqref{SECF.Noninterference.secure f}{\coqdocdefinition{secure\_f}}. \coqdoctac{simpl}. \coqdoctac{reflexivity}. \coqdockw{Qed}.\coqdoceol
\coqdocemptyline
\coqdocnoindent
\coqdockw{Lemma} \coqdef{SECF.Noninterference.interferent insecure f}{interferent\_insecure\_f}{\coqdoclemma{interferent\_insecure\_f}} : \coqexternalref{::type scope:'x7E' x}{http://rocq-prover.org/doc/V9.1.1/stdlib//Corelib.Init.Logic}{\coqdocnotation{\ensuremath{\lnot}}}\coqref{SECF.Noninterference.noninterferent}{\coqdocdefinition{noninterferent}} \coqref{SECF.Noninterference.insecure f}{\coqdocdefinition{insecure\_f}}.\coqdoceol
\coqdocnoindent
\coqdockw{Proof}.\coqdoceol
\coqdocindent{1.00em}
\coqdoctac{unfold} \coqref{SECF.Noninterference.noninterferent}{\coqdocdefinition{noninterferent}}. \coqdoctac{simpl}. \coqdoctac{intros} \coqdocvar{contra}.\coqdoceol
\coqdocindent{1.00em}
\coqdoctac{specialize} (\coqdocvar{contra} 42 0 1). \coqdoctac{simpl} \coqdoctac{in} \coqdocvar{contra}. \coqdoctac{discriminate} \coqdocvar{contra}.\coqdoceol
\coqdocnoindent
\coqdockw{Qed}.\coqdoceol
\coqdocemptyline
\end{coqdoccode}
The \coqref{SECF.Noninterference.secure f}{\coqdocdefinition{secure\_f}} function above is quite obviously noninterferent,
    because the expression \coqdocvar{pi}+1 computing the public output doesn't
    syntactically mention the secret input at all. Since
    noninterference is a semantic property though (not a syntactic
    one), functions where the expression computing the public input
    does syntactically mention the secret input can still be
    noninterferent. Here is a first example: 
\begin{coqdoccode}
\coqdocemptyline
\coqdocnoindent
\coqdockw{Definition} \coqdef{SECF.Noninterference.less obvious f1}{less\_obvious\_f1}{\coqdocdefinition{less\_obvious\_f1}} (\coqdef{SECF.Noninterference.pi:16}{pi}{\coqdocbinder{pi}} \coqdef{SECF.Noninterference.si:17}{si}{\coqdocbinder{si}} : \coqexternalref{nat}{http://rocq-prover.org/doc/V9.1.1/stdlib//Corelib.Init.Datatypes}{\coqdocinductive{nat}}) : \coqexternalref{nat}{http://rocq-prover.org/doc/V9.1.1/stdlib//Corelib.Init.Datatypes}{\coqdocinductive{nat}}\coqexternalref{::type scope:x '*' x}{http://rocq-prover.org/doc/V9.1.1/stdlib//Corelib.Init.Datatypes}{\coqdocnotation{\ensuremath{\times}}}\coqexternalref{nat}{http://rocq-prover.org/doc/V9.1.1/stdlib//Corelib.Init.Datatypes}{\coqdocinductive{nat}} := \coqexternalref{::core scope:'(' x ',' x ',' '..' ',' x ')'}{http://rocq-prover.org/doc/V9.1.1/stdlib//Corelib.Init.Datatypes}{\coqdocnotation{(}}\coqref{SECF.Noninterference.si:17}{\coqdocvariable{si}} \coqexternalref{::nat scope:x '*' x}{http://rocq-prover.org/doc/V9.1.1/stdlib//Corelib.Init.Nat}{\coqdocnotation{\ensuremath{\times}}} 0\coqexternalref{::core scope:'(' x ',' x ',' '..' ',' x ')'}{http://rocq-prover.org/doc/V9.1.1/stdlib//Corelib.Init.Datatypes}{\coqdocnotation{,}} \coqref{SECF.Noninterference.pi:16}{\coqdocvariable{pi}}\coqexternalref{::nat scope:x '+' x}{http://rocq-prover.org/doc/V9.1.1/stdlib//Corelib.Init.Nat}{\coqdocnotation{+}}\coqref{SECF.Noninterference.si:17}{\coqdocvariable{si}}\coqexternalref{::core scope:'(' x ',' x ',' '..' ',' x ')'}{http://rocq-prover.org/doc/V9.1.1/stdlib//Corelib.Init.Datatypes}{\coqdocnotation{)}}.\coqdoceol
\coqdocemptyline
\end{coqdoccode}
This function is noninterferent; since the public output is
    constant 0, so it can't depend on \coqdocvar{si}, even if it syntactically
    mentions it: 
\begin{coqdoccode}
\coqdocemptyline
\coqdocnoindent
\coqdockw{Lemma} \coqdef{SECF.Noninterference.noninterferent less obvious f1}{noninterferent\_less\_obvious\_f1}{\coqdoclemma{noninterferent\_less\_obvious\_f1}} : \coqref{SECF.Noninterference.noninterferent}{\coqdocdefinition{noninterferent}} \coqref{SECF.Noninterference.less obvious f1}{\coqdocdefinition{less\_obvious\_f1}}.\coqdoceol
\coqdocnoindent
\coqdockw{Proof}.\coqdoceol
\coqdocindent{1.00em}
\coqdoctac{unfold} \coqref{SECF.Noninterference.noninterferent}{\coqdocdefinition{noninterferent}}, \coqref{SECF.Noninterference.less obvious f1}{\coqdocdefinition{less\_obvious\_f1}}. \coqdoctac{intros} \coqdocvar{pi} \coqdocvar{si1} \coqdocvar{si2}.\coqdoceol
\coqdocindent{1.00em}
\coqdoctac{simpl}. \coqdoctac{repeat} \coqdoctac{rewrite} \ensuremath{\leftarrow} \coqexternalref{mult n O}{http://rocq-prover.org/doc/V9.1.1/stdlib//Corelib.Init.Peano}{\coqdoclemma{mult\_n\_O}}. \coqdoctac{reflexivity}.\coqdoceol
\coqdocnoindent
\coqdockw{Qed}.\coqdoceol
\coqdocemptyline
\end{coqdoccode}
Here is another example of a function that is noninterferent, even
    if this is not syntactically obvious: 
\begin{coqdoccode}
\coqdocemptyline
\coqdocnoindent
\coqdockw{Definition} \coqdef{SECF.Noninterference.less obvious f2}{less\_obvious\_f2}{\coqdocdefinition{less\_obvious\_f2}} (\coqdef{SECF.Noninterference.pi:18}{pi}{\coqdocbinder{pi}} \coqdef{SECF.Noninterference.si:19}{si}{\coqdocbinder{si}} : \coqexternalref{nat}{http://rocq-prover.org/doc/V9.1.1/stdlib//Corelib.Init.Datatypes}{\coqdocinductive{nat}}) : \coqexternalref{nat}{http://rocq-prover.org/doc/V9.1.1/stdlib//Corelib.Init.Datatypes}{\coqdocinductive{nat}}\coqexternalref{::type scope:x '*' x}{http://rocq-prover.org/doc/V9.1.1/stdlib//Corelib.Init.Datatypes}{\coqdocnotation{\ensuremath{\times}}}\coqexternalref{nat}{http://rocq-prover.org/doc/V9.1.1/stdlib//Corelib.Init.Datatypes}{\coqdocinductive{nat}} :=\coqdoceol
\coqdocindent{1.00em}
\coqexternalref{::core scope:'(' x ',' x ',' '..' ',' x ')'}{http://rocq-prover.org/doc/V9.1.1/stdlib//Corelib.Init.Datatypes}{\coqdocnotation{(}}\coqdockw{if} \coqdocdefinition{Nat.eqb} \coqref{SECF.Noninterference.si:19}{\coqdocvariable{si}} 1 \coqdockw{then} \coqref{SECF.Noninterference.si:19}{\coqdocvariable{si}} \coqexternalref{::nat scope:x '*' x}{http://rocq-prover.org/doc/V9.1.1/stdlib//Corelib.Init.Nat}{\coqdocnotation{\ensuremath{\times}}} \coqref{SECF.Noninterference.pi:18}{\coqdocvariable{pi}} \coqdockw{else} \coqref{SECF.Noninterference.pi:18}{\coqdocvariable{pi}}\coqexternalref{::core scope:'(' x ',' x ',' '..' ',' x ')'}{http://rocq-prover.org/doc/V9.1.1/stdlib//Corelib.Init.Datatypes}{\coqdocnotation{,}} \coqref{SECF.Noninterference.pi:18}{\coqdocvariable{pi}}\coqexternalref{::nat scope:x '+' x}{http://rocq-prover.org/doc/V9.1.1/stdlib//Corelib.Init.Nat}{\coqdocnotation{+}}\coqref{SECF.Noninterference.si:19}{\coqdocvariable{si}}\coqexternalref{::core scope:'(' x ',' x ',' '..' ',' x ')'}{http://rocq-prover.org/doc/V9.1.1/stdlib//Corelib.Init.Datatypes}{\coqdocnotation{)}}.\coqdoceol
\coqdocemptyline
\end{coqdoccode}
For proving this we show that the public output of this function
    is in fact always equal to just its public input: 
\begin{coqdoccode}
\coqdocemptyline
\coqdocnoindent
\coqdockw{Lemma} \coqdef{SECF.Noninterference.aux f2}{aux\_f2}{\coqdoclemma{aux\_f2}} : \coqdockw{\ensuremath{\forall}} \coqdef{SECF.Noninterference.si:20}{si}{\coqdocbinder{si}} \coqdef{SECF.Noninterference.pi:21}{pi}{\coqdocbinder{pi}}, \coqexternalref{::type scope:x '=' x}{http://rocq-prover.org/doc/V9.1.1/stdlib//Corelib.Init.Logic}{\coqdocnotation{(}}\coqdockw{if} \coqdocdefinition{Nat.eqb} \coqref{SECF.Noninterference.si:20}{\coqdocvariable{si}} 1 \coqdockw{then} \coqref{SECF.Noninterference.si:20}{\coqdocvariable{si}} \coqexternalref{::nat scope:x '*' x}{http://rocq-prover.org/doc/V9.1.1/stdlib//Corelib.Init.Nat}{\coqdocnotation{\ensuremath{\times}}} \coqref{SECF.Noninterference.pi:21}{\coqdocvariable{pi}} \coqdockw{else} \coqref{SECF.Noninterference.pi:21}{\coqdocvariable{pi}}\coqexternalref{::type scope:x '=' x}{http://rocq-prover.org/doc/V9.1.1/stdlib//Corelib.Init.Logic}{\coqdocnotation{)}} \coqexternalref{::type scope:x '=' x}{http://rocq-prover.org/doc/V9.1.1/stdlib//Corelib.Init.Logic}{\coqdocnotation{=}} \coqref{SECF.Noninterference.pi:21}{\coqdocvariable{pi}}.\coqdoceol
\coqdocnoindent
\coqdockw{Proof}.\coqdoceol
\coqdocindent{1.00em}
\coqdoctac{intros} \coqdocvar{si} \coqdocvar{pi}. \coqdoctac{destruct} \coqdocvar{si}; \coqdoctac{simpl}.\coqdoceol
\coqdocindent{1.00em}
- \coqdoctac{reflexivity}.\coqdoceol
\coqdocindent{1.00em}
- \coqdoctac{destruct} \coqdocvar{si}.\coqdoceol
\coqdocindent{2.00em}
+ \coqdoctac{simpl}. \coqdoctac{rewrite} \ensuremath{\leftarrow} \coqexternalref{plus n O}{http://rocq-prover.org/doc/V9.1.1/stdlib//Corelib.Init.Peano}{\coqdoclemma{plus\_n\_O}}. \coqdoctac{reflexivity}.\coqdoceol
\coqdocindent{2.00em}
+ \coqdoctac{simpl}. \coqdoctac{reflexivity}.\coqdoceol
\coqdocnoindent
\coqdockw{Qed}.\coqdoceol
\coqdocemptyline
\coqdocnoindent
\coqdockw{Lemma} \coqdef{SECF.Noninterference.noninterferent less obvious f2}{noninterferent\_less\_obvious\_f2}{\coqdoclemma{noninterferent\_less\_obvious\_f2}} : \coqref{SECF.Noninterference.noninterferent}{\coqdocdefinition{noninterferent}} \coqref{SECF.Noninterference.less obvious f2}{\coqdocdefinition{less\_obvious\_f2}}.\coqdoceol
\coqdocnoindent
\coqdockw{Proof}.\coqdoceol
\coqdocindent{1.00em}
\coqdoctac{unfold} \coqref{SECF.Noninterference.noninterferent}{\coqdocdefinition{noninterferent}}, \coqref{SECF.Noninterference.less obvious f2}{\coqdocdefinition{less\_obvious\_f2}}. \coqdoctac{intros} \coqdocvar{pi} \coqdocvar{si1} \coqdocvar{si2}.\coqdoceol
\coqdocindent{1.00em}
\coqdoctac{repeat} \coqdoctac{rewrite} \coqref{SECF.Noninterference.aux f2}{\coqdoclemma{aux\_f2}}. \coqdoctac{simpl}. \coqdoctac{reflexivity}.\coqdoceol
\coqdocnoindent
\coqdockw{Qed}.\coqdoceol
\coqdocemptyline
\end{coqdoccode}
Branching on a secret can, however, be dangerous, since one can
    easily leak the secret this way, even if both the \coqdockw{then} and the
    \coqdockw{else} branches are public. For instance the following function
    leaks whether \coqdocvar{si} is zero or not, so it is not noninterferent. 
\begin{coqdoccode}
\coqdocemptyline
\coqdocnoindent
\coqdockw{Definition} \coqdef{SECF.Noninterference.less obvious f3}{less\_obvious\_f3}{\coqdocdefinition{less\_obvious\_f3}} (\coqdef{SECF.Noninterference.pi:22}{pi}{\coqdocbinder{pi}} \coqdef{SECF.Noninterference.si:23}{si}{\coqdocbinder{si}} : \coqexternalref{nat}{http://rocq-prover.org/doc/V9.1.1/stdlib//Corelib.Init.Datatypes}{\coqdocinductive{nat}}) : \coqexternalref{nat}{http://rocq-prover.org/doc/V9.1.1/stdlib//Corelib.Init.Datatypes}{\coqdocinductive{nat}}\coqexternalref{::type scope:x '*' x}{http://rocq-prover.org/doc/V9.1.1/stdlib//Corelib.Init.Datatypes}{\coqdocnotation{\ensuremath{\times}}}\coqexternalref{nat}{http://rocq-prover.org/doc/V9.1.1/stdlib//Corelib.Init.Datatypes}{\coqdocinductive{nat}} :=\coqdoceol
\coqdocindent{1.00em}
\coqexternalref{::core scope:'(' x ',' x ',' '..' ',' x ')'}{http://rocq-prover.org/doc/V9.1.1/stdlib//Corelib.Init.Datatypes}{\coqdocnotation{(}}\coqdockw{if} \coqdocdefinition{Nat.eqb} \coqref{SECF.Noninterference.si:23}{\coqdocvariable{si}} 0 \coqdockw{then} 1 \coqdockw{else} 0\coqexternalref{::core scope:'(' x ',' x ',' '..' ',' x ')'}{http://rocq-prover.org/doc/V9.1.1/stdlib//Corelib.Init.Datatypes}{\coqdocnotation{,}} \coqref{SECF.Noninterference.pi:22}{\coqdocvariable{pi}}\coqexternalref{::nat scope:x '+' x}{http://rocq-prover.org/doc/V9.1.1/stdlib//Corelib.Init.Nat}{\coqdocnotation{+}}\coqref{SECF.Noninterference.si:23}{\coqdocvariable{si}}\coqexternalref{::core scope:'(' x ',' x ',' '..' ',' x ')'}{http://rocq-prover.org/doc/V9.1.1/stdlib//Corelib.Init.Datatypes}{\coqdocnotation{)}}.\coqdoceol
\coqdocemptyline
\coqdocnoindent
\coqdockw{Lemma} \coqdef{SECF.Noninterference.interferent less obvious f3}{interferent\_less\_obvious\_f3}{\coqdoclemma{interferent\_less\_obvious\_f3}} : \coqexternalref{::type scope:'x7E' x}{http://rocq-prover.org/doc/V9.1.1/stdlib//Corelib.Init.Logic}{\coqdocnotation{\ensuremath{\lnot}}}\coqref{SECF.Noninterference.noninterferent}{\coqdocdefinition{noninterferent}} \coqref{SECF.Noninterference.less obvious f3}{\coqdocdefinition{less\_obvious\_f3}}.\coqdoceol
\coqdocnoindent
\coqdockw{Proof}.\coqdoceol
\coqdocindent{1.00em}
\coqdoctac{unfold} \coqref{SECF.Noninterference.noninterferent}{\coqdocdefinition{noninterferent}}, \coqref{SECF.Noninterference.less obvious f3}{\coqdocdefinition{less\_obvious\_f3}}. \coqdoctac{simpl}. \coqdoctac{intros} \coqdocvar{contra}.\coqdoceol
\coqdocindent{1.00em}
\coqdoctac{specialize} (\coqdocvar{contra} 42 0 1). \coqdoctac{simpl} \coqdoctac{in} \coqdocvar{contra}. \coqdoctac{discriminate} \coqdocvar{contra}.\coqdoceol
\coqdocnoindent
\coqdockw{Qed}.\coqdoceol
\coqdocemptyline
\end{coqdoccode}
\subsection{Noninterference Exercises}



 Let's practice with some ``prove or disprove noninterference''
    exercises, for which you are required to give constructive proofs,
    i.e. the use of classical axioms like excluded middle is not allowed. 

\paragraph{Exercise: 1 star, standard (prove\_or\_disprove\_obvious\_f1)}


\begin{coqdoccode}
\coqdocnoindent
\coqdockw{Definition} \coqdef{SECF.Noninterference.obvious f1}{obvious\_f1}{\coqdocdefinition{obvious\_f1}} (\coqdef{SECF.Noninterference.pi:24}{pi}{\coqdocbinder{pi}} \coqdef{SECF.Noninterference.si:25}{si}{\coqdocbinder{si}} : \coqexternalref{nat}{http://rocq-prover.org/doc/V9.1.1/stdlib//Corelib.Init.Datatypes}{\coqdocinductive{nat}}) : \coqexternalref{nat}{http://rocq-prover.org/doc/V9.1.1/stdlib//Corelib.Init.Datatypes}{\coqdocinductive{nat}}\coqexternalref{::type scope:x '*' x}{http://rocq-prover.org/doc/V9.1.1/stdlib//Corelib.Init.Datatypes}{\coqdocnotation{\ensuremath{\times}}}\coqexternalref{nat}{http://rocq-prover.org/doc/V9.1.1/stdlib//Corelib.Init.Datatypes}{\coqdocinductive{nat}} := \coqexternalref{::core scope:'(' x ',' x ',' '..' ',' x ')'}{http://rocq-prover.org/doc/V9.1.1/stdlib//Corelib.Init.Datatypes}{\coqdocnotation{(}}0\coqexternalref{::core scope:'(' x ',' x ',' '..' ',' x ')'}{http://rocq-prover.org/doc/V9.1.1/stdlib//Corelib.Init.Datatypes}{\coqdocnotation{,}}0\coqexternalref{::core scope:'(' x ',' x ',' '..' ',' x ')'}{http://rocq-prover.org/doc/V9.1.1/stdlib//Corelib.Init.Datatypes}{\coqdocnotation{)}}.\coqdoceol
\coqdocemptyline
\coqdocnoindent
\coqdockw{Lemma} \coqdef{SECF.Noninterference.prove or disprove obvious f1}{prove\_or\_disprove\_obvious\_f1}{\coqdoclemma{prove\_or\_disprove\_obvious\_f1}} :\coqdoceol
\coqdocindent{1.00em}
\coqref{SECF.Noninterference.noninterferent}{\coqdocdefinition{noninterferent}} \coqref{SECF.Noninterference.obvious f1}{\coqdocdefinition{obvious\_f1}} \coqexternalref{::type scope:x 'x5C/' x}{http://rocq-prover.org/doc/V9.1.1/stdlib//Corelib.Init.Logic}{\coqdocnotation{\ensuremath{\lor}}} \coqexternalref{::type scope:'x7E' x}{http://rocq-prover.org/doc/V9.1.1/stdlib//Corelib.Init.Logic}{\coqdocnotation{\ensuremath{\lnot}}}\coqref{SECF.Noninterference.noninterferent}{\coqdocdefinition{noninterferent}} \coqref{SECF.Noninterference.obvious f1}{\coqdocdefinition{obvious\_f1}}.\coqdoceol
\coqdocnoindent
\coqdockw{Proof}.\coqdoceol
 \coqdocvar{Admitted}.\coqdoceol
\end{coqdoccode}
\ensuremath{\Box} 

\paragraph{Exercise: 1 star, standard (prove\_or\_disprove\_obvious\_f2)}


\begin{coqdoccode}
\coqdocnoindent
\coqdockw{Definition} \coqdef{SECF.Noninterference.obvious f2}{obvious\_f2}{\coqdocdefinition{obvious\_f2}} (\coqdef{SECF.Noninterference.pi:26}{pi}{\coqdocbinder{pi}} \coqdef{SECF.Noninterference.si:27}{si}{\coqdocbinder{si}} : \coqexternalref{nat}{http://rocq-prover.org/doc/V9.1.1/stdlib//Corelib.Init.Datatypes}{\coqdocinductive{nat}}) : \coqexternalref{nat}{http://rocq-prover.org/doc/V9.1.1/stdlib//Corelib.Init.Datatypes}{\coqdocinductive{nat}}\coqexternalref{::type scope:x '*' x}{http://rocq-prover.org/doc/V9.1.1/stdlib//Corelib.Init.Datatypes}{\coqdocnotation{\ensuremath{\times}}}\coqexternalref{nat}{http://rocq-prover.org/doc/V9.1.1/stdlib//Corelib.Init.Datatypes}{\coqdocinductive{nat}} := \coqexternalref{::core scope:'(' x ',' x ',' '..' ',' x ')'}{http://rocq-prover.org/doc/V9.1.1/stdlib//Corelib.Init.Datatypes}{\coqdocnotation{(}}\coqref{SECF.Noninterference.pi:26}{\coqdocvariable{pi}}\coqexternalref{::nat scope:x '+' x}{http://rocq-prover.org/doc/V9.1.1/stdlib//Corelib.Init.Nat}{\coqdocnotation{+(}}2\coqexternalref{::nat scope:x '*' x}{http://rocq-prover.org/doc/V9.1.1/stdlib//Corelib.Init.Nat}{\coqdocnotation{\ensuremath{\times}}}\coqref{SECF.Noninterference.si:27}{\coqdocvariable{si}}\coqexternalref{::nat scope:x '+' x}{http://rocq-prover.org/doc/V9.1.1/stdlib//Corelib.Init.Nat}{\coqdocnotation{)}}\coqexternalref{::core scope:'(' x ',' x ',' '..' ',' x ')'}{http://rocq-prover.org/doc/V9.1.1/stdlib//Corelib.Init.Datatypes}{\coqdocnotation{,}}\coqexternalref{::nat scope:x '+' x}{http://rocq-prover.org/doc/V9.1.1/stdlib//Corelib.Init.Nat}{\coqdocnotation{(}}2\coqexternalref{::nat scope:x '*' x}{http://rocq-prover.org/doc/V9.1.1/stdlib//Corelib.Init.Nat}{\coqdocnotation{\ensuremath{\times}}}\coqref{SECF.Noninterference.pi:26}{\coqdocvariable{pi}}\coqexternalref{::nat scope:x '+' x}{http://rocq-prover.org/doc/V9.1.1/stdlib//Corelib.Init.Nat}{\coqdocnotation{)+}}\coqref{SECF.Noninterference.si:27}{\coqdocvariable{si}}\coqexternalref{::core scope:'(' x ',' x ',' '..' ',' x ')'}{http://rocq-prover.org/doc/V9.1.1/stdlib//Corelib.Init.Datatypes}{\coqdocnotation{)}}.\coqdoceol
\coqdocemptyline
\coqdocnoindent
\coqdockw{Lemma} \coqdef{SECF.Noninterference.prove or disprove obvious f2}{prove\_or\_disprove\_obvious\_f2}{\coqdoclemma{prove\_or\_disprove\_obvious\_f2}} :\coqdoceol
\coqdocindent{1.00em}
\coqref{SECF.Noninterference.noninterferent}{\coqdocdefinition{noninterferent}} \coqref{SECF.Noninterference.obvious f2}{\coqdocdefinition{obvious\_f2}} \coqexternalref{::type scope:x 'x5C/' x}{http://rocq-prover.org/doc/V9.1.1/stdlib//Corelib.Init.Logic}{\coqdocnotation{\ensuremath{\lor}}} \coqexternalref{::type scope:'x7E' x}{http://rocq-prover.org/doc/V9.1.1/stdlib//Corelib.Init.Logic}{\coqdocnotation{\ensuremath{\lnot}}}\coqref{SECF.Noninterference.noninterferent}{\coqdocdefinition{noninterferent}} \coqref{SECF.Noninterference.obvious f2}{\coqdocdefinition{obvious\_f2}}.\coqdoceol
\coqdocnoindent
\coqdockw{Proof}.\coqdoceol
 \coqdocvar{Admitted}.\coqdoceol
\end{coqdoccode}
\ensuremath{\Box} 

\paragraph{Exercise: 2 stars, standard (prove\_or\_disprove\_less\_obvious\_f4)}


\begin{coqdoccode}
\coqdocemptyline
\coqdocnoindent
\coqdockw{Definition} \coqdef{SECF.Noninterference.less obvious f4}{less\_obvious\_f4}{\coqdocdefinition{less\_obvious\_f4}} (\coqdef{SECF.Noninterference.pi:28}{pi}{\coqdocbinder{pi}} \coqdef{SECF.Noninterference.si:29}{si}{\coqdocbinder{si}} : \coqexternalref{nat}{http://rocq-prover.org/doc/V9.1.1/stdlib//Corelib.Init.Datatypes}{\coqdocinductive{nat}}) : \coqexternalref{nat}{http://rocq-prover.org/doc/V9.1.1/stdlib//Corelib.Init.Datatypes}{\coqdocinductive{nat}}\coqexternalref{::type scope:x '*' x}{http://rocq-prover.org/doc/V9.1.1/stdlib//Corelib.Init.Datatypes}{\coqdocnotation{\ensuremath{\times}}}\coqexternalref{nat}{http://rocq-prover.org/doc/V9.1.1/stdlib//Corelib.Init.Datatypes}{\coqdocinductive{nat}} :=\coqdoceol
\coqdocindent{1.00em}
\coqexternalref{::core scope:'(' x ',' x ',' '..' ',' x ')'}{http://rocq-prover.org/doc/V9.1.1/stdlib//Corelib.Init.Datatypes}{\coqdocnotation{(}}\coqdockw{if} \coqdocdefinition{Nat.eqb} \coqref{SECF.Noninterference.si:29}{\coqdocvariable{si}} 0 \coqdockw{then} \coqref{SECF.Noninterference.si:29}{\coqdocvariable{si}} \coqexternalref{::nat scope:x '*' x}{http://rocq-prover.org/doc/V9.1.1/stdlib//Corelib.Init.Nat}{\coqdocnotation{\ensuremath{\times}}} \coqref{SECF.Noninterference.pi:28}{\coqdocvariable{pi}} \coqdockw{else} \coqref{SECF.Noninterference.pi:28}{\coqdocvariable{pi}}\coqexternalref{::core scope:'(' x ',' x ',' '..' ',' x ')'}{http://rocq-prover.org/doc/V9.1.1/stdlib//Corelib.Init.Datatypes}{\coqdocnotation{,}} \coqref{SECF.Noninterference.pi:28}{\coqdocvariable{pi}}\coqexternalref{::nat scope:x '+' x}{http://rocq-prover.org/doc/V9.1.1/stdlib//Corelib.Init.Nat}{\coqdocnotation{+}}\coqref{SECF.Noninterference.si:29}{\coqdocvariable{si}}\coqexternalref{::core scope:'(' x ',' x ',' '..' ',' x ')'}{http://rocq-prover.org/doc/V9.1.1/stdlib//Corelib.Init.Datatypes}{\coqdocnotation{)}}.\coqdoceol
\coqdocemptyline
\end{coqdoccode}
Is the \coqref{SECF.Noninterference.less obvious f4}{\coqdocdefinition{less\_obvious\_f4}} function noninterferent or not? 
\begin{coqdoccode}
\coqdocemptyline
\coqdocnoindent
\coqdockw{Lemma} \coqdef{SECF.Noninterference.prove or disprove less obvious f4}{prove\_or\_disprove\_less\_obvious\_f4}{\coqdoclemma{prove\_or\_disprove\_less\_obvious\_f4}} :\coqdoceol
\coqdocindent{1.00em}
\coqref{SECF.Noninterference.noninterferent}{\coqdocdefinition{noninterferent}} \coqref{SECF.Noninterference.less obvious f4}{\coqdocdefinition{less\_obvious\_f4}} \coqexternalref{::type scope:x 'x5C/' x}{http://rocq-prover.org/doc/V9.1.1/stdlib//Corelib.Init.Logic}{\coqdocnotation{\ensuremath{\lor}}} \coqexternalref{::type scope:'x7E' x}{http://rocq-prover.org/doc/V9.1.1/stdlib//Corelib.Init.Logic}{\coqdocnotation{\ensuremath{\lnot}}}\coqref{SECF.Noninterference.noninterferent}{\coqdocdefinition{noninterferent}} \coqref{SECF.Noninterference.less obvious f4}{\coqdocdefinition{less\_obvious\_f4}}.\coqdoceol
\coqdocnoindent
\coqdockw{Proof}.\coqdoceol
 \coqdocvar{Admitted}.\coqdoceol
\end{coqdoccode}
\ensuremath{\Box} 

\paragraph{Exercise: 2 stars, standard (prove\_or\_disprove\_less\_obvious\_f5)}


\begin{coqdoccode}
\coqdocemptyline
\coqdocnoindent
\coqdockw{Definition} \coqdef{SECF.Noninterference.less obvious f5}{less\_obvious\_f5}{\coqdocdefinition{less\_obvious\_f5}} (\coqdef{SECF.Noninterference.pi:30}{pi}{\coqdocbinder{pi}} \coqdef{SECF.Noninterference.si:31}{si}{\coqdocbinder{si}} : \coqexternalref{nat}{http://rocq-prover.org/doc/V9.1.1/stdlib//Corelib.Init.Datatypes}{\coqdocinductive{nat}}) : \coqexternalref{nat}{http://rocq-prover.org/doc/V9.1.1/stdlib//Corelib.Init.Datatypes}{\coqdocinductive{nat}}\coqexternalref{::type scope:x '*' x}{http://rocq-prover.org/doc/V9.1.1/stdlib//Corelib.Init.Datatypes}{\coqdocnotation{\ensuremath{\times}}}\coqexternalref{nat}{http://rocq-prover.org/doc/V9.1.1/stdlib//Corelib.Init.Datatypes}{\coqdocinductive{nat}} :=\coqdoceol
\coqdocindent{1.00em}
\coqexternalref{::core scope:'(' x ',' x ',' '..' ',' x ')'}{http://rocq-prover.org/doc/V9.1.1/stdlib//Corelib.Init.Datatypes}{\coqdocnotation{(}}\coqdockw{if} \coqdocdefinition{Nat.eqb} \coqref{SECF.Noninterference.si:31}{\coqdocvariable{si}} 0 \coqdockw{then} \coqref{SECF.Noninterference.si:31}{\coqdocvariable{si}} \coqexternalref{::nat scope:x '+' x}{http://rocq-prover.org/doc/V9.1.1/stdlib//Corelib.Init.Nat}{\coqdocnotation{+}} \coqref{SECF.Noninterference.pi:30}{\coqdocvariable{pi}} \coqdockw{else} \coqref{SECF.Noninterference.pi:30}{\coqdocvariable{pi}}\coqexternalref{::core scope:'(' x ',' x ',' '..' ',' x ')'}{http://rocq-prover.org/doc/V9.1.1/stdlib//Corelib.Init.Datatypes}{\coqdocnotation{,}} \coqref{SECF.Noninterference.pi:30}{\coqdocvariable{pi}}\coqexternalref{::nat scope:x '+' x}{http://rocq-prover.org/doc/V9.1.1/stdlib//Corelib.Init.Nat}{\coqdocnotation{+}}\coqref{SECF.Noninterference.si:31}{\coqdocvariable{si}}\coqexternalref{::core scope:'(' x ',' x ',' '..' ',' x ')'}{http://rocq-prover.org/doc/V9.1.1/stdlib//Corelib.Init.Datatypes}{\coqdocnotation{)}}.\coqdoceol
\coqdocemptyline
\end{coqdoccode}
Is the \coqref{SECF.Noninterference.less obvious f5}{\coqdocdefinition{less\_obvious\_f5}} function noninterferent or not? 
\begin{coqdoccode}
\coqdocemptyline
\coqdocnoindent
\coqdockw{Lemma} \coqdef{SECF.Noninterference.prove or disprove less obvious f5}{prove\_or\_disprove\_less\_obvious\_f5}{\coqdoclemma{prove\_or\_disprove\_less\_obvious\_f5}} :\coqdoceol
\coqdocindent{1.00em}
\coqref{SECF.Noninterference.noninterferent}{\coqdocdefinition{noninterferent}} \coqref{SECF.Noninterference.less obvious f5}{\coqdocdefinition{less\_obvious\_f5}} \coqexternalref{::type scope:x 'x5C/' x}{http://rocq-prover.org/doc/V9.1.1/stdlib//Corelib.Init.Logic}{\coqdocnotation{\ensuremath{\lor}}} \coqexternalref{::type scope:'x7E' x}{http://rocq-prover.org/doc/V9.1.1/stdlib//Corelib.Init.Logic}{\coqdocnotation{\ensuremath{\lnot}}}\coqref{SECF.Noninterference.noninterferent}{\coqdocdefinition{noninterferent}} \coqref{SECF.Noninterference.less obvious f5}{\coqdocdefinition{less\_obvious\_f5}}.\coqdoceol
\coqdocnoindent
\coqdockw{Proof}.\coqdoceol
 \coqdocvar{Admitted}.\coqdoceol
\end{coqdoccode}
\ensuremath{\Box} 

\paragraph{Exercise: 2 stars, standard (prove\_or\_disprove\_less\_obvious\_f6)}


\begin{coqdoccode}
\coqdocemptyline
\coqdocnoindent
\coqdockw{Definition} \coqdef{SECF.Noninterference.less obvious f6}{less\_obvious\_f6}{\coqdocdefinition{less\_obvious\_f6}} (\coqdef{SECF.Noninterference.pi:32}{pi}{\coqdocbinder{pi}} \coqdef{SECF.Noninterference.si:33}{si}{\coqdocbinder{si}} : \coqexternalref{nat}{http://rocq-prover.org/doc/V9.1.1/stdlib//Corelib.Init.Datatypes}{\coqdocinductive{nat}}): \coqexternalref{nat}{http://rocq-prover.org/doc/V9.1.1/stdlib//Corelib.Init.Datatypes}{\coqdocinductive{nat}}\coqexternalref{::type scope:x '*' x}{http://rocq-prover.org/doc/V9.1.1/stdlib//Corelib.Init.Datatypes}{\coqdocnotation{\ensuremath{\times}}}\coqexternalref{nat}{http://rocq-prover.org/doc/V9.1.1/stdlib//Corelib.Init.Datatypes}{\coqdocinductive{nat}} :=\coqdoceol
\coqdocindent{1.00em}
\coqexternalref{::core scope:'(' x ',' x ',' '..' ',' x ')'}{http://rocq-prover.org/doc/V9.1.1/stdlib//Corelib.Init.Datatypes}{\coqdocnotation{(}}\coqdockw{if} \coqdocdefinition{Nat.ltb} \coqref{SECF.Noninterference.si:33}{\coqdocvariable{si}} \coqref{SECF.Noninterference.pi:32}{\coqdocvariable{pi}} \coqdockw{then} 0 \coqdockw{else} \coqref{SECF.Noninterference.pi:32}{\coqdocvariable{pi}}\coqexternalref{::core scope:'(' x ',' x ',' '..' ',' x ')'}{http://rocq-prover.org/doc/V9.1.1/stdlib//Corelib.Init.Datatypes}{\coqdocnotation{,}} \coqref{SECF.Noninterference.pi:32}{\coqdocvariable{pi}}\coqexternalref{::nat scope:x '+' x}{http://rocq-prover.org/doc/V9.1.1/stdlib//Corelib.Init.Nat}{\coqdocnotation{+}}\coqref{SECF.Noninterference.si:33}{\coqdocvariable{si}}\coqexternalref{::core scope:'(' x ',' x ',' '..' ',' x ')'}{http://rocq-prover.org/doc/V9.1.1/stdlib//Corelib.Init.Datatypes}{\coqdocnotation{)}}.\coqdoceol
\coqdocemptyline
\end{coqdoccode}
Is the \coqref{SECF.Noninterference.less obvious f6}{\coqdocdefinition{less\_obvious\_f6}} function noninterferent or not? 
\begin{coqdoccode}
\coqdocemptyline
\coqdocnoindent
\coqdockw{Lemma} \coqdef{SECF.Noninterference.prove or disprove less obvious f6}{prove\_or\_disprove\_less\_obvious\_f6}{\coqdoclemma{prove\_or\_disprove\_less\_obvious\_f6}} :\coqdoceol
\coqdocindent{1.00em}
\coqref{SECF.Noninterference.noninterferent}{\coqdocdefinition{noninterferent}} \coqref{SECF.Noninterference.less obvious f6}{\coqdocdefinition{less\_obvious\_f6}} \coqexternalref{::type scope:x 'x5C/' x}{http://rocq-prover.org/doc/V9.1.1/stdlib//Corelib.Init.Logic}{\coqdocnotation{\ensuremath{\lor}}} \coqexternalref{::type scope:'x7E' x}{http://rocq-prover.org/doc/V9.1.1/stdlib//Corelib.Init.Logic}{\coqdocnotation{\ensuremath{\lnot}}}\coqref{SECF.Noninterference.noninterferent}{\coqdocdefinition{noninterferent}} \coqref{SECF.Noninterference.less obvious f6}{\coqdocdefinition{less\_obvious\_f6}}.\coqdoceol
\coqdocnoindent
\coqdockw{Proof}.\coqdoceol
 \coqdocvar{Admitted}.\coqdoceol
\end{coqdoccode}
\ensuremath{\Box} 

\paragraph{Exercise: 3 stars, standard, optional (prove\_or\_disprove\_less\_obvious\_f7)}


\begin{coqdoccode}
\coqdocemptyline
\coqdocnoindent
\coqdockw{Definition} \coqdef{SECF.Noninterference.less obvious f7}{less\_obvious\_f7}{\coqdocdefinition{less\_obvious\_f7}} (\coqdef{SECF.Noninterference.pi:34}{pi}{\coqdocbinder{pi}} \coqdef{SECF.Noninterference.si:35}{si}{\coqdocbinder{si}} : \coqexternalref{nat}{http://rocq-prover.org/doc/V9.1.1/stdlib//Corelib.Init.Datatypes}{\coqdocinductive{nat}}): \coqexternalref{nat}{http://rocq-prover.org/doc/V9.1.1/stdlib//Corelib.Init.Datatypes}{\coqdocinductive{nat}}\coqexternalref{::type scope:x '*' x}{http://rocq-prover.org/doc/V9.1.1/stdlib//Corelib.Init.Datatypes}{\coqdocnotation{\ensuremath{\times}}}\coqexternalref{nat}{http://rocq-prover.org/doc/V9.1.1/stdlib//Corelib.Init.Datatypes}{\coqdocinductive{nat}} :=\coqdoceol
\coqdocindent{1.00em}
\coqdockw{if} \coqdocdefinition{Nat.eqb} (\coqref{SECF.Noninterference.si:35}{\coqdocvariable{si}} \coqexternalref{::nat scope:x '+' x}{http://rocq-prover.org/doc/V9.1.1/stdlib//Corelib.Init.Nat}{\coqdocnotation{+}} \coqref{SECF.Noninterference.pi:34}{\coqdocvariable{pi}}) 0 \coqdockw{then} \coqexternalref{::core scope:'(' x ',' x ',' '..' ',' x ')'}{http://rocq-prover.org/doc/V9.1.1/stdlib//Corelib.Init.Datatypes}{\coqdocnotation{(}}\coqref{SECF.Noninterference.si:35}{\coqdocvariable{si}}\coqexternalref{::core scope:'(' x ',' x ',' '..' ',' x ')'}{http://rocq-prover.org/doc/V9.1.1/stdlib//Corelib.Init.Datatypes}{\coqdocnotation{,}}\coqref{SECF.Noninterference.pi:34}{\coqdocvariable{pi}}\coqexternalref{::core scope:'(' x ',' x ',' '..' ',' x ')'}{http://rocq-prover.org/doc/V9.1.1/stdlib//Corelib.Init.Datatypes}{\coqdocnotation{)}} \coqdockw{else} \coqexternalref{::core scope:'(' x ',' x ',' '..' ',' x ')'}{http://rocq-prover.org/doc/V9.1.1/stdlib//Corelib.Init.Datatypes}{\coqdocnotation{(}}\coqref{SECF.Noninterference.pi:34}{\coqdocvariable{pi}}\coqexternalref{::core scope:'(' x ',' x ',' '..' ',' x ')'}{http://rocq-prover.org/doc/V9.1.1/stdlib//Corelib.Init.Datatypes}{\coqdocnotation{,}}\coqref{SECF.Noninterference.si:35}{\coqdocvariable{si}}\coqexternalref{::core scope:'(' x ',' x ',' '..' ',' x ')'}{http://rocq-prover.org/doc/V9.1.1/stdlib//Corelib.Init.Datatypes}{\coqdocnotation{)}}.\coqdoceol
\coqdocemptyline
\coqdocnoindent
\coqdockw{Lemma} \coqdef{SECF.Noninterference.prove or disprove less obvious f7}{prove\_or\_disprove\_less\_obvious\_f7}{\coqdoclemma{prove\_or\_disprove\_less\_obvious\_f7}} :\coqdoceol
\coqdocindent{1.00em}
\coqref{SECF.Noninterference.noninterferent}{\coqdocdefinition{noninterferent}} \coqref{SECF.Noninterference.less obvious f7}{\coqdocdefinition{less\_obvious\_f7}} \coqexternalref{::type scope:x 'x5C/' x}{http://rocq-prover.org/doc/V9.1.1/stdlib//Corelib.Init.Logic}{\coqdocnotation{\ensuremath{\lor}}} \coqexternalref{::type scope:'x7E' x}{http://rocq-prover.org/doc/V9.1.1/stdlib//Corelib.Init.Logic}{\coqdocnotation{\ensuremath{\lnot}}}\coqref{SECF.Noninterference.noninterferent}{\coqdocdefinition{noninterferent}} \coqref{SECF.Noninterference.less obvious f7}{\coqdocdefinition{less\_obvious\_f7}}.\coqdoceol
\coqdocnoindent
\coqdockw{Proof}.\coqdoceol
 \coqdocvar{Admitted}.\coqdoceol
\end{coqdoccode}
\ensuremath{\Box} 
\begin{coqdoccode}
\coqdocemptyline
\end{coqdoccode}
\section{A too-strong secrecy definition}



 In the definition of noninterference above we pass the same public
    inputs to the two executions and this allows public outputs to
    depend on public inputs. To convince ourselves of this, let's look
    at the following overly strong definition of security: 
\begin{coqdoccode}
\coqdocemptyline
\coqdocnoindent
\coqdockw{Definition} \coqdef{SECF.Noninterference.too strong sec def}{too\_strong\_sec\_def}{\coqdocdefinition{too\_strong\_sec\_def}} \{\coqdef{SECF.Noninterference.PI:36}{PI}{\coqdocbinder{PI}} \coqdef{SECF.Noninterference.SI:37}{SI}{\coqdocbinder{SI}} \coqdef{SECF.Noninterference.PO:38}{PO}{\coqdocbinder{PO}} \coqdef{SECF.Noninterference.SO:39}{SO}{\coqdocbinder{SO}} : \coqdockw{Type}\} (\coqdef{SECF.Noninterference.f:40}{f}{\coqdocbinder{f}}:\coqref{SECF.Noninterference.PI:36}{\coqdocvariable{PI}}\coqexternalref{::type scope:x '->' x}{http://rocq-prover.org/doc/V9.1.1/stdlib//Corelib.Init.Logic}{\coqdocnotation{\ensuremath{\rightarrow}}}\coqref{SECF.Noninterference.SI:37}{\coqdocvariable{SI}}\coqexternalref{::type scope:x '->' x}{http://rocq-prover.org/doc/V9.1.1/stdlib//Corelib.Init.Logic}{\coqdocnotation{\ensuremath{\rightarrow}}}\coqref{SECF.Noninterference.PO:38}{\coqdocvariable{PO}}\coqexternalref{::type scope:x '*' x}{http://rocq-prover.org/doc/V9.1.1/stdlib//Corelib.Init.Datatypes}{\coqdocnotation{\ensuremath{\times}}}\coqref{SECF.Noninterference.SO:39}{\coqdocvariable{SO}}) :=\coqdoceol
\coqdocindent{1.00em}
\coqdockw{\ensuremath{\forall}} (\coqdef{SECF.Noninterference.pi1:41}{pi1}{\coqdocbinder{pi1}} \coqdef{SECF.Noninterference.pi2:42}{pi2}{\coqdocbinder{pi2}}:\coqref{SECF.Noninterference.PI:36}{\coqdocvariable{PI}}) (\coqdef{SECF.Noninterference.si1:43}{si1}{\coqdocbinder{si1}} \coqdef{SECF.Noninterference.si2:44}{si2}{\coqdocbinder{si2}}:\coqref{SECF.Noninterference.SI:37}{\coqdocvariable{SI}}), \coqexternalref{fst}{http://rocq-prover.org/doc/V9.1.1/stdlib//Corelib.Init.Datatypes}{\coqdocdefinition{fst}} (\coqref{SECF.Noninterference.f:40}{\coqdocvariable{f}} \coqref{SECF.Noninterference.pi1:41}{\coqdocvariable{pi1}} \coqref{SECF.Noninterference.si1:43}{\coqdocvariable{si1}}) \coqexternalref{::type scope:x '=' x}{http://rocq-prover.org/doc/V9.1.1/stdlib//Corelib.Init.Logic}{\coqdocnotation{=}} \coqexternalref{fst}{http://rocq-prover.org/doc/V9.1.1/stdlib//Corelib.Init.Datatypes}{\coqdocdefinition{fst}} (\coqref{SECF.Noninterference.f:40}{\coqdocvariable{f}} \coqref{SECF.Noninterference.pi2:42}{\coqdocvariable{pi2}} \coqref{SECF.Noninterference.si2:44}{\coqdocvariable{si2}}).\coqdoceol
\coqdocemptyline
\end{coqdoccode}
This basically says that the public output of \coqdocvar{f} can depend
    neither on the public input not on the secret input, so it has to
    be constant, which is not the case for our \coqref{SECF.Noninterference.secure f}{\coqdocdefinition{secure\_f}}. 
\begin{coqdoccode}
\coqdocemptyline
\coqdocnoindent
\coqdockw{Lemma} \coqdef{SECF.Noninterference.secure f rejected again}{secure\_f\_rejected\_again}{\coqdoclemma{secure\_f\_rejected\_again}} : \coqexternalref{::type scope:'x7E' x}{http://rocq-prover.org/doc/V9.1.1/stdlib//Corelib.Init.Logic}{\coqdocnotation{\ensuremath{\lnot}}}\coqref{SECF.Noninterference.too strong sec def}{\coqdocdefinition{too\_strong\_sec\_def}} \coqref{SECF.Noninterference.secure f}{\coqdocdefinition{secure\_f}}.\coqdoceol
\coqdocnoindent
\coqdockw{Proof}.\coqdoceol
\coqdocindent{1.00em}
\coqdoctac{unfold} \coqref{SECF.Noninterference.too strong sec def}{\coqdocdefinition{too\_strong\_sec\_def}}, \coqref{SECF.Noninterference.secure f}{\coqdocdefinition{secure\_f}}. \coqdoctac{simpl}. \coqdoctac{intros} \coqdocvar{contra}.\coqdoceol
\coqdocindent{1.00em}
\coqdoctac{specialize} (\coqdocvar{contra} 0 1 0 0). \coqdoctac{discriminate} \coqdocvar{contra}.\coqdoceol
\coqdocnoindent
\coqdockw{Qed}.\coqdoceol
\coqdocemptyline
\end{coqdoccode}
\section{Noninterferent implies splittable}



 Noninterference is still a very strong property, though. In
    particular, \coqdocvar{f} being noninterferent is equivalent to \coqdocvar{f} being
    splittable into two different functions, one of which doesn't get
    the secret at all. 
\begin{coqdoccode}
\coqdocemptyline
\coqdocnoindent
\coqdockw{Definition} \coqdef{SECF.Noninterference.splittable}{splittable}{\coqdocdefinition{splittable}} \{\coqdef{SECF.Noninterference.PI:45}{PI}{\coqdocbinder{PI}} \coqdef{SECF.Noninterference.SI:46}{SI}{\coqdocbinder{SI}} \coqdef{SECF.Noninterference.PO:47}{PO}{\coqdocbinder{PO}} \coqdef{SECF.Noninterference.SO:48}{SO}{\coqdocbinder{SO}} : \coqdockw{Type}\} (\coqdef{SECF.Noninterference.f:49}{f}{\coqdocbinder{f}}:\coqref{SECF.Noninterference.PI:45}{\coqdocvariable{PI}}\coqexternalref{::type scope:x '->' x}{http://rocq-prover.org/doc/V9.1.1/stdlib//Corelib.Init.Logic}{\coqdocnotation{\ensuremath{\rightarrow}}}\coqref{SECF.Noninterference.SI:46}{\coqdocvariable{SI}}\coqexternalref{::type scope:x '->' x}{http://rocq-prover.org/doc/V9.1.1/stdlib//Corelib.Init.Logic}{\coqdocnotation{\ensuremath{\rightarrow}}}\coqref{SECF.Noninterference.PO:47}{\coqdocvariable{PO}}\coqexternalref{::type scope:x '*' x}{http://rocq-prover.org/doc/V9.1.1/stdlib//Corelib.Init.Datatypes}{\coqdocnotation{\ensuremath{\times}}}\coqref{SECF.Noninterference.SO:48}{\coqdocvariable{SO}}) :=\coqdoceol
\coqdocindent{1.00em}
\coqexternalref{::type scope:'exists' x '..' x ',' x}{http://rocq-prover.org/doc/V9.1.1/stdlib//Corelib.Init.Logic}{\coqdocnotation{\ensuremath{\exists}}} \coqexternalref{::type scope:'exists' x '..' x ',' x}{http://rocq-prover.org/doc/V9.1.1/stdlib//Corelib.Init.Logic}{\coqdocnotation{(}}\coqdef{SECF.Noninterference.pf:50}{pf}{\coqdocbinder{pf}} : \coqref{SECF.Noninterference.PI:45}{\coqdocvariable{PI}} \coqexternalref{::type scope:x '->' x}{http://rocq-prover.org/doc/V9.1.1/stdlib//Corelib.Init.Logic}{\coqdocnotation{\ensuremath{\rightarrow}}} \coqref{SECF.Noninterference.PO:47}{\coqdocvariable{PO}}) (\coqdef{SECF.Noninterference.sf:51}{sf}{\coqdocbinder{sf}} : \coqref{SECF.Noninterference.PI:45}{\coqdocvariable{PI}} \coqexternalref{::type scope:x '->' x}{http://rocq-prover.org/doc/V9.1.1/stdlib//Corelib.Init.Logic}{\coqdocnotation{\ensuremath{\rightarrow}}} \coqref{SECF.Noninterference.SI:46}{\coqdocvariable{SI}} \coqexternalref{::type scope:x '->' x}{http://rocq-prover.org/doc/V9.1.1/stdlib//Corelib.Init.Logic}{\coqdocnotation{\ensuremath{\rightarrow}}} \coqref{SECF.Noninterference.SO:48}{\coqdocvariable{SO}}\coqexternalref{::type scope:'exists' x '..' x ',' x}{http://rocq-prover.org/doc/V9.1.1/stdlib//Corelib.Init.Logic}{\coqdocnotation{),}}\coqdoceol
\coqdocindent{2.00em}
\coqdockw{\ensuremath{\forall}} \coqdef{SECF.Noninterference.pi:52}{pi}{\coqdocbinder{pi}} \coqdef{SECF.Noninterference.si:53}{si}{\coqdocbinder{si}} , \coqref{SECF.Noninterference.f:49}{\coqdocvariable{f}} \coqref{SECF.Noninterference.pi:52}{\coqdocvariable{pi}} \coqref{SECF.Noninterference.si:53}{\coqdocvariable{si}} \coqexternalref{::type scope:x '=' x}{http://rocq-prover.org/doc/V9.1.1/stdlib//Corelib.Init.Logic}{\coqdocnotation{=}} \coqexternalref{::core scope:'(' x ',' x ',' '..' ',' x ')'}{http://rocq-prover.org/doc/V9.1.1/stdlib//Corelib.Init.Datatypes}{\coqdocnotation{(}}\coqref{SECF.Noninterference.pf:50}{\coqdocvariable{pf}} \coqref{SECF.Noninterference.pi:52}{\coqdocvariable{pi}}\coqexternalref{::core scope:'(' x ',' x ',' '..' ',' x ')'}{http://rocq-prover.org/doc/V9.1.1/stdlib//Corelib.Init.Datatypes}{\coqdocnotation{,}} \coqref{SECF.Noninterference.sf:51}{\coqdocvariable{sf}} \coqref{SECF.Noninterference.pi:52}{\coqdocvariable{pi}} \coqref{SECF.Noninterference.si:53}{\coqdocvariable{si}}\coqexternalref{::core scope:'(' x ',' x ',' '..' ',' x ')'}{http://rocq-prover.org/doc/V9.1.1/stdlib//Corelib.Init.Datatypes}{\coqdocnotation{)}}.\coqdoceol
\coqdocemptyline
\coqdocnoindent
\coqdockw{Theorem} \coqdef{SECF.Noninterference.splittable noninterferent}{splittable\_noninterferent}{\coqdoclemma{splittable\_noninterferent}} : \coqdockw{\ensuremath{\forall}} \coqdef{SECF.Noninterference.PI:54}{PI}{\coqdocbinder{PI}} \coqdef{SECF.Noninterference.SI:55}{SI}{\coqdocbinder{SI}} \coqdef{SECF.Noninterference.PO:56}{PO}{\coqdocbinder{PO}} \coqdef{SECF.Noninterference.SO:57}{SO}{\coqdocbinder{SO}} : \coqdockw{Type},\coqdoceol
\coqdocindent{1.00em}
\coqdockw{\ensuremath{\forall}} \coqdef{SECF.Noninterference.f:58}{f}{\coqdocbinder{f}} : \coqref{SECF.Noninterference.PI:54}{\coqdocvariable{PI}} \coqexternalref{::type scope:x '->' x}{http://rocq-prover.org/doc/V9.1.1/stdlib//Corelib.Init.Logic}{\coqdocnotation{\ensuremath{\rightarrow}}} \coqref{SECF.Noninterference.SI:55}{\coqdocvariable{SI}} \coqexternalref{::type scope:x '->' x}{http://rocq-prover.org/doc/V9.1.1/stdlib//Corelib.Init.Logic}{\coqdocnotation{\ensuremath{\rightarrow}}} \coqref{SECF.Noninterference.PO:56}{\coqdocvariable{PO}}\coqexternalref{::type scope:x '*' x}{http://rocq-prover.org/doc/V9.1.1/stdlib//Corelib.Init.Datatypes}{\coqdocnotation{\ensuremath{\times}}}\coqref{SECF.Noninterference.SO:57}{\coqdocvariable{SO}}, \coqref{SECF.Noninterference.splittable}{\coqdocdefinition{splittable}} \coqref{SECF.Noninterference.f:58}{\coqdocvariable{f}} \coqexternalref{::type scope:x '->' x}{http://rocq-prover.org/doc/V9.1.1/stdlib//Corelib.Init.Logic}{\coqdocnotation{\ensuremath{\rightarrow}}} \coqref{SECF.Noninterference.noninterferent}{\coqdocdefinition{noninterferent}} \coqref{SECF.Noninterference.f:58}{\coqdocvariable{f}}.\coqdoceol
\coqdocnoindent
\coqdockw{Proof}.\coqdoceol
\coqdocindent{1.00em}
\coqdoctac{unfold} \coqref{SECF.Noninterference.splittable}{\coqdocdefinition{splittable}}, \coqref{SECF.Noninterference.noninterferent}{\coqdocdefinition{noninterferent}}.\coqdoceol
\coqdocindent{1.00em}
\coqdoctac{intros} \coqdocvar{PI} \coqdocvar{SI} \coqdocvar{PO} \coqdocvar{SO} \coqdocvar{f} [\coqdocvar{pf} [\coqdocvar{sf} \coqdocvar{H}]] \coqdocvar{pi} \coqdocvar{si1} \coqdocvar{si2}.\coqdoceol
\coqdocindent{1.00em}
\coqdoctac{rewrite} \coqdocvar{H}. \coqdoctac{rewrite} \coqdocvar{H}. \coqdoctac{simpl}. \coqdoctac{reflexivity}.\coqdoceol
\coqdocnoindent
\coqdockw{Qed}.\coqdoceol
\coqdocemptyline
\coqdocnoindent
\coqdockw{Theorem} \coqdef{SECF.Noninterference.noninterferent splittable}{noninterferent\_splittable}{\coqdoclemma{noninterferent\_splittable}} : \coqdockw{\ensuremath{\forall}} \coqdef{SECF.Noninterference.PI:59}{PI}{\coqdocbinder{PI}} \coqdef{SECF.Noninterference.SI:60}{SI}{\coqdocbinder{SI}} \coqdef{SECF.Noninterference.PO:61}{PO}{\coqdocbinder{PO}} \coqdef{SECF.Noninterference.SO:62}{SO}{\coqdocbinder{SO}} : \coqdockw{Type},\coqdoceol
\coqdocindent{1.00em}
\coqdockw{\ensuremath{\forall}} \coqdef{SECF.Noninterference.some si:63}{some\_si}{\coqdocbinder{some\_si}} : \coqref{SECF.Noninterference.SI:60}{\coqdocvariable{SI}}, \coqdoceol
\coqdocindent{1.00em}
\coqdockw{\ensuremath{\forall}} \coqdef{SECF.Noninterference.f:64}{f}{\coqdocbinder{f}} : \coqref{SECF.Noninterference.PI:59}{\coqdocvariable{PI}} \coqexternalref{::type scope:x '->' x}{http://rocq-prover.org/doc/V9.1.1/stdlib//Corelib.Init.Logic}{\coqdocnotation{\ensuremath{\rightarrow}}} \coqref{SECF.Noninterference.SI:60}{\coqdocvariable{SI}} \coqexternalref{::type scope:x '->' x}{http://rocq-prover.org/doc/V9.1.1/stdlib//Corelib.Init.Logic}{\coqdocnotation{\ensuremath{\rightarrow}}} \coqref{SECF.Noninterference.PO:61}{\coqdocvariable{PO}}\coqexternalref{::type scope:x '*' x}{http://rocq-prover.org/doc/V9.1.1/stdlib//Corelib.Init.Datatypes}{\coqdocnotation{\ensuremath{\times}}}\coqref{SECF.Noninterference.SO:62}{\coqdocvariable{SO}}, \coqref{SECF.Noninterference.noninterferent}{\coqdocdefinition{noninterferent}} \coqref{SECF.Noninterference.f:64}{\coqdocvariable{f}} \coqexternalref{::type scope:x '->' x}{http://rocq-prover.org/doc/V9.1.1/stdlib//Corelib.Init.Logic}{\coqdocnotation{\ensuremath{\rightarrow}}} \coqref{SECF.Noninterference.splittable}{\coqdocdefinition{splittable}} \coqref{SECF.Noninterference.f:64}{\coqdocvariable{f}}.\coqdoceol
\coqdocnoindent
\coqdockw{Proof}.\coqdoceol
\coqdocindent{1.00em}
\coqdoctac{unfold} \coqref{SECF.Noninterference.splittable}{\coqdocdefinition{splittable}}, \coqref{SECF.Noninterference.noninterferent}{\coqdocdefinition{noninterferent}}.\coqdoceol
\coqdocindent{1.00em}
\coqdoctac{intros} \coqdocvar{PI} \coqdocvar{SI} \coqdocvar{PO} \coqdocvar{SO} \coqdocvar{some\_si} \coqdocvar{f} \coqdocvar{Hni}.\coqdoceol
\coqdocindent{1.00em}
\coqdoctac{\ensuremath{\exists}} (\coqdockw{fun} \coqdef{SECF.Noninterference.pi:65}{pi}{\coqdocbinder{pi}} \ensuremath{\Rightarrow} \coqexternalref{fst}{http://rocq-prover.org/doc/V9.1.1/stdlib//Corelib.Init.Datatypes}{\coqdocdefinition{fst}} (\coqdocvar{f} \coqref{SECF.Noninterference.pi:65}{\coqdocvariable{pi}} \coqdocvar{some\_si})).\coqdoceol
\coqdocindent{1.00em}
\coqdoctac{\ensuremath{\exists}} (\coqdockw{fun} \coqdef{SECF.Noninterference.pi:67}{pi}{\coqdocbinder{pi}} \coqdef{SECF.Noninterference.si:68}{si}{\coqdocbinder{si}} \ensuremath{\Rightarrow} \coqexternalref{snd}{http://rocq-prover.org/doc/V9.1.1/stdlib//Corelib.Init.Datatypes}{\coqdocdefinition{snd}} (\coqdocvar{f} \coqref{SECF.Noninterference.pi:67}{\coqdocvariable{pi}} \coqref{SECF.Noninterference.si:68}{\coqdocvariable{si}})).\coqdoceol
\coqdocindent{1.00em}
\coqdoctac{intros} \coqdocvar{pi} \coqdocvar{si}. \coqdoctac{rewrite} (\coqdocvar{Hni} \coqdocvar{\_} \coqdocvar{\_} \coqdocvar{si}).\coqdoceol
\coqdocindent{1.00em}
\coqdoctac{destruct} (\coqdocvar{f} \coqdocvar{pi} \coqdocvar{si}) \coqdockw{as} [\coqdocvar{po} \coqdocvar{so}]. \coqdoctac{reflexivity}.\coqdoceol
\coqdocnoindent
\coqdockw{Qed}.\coqdoceol
\coqdocemptyline
\end{coqdoccode}
\section{Secure Multi-Execution (SME)}



 The previous proof also captures the key idea behind Secure
    Multi-Execution (SME) \coqdocvar{Devriese} \coqexternalref{and}{http://rocq-prover.org/doc/V9.1.1/stdlib//Corelib.Init.Logic}{\coqdocinductive{and}} \coqdocvar{Piessens} 2010 (in Bib.v), an
    enforcement mechanism that can make \textit{any} function
    noninterferent. To achieve this SME runs the function twice, once
    passing a dummy secret as input to obtain the public output, and
    once using the real secret input to obtain the secret output. 
\begin{coqdoccode}
\coqdocemptyline
\coqdocnoindent
\coqdockw{Definition} \coqdef{SECF.Noninterference.sme}{sme}{\coqdocdefinition{sme}} \{\coqdef{SECF.Noninterference.PI:71}{PI}{\coqdocbinder{PI}} \coqdef{SECF.Noninterference.SI:72}{SI}{\coqdocbinder{SI}} \coqdef{SECF.Noninterference.PO:73}{PO}{\coqdocbinder{PO}} \coqdef{SECF.Noninterference.SO:74}{SO}{\coqdocbinder{SO}} : \coqdockw{Type}\} (\coqdef{SECF.Noninterference.some si:75}{some\_si}{\coqdocbinder{some\_si}} : \coqref{SECF.Noninterference.SI:72}{\coqdocvariable{SI}})\coqdoceol
\coqdocindent{1.00em}
(\coqdef{SECF.Noninterference.f:76}{f}{\coqdocbinder{f}}:\coqref{SECF.Noninterference.PI:71}{\coqdocvariable{PI}}\coqexternalref{::type scope:x '->' x}{http://rocq-prover.org/doc/V9.1.1/stdlib//Corelib.Init.Logic}{\coqdocnotation{\ensuremath{\rightarrow}}}\coqref{SECF.Noninterference.SI:72}{\coqdocvariable{SI}}\coqexternalref{::type scope:x '->' x}{http://rocq-prover.org/doc/V9.1.1/stdlib//Corelib.Init.Logic}{\coqdocnotation{\ensuremath{\rightarrow}}}\coqref{SECF.Noninterference.PO:73}{\coqdocvariable{PO}}\coqexternalref{::type scope:x '*' x}{http://rocq-prover.org/doc/V9.1.1/stdlib//Corelib.Init.Datatypes}{\coqdocnotation{\ensuremath{\times}}}\coqref{SECF.Noninterference.SO:74}{\coqdocvariable{SO}}) : \coqref{SECF.Noninterference.PI:71}{\coqdocvariable{PI}}\coqexternalref{::type scope:x '->' x}{http://rocq-prover.org/doc/V9.1.1/stdlib//Corelib.Init.Logic}{\coqdocnotation{\ensuremath{\rightarrow}}}\coqref{SECF.Noninterference.SI:72}{\coqdocvariable{SI}}\coqexternalref{::type scope:x '->' x}{http://rocq-prover.org/doc/V9.1.1/stdlib//Corelib.Init.Logic}{\coqdocnotation{\ensuremath{\rightarrow}}}\coqref{SECF.Noninterference.PO:73}{\coqdocvariable{PO}}\coqexternalref{::type scope:x '*' x}{http://rocq-prover.org/doc/V9.1.1/stdlib//Corelib.Init.Datatypes}{\coqdocnotation{\ensuremath{\times}}}\coqref{SECF.Noninterference.SO:74}{\coqdocvariable{SO}} :=\coqdoceol
\coqdocindent{2.00em}
\coqdockw{fun} \coqdef{SECF.Noninterference.pi:77}{pi}{\coqdocbinder{pi}} \coqdef{SECF.Noninterference.si:78}{si}{\coqdocbinder{si}} \ensuremath{\Rightarrow} \coqexternalref{::core scope:'(' x ',' x ',' '..' ',' x ')'}{http://rocq-prover.org/doc/V9.1.1/stdlib//Corelib.Init.Datatypes}{\coqdocnotation{(}}\coqexternalref{fst}{http://rocq-prover.org/doc/V9.1.1/stdlib//Corelib.Init.Datatypes}{\coqdocdefinition{fst}} (\coqref{SECF.Noninterference.f:76}{\coqdocvariable{f}} \coqref{SECF.Noninterference.pi:77}{\coqdocvariable{pi}} \coqref{SECF.Noninterference.some si:75}{\coqdocvariable{some\_si}})\coqexternalref{::core scope:'(' x ',' x ',' '..' ',' x ')'}{http://rocq-prover.org/doc/V9.1.1/stdlib//Corelib.Init.Datatypes}{\coqdocnotation{,}} \coqexternalref{snd}{http://rocq-prover.org/doc/V9.1.1/stdlib//Corelib.Init.Datatypes}{\coqdocdefinition{snd}} (\coqref{SECF.Noninterference.f:76}{\coqdocvariable{f}} \coqref{SECF.Noninterference.pi:77}{\coqdocvariable{pi}} \coqref{SECF.Noninterference.si:78}{\coqdocvariable{si}})\coqexternalref{::core scope:'(' x ',' x ',' '..' ',' x ')'}{http://rocq-prover.org/doc/V9.1.1/stdlib//Corelib.Init.Datatypes}{\coqdocnotation{)}}.\coqdoceol
\coqdocemptyline
\end{coqdoccode}
Functions protected by \coqref{SECF.Noninterference.sme}{\coqdocdefinition{sme}} are guaranteed to satisfy noninterference: 
\begin{coqdoccode}
\coqdocemptyline
\coqdocnoindent
\coqdockw{Theorem} \coqdef{SECF.Noninterference.noninterferent sme}{noninterferent\_sme}{\coqdoclemma{noninterferent\_sme}} :  \coqdockw{\ensuremath{\forall}} \coqdef{SECF.Noninterference.PI:79}{PI}{\coqdocbinder{PI}} \coqdef{SECF.Noninterference.SI:80}{SI}{\coqdocbinder{SI}} \coqdef{SECF.Noninterference.PO:81}{PO}{\coqdocbinder{PO}} \coqdef{SECF.Noninterference.SO:82}{SO}{\coqdocbinder{SO}} : \coqdockw{Type},\coqdoceol
\coqdocindent{1.00em}
\coqdockw{\ensuremath{\forall}} \coqdef{SECF.Noninterference.some si:83}{some\_si}{\coqdocbinder{some\_si}} : \coqref{SECF.Noninterference.SI:80}{\coqdocvariable{SI}},\coqdoceol
\coqdocindent{1.00em}
\coqdockw{\ensuremath{\forall}} \coqdef{SECF.Noninterference.f:84}{f}{\coqdocbinder{f}} : \coqref{SECF.Noninterference.PI:79}{\coqdocvariable{PI}} \coqexternalref{::type scope:x '->' x}{http://rocq-prover.org/doc/V9.1.1/stdlib//Corelib.Init.Logic}{\coqdocnotation{\ensuremath{\rightarrow}}} \coqref{SECF.Noninterference.SI:80}{\coqdocvariable{SI}} \coqexternalref{::type scope:x '->' x}{http://rocq-prover.org/doc/V9.1.1/stdlib//Corelib.Init.Logic}{\coqdocnotation{\ensuremath{\rightarrow}}} \coqref{SECF.Noninterference.PO:81}{\coqdocvariable{PO}}\coqexternalref{::type scope:x '*' x}{http://rocq-prover.org/doc/V9.1.1/stdlib//Corelib.Init.Datatypes}{\coqdocnotation{\ensuremath{\times}}}\coqref{SECF.Noninterference.SO:82}{\coqdocvariable{SO}},\coqdoceol
\coqdocindent{2.00em}
\coqref{SECF.Noninterference.noninterferent}{\coqdocdefinition{noninterferent}} (\coqref{SECF.Noninterference.sme}{\coqdocdefinition{sme}} \coqref{SECF.Noninterference.some si:83}{\coqdocvariable{some\_si}} \coqref{SECF.Noninterference.f:84}{\coqdocvariable{f}}).\coqdoceol
\coqdocnoindent
\coqdockw{Proof}. \coqdoctac{intros} \coqdocvar{PI} \coqdocvar{SI} \coqdocvar{PO} \coqdocvar{SO} \coqdocvar{some\_si} \coqdocvar{f} \coqdocvar{pi} \coqdocvar{si1} \coqdocvar{si2}. \coqdoctac{simpl}. \coqdoctac{reflexivity}. \coqdockw{Qed}.\coqdoceol
\coqdocemptyline
\end{coqdoccode}
Moreover, if the function we pass to \coqref{SECF.Noninterference.sme}{\coqdocdefinition{sme}} is already noninterferent,
    then its behavior will not change; so we say that \coqref{SECF.Noninterference.sme}{\coqdocdefinition{sme}} is a \textit{transparent}
    enforcement mechanism for noninterference: 
\begin{coqdoccode}
\coqdocemptyline
\coqdocnoindent
\coqdockw{Theorem} \coqdef{SECF.Noninterference.transparent sme}{transparent\_sme}{\coqdoclemma{transparent\_sme}} : \coqdockw{\ensuremath{\forall}} \coqdef{SECF.Noninterference.PI:85}{PI}{\coqdocbinder{PI}} \coqdef{SECF.Noninterference.SI:86}{SI}{\coqdocbinder{SI}} \coqdef{SECF.Noninterference.PO:87}{PO}{\coqdocbinder{PO}} \coqdef{SECF.Noninterference.SO:88}{SO}{\coqdocbinder{SO}} : \coqdockw{Type},\coqdoceol
\coqdocindent{1.00em}
\coqdockw{\ensuremath{\forall}} \coqdef{SECF.Noninterference.some si:89}{some\_si}{\coqdocbinder{some\_si}} : \coqref{SECF.Noninterference.SI:86}{\coqdocvariable{SI}},\coqdoceol
\coqdocindent{1.00em}
\coqdockw{\ensuremath{\forall}} \coqdef{SECF.Noninterference.f:90}{f}{\coqdocbinder{f}} : \coqref{SECF.Noninterference.PI:85}{\coqdocvariable{PI}} \coqexternalref{::type scope:x '->' x}{http://rocq-prover.org/doc/V9.1.1/stdlib//Corelib.Init.Logic}{\coqdocnotation{\ensuremath{\rightarrow}}} \coqref{SECF.Noninterference.SI:86}{\coqdocvariable{SI}} \coqexternalref{::type scope:x '->' x}{http://rocq-prover.org/doc/V9.1.1/stdlib//Corelib.Init.Logic}{\coqdocnotation{\ensuremath{\rightarrow}}} \coqref{SECF.Noninterference.PO:87}{\coqdocvariable{PO}}\coqexternalref{::type scope:x '*' x}{http://rocq-prover.org/doc/V9.1.1/stdlib//Corelib.Init.Datatypes}{\coqdocnotation{\ensuremath{\times}}}\coqref{SECF.Noninterference.SO:88}{\coqdocvariable{SO}},\coqdoceol
\coqdocindent{2.00em}
\coqref{SECF.Noninterference.noninterferent}{\coqdocdefinition{noninterferent}} \coqref{SECF.Noninterference.f:90}{\coqdocvariable{f}} \coqexternalref{::type scope:x '->' x}{http://rocq-prover.org/doc/V9.1.1/stdlib//Corelib.Init.Logic}{\coqdocnotation{\ensuremath{\rightarrow}}} \coqdockw{\ensuremath{\forall}} \coqdef{SECF.Noninterference.pi:91}{pi}{\coqdocbinder{pi}} \coqdef{SECF.Noninterference.si:92}{si}{\coqdocbinder{si}}, \coqref{SECF.Noninterference.f:90}{\coqdocvariable{f}} \coqref{SECF.Noninterference.pi:91}{\coqdocvariable{pi}} \coqref{SECF.Noninterference.si:92}{\coqdocvariable{si}} \coqexternalref{::type scope:x '=' x}{http://rocq-prover.org/doc/V9.1.1/stdlib//Corelib.Init.Logic}{\coqdocnotation{=}} \coqref{SECF.Noninterference.sme}{\coqdocdefinition{sme}} \coqref{SECF.Noninterference.some si:89}{\coqdocvariable{some\_si}} \coqref{SECF.Noninterference.f:90}{\coqdocvariable{f}} \coqref{SECF.Noninterference.pi:91}{\coqdocvariable{pi}} \coqref{SECF.Noninterference.si:92}{\coqdocvariable{si}}.\coqdoceol
\coqdocnoindent
\coqdockw{Proof}.\coqdoceol
\coqdocindent{1.00em}
\coqdoctac{unfold} \coqref{SECF.Noninterference.noninterferent}{\coqdocdefinition{noninterferent}}, \coqref{SECF.Noninterference.sme}{\coqdocdefinition{sme}}. \coqdoctac{intros} \coqdocvar{PI} \coqdocvar{SI} \coqdocvar{PO} \coqdocvar{SP} \coqdocvar{some\_si} \coqdocvar{f} \coqdocvar{Hni} \coqdocvar{pi} \coqdocvar{si}.\coqdoceol
\coqdocindent{1.00em}
\coqdoctac{rewrite} (\coqdocvar{Hni} \coqdocvar{\_} \coqdocvar{\_} \coqdocvar{si}).\coqdoceol
\coqdocindent{1.00em}
\coqdoctac{destruct} (\coqdocvar{f} \coqdocvar{pi} \coqdocvar{si}) \coqdockw{as} [\coqdocvar{po} \coqdocvar{so}]. \coqdoctac{reflexivity}.\coqdoceol
\coqdocnoindent
\coqdockw{Qed}.\coqdoceol
\coqdocemptyline
\end{coqdoccode}
It is interesting to look at what \coqref{SECF.Noninterference.sme}{\coqdocdefinition{sme}} does for \textit{interferent} functions,
    like \coqref{SECF.Noninterference.insecure f}{\coqdocdefinition{insecure\_f}}, whose public output was one plus its secret input: 
\begin{coqdoccode}
\coqdocemptyline
\coqdocnoindent
\coqdockw{Example} \coqdef{SECF.Noninterference.example1 sme insecure f}{example1\_sme\_insecure\_f}{\coqdocdefinition{example1\_sme\_insecure\_f}}: \coqref{SECF.Noninterference.sme}{\coqdocdefinition{sme}} 0 \coqref{SECF.Noninterference.insecure f}{\coqdocdefinition{insecure\_f}} 0 0 \coqexternalref{::type scope:x '=' x}{http://rocq-prover.org/doc/V9.1.1/stdlib//Corelib.Init.Logic}{\coqdocnotation{=}} \coqexternalref{::core scope:'(' x ',' x ',' '..' ',' x ')'}{http://rocq-prover.org/doc/V9.1.1/stdlib//Corelib.Init.Datatypes}{\coqdocnotation{(}}1\coqexternalref{::core scope:'(' x ',' x ',' '..' ',' x ')'}{http://rocq-prover.org/doc/V9.1.1/stdlib//Corelib.Init.Datatypes}{\coqdocnotation{,}} 0\coqexternalref{::core scope:'(' x ',' x ',' '..' ',' x ')'}{http://rocq-prover.org/doc/V9.1.1/stdlib//Corelib.Init.Datatypes}{\coqdocnotation{)}}.\coqdoceol
\coqdocnoindent
\coqdockw{Proof}. \coqdoctac{reflexivity}. \coqdockw{Qed}.\coqdoceol
\coqdocemptyline
\coqdocnoindent
\coqdockw{Example} \coqdef{SECF.Noninterference.example2 sme insecure f}{example2\_sme\_insecure\_f}{\coqdocdefinition{example2\_sme\_insecure\_f}}: \coqref{SECF.Noninterference.sme}{\coqdocdefinition{sme}} 0 \coqref{SECF.Noninterference.insecure f}{\coqdocdefinition{insecure\_f}} 0 1 \coqexternalref{::type scope:x '=' x}{http://rocq-prover.org/doc/V9.1.1/stdlib//Corelib.Init.Logic}{\coqdocnotation{=}} \coqexternalref{::core scope:'(' x ',' x ',' '..' ',' x ')'}{http://rocq-prover.org/doc/V9.1.1/stdlib//Corelib.Init.Datatypes}{\coqdocnotation{(}}1\coqexternalref{::core scope:'(' x ',' x ',' '..' ',' x ')'}{http://rocq-prover.org/doc/V9.1.1/stdlib//Corelib.Init.Datatypes}{\coqdocnotation{,}} 2\coqexternalref{::core scope:'(' x ',' x ',' '..' ',' x ')'}{http://rocq-prover.org/doc/V9.1.1/stdlib//Corelib.Init.Datatypes}{\coqdocnotation{)}}.\coqdoceol
\coqdocnoindent
\coqdockw{Proof}. \coqdoctac{reflexivity}. \coqdockw{Qed}.\coqdoceol
\coqdocemptyline
\coqdocnoindent
\coqdockw{Example} \coqdef{SECF.Noninterference.example3 sme insecure f}{example3\_sme\_insecure\_f}{\coqdocdefinition{example3\_sme\_insecure\_f}}: \coqref{SECF.Noninterference.sme}{\coqdocdefinition{sme}} 0 \coqref{SECF.Noninterference.insecure f}{\coqdocdefinition{insecure\_f}} 1 1 \coqexternalref{::type scope:x '=' x}{http://rocq-prover.org/doc/V9.1.1/stdlib//Corelib.Init.Logic}{\coqdocnotation{=}} \coqexternalref{::core scope:'(' x ',' x ',' '..' ',' x ')'}{http://rocq-prover.org/doc/V9.1.1/stdlib//Corelib.Init.Datatypes}{\coqdocnotation{(}}1\coqexternalref{::core scope:'(' x ',' x ',' '..' ',' x ')'}{http://rocq-prover.org/doc/V9.1.1/stdlib//Corelib.Init.Datatypes}{\coqdocnotation{,}} 3\coqexternalref{::core scope:'(' x ',' x ',' '..' ',' x ')'}{http://rocq-prover.org/doc/V9.1.1/stdlib//Corelib.Init.Datatypes}{\coqdocnotation{)}}.\coqdoceol
\coqdocnoindent
\coqdockw{Proof}. \coqdoctac{reflexivity}. \coqdockw{Qed}.\coqdoceol
\coqdocemptyline
\end{coqdoccode}
Now the public output of \coqref{SECF.Noninterference.sme}{\coqdocdefinition{sme}} \coqref{SECF.Noninterference.insecure f}{\coqdocdefinition{insecure\_f}} 0 is one plus the dummy
   constant 0, so always the constant 1. 
\begin{coqdoccode}
\coqdocemptyline
\coqdocnoindent
\coqdockw{Lemma} \coqdef{SECF.Noninterference.constant sme insecure f}{constant\_sme\_insecure\_f}{\coqdoclemma{constant\_sme\_insecure\_f}}: \coqdockw{\ensuremath{\forall}} \coqdef{SECF.Noninterference.pi:93}{pi}{\coqdocbinder{pi}} \coqdef{SECF.Noninterference.si:94}{si}{\coqdocbinder{si}},\coqdoceol
\coqdocindent{1.00em}
\coqexternalref{fst}{http://rocq-prover.org/doc/V9.1.1/stdlib//Corelib.Init.Datatypes}{\coqdocdefinition{fst}} (\coqref{SECF.Noninterference.sme}{\coqdocdefinition{sme}} 0 \coqref{SECF.Noninterference.insecure f}{\coqdocdefinition{insecure\_f}} \coqref{SECF.Noninterference.pi:93}{\coqdocvariable{pi}} \coqref{SECF.Noninterference.si:94}{\coqdocvariable{si}}) \coqexternalref{::type scope:x '=' x}{http://rocq-prover.org/doc/V9.1.1/stdlib//Corelib.Init.Logic}{\coqdocnotation{=}} 1.\coqdoceol
\coqdocnoindent
\coqdockw{Proof}. \coqdoctac{reflexivity}. \coqdockw{Qed}.\coqdoceol
\coqdocemptyline
\end{coqdoccode}
This is a secure behavior, but it is different from that of the
    original \coqref{SECF.Noninterference.insecure f}{\coqdocdefinition{insecure\_f}} function. So we are giving up some
    correctness for security. There is no free lunch! 

 Of course the public output of sme does not always become, since
    some functions still use the public input. 
\begin{coqdoccode}
\coqdocemptyline
\coqdocnoindent
\coqdockw{Definition} \coqdef{SECF.Noninterference.another insecure f}{another\_insecure\_f}{\coqdocdefinition{another\_insecure\_f}} (\coqdef{SECF.Noninterference.pi:95}{pi}{\coqdocbinder{pi}} \coqdef{SECF.Noninterference.si:96}{si}{\coqdocbinder{si}} : \coqexternalref{nat}{http://rocq-prover.org/doc/V9.1.1/stdlib//Corelib.Init.Datatypes}{\coqdocinductive{nat}}) : \coqexternalref{nat}{http://rocq-prover.org/doc/V9.1.1/stdlib//Corelib.Init.Datatypes}{\coqdocinductive{nat}}\coqexternalref{::type scope:x '*' x}{http://rocq-prover.org/doc/V9.1.1/stdlib//Corelib.Init.Datatypes}{\coqdocnotation{\ensuremath{\times}}}\coqexternalref{nat}{http://rocq-prover.org/doc/V9.1.1/stdlib//Corelib.Init.Datatypes}{\coqdocinductive{nat}} := \coqexternalref{::core scope:'(' x ',' x ',' '..' ',' x ')'}{http://rocq-prover.org/doc/V9.1.1/stdlib//Corelib.Init.Datatypes}{\coqdocnotation{(}}\coqref{SECF.Noninterference.pi:95}{\coqdocvariable{pi}}\coqexternalref{::nat scope:x '+' x}{http://rocq-prover.org/doc/V9.1.1/stdlib//Corelib.Init.Nat}{\coqdocnotation{+}}\coqref{SECF.Noninterference.si:96}{\coqdocvariable{si}}\coqexternalref{::core scope:'(' x ',' x ',' '..' ',' x ')'}{http://rocq-prover.org/doc/V9.1.1/stdlib//Corelib.Init.Datatypes}{\coqdocnotation{,}} \coqref{SECF.Noninterference.pi:95}{\coqdocvariable{pi}}\coqexternalref{::nat scope:x '+' x}{http://rocq-prover.org/doc/V9.1.1/stdlib//Corelib.Init.Nat}{\coqdocnotation{+}}\coqref{SECF.Noninterference.si:96}{\coqdocvariable{si}}\coqexternalref{::core scope:'(' x ',' x ',' '..' ',' x ')'}{http://rocq-prover.org/doc/V9.1.1/stdlib//Corelib.Init.Datatypes}{\coqdocnotation{)}}.\coqdoceol
\coqdocemptyline
\coqdocnoindent
\coqdockw{Lemma} \coqdef{SECF.Noninterference.sme another insecure f}{sme\_another\_insecure\_f}{\coqdoclemma{sme\_another\_insecure\_f}} : \coqdockw{\ensuremath{\forall}} \coqdef{SECF.Noninterference.pi:97}{pi}{\coqdocbinder{pi}} \coqdef{SECF.Noninterference.si:98}{si}{\coqdocbinder{si}},\coqdoceol
\coqdocindent{1.00em}
\coqref{SECF.Noninterference.sme}{\coqdocdefinition{sme}} 0 (\coqref{SECF.Noninterference.another insecure f}{\coqdocdefinition{another\_insecure\_f}}) \coqref{SECF.Noninterference.pi:97}{\coqdocvariable{pi}} \coqref{SECF.Noninterference.si:98}{\coqdocvariable{si}} \coqexternalref{::type scope:x '=' x}{http://rocq-prover.org/doc/V9.1.1/stdlib//Corelib.Init.Logic}{\coqdocnotation{=}} \coqexternalref{::core scope:'(' x ',' x ',' '..' ',' x ')'}{http://rocq-prover.org/doc/V9.1.1/stdlib//Corelib.Init.Datatypes}{\coqdocnotation{(}}\coqref{SECF.Noninterference.pi:97}{\coqdocvariable{pi}}\coqexternalref{::core scope:'(' x ',' x ',' '..' ',' x ')'}{http://rocq-prover.org/doc/V9.1.1/stdlib//Corelib.Init.Datatypes}{\coqdocnotation{,}}\coqref{SECF.Noninterference.pi:97}{\coqdocvariable{pi}}\coqexternalref{::nat scope:x '+' x}{http://rocq-prover.org/doc/V9.1.1/stdlib//Corelib.Init.Nat}{\coqdocnotation{+}}\coqref{SECF.Noninterference.si:98}{\coqdocvariable{si}}\coqexternalref{::core scope:'(' x ',' x ',' '..' ',' x ')'}{http://rocq-prover.org/doc/V9.1.1/stdlib//Corelib.Init.Datatypes}{\coqdocnotation{)}}.\coqdoceol
\coqdocnoindent
\coqdockw{Proof}. \coqdoctac{unfold} \coqref{SECF.Noninterference.sme}{\coqdocdefinition{sme}}, \coqref{SECF.Noninterference.another insecure f}{\coqdocdefinition{another\_insecure\_f}}.\coqdoceol
\coqdocindent{1.00em}
\coqdoctac{intros} \coqdocvar{pi} \coqdocvar{si}. \coqdoctac{simpl}. \coqdoctac{rewrite} \ensuremath{\leftarrow} \coqexternalref{plus n O}{http://rocq-prover.org/doc/V9.1.1/stdlib//Corelib.Init.Peano}{\coqdoclemma{plus\_n\_O}}. \coqdoctac{reflexivity}. \coqdockw{Qed}.\coqdoceol
\coqdocemptyline
\end{coqdoccode}
\paragraph{Exercise: 1 star, standard (sme\_another\_insecure\_f2)}


\begin{coqdoccode}
\coqdocnoindent
\coqdockw{Definition} \coqdef{SECF.Noninterference.another insecure f2}{another\_insecure\_f2}{\coqdocdefinition{another\_insecure\_f2}} (\coqdef{SECF.Noninterference.pi:99}{pi}{\coqdocbinder{pi}} \coqdef{SECF.Noninterference.si:100}{si}{\coqdocbinder{si}} : \coqexternalref{nat}{http://rocq-prover.org/doc/V9.1.1/stdlib//Corelib.Init.Datatypes}{\coqdocinductive{nat}}) : \coqexternalref{nat}{http://rocq-prover.org/doc/V9.1.1/stdlib//Corelib.Init.Datatypes}{\coqdocinductive{nat}}\coqexternalref{::type scope:x '*' x}{http://rocq-prover.org/doc/V9.1.1/stdlib//Corelib.Init.Datatypes}{\coqdocnotation{\ensuremath{\times}}}\coqexternalref{nat}{http://rocq-prover.org/doc/V9.1.1/stdlib//Corelib.Init.Datatypes}{\coqdocinductive{nat}} :=\coqdoceol
\coqdocindent{1.00em}
\coqexternalref{::core scope:'(' x ',' x ',' '..' ',' x ')'}{http://rocq-prover.org/doc/V9.1.1/stdlib//Corelib.Init.Datatypes}{\coqdocnotation{(}}\coqdockw{if} \coqdocdefinition{Nat.eqb} \coqref{SECF.Noninterference.si:100}{\coqdocvariable{si}} 0 \coqdockw{then} \coqref{SECF.Noninterference.si:100}{\coqdocvariable{si}} \coqexternalref{::nat scope:x '*' x}{http://rocq-prover.org/doc/V9.1.1/stdlib//Corelib.Init.Nat}{\coqdocnotation{\ensuremath{\times}}} \coqref{SECF.Noninterference.pi:99}{\coqdocvariable{pi}} \coqexternalref{::nat scope:x '+' x}{http://rocq-prover.org/doc/V9.1.1/stdlib//Corelib.Init.Nat}{\coqdocnotation{+}} \coqref{SECF.Noninterference.pi:99}{\coqdocvariable{pi}} \coqdockw{else} \coqref{SECF.Noninterference.pi:99}{\coqdocvariable{pi}}\coqexternalref{::core scope:'(' x ',' x ',' '..' ',' x ')'}{http://rocq-prover.org/doc/V9.1.1/stdlib//Corelib.Init.Datatypes}{\coqdocnotation{,}} \coqref{SECF.Noninterference.pi:99}{\coqdocvariable{pi}}\coqexternalref{::nat scope:x '+' x}{http://rocq-prover.org/doc/V9.1.1/stdlib//Corelib.Init.Nat}{\coqdocnotation{+}}\coqref{SECF.Noninterference.si:100}{\coqdocvariable{si}}\coqexternalref{::core scope:'(' x ',' x ',' '..' ',' x ')'}{http://rocq-prover.org/doc/V9.1.1/stdlib//Corelib.Init.Datatypes}{\coqdocnotation{)}}.\coqdoceol
\coqdocemptyline
\coqdocnoindent
\coqdockw{Lemma} \coqdef{SECF.Noninterference.sme another insecure f2}{sme\_another\_insecure\_f2}{\coqdoclemma{sme\_another\_insecure\_f2}} : \coqdockw{\ensuremath{\forall}} \coqdef{SECF.Noninterference.pi:101}{pi}{\coqdocbinder{pi}} \coqdef{SECF.Noninterference.si:102}{si}{\coqdocbinder{si}},\coqdoceol
\coqdocindent{2.00em}
\coqref{SECF.Noninterference.sme}{\coqdocdefinition{sme}} 0 (\coqref{SECF.Noninterference.another insecure f2}{\coqdocdefinition{another\_insecure\_f2}}) \coqref{SECF.Noninterference.pi:101}{\coqdocvariable{pi}} \coqref{SECF.Noninterference.si:102}{\coqdocvariable{si}} \coqexternalref{::type scope:x '=' x}{http://rocq-prover.org/doc/V9.1.1/stdlib//Corelib.Init.Logic}{\coqdocnotation{=}} \coqexternalref{::core scope:'(' x ',' x ',' '..' ',' x ')'}{http://rocq-prover.org/doc/V9.1.1/stdlib//Corelib.Init.Datatypes}{\coqdocnotation{(}}\coqref{SECF.Noninterference.pi:101}{\coqdocvariable{pi}}\coqexternalref{::core scope:'(' x ',' x ',' '..' ',' x ')'}{http://rocq-prover.org/doc/V9.1.1/stdlib//Corelib.Init.Datatypes}{\coqdocnotation{,}} \coqref{SECF.Noninterference.pi:101}{\coqdocvariable{pi}}\coqexternalref{::nat scope:x '+' x}{http://rocq-prover.org/doc/V9.1.1/stdlib//Corelib.Init.Nat}{\coqdocnotation{+}}\coqref{SECF.Noninterference.si:102}{\coqdocvariable{si}}\coqexternalref{::core scope:'(' x ',' x ',' '..' ',' x ')'}{http://rocq-prover.org/doc/V9.1.1/stdlib//Corelib.Init.Datatypes}{\coqdocnotation{)}}.\coqdoceol
\coqdocnoindent
\coqdockw{Proof}.\coqdoceol
 \coqdocvar{Admitted}.\coqdoceol
\end{coqdoccode}
\ensuremath{\Box} 

\paragraph{Exercise: 2 stars, standard (sme\_another\_insecure\_f3)}


\begin{coqdoccode}
\coqdocnoindent
\coqdockw{Definition} \coqdef{SECF.Noninterference.another insecure f3}{another\_insecure\_f3}{\coqdocdefinition{another\_insecure\_f3}} (\coqdef{SECF.Noninterference.pi:103}{pi}{\coqdocbinder{pi}} \coqdef{SECF.Noninterference.si:104}{si}{\coqdocbinder{si}} : \coqexternalref{nat}{http://rocq-prover.org/doc/V9.1.1/stdlib//Corelib.Init.Datatypes}{\coqdocinductive{nat}}) : \coqexternalref{nat}{http://rocq-prover.org/doc/V9.1.1/stdlib//Corelib.Init.Datatypes}{\coqdocinductive{nat}}\coqexternalref{::type scope:x '*' x}{http://rocq-prover.org/doc/V9.1.1/stdlib//Corelib.Init.Datatypes}{\coqdocnotation{\ensuremath{\times}}}\coqexternalref{nat}{http://rocq-prover.org/doc/V9.1.1/stdlib//Corelib.Init.Datatypes}{\coqdocinductive{nat}} :=\coqdoceol
\coqdocindent{1.00em}
\coqexternalref{::core scope:'(' x ',' x ',' '..' ',' x ')'}{http://rocq-prover.org/doc/V9.1.1/stdlib//Corelib.Init.Datatypes}{\coqdocnotation{(}}\coqdockw{if} \coqdocdefinition{Nat.eqb} \coqref{SECF.Noninterference.si:104}{\coqdocvariable{si}} \coqref{SECF.Noninterference.pi:103}{\coqdocvariable{pi}} \coqdockw{then} \coqref{SECF.Noninterference.si:104}{\coqdocvariable{si}} \coqexternalref{::nat scope:x '*' x}{http://rocq-prover.org/doc/V9.1.1/stdlib//Corelib.Init.Nat}{\coqdocnotation{\ensuremath{\times}}} \coqref{SECF.Noninterference.pi:103}{\coqdocvariable{pi}} \coqdockw{else} \coqref{SECF.Noninterference.pi:103}{\coqdocvariable{pi}}\coqexternalref{::core scope:'(' x ',' x ',' '..' ',' x ')'}{http://rocq-prover.org/doc/V9.1.1/stdlib//Corelib.Init.Datatypes}{\coqdocnotation{,}} \coqref{SECF.Noninterference.pi:103}{\coqdocvariable{pi}}\coqexternalref{::nat scope:x '+' x}{http://rocq-prover.org/doc/V9.1.1/stdlib//Corelib.Init.Nat}{\coqdocnotation{+}}\coqref{SECF.Noninterference.si:104}{\coqdocvariable{si}}\coqexternalref{::core scope:'(' x ',' x ',' '..' ',' x ')'}{http://rocq-prover.org/doc/V9.1.1/stdlib//Corelib.Init.Datatypes}{\coqdocnotation{)}}.\coqdoceol
\coqdocemptyline
\coqdocnoindent
\coqdockw{Lemma} \coqdef{SECF.Noninterference.interferent another insecure f3}{interferent\_another\_insecure\_f3}{\coqdoclemma{interferent\_another\_insecure\_f3}} : \coqexternalref{::type scope:'x7E' x}{http://rocq-prover.org/doc/V9.1.1/stdlib//Corelib.Init.Logic}{\coqdocnotation{\ensuremath{\lnot}}} \coqref{SECF.Noninterference.noninterferent}{\coqdocdefinition{noninterferent}} \coqref{SECF.Noninterference.another insecure f3}{\coqdocdefinition{another\_insecure\_f3}}.\coqdoceol
\coqdocnoindent
\coqdockw{Proof}.\coqdoceol
\coqdocindent{1.00em}
\coqdoctac{unfold} \coqref{SECF.Noninterference.noninterferent}{\coqdocdefinition{noninterferent}}, \coqref{SECF.Noninterference.another insecure f3}{\coqdocdefinition{another\_insecure\_f3}}. \coqdoctac{simpl}.\coqdoceol
\coqdocindent{1.00em}
\coqdoctac{intros} \coqdocvar{contra}. \coqdoctac{specialize} (\coqdocvar{contra} 8 2 8). \coqdoctac{simpl} \coqdoctac{in} \coqdocvar{contra}. \coqdoctac{discriminate} \coqdocvar{contra}.\coqdoceol
\coqdocnoindent
\coqdockw{Qed}.\coqdoceol
\coqdocemptyline
\coqdocnoindent
\coqdockw{Lemma} \coqdef{SECF.Noninterference.sme another insecure f3}{sme\_another\_insecure\_f3}{\coqdoclemma{sme\_another\_insecure\_f3}} : \coqdockw{\ensuremath{\forall}} \coqdef{SECF.Noninterference.pi:105}{pi}{\coqdocbinder{pi}} \coqdef{SECF.Noninterference.si:106}{si}{\coqdocbinder{si}},\coqdoceol
\coqdocindent{2.00em}
\coqref{SECF.Noninterference.sme}{\coqdocdefinition{sme}} 0 (\coqref{SECF.Noninterference.another insecure f3}{\coqdocdefinition{another\_insecure\_f3}}) \coqref{SECF.Noninterference.pi:105}{\coqdocvariable{pi}} \coqref{SECF.Noninterference.si:106}{\coqdocvariable{si}} \coqexternalref{::type scope:x '=' x}{http://rocq-prover.org/doc/V9.1.1/stdlib//Corelib.Init.Logic}{\coqdocnotation{=}} \coqexternalref{::core scope:'(' x ',' x ',' '..' ',' x ')'}{http://rocq-prover.org/doc/V9.1.1/stdlib//Corelib.Init.Datatypes}{\coqdocnotation{(}}\coqref{SECF.Noninterference.pi:105}{\coqdocvariable{pi}}\coqexternalref{::core scope:'(' x ',' x ',' '..' ',' x ')'}{http://rocq-prover.org/doc/V9.1.1/stdlib//Corelib.Init.Datatypes}{\coqdocnotation{,}} \coqref{SECF.Noninterference.pi:105}{\coqdocvariable{pi}}\coqexternalref{::nat scope:x '+' x}{http://rocq-prover.org/doc/V9.1.1/stdlib//Corelib.Init.Nat}{\coqdocnotation{+}}\coqref{SECF.Noninterference.si:106}{\coqdocvariable{si}}\coqexternalref{::core scope:'(' x ',' x ',' '..' ',' x ')'}{http://rocq-prover.org/doc/V9.1.1/stdlib//Corelib.Init.Datatypes}{\coqdocnotation{)}}.\coqdoceol
\coqdocnoindent
\coqdockw{Proof}.\coqdoceol
 \coqdocvar{Admitted}.\coqdoceol
\end{coqdoccode}
\ensuremath{\Box} 

 The other downside of \coqref{SECF.Noninterference.sme}{\coqdocdefinition{sme}} is that we have to run the function
    twice for our two security levels, public and secret. In general,
    we need to run the program as many times as we have security
    levels, which is often an exponential number, say if we take our
    security levels to be sets of principals. This is inefficient!


    Other information-flow control mechanisms overcome this downside,
    but have other downsides of their own, for instance:

\begin{itemize}
\item  by requiring nontrivial manual proofs for each individual
          program (e.g., Relational Hoare Logic), or

\item  by using static overapproximations that reject some secure
          programs (security type systems), or

\item  by using dynamic overapproximations that unnecessarily
          change program behavior, for instance forcefully terminating
          even some secure programs to prevent leaks, in which case
          they are not transparent (dynamic information-flow control;
          an extension of dynamic taint tracking to also handle
          implicit flows).

\end{itemize}


    Again, there is no free lunch! 
\begin{coqdoccode}
\coqdocemptyline
\end{coqdoccode}
\section{Noninterference for state transformers}



 The development above is quite easy to adapt to Rocq functions that
    transform states (\coqref{SECF.SpecCT.state}{\coqdocdefinition{state}}\ensuremath{\rightarrow}\coqref{SECF.SpecCT.state}{\coqdocdefinition{state}}), where we label each variable as
    either public or secret using a map of type \coqref{SECF.SpecCT.pub vars}{\coqdocdefinition{pub\_vars}}. 
\begin{coqdoccode}
\coqdocemptyline
\coqdocnoindent
\coqdockw{Print} \coqdocvar{state}. \coqdocemptyline
\coqdocnoindent
\coqdockw{Definition} \coqdef{SECF.Noninterference.pub vars}{pub\_vars}{\coqdocdefinition{pub\_vars}} := \coqref{SECF.Maps.total map}{\coqdocdefinition{total\_map}} \coqexternalref{bool}{http://rocq-prover.org/doc/V9.1.1/stdlib//Corelib.Init.Datatypes}{\coqdocinductive{bool}}. \coqdocemptyline
\end{coqdoccode}
Instead of requiring that the first elements of two pairs are
    equal, we require that the two states have equal values on the
    variables labeled public by the \coqdocvar{pub} map. 
\begin{coqdoccode}
\coqdocemptyline
\coqdocnoindent
\coqdockw{Definition} \coqdef{SECF.Noninterference.pub equiv}{pub\_equiv}{\coqdocdefinition{pub\_equiv}} (\coqdef{SECF.Noninterference.pub:107}{pub}{\coqdocbinder{pub}} : \coqref{SECF.Noninterference.pub vars}{\coqdocdefinition{pub\_vars}}) (\coqdef{SECF.Noninterference.s1:108}{s1}{\coqdocbinder{s1}} \coqdef{SECF.Noninterference.s2:109}{s2}{\coqdocbinder{s2}} : \coqref{SECF.Imp.state}{\coqdocdefinition{state}}) :=\coqdoceol
\coqdocindent{1.00em}
\coqdockw{\ensuremath{\forall}} \coqdef{SECF.Noninterference.x:110}{x}{\coqdocbinder{x}}:\coqdocinductive{string}, \coqref{SECF.Noninterference.pub:107}{\coqdocvariable{pub}} \coqref{SECF.Noninterference.x:110}{\coqdocvariable{x}} \coqexternalref{::type scope:x '=' x}{http://rocq-prover.org/doc/V9.1.1/stdlib//Corelib.Init.Logic}{\coqdocnotation{=}} \coqref{SECF.Imp.:::'true'}{\coqdocnotation{true}} \coqexternalref{::type scope:x '->' x}{http://rocq-prover.org/doc/V9.1.1/stdlib//Corelib.Init.Logic}{\coqdocnotation{\ensuremath{\rightarrow}}} \coqref{SECF.Noninterference.s1:108}{\coqdocvariable{s1}} \coqref{SECF.Noninterference.x:110}{\coqdocvariable{x}} \coqexternalref{::type scope:x '=' x}{http://rocq-prover.org/doc/V9.1.1/stdlib//Corelib.Init.Logic}{\coqdocnotation{=}} \coqref{SECF.Noninterference.s2:109}{\coqdocvariable{s2}} \coqref{SECF.Noninterference.x:110}{\coqdocvariable{x}}.\coqdoceol
\coqdocemptyline
\end{coqdoccode}
This makes the definition more symmetric, since we can use
    \coqref{SECF.SpecCT.pub equiv}{\coqdocdefinition{pub\_equiv}} both for the input states and the output states: 
\begin{coqdoccode}
\coqdocemptyline
\coqdocnoindent
\coqdockw{Definition} \coqdef{SECF.Noninterference.noninterferent state}{noninterferent\_state}{\coqdocdefinition{noninterferent\_state}} \coqdef{SECF.Noninterference.pub:111}{pub}{\coqdocbinder{pub}} (\coqdef{SECF.Noninterference.f:112}{f}{\coqdocbinder{f}} : \coqref{SECF.Imp.state}{\coqdocdefinition{state}} \coqexternalref{::type scope:x '->' x}{http://rocq-prover.org/doc/V9.1.1/stdlib//Corelib.Init.Logic}{\coqdocnotation{\ensuremath{\rightarrow}}} \coqref{SECF.Imp.state}{\coqdocdefinition{state}}) :=\coqdoceol
\coqdocindent{1.00em}
\coqdockw{\ensuremath{\forall}} \coqdef{SECF.Noninterference.s1:113}{s1}{\coqdocbinder{s1}} \coqdef{SECF.Noninterference.s2:114}{s2}{\coqdocbinder{s2}}, \coqref{SECF.Noninterference.pub equiv}{\coqdocdefinition{pub\_equiv}} \coqref{SECF.Noninterference.pub:111}{\coqdocvariable{pub}} \coqref{SECF.Noninterference.s1:113}{\coqdocvariable{s1}} \coqref{SECF.Noninterference.s2:114}{\coqdocvariable{s2}} \coqexternalref{::type scope:x '->' x}{http://rocq-prover.org/doc/V9.1.1/stdlib//Corelib.Init.Logic}{\coqdocnotation{\ensuremath{\rightarrow}}} \coqref{SECF.Noninterference.pub equiv}{\coqdocdefinition{pub\_equiv}} \coqref{SECF.Noninterference.pub:111}{\coqdocvariable{pub}} (\coqref{SECF.Noninterference.f:112}{\coqdocvariable{f}} \coqref{SECF.Noninterference.s1:113}{\coqdocvariable{s1}}) (\coqref{SECF.Noninterference.f:112}{\coqdocvariable{f}} \coqref{SECF.Noninterference.s2:114}{\coqdocvariable{s2}}).\coqdoceol
\coqdocemptyline
\end{coqdoccode}
We can prove an equivalence between \coqref{SECF.Noninterference.noninterferent state}{\coqdocdefinition{noninterferent\_state}} and our original
    \coqref{SECF.StaticIFC.OUTPUT.noninterferent}{\coqdocdefinition{noninterferent}} definition. For this we need to split and merge states.


    We also need a few helper lemmas. 

 The way we define \coqref{SECF.Noninterference.split state}{\coqdocdefinition{split\_state}} and \coqdocvar{merge\_state} is a good example of
    programming with higher-order functions, and there's more of this in
    \coqdoclibrary{Maps}.


    The \coqref{SECF.Noninterference.split state}{\coqdocdefinition{split\_state}} function takes a state \coqdocvar{s} and zeroes out the variables
    \coqdocvar{x} for which \coqdocvar{pub} \coqdocvar{x} is different than an argument bit \coqdocvar{b}. So
    \coqref{SECF.Noninterference.split state}{\coqdocdefinition{split\_state}} \coqdocvar{s} \coqdocvar{pub} \coqexternalref{true}{http://rocq-prover.org/doc/V9.1.1/stdlib//Corelib.Init.Datatypes}{\coqdocconstructor{true}} keeps the public variables, and zeroes out the
    secret ones. Dually, \coqref{SECF.Noninterference.split state}{\coqdocdefinition{split\_state}} \coqdocvar{s} \coqdocvar{pub} \coqexternalref{false}{http://rocq-prover.org/doc/V9.1.1/stdlib//Corelib.Init.Datatypes}{\coqdocconstructor{false}} keeps the secret variables,
    and zeroes out the public ones.  
\begin{coqdoccode}
\coqdocemptyline
\coqdocnoindent
\coqdockw{Definition} \coqdef{SECF.Noninterference.split state}{split\_state}{\coqdocdefinition{split\_state}} (\coqdef{SECF.Noninterference.s:115}{s}{\coqdocbinder{s}}:\coqref{SECF.Imp.state}{\coqdocdefinition{state}}) (\coqdef{SECF.Noninterference.pub:116}{pub}{\coqdocbinder{pub}}:\coqref{SECF.Noninterference.pub vars}{\coqdocdefinition{pub\_vars}}) (\coqdef{SECF.Noninterference.b:117}{b}{\coqdocbinder{b}}:\coqexternalref{bool}{http://rocq-prover.org/doc/V9.1.1/stdlib//Corelib.Init.Datatypes}{\coqdocinductive{bool}}) : \coqref{SECF.Imp.state}{\coqdocdefinition{state}} :=\coqdoceol
\coqdocindent{1.00em}
\coqdockw{fun} \coqdef{SECF.Noninterference.x:118}{x}{\coqdocbinder{x}} : \coqdocinductive{string} \ensuremath{\Rightarrow} \coqdockw{if} \coqdocdefinition{Bool.eqb} (\coqref{SECF.Noninterference.pub:116}{\coqdocvariable{pub}} \coqref{SECF.Noninterference.x:118}{\coqdocvariable{x}}) \coqref{SECF.Noninterference.b:117}{\coqdocvariable{b}} \coqdockw{then} \coqref{SECF.Noninterference.s:115}{\coqdocvariable{s}} \coqref{SECF.Noninterference.x:118}{\coqdocvariable{x}} \coqdockw{else} 0.\coqdoceol
\coqdocemptyline
\end{coqdoccode}
The \coqdocvar{merge\_state} function takes in two states \coqref{SECF.Noninterference.s1}{\coqdocdefinition{s1}} and \coqref{SECF.Noninterference.s2}{\coqdocdefinition{s2}}
    and produces a new state that contains the public variables from
    \coqref{SECF.Noninterference.s1}{\coqdocdefinition{s1}} and the private variables from \coqref{SECF.Noninterference.s2}{\coqdocdefinition{s2}}. 
\begin{coqdoccode}
\coqdocemptyline
\coqdocnoindent
\coqdockw{Definition} \coqdef{SECF.Noninterference.merge states}{merge\_states}{\coqdocdefinition{merge\_states}} (\coqdef{SECF.Noninterference.s1:119}{s1}{\coqdocbinder{s1}} \coqdef{SECF.Noninterference.s2:120}{s2}{\coqdocbinder{s2}}:\coqref{SECF.Imp.state}{\coqdocdefinition{state}}) (\coqdef{SECF.Noninterference.pub:121}{pub}{\coqdocbinder{pub}}:\coqref{SECF.Noninterference.pub vars}{\coqdocdefinition{pub\_vars}}) : \coqref{SECF.Imp.state}{\coqdocdefinition{state}} :=\coqdoceol
\coqdocindent{1.00em}
\coqdockw{fun} \coqdef{SECF.Noninterference.x:122}{x}{\coqdocbinder{x}} : \coqdocinductive{string} \ensuremath{\Rightarrow} \coqdockw{if} \coqref{SECF.Noninterference.pub:121}{\coqdocvariable{pub}} \coqref{SECF.Noninterference.x:122}{\coqdocvariable{x}} \coqdockw{then} \coqref{SECF.Noninterference.s1:119}{\coqdocvariable{s1}} \coqref{SECF.Noninterference.x:122}{\coqdocvariable{x}} \coqdockw{else} \coqref{SECF.Noninterference.s2:120}{\coqdocvariable{s2}} \coqref{SECF.Noninterference.x:122}{\coqdocvariable{x}}.\coqdoceol
\coqdocemptyline
\coqdocnoindent
\coqdockw{Definition} \coqdef{SECF.Noninterference.split state fun}{split\_state\_fun}{\coqdocdefinition{split\_state\_fun}} (\coqdef{SECF.Noninterference.pub:123}{pub}{\coqdocbinder{pub}} : \coqref{SECF.Noninterference.pub vars}{\coqdocdefinition{pub\_vars}}) (\coqdef{SECF.Noninterference.mf:124}{mf}{\coqdocbinder{mf}} : \coqref{SECF.Imp.state}{\coqdocdefinition{state}} \coqexternalref{::type scope:x '->' x}{http://rocq-prover.org/doc/V9.1.1/stdlib//Corelib.Init.Logic}{\coqdocnotation{\ensuremath{\rightarrow}}} \coqref{SECF.Imp.state}{\coqdocdefinition{state}}) :=\coqdoceol
\coqdocindent{1.00em}
\coqdockw{fun} \coqdef{SECF.Noninterference.s1:125}{s1}{\coqdocbinder{s1}} \coqdef{SECF.Noninterference.s2:126}{s2}{\coqdocbinder{s2}} : \coqref{SECF.Imp.state}{\coqdocdefinition{state}} \ensuremath{\Rightarrow}\coqdoceol
\coqdocindent{2.00em}
\coqdockw{let} \coqdef{SECF.Noninterference.ms:127}{ms}{\coqdocbinder{ms}} := \coqref{SECF.Noninterference.mf:124}{\coqdocvariable{mf}} (\coqref{SECF.Noninterference.merge states}{\coqdocdefinition{merge\_states}} \coqref{SECF.Noninterference.s1:125}{\coqdocvariable{s1}} \coqref{SECF.Noninterference.s2:126}{\coqdocvariable{s2}} \coqref{SECF.Noninterference.pub:123}{\coqdocvariable{pub}}) \coqdoctac{in}\coqdoceol
\coqdocindent{2.00em}
\coqexternalref{::core scope:'(' x ',' x ',' '..' ',' x ')'}{http://rocq-prover.org/doc/V9.1.1/stdlib//Corelib.Init.Datatypes}{\coqdocnotation{(}}\coqref{SECF.Noninterference.split state}{\coqdocdefinition{split\_state}} \coqref{SECF.Noninterference.ms:127}{\coqdocvariable{ms}} \coqref{SECF.Noninterference.pub:123}{\coqdocvariable{pub}} \coqref{SECF.Imp.:::'true'}{\coqdocnotation{true}}\coqexternalref{::core scope:'(' x ',' x ',' '..' ',' x ')'}{http://rocq-prover.org/doc/V9.1.1/stdlib//Corelib.Init.Datatypes}{\coqdocnotation{,}} \coqref{SECF.Noninterference.split state}{\coqdocdefinition{split\_state}} \coqref{SECF.Noninterference.ms:127}{\coqdocvariable{ms}} \coqref{SECF.Noninterference.pub:123}{\coqdocvariable{pub}} \coqref{SECF.Imp.:::'false'}{\coqdocnotation{false}}\coqexternalref{::core scope:'(' x ',' x ',' '..' ',' x ')'}{http://rocq-prover.org/doc/V9.1.1/stdlib//Corelib.Init.Datatypes}{\coqdocnotation{)}}.\coqdoceol
\coqdocemptyline
\end{coqdoccode}
The technical development needed for the equivalence proof between
    \coqref{SECF.Noninterference.noninterferent state}{\coqdocdefinition{noninterferent\_state}} and our original \coqref{SECF.StaticIFC.OUTPUT.noninterferent}{\coqdocdefinition{noninterferent}}
    definition is not that interesting though, and one can skip
    directly to the \coqref{SECF.Noninterference.noninterferent state ni}{\coqdoclemma{noninterferent\_state\_ni}} statement on first read. 
\begin{coqdoccode}
\coqdocemptyline
\coqdocnoindent
\coqdockw{Definition} \coqdef{SECF.Noninterference.pub equiv split}{pub\_equiv\_split}{\coqdocdefinition{pub\_equiv\_split}} (\coqdef{SECF.Noninterference.pub:128}{pub}{\coqdocbinder{pub}} : \coqref{SECF.Noninterference.pub vars}{\coqdocdefinition{pub\_vars}}) (\coqdef{SECF.Noninterference.s1:129}{s1}{\coqdocbinder{s1}} \coqdef{SECF.Noninterference.s2:130}{s2}{\coqdocbinder{s2}} : \coqref{SECF.Imp.state}{\coqdocdefinition{state}}) :=\coqdoceol
\coqdocindent{1.00em}
\coqdockw{\ensuremath{\forall}} \coqdef{SECF.Noninterference.x:131}{x}{\coqdocbinder{x}}:\coqdocinductive{string}, (\coqref{SECF.Noninterference.split state}{\coqdocdefinition{split\_state}} \coqref{SECF.Noninterference.s1:129}{\coqdocvariable{s1}} \coqref{SECF.Noninterference.pub:128}{\coqdocvariable{pub}} \coqref{SECF.Imp.:::'true'}{\coqdocnotation{true}}) \coqref{SECF.Noninterference.x:131}{\coqdocvariable{x}} \coqexternalref{::type scope:x '=' x}{http://rocq-prover.org/doc/V9.1.1/stdlib//Corelib.Init.Logic}{\coqdocnotation{=}} (\coqref{SECF.Noninterference.split state}{\coqdocdefinition{split\_state}} \coqref{SECF.Noninterference.s2:130}{\coqdocvariable{s2}} \coqref{SECF.Noninterference.pub:128}{\coqdocvariable{pub}} \coqref{SECF.Imp.:::'true'}{\coqdocnotation{true}}) \coqref{SECF.Noninterference.x:131}{\coqdocvariable{x}}.\coqdoceol
\coqdocemptyline
\coqdocnoindent
\coqdockw{Theorem} \coqdef{SECF.Noninterference.pub equiv split iff}{pub\_equiv\_split\_iff}{\coqdoclemma{pub\_equiv\_split\_iff}} : \coqdockw{\ensuremath{\forall}} \coqdef{SECF.Noninterference.pub:132}{pub}{\coqdocbinder{pub}} \coqdef{SECF.Noninterference.s1:133}{s1}{\coqdocbinder{s1}} \coqdef{SECF.Noninterference.s2:134}{s2}{\coqdocbinder{s2}},\coqdoceol
\coqdocindent{1.00em}
\coqref{SECF.Noninterference.pub equiv}{\coqdocdefinition{pub\_equiv}} \coqref{SECF.Noninterference.pub:132}{\coqdocvariable{pub}} \coqref{SECF.Noninterference.s1:133}{\coqdocvariable{s1}} \coqref{SECF.Noninterference.s2:134}{\coqdocvariable{s2}} \coqexternalref{::type scope:x '<->' x}{http://rocq-prover.org/doc/V9.1.1/stdlib//Corelib.Init.Logic}{\coqdocnotation{\ensuremath{\leftrightarrow}}} \coqref{SECF.Noninterference.pub equiv split}{\coqdocdefinition{pub\_equiv\_split}} \coqref{SECF.Noninterference.pub:132}{\coqdocvariable{pub}} \coqref{SECF.Noninterference.s1:133}{\coqdocvariable{s1}} \coqref{SECF.Noninterference.s2:134}{\coqdocvariable{s2}}.\coqdoceol
\coqdocnoindent
\coqdockw{Proof}.\coqdoceol
\coqdocindent{1.00em}
\coqdoctac{unfold} \coqref{SECF.Noninterference.pub equiv}{\coqdocdefinition{pub\_equiv}}, \coqref{SECF.Noninterference.pub equiv split}{\coqdocdefinition{pub\_equiv\_split}}, \coqref{SECF.Noninterference.split state}{\coqdocdefinition{split\_state}}. \coqdoctac{intros}. \coqdoctac{split}.\coqdoceol
\coqdocindent{1.00em}
- \coqdoctac{intros} \coqdocvar{H} \coqdocvar{x}. \coqdoctac{destruct} (\coqdoclemma{Bool.eqb\_spec} (\coqdocvar{pub} \coqdocvar{x}) \coqref{SECF.Imp.:::'true'}{\coqdocnotation{true}}).\coqdoceol
\coqdocindent{2.00em}
+ \coqdoctac{apply} \coqdocvar{H}. \coqdoctac{apply} \coqdocvar{e}.\coqdoceol
\coqdocindent{2.00em}
+ \coqdoctac{reflexivity}.\coqdoceol
\coqdocindent{1.00em}
- \coqdoctac{intros} \coqdocvar{H} \coqdocvar{x}. \coqdoctac{specialize} (\coqdocvar{H} \coqdocvar{x}). \coqdoctac{destruct} (\coqdoclemma{Bool.eqb\_spec} (\coqdocvar{pub} \coqdocvar{x}) \coqref{SECF.Imp.:::'true'}{\coqdocnotation{true}}).\coqdoceol
\coqdocindent{2.00em}
+ \coqdoctac{intros} \coqdocvar{\_}. \coqdoctac{apply} \coqdocvar{H}.\coqdoceol
\coqdocindent{2.00em}
+ \coqdocvar{contradiction}.\coqdoceol
\coqdocnoindent
\coqdockw{Qed}.\coqdoceol
\coqdocemptyline
\coqdocnoindent
\coqdockw{Theorem} \coqdef{SECF.Noninterference.pub equiv merge states}{pub\_equiv\_merge\_states}{\coqdoclemma{pub\_equiv\_merge\_states}} : \coqdockw{\ensuremath{\forall}} \coqdef{SECF.Noninterference.pub:135}{pub}{\coqdocbinder{pub}} \coqdef{SECF.Noninterference.s:136}{s}{\coqdocbinder{s}} \coqdef{SECF.Noninterference.z1:137}{z1}{\coqdocbinder{z1}} \coqdef{SECF.Noninterference.z2:138}{z2}{\coqdocbinder{z2}},\coqdoceol
\coqdocindent{1.00em}
\coqref{SECF.Noninterference.pub equiv}{\coqdocdefinition{pub\_equiv}} \coqref{SECF.Noninterference.pub:135}{\coqdocvariable{pub}} (\coqref{SECF.Noninterference.merge states}{\coqdocdefinition{merge\_states}} \coqref{SECF.Noninterference.s:136}{\coqdocvariable{s}} \coqref{SECF.Noninterference.z1:137}{\coqdocvariable{z1}} \coqref{SECF.Noninterference.pub:135}{\coqdocvariable{pub}}) (\coqref{SECF.Noninterference.merge states}{\coqdocdefinition{merge\_states}} \coqref{SECF.Noninterference.s:136}{\coqdocvariable{s}} \coqref{SECF.Noninterference.z2:138}{\coqdocvariable{z2}} \coqref{SECF.Noninterference.pub:135}{\coqdocvariable{pub}}).\coqdoceol
\coqdocnoindent
\coqdockw{Proof}.\coqdoceol
\coqdocindent{1.00em}
\coqdoctac{unfold} \coqref{SECF.Noninterference.pub equiv}{\coqdocdefinition{pub\_equiv}}, \coqref{SECF.Noninterference.merge states}{\coqdocdefinition{merge\_states}}. \coqdoctac{intros} \coqdocvar{pub} \coqdocvar{s} \coqdocvar{z1} \coqdocvar{z2} \coqdocvar{x} \coqdocvar{Hx}.\coqdoceol
\coqdocindent{1.00em}
\coqdoctac{rewrite} \coqdocvar{Hx}. \coqdoctac{reflexivity}.\coqdoceol
\coqdocnoindent
\coqdockw{Qed}.\coqdoceol
\coqdocemptyline
\coqdocnoindent
\coqdockw{From} \coqdocvar{Stdlib} \coqdockw{Require} \coqdockw{Import} \coqdoclibrary{FunctionalExtensionality}.\coqdoceol
\coqdocemptyline
\coqdocnoindent
\coqdockw{Theorem} \coqdef{SECF.Noninterference.merge states split state}{merge\_states\_split\_state}{\coqdoclemma{merge\_states\_split\_state}} : \coqdockw{\ensuremath{\forall}} \coqdef{SECF.Noninterference.s:139}{s}{\coqdocbinder{s}} \coqdef{SECF.Noninterference.pub:140}{pub}{\coqdocbinder{pub}},\coqdoceol
\coqdocindent{1.00em}
\coqref{SECF.Noninterference.merge states}{\coqdocdefinition{merge\_states}} (\coqref{SECF.Noninterference.split state}{\coqdocdefinition{split\_state}} \coqref{SECF.Noninterference.s:139}{\coqdocvariable{s}} \coqref{SECF.Noninterference.pub:140}{\coqdocvariable{pub}} \coqref{SECF.Imp.:::'true'}{\coqdocnotation{true}}) (\coqref{SECF.Noninterference.split state}{\coqdocdefinition{split\_state}} \coqref{SECF.Noninterference.s:139}{\coqdocvariable{s}} \coqref{SECF.Noninterference.pub:140}{\coqdocvariable{pub}} \coqref{SECF.Imp.:::'false'}{\coqdocnotation{false}}) \coqref{SECF.Noninterference.pub:140}{\coqdocvariable{pub}} \coqexternalref{::type scope:x '=' x}{http://rocq-prover.org/doc/V9.1.1/stdlib//Corelib.Init.Logic}{\coqdocnotation{=}} \coqref{SECF.Noninterference.s:139}{\coqdocvariable{s}}.\coqdoceol
\coqdocnoindent
\coqdockw{Proof}.\coqdoceol
\coqdocindent{1.00em}
\coqdoctac{unfold} \coqref{SECF.Noninterference.merge states}{\coqdocdefinition{merge\_states}}, \coqref{SECF.Noninterference.split state}{\coqdocdefinition{split\_state}}. \coqdoctac{intros} \coqdocvar{s} \coqdocvar{pub}.\coqdoceol
\coqdocindent{1.00em}
\coqdoctac{apply} \coqdoclemma{functional\_extensionality}. \coqdoctac{intro} \coqdocvar{x}.\coqdoceol
\coqdocindent{1.00em}
\coqdoctac{destruct} (\coqdocvar{pub} \coqdocvar{x}) \coqdocvar{eqn}:\coqdocvar{Heq}; \coqdoctac{reflexivity}.\coqdoceol
\coqdocnoindent
\coqdockw{Qed}.\coqdoceol
\coqdocemptyline
\end{coqdoccode}
Now we can finally state our theorem about the equivalence between
    \coqdocvar{non\_interferent\_state} and \coqref{SECF.StaticIFC.OUTPUT.noninterferent}{\coqdocdefinition{noninterferent}}: 
\begin{coqdoccode}
\coqdocemptyline
\coqdocnoindent
\coqdockw{Theorem} \coqdef{SECF.Noninterference.noninterferent state ni}{noninterferent\_state\_ni}{\coqdoclemma{noninterferent\_state\_ni}} : \coqdockw{\ensuremath{\forall}} \coqdef{SECF.Noninterference.pub:141}{pub}{\coqdocbinder{pub}} \coqdef{SECF.Noninterference.f:142}{f}{\coqdocbinder{f}},\coqdoceol
\coqdocindent{1.00em}
\coqref{SECF.Noninterference.noninterferent state}{\coqdocdefinition{noninterferent\_state}} \coqref{SECF.Noninterference.pub:141}{\coqdocvariable{pub}} \coqref{SECF.Noninterference.f:142}{\coqdocvariable{f}} \coqexternalref{::type scope:x '<->' x}{http://rocq-prover.org/doc/V9.1.1/stdlib//Corelib.Init.Logic}{\coqdocnotation{\ensuremath{\leftrightarrow}}}\coqdoceol
\coqdocindent{1.00em}
\coqref{SECF.Noninterference.noninterferent}{\coqdocdefinition{noninterferent}} (\coqref{SECF.Noninterference.split state fun}{\coqdocdefinition{split\_state\_fun}} \coqref{SECF.Noninterference.pub:141}{\coqdocvariable{pub}} \coqref{SECF.Noninterference.f:142}{\coqdocvariable{f}}).\coqdoceol
\coqdocnoindent
\coqdockw{Proof}.\coqdoceol
\coqdocindent{1.00em}
\coqdoctac{unfold} \coqref{SECF.Noninterference.noninterferent state}{\coqdocdefinition{noninterferent\_state}}, \coqref{SECF.Noninterference.noninterferent}{\coqdocdefinition{noninterferent}}, \coqref{SECF.Noninterference.split state fun}{\coqdocdefinition{split\_state\_fun}}.\coqdoceol
\coqdocindent{1.00em}
\coqdoctac{intros} \coqdocvar{pub} \coqdocvar{f}. \coqdoctac{split}.\coqdoceol
\coqdocindent{1.00em}
- \coqdoctac{intros} \coqdocvar{H} \coqdocvar{s} \coqdocvar{z1} \coqdocvar{z2}. \coqdoctac{simpl}.\coqdoceol
\coqdocindent{2.00em}
\coqdoctac{assert} (\coqdocvar{H'} : \coqref{SECF.Noninterference.pub equiv}{\coqdocdefinition{pub\_equiv}} \coqdocvar{pub} (\coqref{SECF.Noninterference.merge states}{\coqdocdefinition{merge\_states}} \coqdocvar{s} \coqdocvar{z1} \coqdocvar{pub}) (\coqref{SECF.Noninterference.merge states}{\coqdocdefinition{merge\_states}} \coqdocvar{s} \coqdocvar{z2} \coqdocvar{pub})).\coqdoceol
\coqdocindent{3.00em}
\{ \coqdoctac{apply} \coqref{SECF.Noninterference.pub equiv merge states}{\coqdoclemma{pub\_equiv\_merge\_states}}. \}\coqdoceol
\coqdocindent{2.00em}
\coqdoctac{apply} \coqdocvar{H} \coqdoctac{in} \coqdocvar{H'}. \coqdoctac{rewrite} \coqref{SECF.Noninterference.pub equiv split iff}{\coqdoclemma{pub\_equiv\_split\_iff}} \coqdoctac{in} \coqdocvar{H'}.\coqdoceol
\coqdocindent{2.00em}
\coqdoctac{unfold} \coqref{SECF.Noninterference.pub equiv split}{\coqdocdefinition{pub\_equiv\_split}} \coqdoctac{in} \coqdocvar{H'}. \coqdoctac{apply} \coqdoclemma{functional\_extensionality}. \coqdoctac{apply} \coqdocvar{H'}.\coqdoceol
\coqdocindent{1.00em}
- \coqdoctac{intros} \coqdocvar{H} \coqdocvar{s1} \coqdocvar{s2} \coqdocvar{Hequiv}. \coqdoctac{simpl} \coqdoctac{in} \coqdocvar{H}.\coqdoceol
\coqdocindent{2.00em}
\coqdoctac{rewrite} \coqref{SECF.Noninterference.pub equiv split iff}{\coqdoclemma{pub\_equiv\_split\_iff}} \coqdoctac{in} \coqdocvar{Hequiv}. \coqdoctac{unfold} \coqref{SECF.Noninterference.pub equiv split}{\coqdocdefinition{pub\_equiv\_split}} \coqdoctac{in} \coqdocvar{Hequiv}.\coqdoceol
\coqdocindent{2.00em}
\coqdoctac{rewrite} \coqref{SECF.Noninterference.pub equiv split iff}{\coqdoclemma{pub\_equiv\_split\_iff}}. \coqdoctac{unfold} \coqref{SECF.Noninterference.pub equiv split}{\coqdocdefinition{pub\_equiv\_split}}. \coqdoctac{intro} \coqdocvar{x}.\coqdoceol
\coqdocindent{2.00em}
\coqdoctac{specialize} (\coqdocvar{H} (\coqref{SECF.Noninterference.split state}{\coqdocdefinition{split\_state}} \coqdocvar{s1} \coqdocvar{pub} \coqref{SECF.Imp.:::'true'}{\coqdocnotation{true}})\coqdoceol
\coqdocindent{9.00em}
(\coqref{SECF.Noninterference.split state}{\coqdocdefinition{split\_state}} \coqdocvar{s1} \coqdocvar{pub} \coqref{SECF.Imp.:::'false'}{\coqdocnotation{false}})\coqdoceol
\coqdocindent{9.00em}
(\coqref{SECF.Noninterference.split state}{\coqdocdefinition{split\_state}} \coqdocvar{s2} \coqdocvar{pub} \coqref{SECF.Imp.:::'false'}{\coqdocnotation{false}})).\coqdoceol
\coqdocindent{2.00em}
\coqdoctac{rewrite} \coqref{SECF.Noninterference.merge states split state}{\coqdoclemma{merge\_states\_split\_state}} \coqdoctac{in} \coqdocvar{H}.\coqdoceol
\coqdocindent{2.00em}
\coqdoctac{apply} \coqdoclemma{functional\_extensionality} \coqdoctac{in} \coqdocvar{Hequiv}. \coqdoctac{rewrite} \coqdocvar{Hequiv} \coqdoctac{in} \coqdocvar{H}.\coqdoceol
\coqdocindent{2.00em}
\coqdoctac{rewrite} \coqref{SECF.Noninterference.merge states split state}{\coqdoclemma{merge\_states\_split\_state}} \coqdoctac{in} \coqdocvar{H}.\coqdoceol
\coqdocindent{2.00em}
\coqdoctac{rewrite} \coqdocvar{H}. \coqdoctac{reflexivity}.\coqdoceol
\coqdocnoindent
\coqdockw{Qed}.\coqdoceol
\coqdocemptyline
\end{coqdoccode}
\section{SME for state transformers}



 We can use the \coqref{SECF.Noninterference.split state}{\coqdocdefinition{split\_state}} and \coqref{SECF.Noninterference.merge states}{\coqdocdefinition{merge\_states}} functions above to
    also define SME for state transformers. We call the \coqref{SECF.Noninterference.split state}{\coqdocdefinition{split\_state}}
    below to zero out all secret variables before calling \coqdocvar{f} the first
    time to obtain the final value of the public variables. 
\begin{coqdoccode}
\coqdocemptyline
\coqdocnoindent
\coqdockw{Definition} \coqdef{SECF.Noninterference.sme state}{sme\_state}{\coqdocdefinition{sme\_state}} (\coqdef{SECF.Noninterference.f:143}{f}{\coqdocbinder{f}} : \coqref{SECF.Imp.state}{\coqdocdefinition{state}} \coqexternalref{::type scope:x '->' x}{http://rocq-prover.org/doc/V9.1.1/stdlib//Corelib.Init.Logic}{\coqdocnotation{\ensuremath{\rightarrow}}} \coqref{SECF.Imp.state}{\coqdocdefinition{state}}) (\coqdef{SECF.Noninterference.pub:144}{pub}{\coqdocbinder{pub}}:\coqref{SECF.Noninterference.pub vars}{\coqdocdefinition{pub\_vars}}) :=\coqdoceol
\coqdocindent{1.00em}
\coqdockw{fun} \coqdef{SECF.Noninterference.s:145}{s}{\coqdocbinder{s}} \ensuremath{\Rightarrow} \coqref{SECF.Noninterference.merge states}{\coqdocdefinition{merge\_states}} (\coqref{SECF.Noninterference.f:143}{\coqdocvariable{f}} (\coqref{SECF.Noninterference.split state}{\coqdocdefinition{split\_state}} \coqref{SECF.Noninterference.s:145}{\coqdocvariable{s}} \coqref{SECF.Noninterference.pub:144}{\coqdocvariable{pub}} \coqref{SECF.Imp.:::'true'}{\coqdocnotation{true}})) (\coqref{SECF.Noninterference.f:143}{\coqdocvariable{f}} \coqref{SECF.Noninterference.s:145}{\coqdocvariable{s}}) \coqref{SECF.Noninterference.pub:144}{\coqdocvariable{pub}}.\coqdoceol
\coqdocemptyline
\end{coqdoccode}
We will see examples of this in an upcoming section, but for now
    we prove the same two theorems as for \coqref{SECF.Noninterference.sme}{\coqdocdefinition{sme}} above: 
\begin{coqdoccode}
\coqdocemptyline
\coqdocnoindent
\coqdockw{Theorem} \coqdef{SECF.Noninterference.noninterferent sme state}{noninterferent\_sme\_state}{\coqdoclemma{noninterferent\_sme\_state}} : \coqdockw{\ensuremath{\forall}} \coqdef{SECF.Noninterference.pub:146}{pub}{\coqdocbinder{pub}} \coqdef{SECF.Noninterference.f:147}{f}{\coqdocbinder{f}},\coqdoceol
\coqdocindent{1.00em}
\coqref{SECF.Noninterference.noninterferent state}{\coqdocdefinition{noninterferent\_state}} \coqref{SECF.Noninterference.pub:146}{\coqdocvariable{pub}} (\coqref{SECF.Noninterference.sme state}{\coqdocdefinition{sme\_state}} \coqref{SECF.Noninterference.f:147}{\coqdocvariable{f}} \coqref{SECF.Noninterference.pub:146}{\coqdocvariable{pub}}).\coqdoceol
\coqdocnoindent
\coqdockw{Proof}.\coqdoceol
\coqdocindent{1.00em}
\coqdoctac{unfold} \coqref{SECF.Noninterference.noninterferent state}{\coqdocdefinition{noninterferent\_state}}, \coqref{SECF.Noninterference.sme state}{\coqdocdefinition{sme\_state}}.\coqdoceol
\coqdocindent{1.00em}
\coqdoctac{intros} \coqdocvar{pub} \coqdocvar{f} \coqdocvar{s1} \coqdocvar{s2} \coqdocvar{Hequiv}.\coqdoceol
\coqdocindent{1.00em}
\coqdoctac{rewrite} \coqref{SECF.Noninterference.pub equiv split iff}{\coqdoclemma{pub\_equiv\_split\_iff}} \coqdoctac{in} \coqdocvar{Hequiv}.\coqdoceol
\coqdocindent{1.00em}
\coqdoctac{unfold} \coqref{SECF.Noninterference.pub equiv split}{\coqdocdefinition{pub\_equiv\_split}} \coqdoctac{in} \coqdocvar{Hequiv}.\coqdoceol
\coqdocindent{1.00em}
\coqdoctac{apply} \coqdoclemma{functional\_extensionality} \coqdoctac{in} \coqdocvar{Hequiv}. \coqdoctac{rewrite} \coqdocvar{Hequiv}.\coqdoceol
\coqdocindent{1.00em}
\coqdoctac{apply} \coqref{SECF.Noninterference.pub equiv merge states}{\coqdoclemma{pub\_equiv\_merge\_states}}.\coqdoceol
\coqdocnoindent
\coqdockw{Qed}.\coqdoceol
\coqdocemptyline
\coqdocnoindent
\coqdockw{Theorem} \coqdef{SECF.Noninterference.transparent sme state}{transparent\_sme\_state}{\coqdoclemma{transparent\_sme\_state}} : \coqdockw{\ensuremath{\forall}} \coqdef{SECF.Noninterference.f:148}{f}{\coqdocbinder{f}} \coqdef{SECF.Noninterference.pub:149}{pub}{\coqdocbinder{pub}},\coqdoceol
\coqdocindent{1.00em}
\coqref{SECF.Noninterference.noninterferent state}{\coqdocdefinition{noninterferent\_state}} \coqref{SECF.Noninterference.pub:149}{\coqdocvariable{pub}} \coqref{SECF.Noninterference.f:148}{\coqdocvariable{f}} \coqexternalref{::type scope:x '->' x}{http://rocq-prover.org/doc/V9.1.1/stdlib//Corelib.Init.Logic}{\coqdocnotation{\ensuremath{\rightarrow}}} \coqdockw{\ensuremath{\forall}} \coqdef{SECF.Noninterference.s:150}{s}{\coqdocbinder{s}}, \coqref{SECF.Noninterference.f:148}{\coqdocvariable{f}} \coqref{SECF.Noninterference.s:150}{\coqdocvariable{s}} \coqexternalref{::type scope:x '=' x}{http://rocq-prover.org/doc/V9.1.1/stdlib//Corelib.Init.Logic}{\coqdocnotation{=}} \coqref{SECF.Noninterference.sme state}{\coqdocdefinition{sme\_state}} \coqref{SECF.Noninterference.f:148}{\coqdocvariable{f}} \coqref{SECF.Noninterference.pub:149}{\coqdocvariable{pub}} \coqref{SECF.Noninterference.s:150}{\coqdocvariable{s}}.\coqdoceol
\coqdocnoindent
\coqdockw{Proof}.\coqdoceol
\coqdocindent{1.00em}
\coqdoctac{unfold} \coqref{SECF.Noninterference.noninterferent state}{\coqdocdefinition{noninterferent\_state}}, \coqref{SECF.Noninterference.sme state}{\coqdocdefinition{sme\_state}}.\coqdoceol
\coqdocindent{1.00em}
\coqdoctac{intros} \coqdocvar{f} \coqdocvar{pub} \coqdocvar{Hni} \coqdocvar{s}.\coqdoceol
\coqdocindent{1.00em}
\coqdoctac{unfold} \coqref{SECF.Noninterference.merge states}{\coqdocdefinition{merge\_states}}, \coqref{SECF.Noninterference.split state}{\coqdocdefinition{split\_state}}. \coqdoctac{unfold} \coqref{SECF.Noninterference.pub equiv}{\coqdocdefinition{pub\_equiv}} \coqdoctac{in} \coqdocvar{Hni}.\coqdoceol
\coqdocindent{1.00em}
\coqdoctac{apply} \coqdoclemma{functional\_extensionality}. \coqdoctac{intro} \coqdocvar{x}.\coqdoceol
\coqdocindent{1.00em}
\coqdoctac{destruct} (\coqdocvar{pub} \coqdocvar{x}) \coqdocvar{eqn}:\coqdocvar{Eq}.\coqdoceol
\coqdocindent{1.00em}
- \coqdoctac{apply} \coqdocvar{Hni}.\coqdoceol
\coqdocindent{2.00em}
+ \coqdoctac{intros} \coqdocvar{x'} \coqdocvar{Hx'}.\coqdoceol
\coqdocindent{3.00em}
\coqdoctac{destruct} (\coqdoclemma{Bool.eqb\_spec} (\coqdocvar{pub} \coqdocvar{x'}) \coqref{SECF.Imp.:::'true'}{\coqdocnotation{true}}).\coqdoceol
\coqdocindent{3.00em}
\ensuremath{\times} \coqdoctac{reflexivity}.\coqdoceol
\coqdocindent{3.00em}
\ensuremath{\times} \coqdocvar{contradiction}.\coqdoceol
\coqdocindent{2.00em}
+ \coqdoctac{assumption}.\coqdoceol
\coqdocindent{1.00em}
- \coqdoctac{reflexivity}.\coqdoceol
\coqdocnoindent
\coqdockw{Qed}.\coqdoceol
\coqdocemptyline
\end{coqdoccode}
One thing to note in this proof is that we used the lemma
    \coqdocvar{Bool.eqb\_spec} to do case analysis on whether the \coqdocvar{pub} \coqdocvar{x'} is
    equal to \coqexternalref{true}{http://rocq-prover.org/doc/V9.1.1/stdlib//Corelib.Init.Datatypes}{\coqdocconstructor{true}}. For more details on how this works, please check
    out the explanations about the \coqexternalref{reflect}{http://rocq-prover.org/doc/V9.1.1/stdlib//Corelib.Init.Datatypes}{\coqdocinductive{reflect}} inductive predicate in
    \coqdocvar{IndProp}. 
\begin{coqdoccode}
\coqdocemptyline
\end{coqdoccode}
\subsection{Optional: Connection between \texorpdfstring{\protect\coqref{SECF.Noninterference.sme}{\coqdocdefinition{sme}}}{sme} and \texorpdfstring{\protect\coqref{SECF.Noninterference.sme state}{\coqdocdefinition{sme\_state}}}{sme\_state}}



 We can formally relate \coqref{SECF.Noninterference.sme}{\coqdocdefinition{sme}} amd \coqref{SECF.Noninterference.sme state}{\coqdocdefinition{sme\_state}}, but this gets pretty
    technical, so the curious reader can directly skip to the two
    theorems at the end of this subsection. 
\begin{coqdoccode}
\coqdocemptyline
\coqdocnoindent
\coqdockw{Lemma} \coqdef{SECF.Noninterference.split merge public}{split\_merge\_public}{\coqdoclemma{split\_merge\_public}}: \coqdockw{\ensuremath{\forall}} \coqdef{SECF.Noninterference.s:151}{s}{\coqdocbinder{s}} \coqdef{SECF.Noninterference.pub:152}{pub}{\coqdocbinder{pub}},\coqdoceol
\coqdocindent{2.00em}
\coqref{SECF.Noninterference.split state}{\coqdocdefinition{split\_state}} \coqref{SECF.Noninterference.s:151}{\coqdocvariable{s}} \coqref{SECF.Noninterference.pub:152}{\coqdocvariable{pub}} \coqref{SECF.Imp.:::'true'}{\coqdocnotation{true}} \coqexternalref{::type scope:x '=' x}{http://rocq-prover.org/doc/V9.1.1/stdlib//Corelib.Init.Logic}{\coqdocnotation{=}} \coqref{SECF.Noninterference.merge states}{\coqdocdefinition{merge\_states}} \coqref{SECF.Noninterference.s:151}{\coqdocvariable{s}} (\coqdockw{fun} \coqdocvar{\_} \ensuremath{\Rightarrow} 0) \coqref{SECF.Noninterference.pub:152}{\coqdocvariable{pub}}.\coqdoceol
\coqdocnoindent
\coqdockw{Proof}.\coqdoceol
\coqdocindent{1.00em}
\coqdoctac{intros}. \coqdoctac{eapply} \coqdoclemma{functional\_extensionality}. \coqdoctac{intro} \coqdocvar{x}.\coqdoceol
\coqdocindent{1.00em}
\coqdoctac{unfold} \coqref{SECF.Noninterference.split state}{\coqdocdefinition{split\_state}}, \coqref{SECF.Noninterference.merge states}{\coqdocdefinition{merge\_states}}.\coqdoceol
\coqdocindent{1.00em}
\coqdoctac{destruct} (\coqdocvar{pub} \coqdocvar{x}) \coqdocvar{eqn}:\coqdocvar{PUB}; \coqdoctac{simpl}; \coqdoctac{reflexivity}.\coqdoceol
\coqdocnoindent
\coqdockw{Qed}.\coqdoceol
\coqdocemptyline
\coqdocnoindent
\coqdockw{Lemma} \coqdef{SECF.Noninterference.split merge split true}{split\_merge\_split\_true}{\coqdoclemma{split\_merge\_split\_true}}: \coqdockw{\ensuremath{\forall}} \coqdef{SECF.Noninterference.s:153}{s}{\coqdocbinder{s}} \coqdef{SECF.Noninterference.s':154}{s'}{\coqdocbinder{s'}} \coqdef{SECF.Noninterference.pub:155}{pub}{\coqdocbinder{pub}},\coqdoceol
\coqdocindent{2.00em}
\coqref{SECF.Noninterference.split state}{\coqdocdefinition{split\_state}} (\coqref{SECF.Noninterference.merge states}{\coqdocdefinition{merge\_states}} \coqref{SECF.Noninterference.s:153}{\coqdocvariable{s}} \coqref{SECF.Noninterference.s':154}{\coqdocvariable{s'}} \coqref{SECF.Noninterference.pub:155}{\coqdocvariable{pub}}) \coqref{SECF.Noninterference.pub:155}{\coqdocvariable{pub}} \coqref{SECF.Imp.:::'true'}{\coqdocnotation{true}} \coqexternalref{::type scope:x '=' x}{http://rocq-prover.org/doc/V9.1.1/stdlib//Corelib.Init.Logic}{\coqdocnotation{=}} \coqref{SECF.Noninterference.split state}{\coqdocdefinition{split\_state}} \coqref{SECF.Noninterference.s:153}{\coqdocvariable{s}} \coqref{SECF.Noninterference.pub:155}{\coqdocvariable{pub}} \coqref{SECF.Imp.:::'true'}{\coqdocnotation{true}}.\coqdoceol
\coqdocnoindent
\coqdockw{Proof}.\coqdoceol
\coqdocindent{1.00em}
\coqdoctac{intros}. \coqdoctac{eapply} \coqdoclemma{functional\_extensionality}. \coqdoctac{intro} \coqdocvar{x}.\coqdoceol
\coqdocindent{1.00em}
\coqdoctac{unfold} \coqref{SECF.Noninterference.split state}{\coqdocdefinition{split\_state}}, \coqref{SECF.Noninterference.merge states}{\coqdocdefinition{merge\_states}}.\coqdoceol
\coqdocindent{1.00em}
\coqdoctac{destruct} (\coqdocvar{pub} \coqdocvar{x}) \coqdocvar{eqn}:\coqdocvar{PUB}; \coqdoctac{simpl}; \coqdoctac{reflexivity}.\coqdoceol
\coqdocnoindent
\coqdockw{Qed}.\coqdoceol
\coqdocemptyline
\coqdocnoindent
\coqdockw{Lemma} \coqdef{SECF.Noninterference.split merge split false}{split\_merge\_split\_false}{\coqdoclemma{split\_merge\_split\_false}}: \coqdockw{\ensuremath{\forall}} \coqdef{SECF.Noninterference.s:156}{s}{\coqdocbinder{s}} \coqdef{SECF.Noninterference.s':157}{s'}{\coqdocbinder{s'}} \coqdef{SECF.Noninterference.pub:158}{pub}{\coqdocbinder{pub}},\coqdoceol
\coqdocindent{2.00em}
\coqref{SECF.Noninterference.split state}{\coqdocdefinition{split\_state}} (\coqref{SECF.Noninterference.merge states}{\coqdocdefinition{merge\_states}} \coqref{SECF.Noninterference.s:156}{\coqdocvariable{s}} \coqref{SECF.Noninterference.s':157}{\coqdocvariable{s'}} \coqref{SECF.Noninterference.pub:158}{\coqdocvariable{pub}}) \coqref{SECF.Noninterference.pub:158}{\coqdocvariable{pub}} \coqref{SECF.Imp.:::'false'}{\coqdocnotation{false}} \coqexternalref{::type scope:x '=' x}{http://rocq-prover.org/doc/V9.1.1/stdlib//Corelib.Init.Logic}{\coqdocnotation{=}} \coqref{SECF.Noninterference.split state}{\coqdocdefinition{split\_state}} \coqref{SECF.Noninterference.s':157}{\coqdocvariable{s'}} \coqref{SECF.Noninterference.pub:158}{\coqdocvariable{pub}} \coqref{SECF.Imp.:::'false'}{\coqdocnotation{false}}.\coqdoceol
\coqdocnoindent
\coqdockw{Proof}.\coqdoceol
\coqdocindent{1.00em}
\coqdoctac{intros}. \coqdoctac{eapply} \coqdoclemma{functional\_extensionality}. \coqdoctac{intro} \coqdocvar{x}.\coqdoceol
\coqdocindent{1.00em}
\coqdoctac{unfold} \coqref{SECF.Noninterference.split state}{\coqdocdefinition{split\_state}}, \coqref{SECF.Noninterference.merge states}{\coqdocdefinition{merge\_states}}.\coqdoceol
\coqdocindent{1.00em}
\coqdoctac{destruct} (\coqdocvar{pub} \coqdocvar{x}) \coqdocvar{eqn}:\coqdocvar{PUB}; \coqdoctac{simpl}; \coqdoctac{reflexivity}.\coqdoceol
\coqdocnoindent
\coqdockw{Qed}.\coqdoceol
\coqdocemptyline
\coqdocnoindent
\coqdockw{Lemma} \coqdef{SECF.Noninterference.merge states same}{merge\_states\_same}{\coqdoclemma{merge\_states\_same}}: \coqdockw{\ensuremath{\forall}} \coqdef{SECF.Noninterference.s:159}{s}{\coqdocbinder{s}} \coqdef{SECF.Noninterference.pub:160}{pub}{\coqdocbinder{pub}},\coqdoceol
\coqdocindent{2.00em}
\coqref{SECF.Noninterference.merge states}{\coqdocdefinition{merge\_states}} \coqref{SECF.Noninterference.s:159}{\coqdocvariable{s}} \coqref{SECF.Noninterference.s:159}{\coqdocvariable{s}} \coqref{SECF.Noninterference.pub:160}{\coqdocvariable{pub}} \coqexternalref{::type scope:x '=' x}{http://rocq-prover.org/doc/V9.1.1/stdlib//Corelib.Init.Logic}{\coqdocnotation{=}} \coqref{SECF.Noninterference.s:159}{\coqdocvariable{s}}.\coqdoceol
\coqdocnoindent
\coqdockw{Proof}.\coqdoceol
\coqdocindent{1.00em}
\coqdoctac{unfold} \coqref{SECF.Noninterference.merge states}{\coqdocdefinition{merge\_states}}. \coqdoctac{intros}.\coqdoceol
\coqdocindent{1.00em}
\coqdoctac{eapply} \coqdoclemma{functional\_extensionality}. \coqdoctac{intro} \coqdocvar{x}.\coqdoceol
\coqdocindent{1.00em}
\coqdoctac{destruct} (\coqdocvar{pub} \coqdocvar{x}); \coqdoctac{reflexivity}.\coqdoceol
\coqdocnoindent
\coqdockw{Qed}.\coqdoceol
\coqdocemptyline
\coqdocnoindent
\coqdockw{Lemma} \coqdef{SECF.Noninterference.split state idem}{split\_state\_idem}{\coqdoclemma{split\_state\_idem}}: \coqdockw{\ensuremath{\forall}} \coqdef{SECF.Noninterference.s:161}{s}{\coqdocbinder{s}} \coqdef{SECF.Noninterference.pub:162}{pub}{\coqdocbinder{pub}} \coqdef{SECF.Noninterference.b:163}{b}{\coqdocbinder{b}},\coqdoceol
\coqdocindent{2.00em}
\coqref{SECF.Noninterference.split state}{\coqdocdefinition{split\_state}} (\coqref{SECF.Noninterference.split state}{\coqdocdefinition{split\_state}} \coqref{SECF.Noninterference.s:161}{\coqdocvariable{s}} \coqref{SECF.Noninterference.pub:162}{\coqdocvariable{pub}} \coqref{SECF.Noninterference.b:163}{\coqdocvariable{b}}) \coqref{SECF.Noninterference.pub:162}{\coqdocvariable{pub}} \coqref{SECF.Noninterference.b:163}{\coqdocvariable{b}} \coqexternalref{::type scope:x '=' x}{http://rocq-prover.org/doc/V9.1.1/stdlib//Corelib.Init.Logic}{\coqdocnotation{=}} \coqref{SECF.Noninterference.split state}{\coqdocdefinition{split\_state}} \coqref{SECF.Noninterference.s:161}{\coqdocvariable{s}} \coqref{SECF.Noninterference.pub:162}{\coqdocvariable{pub}} \coqref{SECF.Noninterference.b:163}{\coqdocvariable{b}}.\coqdoceol
\coqdocnoindent
\coqdockw{Proof}.\coqdoceol
\coqdocindent{1.00em}
\coqdoctac{unfold} \coqref{SECF.Noninterference.split state}{\coqdocdefinition{split\_state}}. \coqdoctac{intros}.\coqdoceol
\coqdocindent{1.00em}
\coqdoctac{eapply} \coqdoclemma{functional\_extensionality}. \coqdoctac{intro} \coqdocvar{x}.\coqdoceol
\coqdocindent{1.00em}
\coqdoctac{destruct} (\coqdocdefinition{Bool.eqb} (\coqdocvar{pub} \coqdocvar{x}) \coqdocvar{b}); \coqdoctac{reflexivity}.\coqdoceol
\coqdocnoindent
\coqdockw{Qed}.\coqdoceol
\coqdocemptyline
\coqdocnoindent
\coqdockw{Lemma} \coqdef{SECF.Noninterference.eqb neg distr r}{eqb\_neg\_distr\_r}{\coqdoclemma{eqb\_neg\_distr\_r}}: \coqdockw{\ensuremath{\forall}} \coqdef{SECF.Noninterference.b1:164}{b1}{\coqdocbinder{b1}} \coqdef{SECF.Noninterference.b2:165}{b2}{\coqdocbinder{b2}},\coqdoceol
\coqdocindent{2.00em}
\coqdocdefinition{Bool.eqb} \coqref{SECF.Noninterference.b1:164}{\coqdocvariable{b1}} (\coqexternalref{negb}{http://rocq-prover.org/doc/V9.1.1/stdlib//Corelib.Init.Datatypes}{\coqdocdefinition{negb}} \coqref{SECF.Noninterference.b2:165}{\coqdocvariable{b2}}) \coqexternalref{::type scope:x '=' x}{http://rocq-prover.org/doc/V9.1.1/stdlib//Corelib.Init.Logic}{\coqdocnotation{=}} \coqexternalref{negb}{http://rocq-prover.org/doc/V9.1.1/stdlib//Corelib.Init.Datatypes}{\coqdocdefinition{negb}} (\coqdocdefinition{Bool.eqb} \coqref{SECF.Noninterference.b1:164}{\coqdocvariable{b1}} \coqref{SECF.Noninterference.b2:165}{\coqdocvariable{b2}}).\coqdoceol
\coqdocnoindent
\coqdockw{Proof}. \coqdoctac{intros}. \coqdoctac{destruct} \coqdocvar{b1}, \coqdocvar{b2}; \coqdoctac{simpl}; \coqdoctac{reflexivity}. \coqdockw{Qed}.\coqdoceol
\coqdocemptyline
\coqdocnoindent
\coqdockw{Lemma} \coqdef{SECF.Noninterference.split state orthogonal}{split\_state\_orthogonal}{\coqdoclemma{split\_state\_orthogonal}}: \coqdockw{\ensuremath{\forall}} \coqdef{SECF.Noninterference.s:166}{s}{\coqdocbinder{s}} \coqdef{SECF.Noninterference.pub:167}{pub}{\coqdocbinder{pub}} \coqdef{SECF.Noninterference.b:168}{b}{\coqdocbinder{b}},\coqdoceol
\coqdocindent{2.00em}
\coqref{SECF.Noninterference.split state}{\coqdocdefinition{split\_state}} (\coqref{SECF.Noninterference.split state}{\coqdocdefinition{split\_state}} \coqref{SECF.Noninterference.s:166}{\coqdocvariable{s}} \coqref{SECF.Noninterference.pub:167}{\coqdocvariable{pub}} \coqref{SECF.Noninterference.b:168}{\coqdocvariable{b}}) \coqref{SECF.Noninterference.pub:167}{\coqdocvariable{pub}} (\coqexternalref{negb}{http://rocq-prover.org/doc/V9.1.1/stdlib//Corelib.Init.Datatypes}{\coqdocdefinition{negb}} \coqref{SECF.Noninterference.b:168}{\coqdocvariable{b}}) \coqexternalref{::type scope:x '=' x}{http://rocq-prover.org/doc/V9.1.1/stdlib//Corelib.Init.Logic}{\coqdocnotation{=}} \coqdockw{fun} \coqdocvar{\_} \ensuremath{\Rightarrow} 0.\coqdoceol
\coqdocnoindent
\coqdockw{Proof}.\coqdoceol
\coqdocindent{1.00em}
\coqdoctac{unfold} \coqref{SECF.Noninterference.split state}{\coqdocdefinition{split\_state}}. \coqdoctac{intros}.\coqdoceol
\coqdocindent{1.00em}
\coqdoctac{eapply} \coqdoclemma{functional\_extensionality}. \coqdoctac{intro} \coqdocvar{x}.\coqdoceol
\coqdocindent{1.00em}
\coqdoctac{rewrite} \coqref{SECF.Noninterference.eqb neg distr r}{\coqdoclemma{eqb\_neg\_distr\_r}}.\coqdoceol
\coqdocindent{1.00em}
\coqdoctac{destruct} (\coqdocdefinition{Bool.eqb} (\coqdocvar{pub} \coqdocvar{x}) \coqdocvar{b}) \coqdocvar{eqn}:\coqdocvar{BOOL}; \coqdoctac{simpl}; \coqdoctac{reflexivity}.\coqdoceol
\coqdocnoindent
\coqdockw{Qed}.\coqdoceol
\coqdocemptyline
\end{coqdoccode}
First, we show a relationship between \coqref{SECF.Noninterference.sme}{\coqdocdefinition{sme}} and \coqref{SECF.Noninterference.sme state}{\coqdocdefinition{sme\_state}} using \coqref{SECF.Noninterference.split state fun}{\coqdocdefinition{split\_state\_fun}}: 
\begin{coqdoccode}
\coqdocemptyline
\coqdocnoindent
\coqdockw{Theorem} \coqdef{SECF.Noninterference.split sme state sme}{split\_sme\_state\_sme}{\coqdoclemma{split\_sme\_state\_sme}}: \coqdockw{\ensuremath{\forall}} \coqdef{SECF.Noninterference.pub:169}{pub}{\coqdocbinder{pub}} \coqdef{SECF.Noninterference.f:170}{f}{\coqdocbinder{f}},\coqdoceol
\coqdocindent{2.00em}
\coqref{SECF.Noninterference.split state fun}{\coqdocdefinition{split\_state\_fun}} \coqref{SECF.Noninterference.pub:169}{\coqdocvariable{pub}} (\coqref{SECF.Noninterference.sme state}{\coqdocdefinition{sme\_state}} \coqref{SECF.Noninterference.f:170}{\coqdocvariable{f}} \coqref{SECF.Noninterference.pub:169}{\coqdocvariable{pub}}) \coqexternalref{::type scope:x '=' x}{http://rocq-prover.org/doc/V9.1.1/stdlib//Corelib.Init.Logic}{\coqdocnotation{=}} \coqref{SECF.Noninterference.sme}{\coqdocdefinition{sme}} (\coqdockw{fun} \coqdocvar{\_} \ensuremath{\Rightarrow} 0) (\coqref{SECF.Noninterference.split state fun}{\coqdocdefinition{split\_state\_fun}} \coqref{SECF.Noninterference.pub:169}{\coqdocvariable{pub}} \coqref{SECF.Noninterference.f:170}{\coqdocvariable{f}}).\coqdoceol
\coqdocnoindent
\coqdockw{Proof}.\coqdoceol
\coqdocindent{1.00em}
\coqdoctac{intros}.\coqdoceol
\coqdocindent{1.00em}
\coqdoctac{eapply} \coqdoclemma{functional\_extensionality}. \coqdoctac{intro} \coqdocvar{PI}.\coqdoceol
\coqdocindent{1.00em}
\coqdoctac{eapply} \coqdoclemma{functional\_extensionality}. \coqdoctac{intro} \coqdocvar{SI}.\coqdoceol
\coqdocindent{1.00em}
\coqdoctac{unfold} \coqref{SECF.Noninterference.split state fun}{\coqdocdefinition{split\_state\_fun}}, \coqref{SECF.Noninterference.sme}{\coqdocdefinition{sme}}.\coqdoceol
\coqdocindent{1.00em}
\coqdoctac{rewrite} \coqexternalref{pair equal spec}{http://rocq-prover.org/doc/V9.1.1/stdlib//Corelib.Init.Datatypes}{\coqdoclemma{pair\_equal\_spec}}. \coqdoctac{split}.\coqdoceol
\coqdocindent{1.00em}
- \coqdoctac{simpl}. \coqdoctac{unfold} \coqref{SECF.Noninterference.sme state}{\coqdocdefinition{sme\_state}}.\coqdoceol
\coqdocindent{2.00em}
\coqdoctac{rewrite} \ensuremath{\leftarrow} \coqref{SECF.Noninterference.split merge public}{\coqdoclemma{split\_merge\_public}}.\coqdoceol
\coqdocindent{2.00em}
\coqdoctac{repeat} \coqdoctac{rewrite} \coqref{SECF.Noninterference.split merge split true}{\coqdoclemma{split\_merge\_split\_true}}. \coqdoctac{reflexivity}.\coqdoceol
\coqdocindent{1.00em}
- \coqdoctac{simpl}. \coqdoctac{unfold} \coqref{SECF.Noninterference.sme state}{\coqdocdefinition{sme\_state}}.\coqdoceol
\coqdocindent{2.00em}
\coqdoctac{rewrite} \coqref{SECF.Noninterference.split merge split false}{\coqdoclemma{split\_merge\_split\_false}}. \coqdoctac{reflexivity}.\coqdoceol
\coqdocnoindent
\coqdockw{Qed}.\coqdoceol
\coqdocemptyline
\end{coqdoccode}
Second, we also show a relationship between \coqref{SECF.Noninterference.sme}{\coqdocdefinition{sme}} and \coqref{SECF.Noninterference.sme state}{\coqdocdefinition{sme\_state}} using merge\_state\_fun: 
\begin{coqdoccode}
\coqdocemptyline
\coqdocnoindent
\coqdockw{Definition} \coqdef{SECF.Noninterference.merge state fun}{merge\_state\_fun}{\coqdocdefinition{merge\_state\_fun}} (\coqdef{SECF.Noninterference.pub:171}{pub}{\coqdocbinder{pub}} : \coqref{SECF.Noninterference.pub vars}{\coqdocdefinition{pub\_vars}}) (\coqdef{SECF.Noninterference.sf:172}{sf}{\coqdocbinder{sf}} : \coqref{SECF.Imp.state}{\coqdocdefinition{state}} \coqexternalref{::type scope:x '->' x}{http://rocq-prover.org/doc/V9.1.1/stdlib//Corelib.Init.Logic}{\coqdocnotation{\ensuremath{\rightarrow}}} \coqref{SECF.Imp.state}{\coqdocdefinition{state}} \coqexternalref{::type scope:x '->' x}{http://rocq-prover.org/doc/V9.1.1/stdlib//Corelib.Init.Logic}{\coqdocnotation{\ensuremath{\rightarrow}}} \coqref{SECF.Imp.state}{\coqdocdefinition{state}}\coqexternalref{::type scope:x '*' x}{http://rocq-prover.org/doc/V9.1.1/stdlib//Corelib.Init.Datatypes}{\coqdocnotation{\ensuremath{\times}}}\coqref{SECF.Imp.state}{\coqdocdefinition{state}}) :=\coqdoceol
\coqdocindent{1.00em}
\coqdockw{fun} \coqdef{SECF.Noninterference.s:173}{s}{\coqdocbinder{s}} : \coqref{SECF.Imp.state}{\coqdocdefinition{state}} \ensuremath{\Rightarrow}\coqdoceol
\coqdocindent{2.00em}
\coqdockw{let} \coqdef{SECF.Noninterference.ps:174}{ps}{\coqdocbinder{ps}} := \coqref{SECF.Noninterference.sf:172}{\coqdocvariable{sf}} (\coqref{SECF.Noninterference.split state}{\coqdocdefinition{split\_state}} \coqref{SECF.Noninterference.s:173}{\coqdocvariable{s}} \coqref{SECF.Noninterference.pub:171}{\coqdocvariable{pub}} \coqref{SECF.Imp.:::'true'}{\coqdocnotation{true}}) (\coqref{SECF.Noninterference.split state}{\coqdocdefinition{split\_state}} \coqref{SECF.Noninterference.s:173}{\coqdocvariable{s}} \coqref{SECF.Noninterference.pub:171}{\coqdocvariable{pub}} \coqref{SECF.Imp.:::'false'}{\coqdocnotation{false}}) \coqdoctac{in}\coqdoceol
\coqdocindent{2.00em}
\coqref{SECF.Noninterference.merge states}{\coqdocdefinition{merge\_states}} (\coqexternalref{fst}{http://rocq-prover.org/doc/V9.1.1/stdlib//Corelib.Init.Datatypes}{\coqdocdefinition{fst}} \coqref{SECF.Noninterference.ps:174}{\coqdocvariable{ps}}) (\coqexternalref{snd}{http://rocq-prover.org/doc/V9.1.1/stdlib//Corelib.Init.Datatypes}{\coqdocdefinition{snd}} \coqref{SECF.Noninterference.ps:174}{\coqdocvariable{ps}}) \coqref{SECF.Noninterference.pub:171}{\coqdocvariable{pub}}.\coqdoceol
\coqdocemptyline
\coqdocnoindent
\coqdockw{Theorem} \coqdef{SECF.Noninterference.merge sme state sme}{merge\_sme\_state\_sme}{\coqdoclemma{merge\_sme\_state\_sme}}: \coqdockw{\ensuremath{\forall}} \coqdef{SECF.Noninterference.pub:175}{pub}{\coqdocbinder{pub}} \coqdef{SECF.Noninterference.f:176}{f}{\coqdocbinder{f}},\coqdoceol
\coqdocindent{2.00em}
\coqref{SECF.Noninterference.sme state}{\coqdocdefinition{sme\_state}} (\coqref{SECF.Noninterference.merge state fun}{\coqdocdefinition{merge\_state\_fun}} \coqref{SECF.Noninterference.pub:175}{\coqdocvariable{pub}} \coqref{SECF.Noninterference.f:176}{\coqdocvariable{f}}) \coqref{SECF.Noninterference.pub:175}{\coqdocvariable{pub}} \coqexternalref{::type scope:x '=' x}{http://rocq-prover.org/doc/V9.1.1/stdlib//Corelib.Init.Logic}{\coqdocnotation{=}} \coqref{SECF.Noninterference.merge state fun}{\coqdocdefinition{merge\_state\_fun}} \coqref{SECF.Noninterference.pub:175}{\coqdocvariable{pub}} (\coqref{SECF.Noninterference.sme}{\coqdocdefinition{sme}} (\coqdockw{fun} \coqdocvar{\_} \ensuremath{\Rightarrow} 0) \coqref{SECF.Noninterference.f:176}{\coqdocvariable{f}}).\coqdoceol
\coqdocnoindent
\coqdockw{Proof}.\coqdoceol
\coqdocindent{1.00em}
\coqdoctac{intros}.\coqdoceol
\coqdocindent{1.00em}
\coqdoctac{eapply} \coqdoclemma{functional\_extensionality}. \coqdoctac{intro} \coqdocvar{s}.\coqdoceol
\coqdocindent{1.00em}
\coqdoctac{eapply} \coqdoclemma{functional\_extensionality}. \coqdoctac{intro} \coqdocvar{x}.\coqdoceol
\coqdocindent{1.00em}
\coqdoctac{unfold} \coqref{SECF.Noninterference.merge state fun}{\coqdocdefinition{merge\_state\_fun}}. \coqdoctac{simpl}.\coqdoceol
\coqdocindent{1.00em}
\coqdoctac{unfold} \coqref{SECF.Noninterference.sme state}{\coqdocdefinition{sme\_state}}. \coqdoctac{unfold} \coqref{SECF.Noninterference.merge states}{\coqdocdefinition{merge\_states}}.\coqdoceol
\coqdocindent{1.00em}
\coqdoctac{destruct} (\coqdocvar{pub} \coqdocvar{x}) \coqdocvar{eqn}:\coqdocvar{PUB}.\coqdoceol
\coqdocindent{1.00em}
- \coqdoctac{rewrite} \coqref{SECF.Noninterference.split state idem}{\coqdoclemma{split\_state\_idem}}. \coqdoctac{rewrite} \coqref{SECF.Noninterference.split state orthogonal}{\coqdoclemma{split\_state\_orthogonal}}. \coqdoctac{reflexivity}.\coqdoceol
\coqdocindent{1.00em}
- \coqdoctac{reflexivity}.\coqdoceol
\coqdocnoindent
\coqdockw{Qed}.\coqdoceol
\coqdocemptyline
\end{coqdoccode}
\section{Noninterference for Imp programs without loops}



 For programs without loops the ``failed attempt'' evaluation function from
    \coqdoclibrary{Imp} works well and allows us to easily define a state transformer
    function for each Imp command. 
\begin{coqdoccode}
\coqdocemptyline
\coqdocnoindent
\coqdockw{Print} \coqdocvar{ceval\_fun\_no\_while}.\coqdoceol
\coqdocnoindent
\coqdockw{Definition} \coqdef{SECF.Noninterference.flip}{flip}{\coqdocdefinition{flip}} \{\coqdef{SECF.Noninterference.A:177}{A}{\coqdocbinder{A}} \coqdef{SECF.Noninterference.B:178}{B}{\coqdocbinder{B}} \coqdef{SECF.Noninterference.C:179}{C}{\coqdocbinder{C}} : \coqdockw{Type}\} (\coqdef{SECF.Noninterference.f:180}{f}{\coqdocbinder{f}} : \coqref{SECF.Noninterference.A:177}{\coqdocvariable{A}} \coqexternalref{::type scope:x '->' x}{http://rocq-prover.org/doc/V9.1.1/stdlib//Corelib.Init.Logic}{\coqdocnotation{\ensuremath{\rightarrow}}} \coqref{SECF.Noninterference.B:178}{\coqdocvariable{B}} \coqexternalref{::type scope:x '->' x}{http://rocq-prover.org/doc/V9.1.1/stdlib//Corelib.Init.Logic}{\coqdocnotation{\ensuremath{\rightarrow}}} \coqref{SECF.Noninterference.C:179}{\coqdocvariable{C}}) := \coqdockw{fun} \coqdef{SECF.Noninterference.b:181}{b}{\coqdocbinder{b}} \coqdef{SECF.Noninterference.a:182}{a}{\coqdocbinder{a}} \ensuremath{\Rightarrow} \coqref{SECF.Noninterference.f:180}{\coqdocvariable{f}} \coqref{SECF.Noninterference.a:182}{\coqdocvariable{a}} \coqref{SECF.Noninterference.b:181}{\coqdocvariable{b}}.\coqdoceol
\coqdocnoindent
\coqdockw{Definition} \coqdef{SECF.Noninterference.cinterp}{cinterp}{\coqdocdefinition{cinterp}} : \coqref{SECF.Imp.com}{\coqdocinductive{com}} \coqexternalref{::type scope:x '->' x}{http://rocq-prover.org/doc/V9.1.1/stdlib//Corelib.Init.Logic}{\coqdocnotation{\ensuremath{\rightarrow}}} \coqref{SECF.Imp.state}{\coqdocdefinition{state}} \coqexternalref{::type scope:x '->' x}{http://rocq-prover.org/doc/V9.1.1/stdlib//Corelib.Init.Logic}{\coqdocnotation{\ensuremath{\rightarrow}}} \coqref{SECF.Imp.state}{\coqdocdefinition{state}} := \coqref{SECF.Noninterference.flip}{\coqdocdefinition{flip}} \coqref{SECF.Imp.ceval fun no while}{\coqdocdefinition{ceval\_fun\_no\_while}}.\coqdoceol
\coqdocemptyline
\coqdocnoindent
\coqdockw{Definition} \coqdef{SECF.Noninterference.noninterferent no while}{noninterferent\_no\_while}{\coqdocdefinition{noninterferent\_no\_while}} \coqdef{SECF.Noninterference.pub:183}{pub}{\coqdocbinder{pub}} \coqdef{SECF.Noninterference.c:184}{c}{\coqdocbinder{c}} : \coqdockw{Prop} :=\coqdoceol
\coqdocindent{1.00em}
\coqref{SECF.Noninterference.noninterferent state}{\coqdocdefinition{noninterferent\_state}} \coqref{SECF.Noninterference.pub:183}{\coqdocvariable{pub}} (\coqref{SECF.Noninterference.cinterp}{\coqdocdefinition{cinterp}} \coqref{SECF.Noninterference.c:184}{\coqdocvariable{c}}).\coqdoceol
\coqdocemptyline
\end{coqdoccode}
A command \coqdocvar{c} without loops is noninterferent if the state
    transformer obtained by interpreting the command with \coqref{SECF.Noninterference.cinterp}{\coqdocdefinition{cinterp}}
    maps public-equivalent States to public-equivalent states.


    Let's use this definition to prove that the following command is
    noninterferent: 
\begin{coqdoccode}
\coqdocemptyline
\coqdocnoindent
\coqdockw{Definition} \coqdef{SECF.Noninterference.xpub}{xpub}{\coqdocdefinition{xpub}} : \coqref{SECF.Noninterference.pub vars}{\coqdocdefinition{pub\_vars}} := (\coqref{SECF.Imp.X}{\coqdocdefinition{X}} \coqref{SECF.Maps.:::x '!->' x ';' x}{\coqdocnotation{!->}} \coqref{SECF.Imp.:::'true'}{\coqdocnotation{true}}\coqref{SECF.Maps.:::x '!->' x ';' x}{\coqdocnotation{;}} \coqref{SECF.Maps.:::'  ' '!->' x}{\coqdocnotation{\_\_}} \coqref{SECF.Maps.:::'  ' '!->' x}{\coqdocnotation{!->}} \coqref{SECF.Imp.:::'false'}{\coqdocnotation{false}}).\coqdoceol
\coqdocemptyline
\coqdocnoindent
\coqdockw{Definition} \coqdef{SECF.Noninterference.secure com}{secure\_com}{\coqdocdefinition{secure\_com}} : \coqref{SECF.Imp.com}{\coqdocinductive{com}} :=\coqdoceol
\coqdocindent{1.00em}
\coqref{SECF.Imp.::com scope:'<x7B' x 'x7D>'}{\coqdocnotation{<\{}} \coqref{SECF.Imp.X}{\coqdocdefinition{X}} \coqref{SECF.Imp.:com:com scope:x ':=' x}{\coqdocnotation{:=}} \coqref{SECF.Imp.X}{\coqdocdefinition{X}}\coqref{SECF.Imp.:com::x '+' x}{\coqdocnotation{+}}1\coqref{SECF.Imp.:com:com scope:x ';' x}{\coqdocnotation{;}}\coqdoceol
\coqdocindent{2.50em}
\coqref{SECF.Imp.Y}{\coqdocdefinition{Y}} \coqref{SECF.Imp.:com:com scope:x ':=' x}{\coqdocnotation{:=}} \coqref{SECF.Imp.:com::'(' x ')'}{\coqdocnotation{(}}\coqref{SECF.Imp.X}{\coqdocdefinition{X}}\coqref{SECF.Imp.:com::x '-' x}{\coqdocnotation{-}}1\coqref{SECF.Imp.:com::'(' x ')'}{\coqdocnotation{)}}\coqref{SECF.Imp.:com::x '+' x}{\coqdocnotation{+}}\coqref{SECF.Imp.Y}{\coqdocdefinition{Y}}\coqref{SECF.Imp.:com::x '*' x}{\coqdocnotation{\ensuremath{\times}}}2 \coqref{SECF.Imp.::com scope:'<x7B' x 'x7D>'}{\coqdocnotation{\}>}}.\coqdoceol
\coqdocemptyline
\end{coqdoccode}
For proving \coqref{SECF.StaticIFC.secure com}{\coqdocdefinition{secure\_com}} noninterferent we first prove a few
    helper lemmas. 
\begin{coqdoccode}
\coqdocemptyline
\coqdocnoindent
\coqdockw{Lemma} \coqdef{SECF.Noninterference.xpub true}{xpub\_true}{\coqdoclemma{xpub\_true}} : \coqdockw{\ensuremath{\forall}} \coqdef{SECF.Noninterference.x:185}{x}{\coqdocbinder{x}}, \coqref{SECF.Noninterference.xpub}{\coqdocdefinition{xpub}} \coqref{SECF.Noninterference.x:185}{\coqdocvariable{x}} \coqexternalref{::type scope:x '=' x}{http://rocq-prover.org/doc/V9.1.1/stdlib//Corelib.Init.Logic}{\coqdocnotation{=}} \coqref{SECF.Imp.:::'true'}{\coqdocnotation{true}} \coqexternalref{::type scope:x '->' x}{http://rocq-prover.org/doc/V9.1.1/stdlib//Corelib.Init.Logic}{\coqdocnotation{\ensuremath{\rightarrow}}} \coqref{SECF.Noninterference.x:185}{\coqdocvariable{x}} \coqexternalref{::type scope:x '=' x}{http://rocq-prover.org/doc/V9.1.1/stdlib//Corelib.Init.Logic}{\coqdocnotation{=}} \coqref{SECF.Imp.X}{\coqdocdefinition{X}}.\coqdoceol
\coqdocnoindent
\coqdockw{Proof}.\coqdoceol
\coqdocindent{1.00em}
\coqdoctac{unfold} \coqref{SECF.Noninterference.xpub}{\coqdocdefinition{xpub}}. \coqdoctac{intros} \coqdocvar{x} \coqdocvar{Hx}.\coqdoceol
\coqdocindent{1.00em}
\coqdoctac{destruct} (\coqdoclemma{eqb\_spec} \coqdocvar{x} \coqref{SECF.Imp.X}{\coqdocdefinition{X}}).\coqdoceol
\coqdocindent{1.00em}
- \coqdoctac{subst}. \coqdoctac{reflexivity}.\coqdoceol
\coqdocindent{1.00em}
- \coqdoctac{rewrite} \coqref{SECF.Maps.t update neq}{\coqdocaxiom{t\_update\_neq}} \coqdoctac{in} \coqdocvar{Hx}.\coqdoceol
\coqdocindent{2.00em}
+ \coqdoctac{rewrite} \coqref{SECF.Maps.t apply empty}{\coqdocaxiom{t\_apply\_empty}} \coqdoctac{in} \coqdocvar{Hx}. \coqdoctac{discriminate}.\coqdoceol
\coqdocindent{2.00em}
+ \coqdoctac{intro} \coqdocvar{contra}. \coqdoctac{subst}. \coqdocvar{contradiction}.\coqdoceol
\coqdocnoindent
\coqdockw{Qed}.\coqdoceol
\coqdocemptyline
\end{coqdoccode}
Here we are using the \coqref{SECF.Maps.t update neq}{\coqdoclemma{t\_update\_neq}} and \coqref{SECF.Maps.t apply empty}{\coqdoclemma{t\_apply\_empty}} lemmas from \coqdoclibrary{Maps} 
\begin{coqdoccode}
\coqdocemptyline
\coqdocnoindent
\coqdockw{Lemma} \coqdef{SECF.Noninterference.xpubX}{xpubX}{\coqdoclemma{xpubX}} : \coqref{SECF.Noninterference.xpub}{\coqdocdefinition{xpub}} \coqref{SECF.Imp.X}{\coqdocdefinition{X}} \coqexternalref{::type scope:x '=' x}{http://rocq-prover.org/doc/V9.1.1/stdlib//Corelib.Init.Logic}{\coqdocnotation{=}} \coqref{SECF.Imp.:::'true'}{\coqdocnotation{true}}.\coqdoceol
\coqdocnoindent
\coqdockw{Proof}. \coqdoctac{reflexivity}. \coqdockw{Qed}.\coqdoceol
\coqdocemptyline
\end{coqdoccode}
Using these lemmas the noninterference proof for \coqref{SECF.StaticIFC.secure com}{\coqdocdefinition{secure\_com}} is easy: 
\begin{coqdoccode}
\coqdocemptyline
\coqdocnoindent
\coqdockw{Lemma} \coqdef{SECF.Noninterference.noninterferent secure com}{noninterferent\_secure\_com}{\coqdoclemma{noninterferent\_secure\_com}} :\coqdoceol
\coqdocindent{1.00em}
\coqref{SECF.Noninterference.noninterferent no while}{\coqdocdefinition{noninterferent\_no\_while}} \coqref{SECF.Noninterference.xpub}{\coqdocdefinition{xpub}} \coqref{SECF.Noninterference.secure com}{\coqdocdefinition{secure\_com}}.\coqdoceol
\coqdocnoindent
\coqdockw{Proof}.\coqdoceol
\coqdocindent{1.00em}
\coqdoctac{unfold} \coqref{SECF.Noninterference.noninterferent no while}{\coqdocdefinition{noninterferent\_no\_while}}, \coqref{SECF.Noninterference.noninterferent state}{\coqdocdefinition{noninterferent\_state}}, \coqref{SECF.Noninterference.secure com}{\coqdocdefinition{secure\_com}}.\coqdoceol
\coqdocindent{1.00em}
\coqdoctac{intros} \coqdocvar{s1} \coqdocvar{s2} \coqdocvar{PEQUIV} \coqdocvar{x} \coqdocvar{Hx}.\coqdoceol
\coqdocemptyline
\coqdocindent{1.00em}
\coqdoctac{apply} \coqref{SECF.Noninterference.xpub true}{\coqdoclemma{xpub\_true}} \coqdoctac{in} \coqdocvar{Hx}. \coqdoctac{subst}.\coqdoceol
\coqdocemptyline
\coqdocindent{1.00em}
\coqdoctac{specialize} (\coqdocvar{PEQUIV} \coqref{SECF.Imp.X}{\coqdocdefinition{X}} \coqref{SECF.Noninterference.xpubX}{\coqdoclemma{xpubX}}).\coqdoceol
\coqdocemptyline
\coqdocindent{1.00em}
\coqdoctac{simpl}. \coqdoctac{rewrite} \coqdocvar{PEQUIV}. \coqdoctac{reflexivity}.\coqdoceol
\coqdocnoindent
\coqdockw{Qed}.\coqdoceol
\coqdocemptyline
\end{coqdoccode}
\paragraph{Exercise: 2 stars, standard (noninterferent\_secure\_ex1)}


\begin{coqdoccode}
\coqdocnoindent
\coqdockw{Definition} \coqdef{SECF.Noninterference.secure ex1}{secure\_ex1}{\coqdocdefinition{secure\_ex1}} :=\coqdoceol
\coqdocindent{1.00em}
\coqref{SECF.Imp.::com scope:'<x7B' x 'x7D>'}{\coqdocnotation{<\{}} \coqref{SECF.Imp.Y}{\coqdocdefinition{Y}} \coqref{SECF.Imp.:com:com scope:x ':=' x}{\coqdocnotation{:=}} \coqref{SECF.Imp.Y}{\coqdocdefinition{Y}} \coqref{SECF.Imp.:com::x '-' x}{\coqdocnotation{-}} 1\coqref{SECF.Imp.:com:com scope:x ';' x}{\coqdocnotation{;}}\coqdoceol
\coqdocindent{2.50em}
\coqref{SECF.Imp.X}{\coqdocdefinition{X}} \coqref{SECF.Imp.:com:com scope:x ':=' x}{\coqdocnotation{:=}} 1 \coqref{SECF.Imp.::com scope:'<x7B' x 'x7D>'}{\coqdocnotation{\}>}}.\coqdoceol
\coqdocemptyline
\coqdocnoindent
\coqdockw{Lemma} \coqdef{SECF.Noninterference.noninterferent secure ex1}{noninterferent\_secure\_ex1}{\coqdoclemma{noninterferent\_secure\_ex1}} :\coqdoceol
\coqdocindent{1.00em}
\coqref{SECF.Noninterference.noninterferent no while}{\coqdocdefinition{noninterferent\_no\_while}} \coqref{SECF.Noninterference.xpub}{\coqdocdefinition{xpub}} \coqref{SECF.Noninterference.secure ex1}{\coqdocdefinition{secure\_ex1}}.\coqdoceol
\coqdocnoindent
\coqdockw{Proof}.\coqdoceol
 \coqdocvar{Admitted}.\coqdoceol
\end{coqdoccode}
\ensuremath{\Box} 

\paragraph{Exercise: 3 stars, standard, optional (noninterferent\_secure\_ex2)}


\begin{coqdoccode}
\coqdocnoindent
\coqdockw{Definition} \coqdef{SECF.Noninterference.secure ex2}{secure\_ex2}{\coqdocdefinition{secure\_ex2}} :=\coqdoceol
\coqdocindent{1.00em}
\coqref{SECF.Imp.::com scope:'<x7B' x 'x7D>'}{\coqdocnotation{<\{}} \coqref{SECF.Imp.:com:com scope:'if' x 'then' x 'else' x 'end'}{\coqdocnotation{if}} \coqref{SECF.Imp.X}{\coqdocdefinition{X}} \coqref{SECF.Imp.:com::x '=' x}{\coqdocnotation{=}} 0 \coqref{SECF.Imp.:com:com scope:'if' x 'then' x 'else' x 'end'}{\coqdocnotation{then}}\coqdoceol
\coqdocindent{3.50em}
\coqref{SECF.Imp.X}{\coqdocdefinition{X}} \coqref{SECF.Imp.:com:com scope:x ':=' x}{\coqdocnotation{:=}} \coqref{SECF.Imp.X}{\coqdocdefinition{X}} \coqref{SECF.Imp.:com::x '+' x}{\coqdocnotation{+}} 5\coqdoceol
\coqdocindent{2.50em}
\coqref{SECF.Imp.:com:com scope:'if' x 'then' x 'else' x 'end'}{\coqdocnotation{else}}\coqdoceol
\coqdocindent{3.50em}
\coqref{SECF.Imp.Y}{\coqdocdefinition{Y}} \coqref{SECF.Imp.:com:com scope:x ':=' x}{\coqdocnotation{:=}} \coqref{SECF.Imp.X}{\coqdocdefinition{X}}\coqdoceol
\coqdocindent{2.50em}
\coqref{SECF.Imp.:com:com scope:'if' x 'then' x 'else' x 'end'}{\coqdocnotation{end}} \coqref{SECF.Imp.::com scope:'<x7B' x 'x7D>'}{\coqdocnotation{\}>}}.\coqdoceol
\coqdocemptyline
\coqdocnoindent
\coqdockw{Lemma} \coqdef{SECF.Noninterference.noninterferent secure ex2}{noninterferent\_secure\_ex2}{\coqdoclemma{noninterferent\_secure\_ex2}} :\coqdoceol
\coqdocindent{1.00em}
\coqref{SECF.Noninterference.noninterferent no while}{\coqdocdefinition{noninterferent\_no\_while}} \coqref{SECF.Noninterference.xpub}{\coqdocdefinition{xpub}} \coqref{SECF.Noninterference.secure ex2}{\coqdocdefinition{secure\_ex2}}.\coqdoceol
\coqdocnoindent
\coqdockw{Proof}.\coqdoceol
 \coqdocvar{Admitted}.\coqdoceol
\end{coqdoccode}
\ensuremath{\Box} 

 Now let's look at a couple of insecure commands: 
\begin{coqdoccode}
\coqdocemptyline
\coqdocnoindent
\coqdockw{Definition} \coqdef{SECF.Noninterference.insecure com1}{insecure\_com1}{\coqdocdefinition{insecure\_com1}} : \coqref{SECF.Imp.com}{\coqdocinductive{com}} :=\coqdoceol
\coqdocindent{1.00em}
\coqref{SECF.Imp.::com scope:'<x7B' x 'x7D>'}{\coqdocnotation{<\{}} \coqref{SECF.Imp.X}{\coqdocdefinition{X}} \coqref{SECF.Imp.:com:com scope:x ':=' x}{\coqdocnotation{:=}} \coqref{SECF.Imp.Y}{\coqdocdefinition{Y}}\coqref{SECF.Imp.:com::x '+' x}{\coqdocnotation{+}}1\coqref{SECF.Imp.:com:com scope:x ';' x}{\coqdocnotation{;}} \coqdoceol
\coqdocindent{2.50em}
\coqref{SECF.Imp.Y}{\coqdocdefinition{Y}} \coqref{SECF.Imp.:com:com scope:x ':=' x}{\coqdocnotation{:=}} \coqref{SECF.Imp.:com::'(' x ')'}{\coqdocnotation{(}}\coqref{SECF.Imp.X}{\coqdocdefinition{X}}\coqref{SECF.Imp.:com::x '-' x}{\coqdocnotation{-}}1\coqref{SECF.Imp.:com::'(' x ')'}{\coqdocnotation{)}}\coqref{SECF.Imp.:com::x '+' x}{\coqdocnotation{+}}\coqref{SECF.Imp.Y}{\coqdocdefinition{Y}}\coqref{SECF.Imp.:com::x '*' x}{\coqdocnotation{\ensuremath{\times}}}2 \coqref{SECF.Imp.::com scope:'<x7B' x 'x7D>'}{\coqdocnotation{\}>}}.\coqdoceol
\coqdocemptyline
\end{coqdoccode}
An \textit{explicit flow} is when a command directly assigns an expression
    depending on secret variables to a public variable, like the \coqref{SECF.SpecCT.X}{\coqdocdefinition{X}} := \coqref{SECF.SpecCT.Y}{\coqdocdefinition{Y}}+1
    assignment above. Explicit flows are easier to find automatically
    and even simple taint-tracking would be enough for discovering this.


    We prove that \coqref{SECF.StaticIFC.insecure com1}{\coqdocdefinition{insecure\_com1}} is interferent as follows: 
\begin{coqdoccode}
\coqdocemptyline
\coqdocnoindent
\coqdockw{Lemma} \coqdef{SECF.Noninterference.interferent insecure com1}{interferent\_insecure\_com1}{\coqdoclemma{interferent\_insecure\_com1}} :\coqdoceol
\coqdocindent{1.00em}
\coqexternalref{::type scope:'x7E' x}{http://rocq-prover.org/doc/V9.1.1/stdlib//Corelib.Init.Logic}{\coqdocnotation{\ensuremath{\lnot}}}\coqref{SECF.Noninterference.noninterferent no while}{\coqdocdefinition{noninterferent\_no\_while}} \coqref{SECF.Noninterference.xpub}{\coqdocdefinition{xpub}} \coqref{SECF.Noninterference.insecure com1}{\coqdocdefinition{insecure\_com1}}.\coqdoceol
\coqdocnoindent
\coqdockw{Proof}.\coqdoceol
\coqdocindent{1.00em}
\coqdoctac{unfold} \coqref{SECF.Noninterference.noninterferent no while}{\coqdocdefinition{noninterferent\_no\_while}}, \coqref{SECF.Noninterference.noninterferent state}{\coqdocdefinition{noninterferent\_state}}, \coqref{SECF.Noninterference.insecure com1}{\coqdocdefinition{insecure\_com1}}.\coqdoceol
\coqdocindent{1.00em}
\coqdoctac{intro} \coqdocvar{Hc}.\coqdoceol
\coqdocemptyline
\coqdocindent{1.00em}
\coqdoctac{set} (\coqdocvar{s1} := (\coqref{SECF.Imp.X}{\coqdocdefinition{X}} \coqref{SECF.Maps.:::x '!->' x ';' x}{\coqdocnotation{!->}} 0 \coqref{SECF.Maps.:::x '!->' x ';' x}{\coqdocnotation{;}} \coqref{SECF.Imp.Y}{\coqdocdefinition{Y}} \coqref{SECF.Imp.:::x '!->' x}{\coqdocnotation{!->}} 0)).\coqdoceol
\coqdocindent{1.00em}
\coqdoctac{set} (\coqdocvar{s2} := (\coqref{SECF.Imp.X}{\coqdocdefinition{X}} \coqref{SECF.Maps.:::x '!->' x ';' x}{\coqdocnotation{!->}} 0 \coqref{SECF.Maps.:::x '!->' x ';' x}{\coqdocnotation{;}} \coqref{SECF.Imp.Y}{\coqdocdefinition{Y}} \coqref{SECF.Imp.:::x '!->' x}{\coqdocnotation{!->}} 1)).\coqdoceol
\coqdocindent{1.00em}
\coqdoctac{specialize} (\coqdocvar{Hc} \coqdocvar{s1} \coqdocvar{s2}).\coqdoceol
\coqdocemptyline
\coqdocindent{1.00em}
\coqdoctac{assert} (\coqdocvar{PEQUIV}: \coqref{SECF.Noninterference.pub equiv}{\coqdocdefinition{pub\_equiv}} \coqref{SECF.Noninterference.xpub}{\coqdocdefinition{xpub}} \coqdocvar{s1} \coqdocvar{s2}).\coqdoceol
\coqdocindent{1.00em}
\{ \coqdoctac{clear} \coqdocvar{Hc}. \coqdoctac{intros} \coqdocvar{x} \coqdocvar{H}. \coqdoctac{apply} \coqref{SECF.Noninterference.xpub true}{\coqdoclemma{xpub\_true}} \coqdoctac{in} \coqdocvar{H}. \coqdoctac{subst}. \coqdoctac{reflexivity}. \}\coqdoceol
\coqdocnoindent
\coqdoceol
\coqdocindent{1.00em}
\coqdoctac{specialize} (\coqdocvar{Hc} \coqdocvar{PEQUIV} \coqref{SECF.Imp.X}{\coqdocdefinition{X}} \coqref{SECF.Noninterference.xpubX}{\coqdoclemma{xpubX}}).\coqdoceol
\coqdocemptyline
\coqdocindent{1.00em}
\coqdoctac{simpl} \coqdoctac{in} \coqdocvar{Hc}. \coqdoctac{unfold} \coqdocvar{s1}, \coqdocvar{s2}, \coqref{SECF.Maps.t update}{\coqdocdefinition{t\_update}} \coqdoctac{in} \coqdocvar{Hc}. \coqdoctac{simpl} \coqdoctac{in} \coqdocvar{Hc}.\coqdoceol
\coqdocemptyline
\coqdocindent{1.00em}
\coqdoctac{discriminate} \coqdocvar{Hc}.\coqdoceol
\coqdocnoindent
\coqdockw{Qed}.\coqdoceol
\coqdocemptyline
\end{coqdoccode}
As we saw above, the \coqdoctac{set} tactic allows us to give names to
    complex expression, making proofs more readable and
    manageable. It's particularly useful when constructing concrete
    counterexamples where one needs to work with specific values. 

\paragraph{Exercise: 2 stars, standard (interferent\_insecure\_com\_explicit)}


\begin{coqdoccode}
\coqdocnoindent
\coqdockw{Definition} \coqdef{SECF.Noninterference.insecure com explicit}{insecure\_com\_explicit}{\coqdocdefinition{insecure\_com\_explicit}} :=\coqdoceol
\coqdocindent{1.00em}
\coqref{SECF.Imp.::com scope:'<x7B' x 'x7D>'}{\coqdocnotation{<\{}} \coqref{SECF.Imp.X}{\coqdocdefinition{X}} \coqref{SECF.Imp.:com:com scope:x ':=' x}{\coqdocnotation{:=}} \coqref{SECF.Imp.Y}{\coqdocdefinition{Y}} \coqref{SECF.Imp.:com::x '+' x}{\coqdocnotation{+}} \coqref{SECF.Imp.X}{\coqdocdefinition{X}}\coqref{SECF.Imp.:com:com scope:x ';' x}{\coqdocnotation{;}} \coqdoceol
\coqdocindent{2.50em}
\coqref{SECF.Imp.Y}{\coqdocdefinition{Y}} \coqref{SECF.Imp.:com:com scope:x ':=' x}{\coqdocnotation{:=}} \coqref{SECF.Imp.Y}{\coqdocdefinition{Y}} \coqref{SECF.Imp.:com::x '-' x}{\coqdocnotation{-}} 1 \coqref{SECF.Imp.::com scope:'<x7B' x 'x7D>'}{\coqdocnotation{\}>}}.\coqdoceol
\coqdocemptyline
\coqdocnoindent
\coqdockw{Lemma} \coqdef{SECF.Noninterference.interferent insecure com explicit}{interferent\_insecure\_com\_explicit}{\coqdoclemma{interferent\_insecure\_com\_explicit}} :\coqdoceol
\coqdocindent{1.00em}
\coqexternalref{::type scope:'x7E' x}{http://rocq-prover.org/doc/V9.1.1/stdlib//Corelib.Init.Logic}{\coqdocnotation{\ensuremath{\lnot}}}\coqref{SECF.Noninterference.noninterferent no while}{\coqdocdefinition{noninterferent\_no\_while}} \coqref{SECF.Noninterference.xpub}{\coqdocdefinition{xpub}} \coqref{SECF.Noninterference.insecure com explicit}{\coqdocdefinition{insecure\_com\_explicit}}.\coqdoceol
\coqdocnoindent
\coqdockw{Proof}.\coqdoceol
 \coqdocvar{Admitted}.\coqdoceol
\end{coqdoccode}
\ensuremath{\Box} 

 Noninterference can be violated not only by explicit flows, but also by
    \textit{implicit flows}, which leak secret information via the control-flow of the
    program. Here is a simple example: 
\begin{coqdoccode}
\coqdocemptyline
\coqdocnoindent
\coqdockw{Definition} \coqdef{SECF.Noninterference.insecure com2}{insecure\_com2}{\coqdocdefinition{insecure\_com2}} : \coqref{SECF.Imp.com}{\coqdocinductive{com}} :=\coqdoceol
\coqdocindent{1.00em}
\coqref{SECF.Imp.::com scope:'<x7B' x 'x7D>'}{\coqdocnotation{<\{}} \coqref{SECF.Imp.:com:com scope:'if' x 'then' x 'else' x 'end'}{\coqdocnotation{if}} \coqref{SECF.Imp.Y}{\coqdocdefinition{Y}} \coqref{SECF.Imp.:com::x '=' x}{\coqdocnotation{=}} 0 \coqref{SECF.Imp.:com:com scope:'if' x 'then' x 'else' x 'end'}{\coqdocnotation{then}}\coqdoceol
\coqdocindent{3.50em}
\coqref{SECF.Imp.Y}{\coqdocdefinition{Y}} \coqref{SECF.Imp.:com:com scope:x ':=' x}{\coqdocnotation{:=}} 42\coqdoceol
\coqdocindent{2.50em}
\coqref{SECF.Imp.:com:com scope:'if' x 'then' x 'else' x 'end'}{\coqdocnotation{else}}\coqdoceol
\coqdocindent{3.50em}
\coqref{SECF.Imp.X}{\coqdocdefinition{X}} \coqref{SECF.Imp.:com:com scope:x ':=' x}{\coqdocnotation{:=}} \coqref{SECF.Imp.X}{\coqdocdefinition{X}}\coqref{SECF.Imp.:com::x '+' x}{\coqdocnotation{+}}1 \coqdoceol
\coqdocindent{2.50em}
\coqref{SECF.Imp.:com:com scope:'if' x 'then' x 'else' x 'end'}{\coqdocnotation{end}} \coqref{SECF.Imp.::com scope:'<x7B' x 'x7D>'}{\coqdocnotation{\}>}}.\coqdoceol
\coqdocemptyline
\end{coqdoccode}
Here the expression \coqref{SECF.SpecCT.X}{\coqdocdefinition{X}}+1 we are assigning to \coqref{SECF.SpecCT.X}{\coqdocdefinition{X}} is public information, but
    we are doing this assignment after we branched on a secret condition \coqref{SECF.SpecCT.Y}{\coqdocdefinition{Y}} =
    0, so we are indirectly leaking information about the value of \coqref{SECF.SpecCT.Y}{\coqdocdefinition{Y}}. In this
    case we can infer that if \coqref{SECF.SpecCT.X}{\coqdocdefinition{X}} gets incremented the value of \coqref{SECF.SpecCT.Y}{\coqdocdefinition{Y}} is not 0. 
\begin{coqdoccode}
\coqdocemptyline
\coqdocnoindent
\coqdockw{Lemma} \coqdef{SECF.Noninterference.interferent insecure com2}{interferent\_insecure\_com2}{\coqdoclemma{interferent\_insecure\_com2}} :\coqdoceol
\coqdocindent{1.00em}
\coqexternalref{::type scope:'x7E' x}{http://rocq-prover.org/doc/V9.1.1/stdlib//Corelib.Init.Logic}{\coqdocnotation{\ensuremath{\lnot}}}\coqref{SECF.Noninterference.noninterferent no while}{\coqdocdefinition{noninterferent\_no\_while}} \coqref{SECF.Noninterference.xpub}{\coqdocdefinition{xpub}} \coqref{SECF.Noninterference.insecure com2}{\coqdocdefinition{insecure\_com2}}.\coqdoceol
\coqdocnoindent
\coqdockw{Proof}.\coqdoceol
\coqdocindent{1.00em}
\coqdoctac{unfold} \coqref{SECF.Noninterference.noninterferent no while}{\coqdocdefinition{noninterferent\_no\_while}}, \coqref{SECF.Noninterference.noninterferent state}{\coqdocdefinition{noninterferent\_state}}, \coqref{SECF.Noninterference.insecure com1}{\coqdocdefinition{insecure\_com1}}.\coqdoceol
\coqdocindent{1.00em}
\coqdoctac{intro} \coqdocvar{Hc}.\coqdoceol
\coqdocemptyline
\coqdocindent{1.00em}
\coqdoctac{set} (\coqdocvar{s1} := (\coqref{SECF.Imp.X}{\coqdocdefinition{X}} \coqref{SECF.Maps.:::x '!->' x ';' x}{\coqdocnotation{!->}} 0 \coqref{SECF.Maps.:::x '!->' x ';' x}{\coqdocnotation{;}} \coqref{SECF.Imp.Y}{\coqdocdefinition{Y}} \coqref{SECF.Imp.:::x '!->' x}{\coqdocnotation{!->}} 0)).\coqdoceol
\coqdocindent{1.00em}
\coqdoctac{set} (\coqdocvar{s2} := (\coqref{SECF.Imp.X}{\coqdocdefinition{X}} \coqref{SECF.Maps.:::x '!->' x ';' x}{\coqdocnotation{!->}} 0 \coqref{SECF.Maps.:::x '!->' x ';' x}{\coqdocnotation{;}} \coqref{SECF.Imp.Y}{\coqdocdefinition{Y}} \coqref{SECF.Imp.:::x '!->' x}{\coqdocnotation{!->}} 1)).\coqdoceol
\coqdocindent{1.00em}
\coqdoctac{specialize} (\coqdocvar{Hc} \coqdocvar{s1} \coqdocvar{s2}).\coqdoceol
\coqdocemptyline
\coqdocindent{1.00em}
\coqdoctac{assert} (\coqdocvar{PEQUIV}: \coqref{SECF.Noninterference.pub equiv}{\coqdocdefinition{pub\_equiv}} \coqref{SECF.Noninterference.xpub}{\coqdocdefinition{xpub}} \coqdocvar{s1} \coqdocvar{s2}).\coqdoceol
\coqdocindent{1.00em}
\{ \coqdoctac{clear} \coqdocvar{Hc}. \coqdoctac{intros} \coqdocvar{x} \coqdocvar{H}. \coqdoctac{apply} \coqref{SECF.Noninterference.xpub true}{\coqdoclemma{xpub\_true}} \coqdoctac{in} \coqdocvar{H}. \coqdoctac{subst}. \coqdoctac{reflexivity}. \}\coqdoceol
\coqdocnoindent
\coqdoceol
\coqdocindent{1.00em}
\coqdoctac{specialize} (\coqdocvar{Hc} \coqdocvar{PEQUIV} \coqref{SECF.Imp.X}{\coqdocdefinition{X}} \coqref{SECF.Noninterference.xpubX}{\coqdoclemma{xpubX}}).\coqdoceol
\coqdocemptyline
\coqdocindent{1.00em}
\coqdoctac{simpl} \coqdoctac{in} \coqdocvar{Hc}. \coqdoctac{unfold} \coqdocvar{s1}, \coqdocvar{s2}, \coqref{SECF.Maps.t update}{\coqdocdefinition{t\_update}} \coqdoctac{in} \coqdocvar{Hc}. \coqdoctac{simpl} \coqdoctac{in} \coqdocvar{Hc}.\coqdoceol
\coqdocemptyline
\coqdocindent{1.00em}
\coqdoctac{discriminate} \coqdocvar{Hc}.\coqdoceol
\coqdocnoindent
\coqdockw{Qed}.\coqdoceol
\coqdocemptyline
\end{coqdoccode}
\paragraph{Exercise: 3 stars, standard (interferent\_insecure\_com\_implicit)}


\begin{coqdoccode}
\coqdocnoindent
\coqdockw{Definition} \coqdef{SECF.Noninterference.insecure com implicit}{insecure\_com\_implicit}{\coqdocdefinition{insecure\_com\_implicit}} :=\coqdoceol
\coqdocindent{1.00em}
\coqref{SECF.Imp.::com scope:'<x7B' x 'x7D>'}{\coqdocnotation{<\{}} \coqref{SECF.Imp.:com:com scope:'if' x 'then' x 'else' x 'end'}{\coqdocnotation{if}} \coqref{SECF.Imp.Y}{\coqdocdefinition{Y}} \coqref{SECF.Imp.:com::x '=' x}{\coqdocnotation{=}} 42 \coqref{SECF.Imp.:com:com scope:'if' x 'then' x 'else' x 'end'}{\coqdocnotation{then}}\coqdoceol
\coqdocindent{3.50em}
\coqref{SECF.Imp.X}{\coqdocdefinition{X}} \coqref{SECF.Imp.:com:com scope:x ':=' x}{\coqdocnotation{:=}} \coqref{SECF.Imp.X}{\coqdocdefinition{X}} \coqref{SECF.Imp.:com::x '-' x}{\coqdocnotation{-}} 1 \coqdoceol
\coqdocindent{2.50em}
\coqref{SECF.Imp.:com:com scope:'if' x 'then' x 'else' x 'end'}{\coqdocnotation{else}}\coqdoceol
\coqdocindent{3.50em}
\coqref{SECF.Imp.Y}{\coqdocdefinition{Y}} \coqref{SECF.Imp.:com:com scope:x ':=' x}{\coqdocnotation{:=}} 2 \coqref{SECF.Imp.:com::x '*' x}{\coqdocnotation{\ensuremath{\times}}} \coqref{SECF.Imp.Y}{\coqdocdefinition{Y}}\coqdoceol
\coqdocindent{2.50em}
\coqref{SECF.Imp.:com:com scope:'if' x 'then' x 'else' x 'end'}{\coqdocnotation{end}} \coqref{SECF.Imp.::com scope:'<x7B' x 'x7D>'}{\coqdocnotation{\}>}}.\coqdoceol
\coqdocemptyline
\coqdocnoindent
\coqdockw{Lemma} \coqdef{SECF.Noninterference.interferent insecure com implicit}{interferent\_insecure\_com\_implicit}{\coqdoclemma{interferent\_insecure\_com\_implicit}} :\coqdoceol
\coqdocindent{1.00em}
\coqexternalref{::type scope:'x7E' x}{http://rocq-prover.org/doc/V9.1.1/stdlib//Corelib.Init.Logic}{\coqdocnotation{\ensuremath{\lnot}}}\coqref{SECF.Noninterference.noninterferent no while}{\coqdocdefinition{noninterferent\_no\_while}} \coqref{SECF.Noninterference.xpub}{\coqdocdefinition{xpub}} \coqref{SECF.Noninterference.insecure com implicit}{\coqdocdefinition{insecure\_com\_implicit}}.\coqdoceol
\coqdocnoindent
\coqdockw{Proof}.\coqdoceol
 \coqdocvar{Admitted}.\coqdoceol
\end{coqdoccode}
\ensuremath{\Box} 

 We will return to explicit and implicit flows in the \coqdoclibrary{StaticIFC} chapter. 
\begin{coqdoccode}
\coqdocemptyline
\end{coqdoccode}
\section{SME for Imp programs without loops}



 We can use \coqref{SECF.Noninterference.sme state}{\coqdocdefinition{sme\_state}} to execute such programs to obtain a
    noninterferent state transformer by running programs 2 times, once
    on a state where the secrets were zeroed out and once on the
    original input state, and then merging the final states. 
\begin{coqdoccode}
\coqdocemptyline
\coqdocnoindent
\coqdockw{Print} \coqdocvar{sme\_state}.\coqdoceol
\coqdocemptyline
\coqdocnoindent
\coqdockw{Definition} \coqdef{SECF.Noninterference.sme cmd}{sme\_cmd}{\coqdocdefinition{sme\_cmd}} \coqdef{SECF.Noninterference.c:186}{c}{\coqdocbinder{c}} : \coqref{SECF.Noninterference.pub vars}{\coqdocdefinition{pub\_vars}}\coqexternalref{::type scope:x '->' x}{http://rocq-prover.org/doc/V9.1.1/stdlib//Corelib.Init.Logic}{\coqdocnotation{\ensuremath{\rightarrow}}}\coqref{SECF.Imp.state}{\coqdocdefinition{state}}\coqexternalref{::type scope:x '->' x}{http://rocq-prover.org/doc/V9.1.1/stdlib//Corelib.Init.Logic}{\coqdocnotation{\ensuremath{\rightarrow}}}\coqref{SECF.Imp.state}{\coqdocdefinition{state}} := \coqref{SECF.Noninterference.sme state}{\coqdocdefinition{sme\_state}} (\coqref{SECF.Noninterference.cinterp}{\coqdocdefinition{cinterp}} \coqref{SECF.Noninterference.c:186}{\coqdocvariable{c}}).\coqdoceol
\coqdocemptyline
\end{coqdoccode}
The result of applying \coqref{SECF.Noninterference.sme cmd}{\coqdocdefinition{sme\_cmd}} to a program is not a program,
    but a state transformer. We prove noninterference and transparency
    for the state transformers obtained by \coqref{SECF.Noninterference.sme cmd}{\coqdocdefinition{sme\_cmd}} using our
    noninterference and transparency theorems about \coqref{SECF.Noninterference.sme state}{\coqdocdefinition{sme\_state}}: 
\begin{coqdoccode}
\coqdocemptyline
\coqdocnoindent
\coqdockw{Theorem} \coqdef{SECF.Noninterference.noninterferent sme cmd}{noninterferent\_sme\_cmd}{\coqdoclemma{noninterferent\_sme\_cmd}} : \coqdockw{\ensuremath{\forall}} \coqdef{SECF.Noninterference.c:187}{c}{\coqdocbinder{c}} \coqdef{SECF.Noninterference.pub:188}{pub}{\coqdocbinder{pub}},\coqdoceol
\coqdocindent{1.00em}
\coqref{SECF.Noninterference.noninterferent state}{\coqdocdefinition{noninterferent\_state}} \coqref{SECF.Noninterference.pub:188}{\coqdocvariable{pub}} (\coqref{SECF.Noninterference.sme cmd}{\coqdocdefinition{sme\_cmd}} \coqref{SECF.Noninterference.c:187}{\coqdocvariable{c}} \coqref{SECF.Noninterference.pub:188}{\coqdocvariable{pub}}).\coqdoceol
\coqdocnoindent
\coqdockw{Proof}. \coqdoctac{intros} \coqdocvar{c} \coqdocvar{p}. \coqdoctac{apply} \coqref{SECF.Noninterference.noninterferent sme state}{\coqdoclemma{noninterferent\_sme\_state}}. \coqdockw{Qed}.\coqdoceol
\coqdocemptyline
\coqdocnoindent
\coqdockw{Theorem} \coqdef{SECF.Noninterference.transparent sme cmd}{transparent\_sme\_cmd}{\coqdoclemma{transparent\_sme\_cmd}} : \coqdockw{\ensuremath{\forall}} \coqdef{SECF.Noninterference.c:189}{c}{\coqdocbinder{c}} \coqdef{SECF.Noninterference.pub:190}{pub}{\coqdocbinder{pub}},\coqdoceol
\coqdocindent{2.00em}
\coqref{SECF.Noninterference.noninterferent state}{\coqdocdefinition{noninterferent\_state}} \coqref{SECF.Noninterference.pub:190}{\coqdocvariable{pub}} (\coqdockw{fun} \coqdef{SECF.Noninterference.s:191}{s}{\coqdocbinder{s}} \ensuremath{\Rightarrow} \coqref{SECF.Imp.ceval fun no while}{\coqdocdefinition{ceval\_fun\_no\_while}} \coqref{SECF.Noninterference.s:191}{\coqdocvariable{s}} \coqref{SECF.Noninterference.c:189}{\coqdocvariable{c}}) \coqexternalref{::type scope:x '->' x}{http://rocq-prover.org/doc/V9.1.1/stdlib//Corelib.Init.Logic}{\coqdocnotation{\ensuremath{\rightarrow}}}\coqdoceol
\coqdocindent{2.00em}
\coqdockw{\ensuremath{\forall}} \coqdef{SECF.Noninterference.s:192}{s}{\coqdocbinder{s}}, \coqref{SECF.Noninterference.cinterp}{\coqdocdefinition{cinterp}} \coqref{SECF.Noninterference.c:189}{\coqdocvariable{c}} \coqref{SECF.Noninterference.s:192}{\coqdocvariable{s}} \coqexternalref{::type scope:x '=' x}{http://rocq-prover.org/doc/V9.1.1/stdlib//Corelib.Init.Logic}{\coqdocnotation{=}} \coqref{SECF.Noninterference.sme cmd}{\coqdocdefinition{sme\_cmd}} \coqref{SECF.Noninterference.c:189}{\coqdocvariable{c}} \coqref{SECF.Noninterference.pub:190}{\coqdocvariable{pub}} \coqref{SECF.Noninterference.s:192}{\coqdocvariable{s}}.\coqdoceol
\coqdocnoindent
\coqdockw{Proof}.\coqdoceol
\coqdocindent{1.00em}
\coqdoctac{unfold} \coqref{SECF.Noninterference.sme cmd}{\coqdocdefinition{sme\_cmd}}. \coqdoctac{intros} \coqdocvar{c} \coqdocvar{pub} \coqdocvar{NI}. \coqdoctac{apply} \coqref{SECF.Noninterference.transparent sme state}{\coqdoclemma{transparent\_sme\_state}}. \coqdoctac{apply} \coqdocvar{NI}.\coqdoceol
\coqdocnoindent
\coqdockw{Qed}.\coqdoceol
\coqdocemptyline
\end{coqdoccode}
Perhaps more interesting is to look at how \coqref{SECF.Noninterference.sme cmd}{\coqdocdefinition{sme\_cmd}}
    changes the behavior of some insecure commands: 
\begin{coqdoccode}
\coqdocemptyline
\coqdocnoindent
\coqdockw{Print} \coqdocvar{insecure\_com1}. \coqdocnoindent
\coqdockw{Definition} \coqdef{SECF.Noninterference.secure com1}{secure\_com1}{\coqdocdefinition{secure\_com1}} : \coqref{SECF.Imp.com}{\coqdocinductive{com}} :=\coqdoceol
\coqdocindent{1.00em}
\coqref{SECF.Imp.::com scope:'<x7B' x 'x7D>'}{\coqdocnotation{<\{}} \coqref{SECF.Imp.X}{\coqdocdefinition{X}} \coqref{SECF.Imp.:com:com scope:x ':=' x}{\coqdocnotation{:=}} 1\coqref{SECF.Imp.:com:com scope:x ';' x}{\coqdocnotation{;}} \coqdoceol
\coqdocindent{2.50em}
\coqref{SECF.Imp.Y}{\coqdocdefinition{Y}} \coqref{SECF.Imp.:com:com scope:x ':=' x}{\coqdocnotation{:=}} \coqref{SECF.Imp.Y}{\coqdocdefinition{Y}}\coqref{SECF.Imp.:com::x '*' x}{\coqdocnotation{\ensuremath{\times}}}3  \coqref{SECF.Imp.::com scope:'<x7B' x 'x7D>'}{\coqdocnotation{\}>}}.\coqdoceol
\coqdocemptyline
\coqdocnoindent
\coqdockw{Lemma} \coqdef{SECF.Noninterference.sme insecure com1}{sme\_insecure\_com1}{\coqdoclemma{sme\_insecure\_com1}} : \coqref{SECF.Noninterference.sme cmd}{\coqdocdefinition{sme\_cmd}} \coqref{SECF.Noninterference.insecure com1}{\coqdocdefinition{insecure\_com1}} \coqref{SECF.Noninterference.xpub}{\coqdocdefinition{xpub}} \coqexternalref{::type scope:x '=' x}{http://rocq-prover.org/doc/V9.1.1/stdlib//Corelib.Init.Logic}{\coqdocnotation{=}} \coqref{SECF.Noninterference.cinterp}{\coqdocdefinition{cinterp}} \coqref{SECF.Noninterference.secure com1}{\coqdocdefinition{secure\_com1}}.\coqdoceol
\coqdocnoindent
\coqdockw{Proof}.\coqdoceol
\coqdocindent{1.00em}
\coqdoctac{eapply} \coqdoclemma{functional\_extensionality}. \coqdoctac{intros} \coqdocvar{st}.\coqdoceol
\coqdocindent{1.00em}
\coqdoctac{unfold} \coqref{SECF.Noninterference.sme cmd}{\coqdocdefinition{sme\_cmd}}, \coqref{SECF.Noninterference.sme state}{\coqdocdefinition{sme\_state}}, \coqref{SECF.Noninterference.insecure com1}{\coqdocdefinition{insecure\_com1}}. \coqdoctac{simpl}.\coqdoceol
\coqdocindent{1.00em}
\coqdoctac{eapply} \coqdoclemma{functional\_extensionality}. \coqdoctac{intros} \coqdocvar{x}.\coqdoceol
\coqdocindent{1.00em}
\coqdoctac{unfold} \coqref{SECF.Noninterference.merge states}{\coqdocdefinition{merge\_states}}. \coqdoctac{simpl}.\coqdoceol
\coqdocemptyline
\coqdocindent{1.00em}
\coqdoctac{destruct} (\coqref{SECF.Noninterference.xpub}{\coqdocdefinition{xpub}} \coqdocvar{x}) \coqdocvar{eqn}:\coqdocvar{HXP}.\coqdoceol
\coqdocindent{1.00em}
\{ \coqdoctac{eapply} \coqref{SECF.Noninterference.xpub true}{\coqdoclemma{xpub\_true}} \coqdoctac{in} \coqdocvar{HXP}. \coqdoctac{subst}.\coqdoceol
\coqdocindent{2.00em}
\coqdoctac{rewrite} \coqref{SECF.Maps.t update neq}{\coqdocaxiom{t\_update\_neq}}; \coqdoctac{try} (\coqdoctac{intros} \coqdocvar{Hcontra}; \coqdoctac{discriminate}).\coqdoceol
\coqdocindent{2.00em}
\coqdoctac{rewrite} \coqref{SECF.Maps.t update eq}{\coqdocaxiom{t\_update\_eq}}; \coqdoctac{simpl}; \coqdoctac{auto}. \}\coqdoceol
\coqdocnoindent
\coqdoceol
\coqdocindent{1.00em}
\coqdoctac{destruct} (\coqdocdefinition{eqb} \coqdocvar{x} \coqref{SECF.Imp.Y}{\coqdocdefinition{Y}}) \coqdocvar{eqn}:\coqdocvar{HY}.\coqdoceol
\coqdocindent{1.00em}
\{ \coqdoctac{rewrite} \coqdoclemma{eqb\_eq} \coqdoctac{in} \coqdocvar{HY}. \coqdoctac{subst}.\coqdoceol
\coqdocindent{2.00em}
\coqdoctac{repeat} \coqdoctac{rewrite} \coqref{SECF.Maps.t update eq}{\coqdocaxiom{t\_update\_eq}}.\coqdoceol
\coqdocindent{2.00em}
\coqdoctac{repeat} \coqdoctac{rewrite} \coqref{SECF.Maps.t update neq}{\coqdocaxiom{t\_update\_neq}}; \coqdoctac{try} (\coqdoctac{intros} \coqdocvar{Hcontra}; \coqdoctac{discriminate}).\coqdoceol
\coqdocindent{2.00em}
\coqdocvar{lia}. \}\coqdoceol
\coqdocnoindent
\coqdoceol
\coqdocindent{1.00em}
\coqdoctac{assert} (\coqdocvar{HX}: \coqdocvar{x} \coqexternalref{::type scope:x '<>' x}{http://rocq-prover.org/doc/V9.1.1/stdlib//Corelib.Init.Logic}{\coqdocnotation{\ensuremath{\not=}}} \coqref{SECF.Imp.X}{\coqdocdefinition{X}}).\coqdoceol
\coqdocindent{1.00em}
\{ \coqdoctac{intros} \coqdocvar{Hx}. \coqdoctac{subst}. \coqdoctac{rewrite} \coqref{SECF.Noninterference.xpubX}{\coqdoclemma{xpubX}} \coqdoctac{in} \coqdocvar{HXP}. \coqdoctac{discriminate}. \}\coqdoceol
\coqdocnoindent
\coqdoceol
\coqdocindent{1.00em}
\coqdoctac{rewrite} \coqdoclemma{eqb\_neq} \coqdoctac{in} \coqdocvar{HY}.\coqdoceol
\coqdocindent{1.00em}
\coqdoctac{repeat} (\coqdoctac{rewrite} \coqref{SECF.Maps.t update neq}{\coqdocaxiom{t\_update\_neq}}; \coqdoctac{auto}).\coqdoceol
\coqdocnoindent
\coqdockw{Qed}.\coqdoceol
\coqdocemptyline
\end{coqdoccode}
The example above shows that the effect of applying \coqref{SECF.Noninterference.sme cmd}{\coqdocdefinition{sme\_cmd}} is
    hard to predict statically and it is not just a simple syntactic
    transformation of the original command. Here is another example of that: 
\begin{coqdoccode}
\coqdocemptyline
\coqdocnoindent
\coqdockw{Definition} \coqdef{SECF.Noninterference.insecure com2'}{insecure\_com2'}{\coqdocdefinition{insecure\_com2'}} : \coqref{SECF.Imp.com}{\coqdocinductive{com}} :=\coqdoceol
\coqdocindent{1.00em}
\coqref{SECF.Imp.::com scope:'<x7B' x 'x7D>'}{\coqdocnotation{<\{}} \coqref{SECF.Imp.:com:com scope:'if' x 'then' x 'else' x 'end'}{\coqdocnotation{if}} \coqref{SECF.Imp.Y}{\coqdocdefinition{Y}} \coqref{SECF.Imp.:com::x '=' x}{\coqdocnotation{=}} 0 \coqref{SECF.Imp.:com:com scope:'if' x 'then' x 'else' x 'end'}{\coqdocnotation{then}}\coqdoceol
\coqdocindent{3.50em}
\coqref{SECF.Imp.X}{\coqdocdefinition{X}} \coqref{SECF.Imp.:com:com scope:x ':=' x}{\coqdocnotation{:=}} 42  \coqdoceol
\coqdocindent{2.50em}
\coqref{SECF.Imp.:com:com scope:'if' x 'then' x 'else' x 'end'}{\coqdocnotation{else}}\coqdoceol
\coqdocindent{3.50em}
\coqref{SECF.Imp.X}{\coqdocdefinition{X}} \coqref{SECF.Imp.:com:com scope:x ':=' x}{\coqdocnotation{:=}} \coqref{SECF.Imp.X}{\coqdocdefinition{X}} \coqref{SECF.Imp.:com::x '+' x}{\coqdocnotation{+}} 1 \coqdoceol
\coqdocindent{2.50em}
\coqref{SECF.Imp.:com:com scope:'if' x 'then' x 'else' x 'end'}{\coqdocnotation{end}} \coqref{SECF.Imp.::com scope:'<x7B' x 'x7D>'}{\coqdocnotation{\}>}}.\coqdoceol
\coqdocemptyline
\coqdocnoindent
\coqdockw{Definition} \coqdef{SECF.Noninterference.secure com2'}{secure\_com2'}{\coqdocdefinition{secure\_com2'}} : \coqref{SECF.Imp.com}{\coqdocinductive{com}} :=\coqdoceol
\coqdocindent{1.00em}
\coqref{SECF.Imp.::com scope:'<x7B' x 'x7D>'}{\coqdocnotation{<\{}} \coqref{SECF.Imp.X}{\coqdocdefinition{X}} \coqref{SECF.Imp.:com:com scope:x ':=' x}{\coqdocnotation{:=}} 42  \coqref{SECF.Imp.::com scope:'<x7B' x 'x7D>'}{\coqdocnotation{\}>}}.\coqdoceol
\coqdocemptyline
\coqdocnoindent
\coqdockw{Lemma} \coqdef{SECF.Noninterference.sme insecure com2'}{sme\_insecure\_com2'}{\coqdoclemma{sme\_insecure\_com2'}} : \coqref{SECF.Noninterference.sme cmd}{\coqdocdefinition{sme\_cmd}} \coqref{SECF.Noninterference.insecure com2'}{\coqdocdefinition{insecure\_com2'}} \coqref{SECF.Noninterference.xpub}{\coqdocdefinition{xpub}} \coqexternalref{::type scope:x '=' x}{http://rocq-prover.org/doc/V9.1.1/stdlib//Corelib.Init.Logic}{\coqdocnotation{=}} \coqref{SECF.Noninterference.cinterp}{\coqdocdefinition{cinterp}} \coqref{SECF.Noninterference.secure com2'}{\coqdocdefinition{secure\_com2'}}.\coqdoceol
\coqdocnoindent
\coqdockw{Proof}.\coqdoceol
\coqdocindent{1.00em}
\coqdoctac{eapply} \coqdoclemma{functional\_extensionality}. \coqdoctac{intros} \coqdocvar{st}.\coqdoceol
\coqdocindent{1.00em}
\coqdoctac{unfold} \coqref{SECF.Noninterference.sme cmd}{\coqdocdefinition{sme\_cmd}}. \coqdoctac{unfold} \coqref{SECF.Noninterference.sme state}{\coqdocdefinition{sme\_state}}. \coqdoctac{simpl}.\coqdoceol
\coqdocindent{1.00em}
\coqdoctac{eapply} \coqdoclemma{functional\_extensionality}. \coqdoctac{intros} \coqdocvar{x}.\coqdoceol
\coqdocindent{1.00em}
\coqdoctac{unfold} \coqref{SECF.Noninterference.merge states}{\coqdocdefinition{merge\_states}}. \coqdoctac{simpl}.\coqdoceol
\coqdocemptyline
\coqdocindent{1.00em}
\coqdoctac{destruct} (\coqref{SECF.Noninterference.xpub}{\coqdocdefinition{xpub}} \coqdocvar{x}) \coqdocvar{eqn}:\coqdocvar{HXP}.\coqdoceol
\coqdocindent{1.00em}
\{ \coqdoctac{eapply} \coqref{SECF.Noninterference.xpub true}{\coqdoclemma{xpub\_true}} \coqdoctac{in} \coqdocvar{HXP}. \coqdoctac{subst}.\coqdoceol
\coqdocindent{2.00em}
\coqdoctac{rewrite} \coqref{SECF.Maps.t update eq}{\coqdocaxiom{t\_update\_eq}}; \coqdoctac{simpl}; \coqdoctac{auto}. \}\coqdoceol
\coqdocnoindent
\coqdoceol
\coqdocindent{1.00em}
\coqdoctac{destruct} (\coqdocdefinition{eqb} \coqdocvar{x} \coqref{SECF.Imp.Y}{\coqdocdefinition{Y}}) \coqdocvar{eqn}:\coqdocvar{HY}.\coqdoceol
\coqdocindent{1.00em}
\{ \coqdoctac{rewrite} \coqdoclemma{eqb\_eq} \coqdoctac{in} \coqdocvar{HY}. \coqdoctac{subst}.\coqdoceol
\coqdocindent{2.00em}
\coqdoctac{destruct} (\coqdocvar{st} \coqref{SECF.Imp.Y}{\coqdocdefinition{Y}} \coqexternalref{::nat scope:x '=?' x}{http://rocq-prover.org/doc/V9.1.1/stdlib//Corelib.Init.Nat}{\coqdocnotation{=?}} 0);\coqdoceol
\coqdocindent{3.00em}
\coqdoctac{repeat} \coqdoctac{rewrite} \coqref{SECF.Maps.t update neq}{\coqdocaxiom{t\_update\_neq}};\coqdoceol
\coqdocindent{3.00em}
\coqdoctac{try} (\coqdoctac{intros} \coqdocvar{Hcontra}; \coqdoctac{discriminate}); \coqdoctac{reflexivity}. \}\coqdoceol
\coqdocnoindent
\coqdoceol
\coqdocindent{1.00em}
\coqdoctac{assert} (\coqdocvar{HX}: \coqdocvar{x} \coqexternalref{::type scope:x '<>' x}{http://rocq-prover.org/doc/V9.1.1/stdlib//Corelib.Init.Logic}{\coqdocnotation{\ensuremath{\not=}}} \coqref{SECF.Imp.X}{\coqdocdefinition{X}}).\coqdoceol
\coqdocindent{1.00em}
\{ \coqdoctac{intros} \coqdocvar{Hx}. \coqdoctac{subst}. \coqdoctac{rewrite} \coqref{SECF.Noninterference.xpubX}{\coqdoclemma{xpubX}} \coqdoctac{in} \coqdocvar{HXP}. \coqdoctac{discriminate}. \}\coqdoceol
\coqdocnoindent
\coqdoceol
\coqdocindent{1.00em}
\coqdoctac{rewrite} \coqdoclemma{eqb\_neq} \coqdoctac{in} \coqdocvar{HY}.\coqdoceol
\coqdocindent{1.00em}
\coqdoctac{destruct} (\coqdocvar{st} \coqref{SECF.Imp.Y}{\coqdocdefinition{Y}} \coqexternalref{::nat scope:x '=?' x}{http://rocq-prover.org/doc/V9.1.1/stdlib//Corelib.Init.Nat}{\coqdocnotation{=?}} 0).\coqdoceol
\coqdocindent{1.00em}
- \coqdoctac{rewrite} \coqref{SECF.Maps.t update neq}{\coqdocaxiom{t\_update\_neq}}; \coqdoctac{auto}.\coqdoceol
\coqdocindent{1.00em}
- \coqdoctac{repeat} \coqdoctac{rewrite} \coqref{SECF.Maps.t update neq}{\coqdocaxiom{t\_update\_neq}}; \coqdoctac{auto}.\coqdoceol
\coqdocnoindent
\coqdockw{Qed}.\coqdoceol
\coqdocemptyline
\end{coqdoccode}
For simplicity, above we looked at a modified \coqref{SECF.Noninterference.insecure com2'}{\coqdocdefinition{insecure\_com2'}}.
    What about the effect of \coqref{SECF.Noninterference.sme cmd}{\coqdocdefinition{sme\_cmd}} on the \textit{original} \coqref{SECF.StaticIFC.insecure com2}{\coqdocdefinition{insecure\_com2}}? 
\begin{coqdoccode}
\coqdocnoindent
\coqdockw{Print} \coqdocvar{insecure\_com2}.\coqdoceol
\coqdocemptyline
\end{coqdoccode}
This is more challenging, but it turns out there is a general and
    systematic way to characterize the effect of \coqref{SECF.Noninterference.sme cmd}{\coqdocdefinition{sme\_cmd}} as a single
    program. This program is called a \textit{self-composition} and it
    captures two executions of the original program (in this case the
    two executions performed by \coqref{SECF.Noninterference.sme cmd}{\coqdocdefinition{sme\_cmd}}): 
\begin{coqdoccode}
\coqdocemptyline
\coqdocnoindent
\coqdockw{Definition} \coqdef{SECF.Noninterference.pX}{pX}{\coqdocdefinition{pX}} := "pX"\%\coqdocvar{string}.\coqdoceol
\coqdocnoindent
\coqdockw{Definition} \coqdef{SECF.Noninterference.pY}{pY}{\coqdocdefinition{pY}} := "pY"\%\coqdocvar{string}.\coqdoceol
\coqdocnoindent
\coqdockw{Definition} \coqdef{SECF.Noninterference.secure com2}{secure\_com2}{\coqdocdefinition{secure\_com2}} :=\coqdoceol
\coqdocindent{1.00em}
\coqref{SECF.Imp.::com scope:'<x7B' x 'x7D>'}{\coqdocnotation{<\{}} \coqdoceol
\coqdocindent{2.50em}
\coqref{SECF.Noninterference.pX}{\coqdocdefinition{pX}} \coqref{SECF.Imp.:com:com scope:x ':=' x}{\coqdocnotation{:=}} \coqref{SECF.Imp.X}{\coqdocdefinition{X}}\coqref{SECF.Imp.:com:com scope:x ';' x}{\coqdocnotation{;}}\coqdoceol
\coqdocindent{2.50em}
\coqdoceol
\coqdocindent{2.50em}
\coqref{SECF.Imp.:com:com scope:'if' x 'then' x 'else' x 'end'}{\coqdocnotation{if}} \coqref{SECF.Imp.Y}{\coqdocdefinition{Y}} \coqref{SECF.Imp.:com::x '=' x}{\coqdocnotation{=}} 0 \coqref{SECF.Imp.:com:com scope:'if' x 'then' x 'else' x 'end'}{\coqdocnotation{then}} \coqref{SECF.Imp.Y}{\coqdocdefinition{Y}} \coqref{SECF.Imp.:com:com scope:x ':=' x}{\coqdocnotation{:=}} 42\coqdoceol
\coqdocindent{7.00em}
\coqref{SECF.Imp.:com:com scope:'if' x 'then' x 'else' x 'end'}{\coqdocnotation{else}} \coqref{SECF.Imp.X}{\coqdocdefinition{X}} \coqref{SECF.Imp.:com:com scope:x ':=' x}{\coqdocnotation{:=}} \coqref{SECF.Imp.X}{\coqdocdefinition{X}}\coqref{SECF.Imp.:com::x '+' x}{\coqdocnotation{+}}1 \coqref{SECF.Imp.:com:com scope:'if' x 'then' x 'else' x 'end'}{\coqdocnotation{end}}\coqref{SECF.Imp.:com:com scope:x ';' x}{\coqdocnotation{;}} \coqdoceol
\coqdocindent{2.50em}
\coqdoceol
\coqdocindent{2.50em}
\coqref{SECF.Noninterference.pY}{\coqdocdefinition{pY}} \coqref{SECF.Imp.:com:com scope:x ':=' x}{\coqdocnotation{:=}} 0\coqref{SECF.Imp.:com:com scope:x ';' x}{\coqdocnotation{;}}\coqdoceol
\coqdocindent{2.50em}
\coqdoceol
\coqdocindent{2.50em}
\coqref{SECF.Imp.:com:com scope:'if' x 'then' x 'else' x 'end'}{\coqdocnotation{if}} \coqref{SECF.Noninterference.pY}{\coqdocdefinition{pY}} \coqref{SECF.Imp.:com::x '=' x}{\coqdocnotation{=}} 0 \coqref{SECF.Imp.:com:com scope:'if' x 'then' x 'else' x 'end'}{\coqdocnotation{then}} \coqref{SECF.Noninterference.pY}{\coqdocdefinition{pY}} \coqref{SECF.Imp.:com:com scope:x ':=' x}{\coqdocnotation{:=}} 42\coqdoceol
\coqdocindent{7.50em}
\coqref{SECF.Imp.:com:com scope:'if' x 'then' x 'else' x 'end'}{\coqdocnotation{else}} \coqref{SECF.Noninterference.pX}{\coqdocdefinition{pX}} \coqref{SECF.Imp.:com:com scope:x ':=' x}{\coqdocnotation{:=}} \coqref{SECF.Noninterference.pX}{\coqdocdefinition{pX}}\coqref{SECF.Imp.:com::x '+' x}{\coqdocnotation{+}}1 \coqref{SECF.Imp.:com:com scope:'if' x 'then' x 'else' x 'end'}{\coqdocnotation{end}}\coqref{SECF.Imp.:com:com scope:x ';' x}{\coqdocnotation{;}} \coqdoceol
\coqdocindent{2.50em}
\coqdoceol
\coqdocindent{2.50em}
\coqref{SECF.Imp.X}{\coqdocdefinition{X}} \coqref{SECF.Imp.:com:com scope:x ':=' x}{\coqdocnotation{:=}} \coqref{SECF.Noninterference.pX}{\coqdocdefinition{pX}}\coqdoceol
\coqdocnoindent
\coqref{SECF.Imp.::com scope:'<x7B' x 'x7D>'}{\coqdocnotation{\}>}}.\coqdoceol
\coqdocemptyline
\end{coqdoccode}
Because in our simple Imp language we have no way to restore the
    \coqref{SECF.Noninterference.pX}{\coqdocdefinition{pX}} and \coqref{SECF.Noninterference.pY}{\coqdocdefinition{pY}} variables to their original state, the equivalence
    lemma below needs to account for the fact that their values will be
    different. We do this by reusing our old friend \coqref{SECF.SpecCT.pub equiv}{\coqdocdefinition{pub\_equiv}}: 
\begin{coqdoccode}
\coqdocemptyline
\coqdocnoindent
\coqdockw{Definition} \coqdef{SECF.Noninterference.psecret}{psecret}{\coqdocdefinition{psecret}} := (\coqref{SECF.Noninterference.pX}{\coqdocdefinition{pX}} \coqref{SECF.Maps.:::x '!->' x ';' x}{\coqdocnotation{!->}} \coqref{SECF.Imp.:::'false'}{\coqdocnotation{false}}\coqref{SECF.Maps.:::x '!->' x ';' x}{\coqdocnotation{;}} \coqref{SECF.Noninterference.pY}{\coqdocdefinition{pY}} \coqref{SECF.Maps.:::x '!->' x ';' x}{\coqdocnotation{!->}} \coqref{SECF.Imp.:::'false'}{\coqdocnotation{false}}\coqref{SECF.Maps.:::x '!->' x ';' x}{\coqdocnotation{;}} \coqref{SECF.Maps.:::'  ' '!->' x}{\coqdocnotation{\_\_}} \coqref{SECF.Maps.:::'  ' '!->' x}{\coqdocnotation{!->}} \coqref{SECF.Imp.:::'true'}{\coqdocnotation{true}}).\coqdoceol
\coqdocemptyline
\coqdocnoindent
\coqdockw{Lemma} \coqdef{SECF.Noninterference.sme insecure com2}{sme\_insecure\_com2}{\coqdoclemma{sme\_insecure\_com2}} : \coqdockw{\ensuremath{\forall}} \coqdef{SECF.Noninterference.st:193}{st}{\coqdocbinder{st}},\coqdoceol
\coqdocindent{2.00em}
\coqref{SECF.Noninterference.pub equiv}{\coqdocdefinition{pub\_equiv}} \coqref{SECF.Noninterference.psecret}{\coqdocdefinition{psecret}} (\coqref{SECF.Noninterference.sme cmd}{\coqdocdefinition{sme\_cmd}} \coqref{SECF.Noninterference.insecure com2}{\coqdocdefinition{insecure\_com2}} \coqref{SECF.Noninterference.xpub}{\coqdocdefinition{xpub}} \coqref{SECF.Noninterference.st:193}{\coqdocvariable{st}})\coqdoceol
\coqdocindent{11.00em}
(\coqref{SECF.Noninterference.cinterp}{\coqdocdefinition{cinterp}} \coqref{SECF.Noninterference.secure com2}{\coqdocdefinition{secure\_com2}} \coqref{SECF.Noninterference.st:193}{\coqdocvariable{st}}).\coqdoceol
\coqdocnoindent
\coqdockw{Proof}.\coqdoceol
\coqdocindent{1.00em}
\coqdoctac{unfold} \coqref{SECF.Noninterference.pub equiv}{\coqdocdefinition{pub\_equiv}}. \coqdoctac{intros} \coqdocvar{st} \coqdocvar{x} \coqdocvar{PSEC}.\coqdoceol
\coqdocemptyline
\coqdocindent{1.00em}
\coqdoctac{unfold} \coqref{SECF.Noninterference.sme cmd}{\coqdocdefinition{sme\_cmd}}, \coqref{SECF.Noninterference.sme state}{\coqdocdefinition{sme\_state}}, \coqref{SECF.Noninterference.insecure com2}{\coqdocdefinition{insecure\_com2}}. \coqdoctac{simpl}.\coqdoceol
\coqdocindent{1.00em}
\coqdoctac{unfold} \coqref{SECF.Noninterference.merge states}{\coqdocdefinition{merge\_states}}. \coqdoctac{simpl}.\coqdoceol
\coqdocemptyline
\coqdocindent{1.00em}
\coqdoctac{destruct} (\coqref{SECF.Noninterference.xpub}{\coqdocdefinition{xpub}} \coqdocvar{x}) \coqdocvar{eqn}:\coqdocvar{HXP}.\coqdoceol
\coqdocindent{1.00em}
\{ \coqdoctac{eapply} \coqref{SECF.Noninterference.xpub true}{\coqdoclemma{xpub\_true}} \coqdoctac{in} \coqdocvar{HXP}. \coqdoctac{subst}.\coqdoceol
\coqdocindent{2.00em}
\coqdoctac{rewrite} \coqref{SECF.Maps.t update neq}{\coqdocaxiom{t\_update\_neq}}; \coqdoctac{try} \coqdoctac{discriminate}.\coqdoceol
\coqdocindent{2.00em}
\coqdoctac{rewrite} \coqref{SECF.Maps.t update eq}{\coqdocaxiom{t\_update\_eq}}.\coqdoceol
\coqdocindent{2.00em}
\coqdoctac{unfold} \coqref{SECF.Noninterference.split state}{\coqdocdefinition{split\_state}}. \coqdoctac{simpl}.\coqdoceol
\coqdocindent{2.00em}
\coqdoctac{repeat} (\coqdoctac{rewrite} \coqref{SECF.Maps.t update neq}{\coqdocaxiom{t\_update\_neq}}; \coqdoctac{try} \coqdoctac{discriminate}).\coqdoceol
\coqdocindent{2.00em}
\coqdoctac{destruct} (\coqdocvar{st} \coqref{SECF.Imp.Y}{\coqdocdefinition{Y}} \coqexternalref{::nat scope:x '=?' x}{http://rocq-prover.org/doc/V9.1.1/stdlib//Corelib.Init.Nat}{\coqdocnotation{=?}} 0) \coqdocvar{eqn}:\coqdocvar{HY0}.\coqdoceol
\coqdocindent{2.00em}
- \coqdoctac{rewrite} \coqref{SECF.Maps.t update neq}{\coqdocaxiom{t\_update\_neq}}; \coqdoctac{try} \coqdoctac{discriminate}.\coqdoceol
\coqdocindent{3.00em}
\coqdoctac{rewrite} \coqref{SECF.Maps.t update eq}{\coqdocaxiom{t\_update\_eq}}. \coqdoctac{reflexivity}.\coqdoceol
\coqdocindent{2.00em}
- \coqdoctac{rewrite} \coqref{SECF.Maps.t update neq}{\coqdocaxiom{t\_update\_neq}}; \coqdoctac{try} \coqdoctac{discriminate}.\coqdoceol
\coqdocindent{3.00em}
\coqdoctac{rewrite} \coqref{SECF.Maps.t update eq}{\coqdocaxiom{t\_update\_eq}}. \coqdoctac{reflexivity}. \}\coqdoceol
\coqdocnoindent
\coqdoceol
\coqdocindent{1.00em}
\coqdoctac{destruct} (\coqdocdefinition{eqb} \coqdocvar{x} \coqref{SECF.Imp.Y}{\coqdocdefinition{Y}}) \coqdocvar{eqn}:\coqdocvar{HY}.\coqdoceol
\coqdocindent{1.00em}
\{ \coqdoctac{rewrite} \coqdoclemma{eqb\_eq} \coqdoctac{in} \coqdocvar{HY}. \coqdoctac{subst}.\coqdoceol
\coqdocindent{2.00em}
\coqdoctac{destruct} (\coqdocvar{st} \coqref{SECF.Imp.Y}{\coqdocdefinition{Y}} \coqexternalref{::nat scope:x '=?' x}{http://rocq-prover.org/doc/V9.1.1/stdlib//Corelib.Init.Nat}{\coqdocnotation{=?}} 0) \coqdocvar{eqn}:\coqdocvar{HY0}.\coqdoceol
\coqdocindent{2.00em}
- \coqdoctac{rewrite} \coqref{SECF.Maps.t update eq}{\coqdocaxiom{t\_update\_eq}}.\coqdoceol
\coqdocindent{3.00em}
\coqdoctac{repeat} (\coqdoctac{rewrite} \coqref{SECF.Maps.t update neq}{\coqdocaxiom{t\_update\_neq}}; \coqdoctac{try} \coqdoctac{discriminate}).\coqdoceol
\coqdocindent{3.00em}
\coqdoctac{rewrite} \coqdocvar{HY0}.\coqdoceol
\coqdocindent{3.00em}
\coqdoctac{rewrite} \coqref{SECF.Maps.t update eq}{\coqdocaxiom{t\_update\_eq}}. \coqdoctac{reflexivity}.\coqdoceol
\coqdocindent{2.00em}
- \coqdoctac{repeat} (\coqdoctac{rewrite} \coqref{SECF.Maps.t update neq}{\coqdocaxiom{t\_update\_neq}}; \coqdoctac{try} \coqdoctac{discriminate}).\coqdoceol
\coqdocindent{3.00em}
\coqdoctac{rewrite} \coqdocvar{HY0}.\coqdoceol
\coqdocindent{3.00em}
\coqdoctac{repeat} (\coqdoctac{rewrite} \coqref{SECF.Maps.t update neq}{\coqdocaxiom{t\_update\_neq}}; \coqdoctac{try} \coqdoctac{discriminate}).\coqdoceol
\coqdocindent{3.00em}
\coqdoctac{reflexivity}. \}\coqdoceol
\coqdocnoindent
\coqdoceol
\coqdocindent{1.00em}
\coqdoctac{rewrite} \coqdoclemma{eqb\_neq} \coqdoctac{in} \coqdocvar{HY}.\coqdoceol
\coqdocemptyline
\coqdocindent{1.00em}
\coqdoctac{assert} (\coqdocvar{HX}: \coqdocvar{x} \coqexternalref{::type scope:x '<>' x}{http://rocq-prover.org/doc/V9.1.1/stdlib//Corelib.Init.Logic}{\coqdocnotation{\ensuremath{\not=}}} \coqref{SECF.Imp.X}{\coqdocdefinition{X}}).\coqdoceol
\coqdocindent{1.00em}
\{ \coqdoctac{intros} \coqdocvar{Hx}. \coqdoctac{subst}. \coqdoctac{rewrite} \coqref{SECF.Noninterference.xpubX}{\coqdoclemma{xpubX}} \coqdoctac{in} \coqdocvar{HXP}. \coqdoctac{discriminate}. \}\coqdoceol
\coqdocnoindent
\coqdoceol
\coqdocindent{1.00em}
\coqdoctac{assert} (\coqdocvar{HpXY}: \coqdocvar{x} \coqexternalref{::type scope:x '<>' x}{http://rocq-prover.org/doc/V9.1.1/stdlib//Corelib.Init.Logic}{\coqdocnotation{\ensuremath{\not=}}} \coqref{SECF.Noninterference.pX}{\coqdocdefinition{pX}} \coqexternalref{::type scope:x '/x5C' x}{http://rocq-prover.org/doc/V9.1.1/stdlib//Corelib.Init.Logic}{\coqdocnotation{\ensuremath{\land}}} \coqdocvar{x} \coqexternalref{::type scope:x '<>' x}{http://rocq-prover.org/doc/V9.1.1/stdlib//Corelib.Init.Logic}{\coqdocnotation{\ensuremath{\not=}}} \coqref{SECF.Noninterference.pY}{\coqdocdefinition{pY}}).\coqdoceol
\coqdocindent{1.00em}
\{ \coqdoctac{clear} - \coqdocvar{PSEC}.\coqdoceol
\coqdocindent{2.00em}
\coqdoctac{unfold} \coqref{SECF.Noninterference.psecret}{\coqdocdefinition{psecret}} \coqdoctac{in} \coqdocvar{PSEC}.\coqdoceol
\coqdocindent{2.00em}
\coqdoctac{destruct} (\coqdocdefinition{eqb} \coqdocvar{x} \coqref{SECF.Noninterference.pX}{\coqdocdefinition{pX}}) \coqdocvar{eqn}:\coqdocvar{HpX}.\coqdoceol
\coqdocindent{2.00em}
- \coqdoctac{rewrite} \coqdoclemma{eqb\_eq} \coqdoctac{in} \coqdocvar{HpX}. \coqdoctac{subst}.\coqdoceol
\coqdocindent{3.00em}
\coqdoctac{rewrite} \coqref{SECF.Maps.t update eq}{\coqdocaxiom{t\_update\_eq}} \coqdoctac{in} \coqdocvar{PSEC}. \coqdoctac{discriminate}.\coqdoceol
\coqdocindent{2.00em}
- \coqdoctac{destruct} (\coqdocdefinition{eqb} \coqdocvar{x} \coqref{SECF.Noninterference.pY}{\coqdocdefinition{pY}}) \coqdocvar{eqn}:\coqdocvar{HpY}.\coqdoceol
\coqdocindent{3.00em}
+ \coqdoctac{rewrite} \coqdoclemma{eqb\_eq} \coqdoctac{in} \coqdocvar{HpY}. \coqdoctac{subst}.\coqdoceol
\coqdocindent{4.00em}
\coqdoctac{rewrite} \coqref{SECF.Maps.t update neq}{\coqdocaxiom{t\_update\_neq}} \coqdoctac{in} \coqdocvar{PSEC}; \coqdoctac{discriminate}.\coqdoceol
\coqdocindent{3.00em}
+ \coqdoctac{rewrite} \coqdoclemma{eqb\_neq} \coqdoctac{in} \coqdocvar{HpX}. \coqdoctac{subst}.\coqdoceol
\coqdocindent{4.00em}
\coqdoctac{rewrite} \coqdoclemma{eqb\_neq} \coqdoctac{in} \coqdocvar{HpY}. \coqdoctac{subst}. \coqdoctac{auto}. \}\coqdoceol
\coqdocnoindent
\coqdoceol
\coqdocindent{1.00em}
\coqdoctac{destruct} \coqdocvar{HpXY} \coqdockw{as} [\coqdocvar{HpX} \coqdocvar{HpY}].\coqdoceol
\coqdocemptyline
\coqdocindent{1.00em}
\coqdoctac{repeat} (\coqdoctac{rewrite} \coqref{SECF.Maps.t update neq}{\coqdocaxiom{t\_update\_neq}}; \coqdoctac{auto}); \coqdoctac{try} \coqdoctac{discriminate}.\coqdoceol
\coqdocindent{1.00em}
\coqdoctac{destruct} (\coqdocvar{st} \coqref{SECF.Imp.Y}{\coqdocdefinition{Y}} \coqexternalref{::nat scope:x '=?' x}{http://rocq-prover.org/doc/V9.1.1/stdlib//Corelib.Init.Nat}{\coqdocnotation{=?}} 0) \coqdocvar{eqn}:\coqdocvar{HY0}.\coqdoceol
\coqdocindent{1.00em}
- \coqdoctac{repeat} (\coqdoctac{rewrite} \coqref{SECF.Maps.t update neq}{\coqdocaxiom{t\_update\_neq}}; \coqdoctac{auto}).\coqdoceol
\coqdocindent{1.00em}
- \coqdoctac{repeat} (\coqdoctac{rewrite} \coqref{SECF.Maps.t update neq}{\coqdocaxiom{t\_update\_neq}}; \coqdoctac{auto}).\coqdoceol
\coqdocnoindent
\coqdockw{Qed}.\coqdoceol
\coqdocemptyline
\end{coqdoccode}
By optimizing the self-composition program above quite a bit we
    can finally figure out what \coqref{SECF.Noninterference.sme cmd}{\coqdocdefinition{sme\_cmd}} does for \coqref{SECF.StaticIFC.insecure com2}{\coqdocdefinition{insecure\_com2}}: 
\begin{coqdoccode}
\coqdocemptyline
\coqdocnoindent
\coqdockw{Definition} \coqdef{SECF.Noninterference.secure com2 simple}{secure\_com2\_simple}{\coqdocdefinition{secure\_com2\_simple}} :=\coqdoceol
\coqdocindent{1.00em}
\coqref{SECF.Imp.::com scope:'<x7B' x 'x7D>'}{\coqdocnotation{<\{}} \coqref{SECF.Imp.:com:com scope:'if' x 'then' x 'else' x 'end'}{\coqdocnotation{if}} \coqref{SECF.Imp.Y}{\coqdocdefinition{Y}} \coqref{SECF.Imp.:com::x '=' x}{\coqdocnotation{=}} 0 \coqref{SECF.Imp.:com:com scope:'if' x 'then' x 'else' x 'end'}{\coqdocnotation{then}}\coqdoceol
\coqdocindent{3.50em}
\coqref{SECF.Imp.Y}{\coqdocdefinition{Y}} \coqref{SECF.Imp.:com:com scope:x ':=' x}{\coqdocnotation{:=}} 42\coqdoceol
\coqdocindent{2.50em}
\coqref{SECF.Imp.:com:com scope:'if' x 'then' x 'else' x 'end'}{\coqdocnotation{else}}\coqdoceol
\coqdocindent{3.50em}
\coqref{SECF.Imp.:com:com scope:'skip'}{\coqdocnotation{skip}} \coqdoceol
\coqdocindent{2.50em}
\coqref{SECF.Imp.:com:com scope:'if' x 'then' x 'else' x 'end'}{\coqdocnotation{end}}\coqdoceol
\coqdocnoindent
\coqref{SECF.Imp.::com scope:'<x7B' x 'x7D>'}{\coqdocnotation{\}>}}.\coqdoceol
\coqdocemptyline
\coqdocnoindent
\coqdockw{Lemma} \coqdef{SECF.Noninterference.sme insecure com2 simple}{sme\_insecure\_com2\_simple}{\coqdoclemma{sme\_insecure\_com2\_simple}} :\coqdoceol
\coqdocindent{1.00em}
\coqref{SECF.Noninterference.sme cmd}{\coqdocdefinition{sme\_cmd}} \coqref{SECF.Noninterference.insecure com2}{\coqdocdefinition{insecure\_com2}} \coqref{SECF.Noninterference.xpub}{\coqdocdefinition{xpub}} \coqexternalref{::type scope:x '=' x}{http://rocq-prover.org/doc/V9.1.1/stdlib//Corelib.Init.Logic}{\coqdocnotation{=}} \coqref{SECF.Noninterference.cinterp}{\coqdocdefinition{cinterp}} \coqref{SECF.Noninterference.secure com2 simple}{\coqdocdefinition{secure\_com2\_simple}}.\coqdoceol
\coqdocnoindent
\coqdockw{Proof}.\coqdoceol
\coqdocindent{1.00em}
\coqdoctac{eapply} \coqdoclemma{functional\_extensionality}. \coqdoctac{intros} \coqdocvar{st}.\coqdoceol
\coqdocindent{1.00em}
\coqdoctac{unfold} \coqref{SECF.Noninterference.sme cmd}{\coqdocdefinition{sme\_cmd}}. \coqdoctac{unfold} \coqref{SECF.Noninterference.sme state}{\coqdocdefinition{sme\_state}}. \coqdoctac{simpl}.\coqdoceol
\coqdocindent{1.00em}
\coqdoctac{eapply} \coqdoclemma{functional\_extensionality}. \coqdoctac{intros} \coqdocvar{x}.\coqdoceol
\coqdocindent{1.00em}
\coqdoctac{unfold} \coqref{SECF.Noninterference.merge states}{\coqdocdefinition{merge\_states}}. \coqdoctac{simpl}.\coqdoceol
\coqdocemptyline
\coqdocindent{1.00em}
\coqdoctac{destruct} (\coqref{SECF.Noninterference.xpub}{\coqdocdefinition{xpub}} \coqdocvar{x}) \coqdocvar{eqn}:\coqdocvar{HXP}.\coqdoceol
\coqdocindent{1.00em}
\{ \coqdoctac{eapply} \coqref{SECF.Noninterference.xpub true}{\coqdoclemma{xpub\_true}} \coqdoctac{in} \coqdocvar{HXP}. \coqdoctac{subst}.\coqdoceol
\coqdocindent{2.00em}
\coqdoctac{rewrite} \coqref{SECF.Maps.t update neq}{\coqdocaxiom{t\_update\_neq}}; \coqdoctac{try} \coqdoctac{discriminate}.\coqdoceol
\coqdocindent{2.00em}
\coqdoctac{unfold} \coqref{SECF.Noninterference.split state}{\coqdocdefinition{split\_state}}. \coqdoctac{simpl}.\coqdoceol
\coqdocindent{2.00em}
\coqdoctac{destruct} (\coqdocvar{st} \coqref{SECF.Imp.Y}{\coqdocdefinition{Y}} \coqexternalref{::nat scope:x '=?' x}{http://rocq-prover.org/doc/V9.1.1/stdlib//Corelib.Init.Nat}{\coqdocnotation{=?}} 0).\coqdoceol
\coqdocindent{2.00em}
- \coqdoctac{rewrite} \coqref{SECF.Maps.t update neq}{\coqdocaxiom{t\_update\_neq}}; \coqdoctac{try} \coqdoctac{discriminate}.\coqdoceol
\coqdocindent{3.00em}
\coqdoctac{reflexivity}.\coqdoceol
\coqdocindent{2.00em}
- \coqdoctac{reflexivity}. \}\coqdoceol
\coqdocnoindent
\coqdoceol
\coqdocindent{1.00em}
\coqdoctac{destruct} (\coqdocdefinition{eqb} \coqdocvar{x} \coqref{SECF.Imp.Y}{\coqdocdefinition{Y}}) \coqdocvar{eqn}:\coqdocvar{HY}.\coqdoceol
\coqdocindent{1.00em}
\{ \coqdoctac{rewrite} \coqdoclemma{eqb\_eq} \coqdoctac{in} \coqdocvar{HY}. \coqdoctac{subst}.\coqdoceol
\coqdocindent{2.00em}
\coqdoctac{destruct} (\coqdocvar{st} \coqref{SECF.Imp.Y}{\coqdocdefinition{Y}} \coqexternalref{::nat scope:x '=?' x}{http://rocq-prover.org/doc/V9.1.1/stdlib//Corelib.Init.Nat}{\coqdocnotation{=?}} 0).\coqdoceol
\coqdocindent{2.00em}
- \coqdoctac{repeat} \coqdoctac{rewrite} \coqref{SECF.Maps.t update eq}{\coqdocaxiom{t\_update\_eq}}. \coqdoctac{reflexivity}.\coqdoceol
\coqdocindent{2.00em}
- \coqdoctac{rewrite} \coqref{SECF.Maps.t update neq}{\coqdocaxiom{t\_update\_neq}}; \coqdoctac{try} \coqdoctac{discriminate}.\coqdoceol
\coqdocindent{3.00em}
\coqdoctac{reflexivity}. \}\coqdoceol
\coqdocnoindent
\coqdoceol
\coqdocindent{1.00em}
\coqdoctac{assert} (\coqdocvar{HX}: \coqdocvar{x} \coqexternalref{::type scope:x '<>' x}{http://rocq-prover.org/doc/V9.1.1/stdlib//Corelib.Init.Logic}{\coqdocnotation{\ensuremath{\not=}}} \coqref{SECF.Imp.X}{\coqdocdefinition{X}}).\coqdoceol
\coqdocindent{1.00em}
\{ \coqdoctac{intros} \coqdocvar{Hx}. \coqdoctac{subst}. \coqdoctac{rewrite} \coqref{SECF.Noninterference.xpubX}{\coqdoclemma{xpubX}} \coqdoctac{in} \coqdocvar{HXP}. \coqdoctac{discriminate}. \}\coqdoceol
\coqdocnoindent
\coqdoceol
\coqdocindent{1.00em}
\coqdoctac{rewrite} \coqdoclemma{eqb\_neq} \coqdoctac{in} \coqdocvar{HY}.\coqdoceol
\coqdocindent{1.00em}
\coqdoctac{destruct} (\coqdocvar{st} \coqref{SECF.Imp.Y}{\coqdocdefinition{Y}} \coqexternalref{::nat scope:x '=?' x}{http://rocq-prover.org/doc/V9.1.1/stdlib//Corelib.Init.Nat}{\coqdocnotation{=?}} 0).\coqdoceol
\coqdocindent{1.00em}
- \coqdoctac{reflexivity}.\coqdoceol
\coqdocindent{1.00em}
- \coqdoctac{rewrite} \coqref{SECF.Maps.t update neq}{\coqdocaxiom{t\_update\_neq}}; \coqdoctac{auto}.\coqdoceol
\coqdocnoindent
\coqdockw{Qed}.\coqdoceol
\coqdocemptyline
\end{coqdoccode}
Self-composition and the more general concept of a \textit{product program}
    are generally useful techniques of their own (e.g., for reducing
    relational properties proved by Relational Hoare Logic to regular
    properties proved by standard Hoare Logic), but we will not
    discuss them here any further. 
\begin{coqdoccode}
\coqdocemptyline
\end{coqdoccode}
\section{Noninterference for Imp programs with loops}



 In the presence of loops, we need to define noninterference using the
    evaluation relation (\coqref{SECF.Hoare.HoareAssertAssume.ceval}{\coqdocinductive{ceval}}) of Imp: 
\begin{coqdoccode}
\coqdocemptyline
\coqdocnoindent
\coqdockw{Definition} \coqdef{SECF.Noninterference.noninterferent while}{noninterferent\_while}{\coqdocdefinition{noninterferent\_while}} \coqdef{SECF.Noninterference.pub:194}{pub}{\coqdocbinder{pub}} \coqdef{SECF.Noninterference.c:195}{c}{\coqdocbinder{c}} := \coqdockw{\ensuremath{\forall}} \coqdef{SECF.Noninterference.s1:196}{s1}{\coqdocbinder{s1}} \coqdef{SECF.Noninterference.s2:197}{s2}{\coqdocbinder{s2}} \coqdef{SECF.Noninterference.s1':198}{s1'}{\coqdocbinder{s1'}} \coqdef{SECF.Noninterference.s2':199}{s2'}{\coqdocbinder{s2'}},\coqdoceol
\coqdocindent{1.00em}
\coqref{SECF.Noninterference.pub equiv}{\coqdocdefinition{pub\_equiv}} \coqref{SECF.Noninterference.pub:194}{\coqdocvariable{pub}} \coqref{SECF.Noninterference.s1:196}{\coqdocvariable{s1}} \coqref{SECF.Noninterference.s2:197}{\coqdocvariable{s2}} \coqexternalref{::type scope:x '->' x}{http://rocq-prover.org/doc/V9.1.1/stdlib//Corelib.Init.Logic}{\coqdocnotation{\ensuremath{\rightarrow}}}\coqdoceol
\coqdocindent{1.00em}
\coqref{SECF.Noninterference.s1:196}{\coqdocvariable{s1}} \coqref{SECF.Imp.:::x '=[' x ']=>' x}{\coqdocnotation{=[}} \coqref{SECF.Noninterference.c:195}{\coqdocvariable{c}} \coqref{SECF.Imp.:::x '=[' x ']=>' x}{\coqdocnotation{]=>}} \coqref{SECF.Noninterference.s1':198}{\coqdocvariable{s1'}} \coqexternalref{::type scope:x '->' x}{http://rocq-prover.org/doc/V9.1.1/stdlib//Corelib.Init.Logic}{\coqdocnotation{\ensuremath{\rightarrow}}}\coqdoceol
\coqdocindent{1.00em}
\coqref{SECF.Noninterference.s2:197}{\coqdocvariable{s2}} \coqref{SECF.Imp.:::x '=[' x ']=>' x}{\coqdocnotation{=[}} \coqref{SECF.Noninterference.c:195}{\coqdocvariable{c}} \coqref{SECF.Imp.:::x '=[' x ']=>' x}{\coqdocnotation{]=>}} \coqref{SECF.Noninterference.s2':199}{\coqdocvariable{s2'}} \coqexternalref{::type scope:x '->' x}{http://rocq-prover.org/doc/V9.1.1/stdlib//Corelib.Init.Logic}{\coqdocnotation{\ensuremath{\rightarrow}}}\coqdoceol
\coqdocindent{1.00em}
\coqref{SECF.Noninterference.pub equiv}{\coqdocdefinition{pub\_equiv}} \coqref{SECF.Noninterference.pub:194}{\coqdocvariable{pub}} \coqref{SECF.Noninterference.s1':198}{\coqdocvariable{s1'}} \coqref{SECF.Noninterference.s2':199}{\coqdocvariable{s2'}}.\coqdoceol
\coqdocemptyline
\coqdocnoindent
\coqdockw{Ltac} \coqdocvar{invert} \coqdocvar{H} := \coqdoctac{inversion} \coqdocvar{H}; \coqdoctac{subst}; \coqdoctac{clear} \coqdocvar{H}.\coqdoceol
\coqdocemptyline
\end{coqdoccode}
We re-prove noninterference of \coqref{SECF.StaticIFC.secure com}{\coqdocdefinition{secure\_com}} for this new definition: 
\begin{coqdoccode}
\coqdocemptyline
\coqdocnoindent
\coqdockw{Lemma} \coqdef{SECF.Noninterference.noninterferent secure com a bit harder}{noninterferent\_secure\_com\_a\_bit\_harder}{\coqdoclemma{noninterferent\_secure\_com\_a\_bit\_harder}} :\coqdoceol
\coqdocindent{1.00em}
\coqref{SECF.Noninterference.noninterferent while}{\coqdocdefinition{noninterferent\_while}} \coqref{SECF.Noninterference.xpub}{\coqdocdefinition{xpub}} \coqref{SECF.Noninterference.secure com}{\coqdocdefinition{secure\_com}}.\coqdoceol
\coqdocnoindent
\coqdockw{Proof}.\coqdoceol
\coqdocindent{1.00em}
\coqdoctac{unfold} \coqref{SECF.Noninterference.noninterferent while}{\coqdocdefinition{noninterferent\_while}}, \coqref{SECF.Noninterference.secure com}{\coqdocdefinition{secure\_com}}, \coqref{SECF.Noninterference.pub equiv}{\coqdocdefinition{pub\_equiv}}.\coqdoceol
\coqdocindent{1.00em}
\coqdoctac{intros} \coqdocvar{s1} \coqdocvar{s2} \coqdocvar{s1'} \coqdocvar{s2'} \coqdocvar{H} \coqdocvar{H1} \coqdocvar{H2} \coqdocvar{x} \coqdocvar{Hx}.\coqdoceol
\coqdocindent{1.00em}
\coqdoctac{apply} \coqref{SECF.Noninterference.xpub true}{\coqdoclemma{xpub\_true}} \coqdoctac{in} \coqdocvar{Hx}. \coqdoctac{subst}.\coqdoceol
\coqdocindent{1.00em}
\coqdocvar{invert} \coqdocvar{H1}. \coqdocvar{invert} \coqdocvar{H4}. \coqdocvar{invert} \coqdocvar{H7}.\coqdoceol
\coqdocindent{1.00em}
\coqdocvar{invert} \coqdocvar{H2}. \coqdocvar{invert} \coqdocvar{H3}. \coqdocvar{invert} \coqdocvar{H6}. \coqdoctac{simpl}.\coqdoceol
\coqdocindent{1.00em}
\coqdoctac{rewrite} (\coqdocvar{H} \coqref{SECF.Imp.X}{\coqdocdefinition{X}} \coqref{SECF.Noninterference.xpubX}{\coqdoclemma{xpubX}}). \coqdoctac{reflexivity}.\coqdoceol
\coqdocnoindent
\coqdockw{Qed}.\coqdoceol
\coqdocemptyline
\end{coqdoccode}
The advantage of the new definition is that it also says something
    meaningful about programs with while loops. 

 For instance, we can prove that \coqref{SECF.Imp.fact in coq}{\coqdocdefinition{fact\_in\_coq}} from \coqdoclibrary{Imp} does
    not leak the old value of \coqref{SECF.SpecCT.Y}{\coqdocdefinition{Y}} and \coqref{SECF.SpecCT.Z}{\coqdocdefinition{Z}} to \coqref{SECF.SpecCT.X}{\coqdocdefinition{X}}: 
\begin{coqdoccode}
\coqdocemptyline
\coqdocnoindent
\coqdockw{Print} \coqdocvar{fact\_in\_coq}.\coqdoceol
\coqdocemptyline
\coqdocnoindent
\coqdockw{Lemma} \coqdef{SECF.Noninterference.noninterferent fact in coq}{noninterferent\_fact\_in\_coq}{\coqdoclemma{noninterferent\_fact\_in\_coq}} :\coqdoceol
\coqdocindent{1.00em}
\coqref{SECF.Noninterference.noninterferent while}{\coqdocdefinition{noninterferent\_while}} \coqref{SECF.Noninterference.xpub}{\coqdocdefinition{xpub}} \coqref{SECF.Imp.fact in coq}{\coqdocdefinition{fact\_in\_coq}}.\coqdoceol
\coqdocnoindent
\coqdockw{Proof}.\coqdoceol
\coqdocindent{1.00em}
\coqdoctac{unfold} \coqref{SECF.Noninterference.noninterferent while}{\coqdocdefinition{noninterferent\_while}}, \coqref{SECF.Imp.fact in coq}{\coqdocdefinition{fact\_in\_coq}}, \coqref{SECF.Noninterference.pub equiv}{\coqdocdefinition{pub\_equiv}}.\coqdoceol
\coqdocindent{1.00em}
\coqdoctac{intros} \coqdocvar{s1} \coqdocvar{s2} \coqdocvar{s1'} \coqdocvar{s2'} \coqdocvar{H} \coqdocvar{H1} \coqdocvar{H2} \coqdocvar{x} \coqdocvar{Hx}.\coqdoceol
\coqdocindent{1.00em}
\coqdoctac{apply} \coqref{SECF.Noninterference.xpub true}{\coqdoclemma{xpub\_true}} \coqdoctac{in} \coqdocvar{Hx}. \coqdoctac{subst}.\coqdoceol
\coqdocindent{1.00em}
\coqdoctac{assert} (\coqdocvar{Hs}: \coqdockw{\ensuremath{\forall}} \coqdef{SECF.Noninterference.s:200}{s}{\coqdocbinder{s}} \coqdef{SECF.Noninterference.s':201}{s'}{\coqdocbinder{s'}}, \coqref{SECF.Noninterference.s:200}{\coqdocvariable{s}} \coqref{SECF.Imp.:::x '=[' x ']=>' x}{\coqdocnotation{=[}} \coqref{SECF.Imp.Z}{\coqdocdefinition{Z}} \coqref{SECF.Imp.:com:com scope:x ':=' x}{\coqdocnotation{:=}} \coqref{SECF.Imp.X}{\coqdocdefinition{X}}\coqref{SECF.Imp.:com:com scope:x ';' x}{\coqdocnotation{;}} \coqref{SECF.Imp.Y}{\coqdocdefinition{Y}} \coqref{SECF.Imp.:com:com scope:x ':=' x}{\coqdocnotation{:=}} 1\coqref{SECF.Imp.:com:com scope:x ';' x}{\coqdocnotation{;}} \coqref{SECF.Imp.:com:com scope:'while' x 'do' x 'end'}{\coqdocnotation{while}} \coqref{SECF.Imp.Z}{\coqdocdefinition{Z}} \coqref{SECF.Imp.:com::x '<>' x}{\coqdocnotation{\ensuremath{\not=}}} 0 \coqref{SECF.Imp.:com:com scope:'while' x 'do' x 'end'}{\coqdocnotation{do}} \coqref{SECF.Imp.Y}{\coqdocdefinition{Y}} \coqref{SECF.Imp.:com:com scope:x ':=' x}{\coqdocnotation{:=}} \coqref{SECF.Imp.Y}{\coqdocdefinition{Y}} \coqref{SECF.Imp.:com::x '*' x}{\coqdocnotation{\ensuremath{\times}}} \coqref{SECF.Imp.Z}{\coqdocdefinition{Z}}\coqref{SECF.Imp.:com:com scope:x ';' x}{\coqdocnotation{;}} \coqref{SECF.Imp.Z}{\coqdocdefinition{Z}} \coqref{SECF.Imp.:com:com scope:x ':=' x}{\coqdocnotation{:=}} \coqref{SECF.Imp.Z}{\coqdocdefinition{Z}} \coqref{SECF.Imp.:com::x '-' x}{\coqdocnotation{-}} 1 \coqref{SECF.Imp.:com:com scope:'while' x 'do' x 'end'}{\coqdocnotation{end}} \coqref{SECF.Imp.:::x '=[' x ']=>' x}{\coqdocnotation{]=>}} \coqref{SECF.Noninterference.s':201}{\coqdocvariable{s'}} \coqexternalref{::type scope:x '->' x}{http://rocq-prover.org/doc/V9.1.1/stdlib//Corelib.Init.Logic}{\coqdocnotation{\ensuremath{\rightarrow}}}\coqdoceol
\coqdocindent{11.00em}
\coqref{SECF.Noninterference.s:200}{\coqdocvariable{s}} \coqref{SECF.Imp.X}{\coqdocdefinition{X}} \coqexternalref{::type scope:x '=' x}{http://rocq-prover.org/doc/V9.1.1/stdlib//Corelib.Init.Logic}{\coqdocnotation{=}} \coqref{SECF.Noninterference.s':201}{\coqdocvariable{s'}} \coqref{SECF.Imp.X}{\coqdocdefinition{X}}).\coqdoceol
\coqdocindent{1.00em}
\{ \coqdoctac{intros}. \coqdoctac{clear} -\coqdocvar{H0}. \coqdocvar{invert} \coqdocvar{H0}. \coqdocvar{invert} \coqdocvar{H5}.\coqdoceol
\coqdocindent{2.00em}
\coqdocvar{invert} \coqdocvar{H2}. \coqdocvar{invert} \coqdocvar{H1}. \coqdoctac{simpl} \coqdoctac{in} \coqdocvar{H6}.\coqdoceol
\coqdocindent{2.00em}
\coqdocvar{remember} (\coqref{SECF.Imp.Y}{\coqdocdefinition{Y}} \coqref{SECF.Maps.:::x '!->' x ';' x}{\coqdocnotation{!->}} 1\coqref{SECF.Maps.:::x '!->' x ';' x}{\coqdocnotation{;}} \coqref{SECF.Imp.Z}{\coqdocdefinition{Z}} \coqref{SECF.Maps.:::x '!->' x ';' x}{\coqdocnotation{!->}} \coqdocvar{s} \coqref{SECF.Imp.X}{\coqdocdefinition{X}}\coqref{SECF.Maps.:::x '!->' x ';' x}{\coqdocnotation{;}} \coqdocvar{s}) \coqdockw{as} \coqdocvar{st}.\coqdoceol
\coqdocindent{2.00em}
\coqdoctac{replace} (\coqdocvar{s} \coqref{SECF.Imp.X}{\coqdocdefinition{X}}) \coqdockw{with} (\coqdocvar{st} \coqref{SECF.Imp.X}{\coqdocdefinition{X}}); \coqdocvar{cycle} 1.\coqdoceol
\coqdocindent{2.00em}
\{ \coqdocvar{invert} \coqdocvar{Heqst}. \coqdoctac{rewrite} \coqref{SECF.Maps.t update neq}{\coqdocaxiom{t\_update\_neq}}; \coqdoctac{eauto}. \coqdoctac{intros} \coqdocvar{contra}.\coqdoceol
\coqdocindent{3.00em}
\coqdoctac{discriminate} \coqdocvar{contra}. \}\coqdoceol
\coqdocindent{2.00em}
\coqdoctac{clear} -\coqdocvar{H6}.\coqdoceol
\coqdocindent{2.00em}
\coqdocvar{remember} \coqref{SECF.Imp.::com scope:'<x7B' x 'x7D>'}{\coqdocnotation{<\{}} \coqref{SECF.Imp.:com:com scope:'while' x 'do' x 'end'}{\coqdocnotation{while}} \coqref{SECF.Imp.Z}{\coqdocdefinition{Z}} \coqref{SECF.Imp.:com::x '<>' x}{\coqdocnotation{\ensuremath{\not=}}} 0 \coqref{SECF.Imp.:com:com scope:'while' x 'do' x 'end'}{\coqdocnotation{do}} \coqref{SECF.Imp.Y}{\coqdocdefinition{Y}} \coqref{SECF.Imp.:com:com scope:x ':=' x}{\coqdocnotation{:=}} \coqref{SECF.Imp.Y}{\coqdocdefinition{Y}} \coqref{SECF.Imp.:com::x '*' x}{\coqdocnotation{\ensuremath{\times}}} \coqref{SECF.Imp.Z}{\coqdocdefinition{Z}}\coqref{SECF.Imp.:com:com scope:x ';' x}{\coqdocnotation{;}} \coqref{SECF.Imp.Z}{\coqdocdefinition{Z}} \coqref{SECF.Imp.:com:com scope:x ':=' x}{\coqdocnotation{:=}} \coqref{SECF.Imp.Z}{\coqdocdefinition{Z}} \coqref{SECF.Imp.:com::x '-' x}{\coqdocnotation{-}} 1 \coqref{SECF.Imp.:com:com scope:'while' x 'do' x 'end'}{\coqdocnotation{end}} \coqref{SECF.Imp.::com scope:'<x7B' x 'x7D>'}{\coqdocnotation{\}>}} \coqdockw{as} \coqdocvar{loopdef}\coqdoceol
\coqdocindent{3.00em}
\coqdocvar{eqn}:\coqdocvar{Heqloopdef}.\coqdoceol
\coqdocindent{2.00em}
\coqdocvar{revert} \coqdocvar{Heqloopdef}.\coqdoceol
\coqdocindent{2.00em}
\coqdoctac{induction} \coqdocvar{H6}; \coqdoctac{intros}.\coqdoceol
\coqdocindent{2.00em}
- \coqdoctac{discriminate} \coqdocvar{Heqloopdef}.\coqdoceol
\coqdocindent{2.00em}
- \coqdoctac{discriminate} \coqdocvar{Heqloopdef}.\coqdoceol
\coqdocindent{2.00em}
- \coqdoctac{discriminate} \coqdocvar{Heqloopdef}.\coqdoceol
\coqdocindent{2.00em}
- \coqdoctac{discriminate} \coqdocvar{Heqloopdef}.\coqdoceol
\coqdocindent{2.00em}
- \coqdoctac{discriminate} \coqdocvar{Heqloopdef}.\coqdoceol
\coqdocindent{2.00em}
- \coqdoctac{reflexivity}.\coqdoceol
\coqdocindent{2.00em}
- \coqdocvar{invert} \coqdocvar{Heqloopdef}.\coqdoceol
\coqdocindent{3.00em}
\coqdoctac{rewrite} \ensuremath{\leftarrow} \coqdocvar{IHceval2}; \coqdoctac{eauto}.\coqdoceol
\coqdocindent{3.00em}
\coqdocvar{invert} \coqdocvar{H6\_}. \coqdocvar{invert} \coqdocvar{H5}. \coqdocvar{invert} \coqdocvar{H2}. \coqdoctac{simpl}.\coqdoceol
\coqdocindent{3.00em}
\coqdoctac{rewrite} \coqref{SECF.Maps.t update neq}{\coqdocaxiom{t\_update\_neq}}; \coqdocvar{cycle} 1.\coqdoceol
\coqdocindent{3.00em}
\{ \coqdoctac{intros} \coqdocvar{contra}. \coqdoctac{discriminate} \coqdocvar{contra}. \}\coqdoceol
\coqdocindent{3.00em}
\coqdoctac{rewrite} \coqref{SECF.Maps.t update neq}{\coqdocaxiom{t\_update\_neq}}; \coqdoctac{eauto}.\coqdoceol
\coqdocindent{3.00em}
\coqdoctac{intros} \coqdocvar{contra}. \coqdoctac{discriminate} \coqdocvar{contra}. \}\coqdoceol
\coqdocindent{1.00em}
\coqdoctac{eapply} \coqdocvar{Hs} \coqdoctac{in} \coqdocvar{H1}. \coqdoctac{eapply} \coqdocvar{Hs} \coqdoctac{in} \coqdocvar{H2}. \coqdoctac{rewrite} \ensuremath{\leftarrow} \coqdocvar{H1}, \ensuremath{\leftarrow} \coqdocvar{H2}.\coqdoceol
\coqdocindent{1.00em}
\coqdoctac{rewrite} (\coqdocvar{H} \coqref{SECF.Imp.X}{\coqdocdefinition{X}} \coqref{SECF.Noninterference.xpubX}{\coqdoclemma{xpubX}}). \coqdoctac{reflexivity}.\coqdoceol
\coqdocnoindent
\coqdockw{Qed}.\coqdoceol
\coqdocemptyline
\end{coqdoccode}
\section{SME for Imp programs with loops}



 To define SME in the presence of while loops we also need to use a
    relation, of a similar type to \coqref{SECF.Hoare.HoareAssertAssume.ceval}{\coqdocinductive{ceval}}: 
\begin{coqdoccode}
\coqdocemptyline
\coqdocnoindent
\coqdockw{Check} \coqref{SECF.Imp.ceval}{\coqdocinductive{ceval}} : \coqref{SECF.Imp.com}{\coqdocinductive{com}} \coqexternalref{::type scope:x '->' x}{http://rocq-prover.org/doc/V9.1.1/stdlib//Corelib.Init.Logic}{\coqdocnotation{\ensuremath{\rightarrow}}} \coqref{SECF.Imp.state}{\coqdocdefinition{state}} \coqexternalref{::type scope:x '->' x}{http://rocq-prover.org/doc/V9.1.1/stdlib//Corelib.Init.Logic}{\coqdocnotation{\ensuremath{\rightarrow}}} \coqref{SECF.Imp.state}{\coqdocdefinition{state}} \coqexternalref{::type scope:x '->' x}{http://rocq-prover.org/doc/V9.1.1/stdlib//Corelib.Init.Logic}{\coqdocnotation{\ensuremath{\rightarrow}}} \coqdockw{Prop}.\coqdoceol
\coqdocemptyline
\coqdocnoindent
\coqdockw{Definition} \coqdef{SECF.Noninterference.sme while}{sme\_while}{\coqdocdefinition{sme\_while}} (\coqdef{SECF.Noninterference.pub:204}{pub}{\coqdocbinder{pub}}:\coqref{SECF.Noninterference.pub vars}{\coqdocdefinition{pub\_vars}}) (\coqdef{SECF.Noninterference.c:205}{c}{\coqdocbinder{c}}:\coqref{SECF.Imp.com}{\coqdocinductive{com}}) (\coqdef{SECF.Noninterference.s:206}{s}{\coqdocbinder{s}} \coqdef{SECF.Noninterference.s':207}{s'}{\coqdocbinder{s'}}:\coqref{SECF.Imp.state}{\coqdocdefinition{state}}) : \coqdockw{Prop} :=\coqdoceol
\coqdocindent{1.00em}
\coqexternalref{::type scope:'exists' x '..' x ',' x}{http://rocq-prover.org/doc/V9.1.1/stdlib//Corelib.Init.Logic}{\coqdocnotation{\ensuremath{\exists}}} \coqdef{SECF.Noninterference.ps:208}{ps}{\coqdocbinder{ps}} \coqdef{SECF.Noninterference.ss:209}{ss}{\coqdocbinder{ss}}\coqexternalref{::type scope:'exists' x '..' x ',' x}{http://rocq-prover.org/doc/V9.1.1/stdlib//Corelib.Init.Logic}{\coqdocnotation{,}} \coqref{SECF.Noninterference.split state}{\coqdocdefinition{split\_state}} \coqref{SECF.Noninterference.s:206}{\coqdocvariable{s}} \coqref{SECF.Noninterference.pub:204}{\coqdocvariable{pub}} \coqref{SECF.Imp.:::'true'}{\coqdocnotation{true}} \coqref{SECF.Imp.:::x '=[' x ']=>' x}{\coqdocnotation{=[}} \coqref{SECF.Noninterference.c:205}{\coqdocvariable{c}} \coqref{SECF.Imp.:::x '=[' x ']=>' x}{\coqdocnotation{]=>}} \coqref{SECF.Noninterference.ps:208}{\coqdocvariable{ps}} \coqexternalref{::type scope:x '/x5C' x}{http://rocq-prover.org/doc/V9.1.1/stdlib//Corelib.Init.Logic}{\coqdocnotation{\ensuremath{\land}}}\coqdoceol
\coqdocindent{2.00em}
\coqref{SECF.Noninterference.s:206}{\coqdocvariable{s}} \coqref{SECF.Imp.:::x '=[' x ']=>' x}{\coqdocnotation{=[}} \coqref{SECF.Noninterference.c:205}{\coqdocvariable{c}} \coqref{SECF.Imp.:::x '=[' x ']=>' x}{\coqdocnotation{]=>}} \coqref{SECF.Noninterference.ss:209}{\coqdocvariable{ss}} \coqexternalref{::type scope:x '/x5C' x}{http://rocq-prover.org/doc/V9.1.1/stdlib//Corelib.Init.Logic}{\coqdocnotation{\ensuremath{\land}}}\coqdoceol
\coqdocindent{2.00em}
\coqref{SECF.Noninterference.merge states}{\coqdocdefinition{merge\_states}} \coqref{SECF.Noninterference.ps:208}{\coqdocvariable{ps}} \coqref{SECF.Noninterference.ss:209}{\coqdocvariable{ss}} \coqref{SECF.Noninterference.pub:204}{\coqdocvariable{pub}} \coqexternalref{::type scope:x '=' x}{http://rocq-prover.org/doc/V9.1.1/stdlib//Corelib.Init.Logic}{\coqdocnotation{=}} \coqref{SECF.Noninterference.s':207}{\coqdocvariable{s'}}.\coqdoceol
\coqdocemptyline
\end{coqdoccode}
To state that sme\_eval is secure, we need to generalize our noninterference
    definition, so that we can apply it not only to \coqref{SECF.Hoare.HoareAssertAssume.ceval}{\coqdocinductive{ceval}}, but with
    any evaluation relation, including \coqref{SECF.Noninterference.sme while}{\coqdocdefinition{sme\_while}} \coqdocvar{pub}. 
\begin{coqdoccode}
\coqdocemptyline
\coqdocnoindent
\coqdockw{Definition} \coqdef{SECF.Noninterference.noninterferent while R}{noninterferent\_while\_R}{\coqdocdefinition{noninterferent\_while\_R}} (\coqdef{SECF.Noninterference.R:210}{R}{\coqdocbinder{R}}:\coqref{SECF.Imp.com}{\coqdocinductive{com}}\coqexternalref{::type scope:x '->' x}{http://rocq-prover.org/doc/V9.1.1/stdlib//Corelib.Init.Logic}{\coqdocnotation{\ensuremath{\rightarrow}}}\coqref{SECF.Imp.state}{\coqdocdefinition{state}}\coqexternalref{::type scope:x '->' x}{http://rocq-prover.org/doc/V9.1.1/stdlib//Corelib.Init.Logic}{\coqdocnotation{\ensuremath{\rightarrow}}}\coqref{SECF.Imp.state}{\coqdocdefinition{state}}\coqexternalref{::type scope:x '->' x}{http://rocq-prover.org/doc/V9.1.1/stdlib//Corelib.Init.Logic}{\coqdocnotation{\ensuremath{\rightarrow}}}\coqdockw{Prop}) \coqdef{SECF.Noninterference.pub:211}{pub}{\coqdocbinder{pub}} \coqdef{SECF.Noninterference.c:212}{c}{\coqdocbinder{c}} :=\coqdoceol
\coqdocindent{1.00em}
\coqdockw{\ensuremath{\forall}} \coqdef{SECF.Noninterference.s1:213}{s1}{\coqdocbinder{s1}} \coqdef{SECF.Noninterference.s2:214}{s2}{\coqdocbinder{s2}} \coqdef{SECF.Noninterference.s1':215}{s1'}{\coqdocbinder{s1'}} \coqdef{SECF.Noninterference.s2':216}{s2'}{\coqdocbinder{s2'}},\coqdoceol
\coqdocindent{1.00em}
\coqref{SECF.Noninterference.pub equiv}{\coqdocdefinition{pub\_equiv}} \coqref{SECF.Noninterference.pub:211}{\coqdocvariable{pub}} \coqref{SECF.Noninterference.s1:213}{\coqdocvariable{s1}} \coqref{SECF.Noninterference.s2:214}{\coqdocvariable{s2}} \coqexternalref{::type scope:x '->' x}{http://rocq-prover.org/doc/V9.1.1/stdlib//Corelib.Init.Logic}{\coqdocnotation{\ensuremath{\rightarrow}}}\coqdoceol
\coqdocindent{1.00em}
\coqref{SECF.Noninterference.R:210}{\coqdocvariable{R}} \coqref{SECF.Noninterference.c:212}{\coqdocvariable{c}} \coqref{SECF.Noninterference.s1:213}{\coqdocvariable{s1}} \coqref{SECF.Noninterference.s1':215}{\coqdocvariable{s1'}} \coqexternalref{::type scope:x '->' x}{http://rocq-prover.org/doc/V9.1.1/stdlib//Corelib.Init.Logic}{\coqdocnotation{\ensuremath{\rightarrow}}}\coqdoceol
\coqdocindent{1.00em}
\coqref{SECF.Noninterference.R:210}{\coqdocvariable{R}} \coqref{SECF.Noninterference.c:212}{\coqdocvariable{c}} \coqref{SECF.Noninterference.s2:214}{\coqdocvariable{s2}} \coqref{SECF.Noninterference.s2':216}{\coqdocvariable{s2'}} \coqexternalref{::type scope:x '->' x}{http://rocq-prover.org/doc/V9.1.1/stdlib//Corelib.Init.Logic}{\coqdocnotation{\ensuremath{\rightarrow}}}\coqdoceol
\coqdocindent{1.00em}
\coqref{SECF.Noninterference.pub equiv}{\coqdocdefinition{pub\_equiv}} \coqref{SECF.Noninterference.pub:211}{\coqdocvariable{pub}} \coqref{SECF.Noninterference.s1':215}{\coqdocvariable{s1'}} \coqref{SECF.Noninterference.s2':216}{\coqdocvariable{s2'}}.\coqdoceol
\coqdocemptyline
\end{coqdoccode}
The proof that \coqdocvar{while\_sme} is noninterferent is as before, but now it relies
    on the determinism of \coqref{SECF.Hoare.HoareAssertAssume.ceval}{\coqdocinductive{ceval}}, which was obvious for state transformer
    functions, but is not obvious for evaluation relations. 
\begin{coqdoccode}
\coqdocemptyline
\coqdocnoindent
\coqdockw{Check} \coqref{SECF.Imp.ceval deterministic}{\coqdoclemma{ceval\_deterministic}} : \coqdockw{\ensuremath{\forall}} (\coqdef{SECF.Noninterference.c:217}{c}{\coqdocbinder{c}} : \coqref{SECF.Imp.com}{\coqdocinductive{com}}) (\coqdef{SECF.Noninterference.st:218}{st}{\coqdocbinder{st}} \coqdef{SECF.Noninterference.st1:219}{st1}{\coqdocbinder{st1}} \coqdef{SECF.Noninterference.st2:220}{st2}{\coqdocbinder{st2}} : \coqref{SECF.Imp.state}{\coqdocdefinition{state}}),\coqdoceol
\coqdocindent{2.00em}
\coqref{SECF.Noninterference.st:218}{\coqdocvariable{st}} \coqref{SECF.Imp.:::x '=[' x ']=>' x}{\coqdocnotation{=[}} \coqref{SECF.Noninterference.c:217}{\coqdocvariable{c}} \coqref{SECF.Imp.:::x '=[' x ']=>' x}{\coqdocnotation{]=>}} \coqref{SECF.Noninterference.st1:219}{\coqdocvariable{st1}} \coqexternalref{::type scope:x '->' x}{http://rocq-prover.org/doc/V9.1.1/stdlib//Corelib.Init.Logic}{\coqdocnotation{\ensuremath{\rightarrow}}}\coqdoceol
\coqdocindent{2.00em}
\coqref{SECF.Noninterference.st:218}{\coqdocvariable{st}} \coqref{SECF.Imp.:::x '=[' x ']=>' x}{\coqdocnotation{=[}} \coqref{SECF.Noninterference.c:217}{\coqdocvariable{c}} \coqref{SECF.Imp.:::x '=[' x ']=>' x}{\coqdocnotation{]=>}} \coqref{SECF.Noninterference.st2:220}{\coqdocvariable{st2}} \coqexternalref{::type scope:x '->' x}{http://rocq-prover.org/doc/V9.1.1/stdlib//Corelib.Init.Logic}{\coqdocnotation{\ensuremath{\rightarrow}}}\coqdoceol
\coqdocindent{2.00em}
\coqref{SECF.Noninterference.st1:219}{\coqdocvariable{st1}} \coqexternalref{::type scope:x '=' x}{http://rocq-prover.org/doc/V9.1.1/stdlib//Corelib.Init.Logic}{\coqdocnotation{=}} \coqref{SECF.Noninterference.st2:220}{\coqdocvariable{st2}}.\coqdoceol
\coqdocemptyline
\coqdocnoindent
\coqdockw{Theorem} \coqdef{SECF.Noninterference.noninterferent while sme}{noninterferent\_while\_sme}{\coqdoclemma{noninterferent\_while\_sme}} : \coqdockw{\ensuremath{\forall}} \coqdef{SECF.Noninterference.pub:221}{pub}{\coqdocbinder{pub}} \coqdef{SECF.Noninterference.c:222}{c}{\coqdocbinder{c}},\coqdoceol
\coqdocindent{1.00em}
\coqref{SECF.Noninterference.noninterferent while R}{\coqdocdefinition{noninterferent\_while\_R}} (\coqref{SECF.Noninterference.sme while}{\coqdocdefinition{sme\_while}} \coqref{SECF.Noninterference.pub:221}{\coqdocvariable{pub}}) \coqref{SECF.Noninterference.pub:221}{\coqdocvariable{pub}} \coqref{SECF.Noninterference.c:222}{\coqdocvariable{c}}.\coqdoceol
\coqdocnoindent
\coqdockw{Proof}.\coqdoceol
\coqdocindent{1.00em}
\coqdoctac{unfold} \coqref{SECF.Noninterference.noninterferent while R}{\coqdocdefinition{noninterferent\_while\_R}}, \coqref{SECF.Noninterference.sme while}{\coqdocdefinition{sme\_while}}.\coqdoceol
\coqdocindent{1.00em}
\coqdoctac{intros} \coqdocvar{pub} \coqdocvar{c} \coqdocvar{s1} \coqdocvar{s2} \coqdocvar{s1'} \coqdocvar{s2'} \coqdocvar{H} [\coqdocvar{ps1} [\coqdocvar{ss1} [\coqdocvar{H1p} [\coqdocvar{H1s} \coqdocvar{H1m}]]]]\coqdoceol
                               [\coqdocvar{ps2} [\coqdocvar{ss2} [\coqdocvar{H2p} [\coqdocvar{H2s} \coqdocvar{H2m}]]]].\coqdoceol
\coqdocindent{1.00em}
\coqdoctac{subst}. \coqdoctac{rewrite} \coqref{SECF.Noninterference.pub equiv split iff}{\coqdoclemma{pub\_equiv\_split\_iff}} \coqdoctac{in} \coqdocvar{H}. \coqdoctac{unfold} \coqref{SECF.Noninterference.pub equiv split}{\coqdocdefinition{pub\_equiv\_split}} \coqdoctac{in} \coqdocvar{H}.\coqdoceol
\coqdocindent{1.00em}
\coqdoctac{apply} \coqdoclemma{functional\_extensionality} \coqdoctac{in} \coqdocvar{H}. \coqdoctac{rewrite} \coqdocvar{H} \coqdoctac{in} \coqdocvar{H1p}.\coqdoceol
\coqdocindent{1.00em}
\coqdoctac{rewrite} (\coqref{SECF.Imp.ceval deterministic}{\coqdoclemma{ceval\_deterministic}} \coqdocvar{\_} \coqdocvar{\_} \coqdocvar{\_} \coqdocvar{\_} \coqdocvar{H1p} \coqdocvar{H2p}).\coqdoceol
\coqdocindent{1.00em}
\coqdoctac{apply} \coqref{SECF.Noninterference.pub equiv merge states}{\coqdoclemma{pub\_equiv\_merge\_states}}.\coqdoceol
\coqdocnoindent
\coqdockw{Qed}.\coqdoceol
\coqdocemptyline
\end{coqdoccode}
Turns out we can only prove a weak version of transparency for
    noninterferent programs, and this has to do with nontermination
    (more later). 

 But first we need a few lemmas:  
\begin{coqdoccode}
\coqdocemptyline
\coqdocnoindent
\coqdockw{Lemma} \coqdef{SECF.Noninterference.pub equiv split state}{pub\_equiv\_split\_state}{\coqdoclemma{pub\_equiv\_split\_state}} : \coqdockw{\ensuremath{\forall}} (\coqdef{SECF.Noninterference.pub:223}{pub}{\coqdocbinder{pub}}:\coqref{SECF.Noninterference.pub vars}{\coqdocdefinition{pub\_vars}}) \coqdef{SECF.Noninterference.s:224}{s}{\coqdocbinder{s}},\coqdoceol
\coqdocindent{1.00em}
\coqref{SECF.Noninterference.pub equiv}{\coqdocdefinition{pub\_equiv}} \coqref{SECF.Noninterference.pub:223}{\coqdocvariable{pub}} (\coqref{SECF.Noninterference.split state}{\coqdocdefinition{split\_state}} \coqref{SECF.Noninterference.s:224}{\coqdocvariable{s}} \coqref{SECF.Noninterference.pub:223}{\coqdocvariable{pub}} \coqref{SECF.Imp.:::'true'}{\coqdocnotation{true}}) \coqref{SECF.Noninterference.s:224}{\coqdocvariable{s}}.\coqdoceol
\coqdocnoindent
\coqdockw{Proof}.\coqdoceol
\coqdocindent{1.00em}
\coqdoctac{unfold} \coqref{SECF.Noninterference.pub equiv}{\coqdocdefinition{pub\_equiv}}, \coqref{SECF.Noninterference.split state}{\coqdocdefinition{split\_state}}.\coqdoceol
\coqdocindent{1.00em}
\coqdoctac{intros} \coqdocvar{pub} \coqdocvar{s} \coqdocvar{x} \coqdocvar{Hx}. \coqdoctac{destruct} (\coqdoclemma{Bool.eqb\_spec} (\coqdocvar{pub} \coqdocvar{x}) \coqref{SECF.Imp.:::'true'}{\coqdocnotation{true}}).\coqdoceol
\coqdocindent{1.00em}
- \coqdoctac{reflexivity}.\coqdoceol
\coqdocindent{1.00em}
- \coqdocvar{contradiction}.\coqdoceol
\coqdocnoindent
\coqdockw{Qed}.\coqdoceol
\coqdocemptyline
\coqdocnoindent
\coqdockw{Lemma} \coqdef{SECF.Noninterference.pub equiv sym}{pub\_equiv\_sym}{\coqdoclemma{pub\_equiv\_sym}} : \coqdockw{\ensuremath{\forall}} (\coqdef{SECF.Noninterference.pub:225}{pub}{\coqdocbinder{pub}}:\coqref{SECF.Noninterference.pub vars}{\coqdocdefinition{pub\_vars}}) \coqdef{SECF.Noninterference.s1:226}{s1}{\coqdocbinder{s1}} \coqdef{SECF.Noninterference.s2:227}{s2}{\coqdocbinder{s2}},\coqdoceol
\coqdocindent{1.00em}
\coqref{SECF.Noninterference.pub equiv}{\coqdocdefinition{pub\_equiv}} \coqref{SECF.Noninterference.pub:225}{\coqdocvariable{pub}} \coqref{SECF.Noninterference.s1:226}{\coqdocvariable{s1}} \coqref{SECF.Noninterference.s2:227}{\coqdocvariable{s2}} \coqexternalref{::type scope:x '->' x}{http://rocq-prover.org/doc/V9.1.1/stdlib//Corelib.Init.Logic}{\coqdocnotation{\ensuremath{\rightarrow}}}\coqdoceol
\coqdocindent{1.00em}
\coqref{SECF.Noninterference.pub equiv}{\coqdocdefinition{pub\_equiv}} \coqref{SECF.Noninterference.pub:225}{\coqdocvariable{pub}} \coqref{SECF.Noninterference.s2:227}{\coqdocvariable{s2}} \coqref{SECF.Noninterference.s1:226}{\coqdocvariable{s1}}.\coqdoceol
\coqdocnoindent
\coqdockw{Proof}.\coqdoceol
\coqdocindent{1.00em}
\coqdoctac{unfold} \coqref{SECF.Noninterference.pub equiv}{\coqdocdefinition{pub\_equiv}}. \coqdoctac{intros} \coqdocvar{pub} \coqdocvar{s1} \coqdocvar{s2} \coqdocvar{H} \coqdocvar{x} \coqdocvar{Hx}.\coqdoceol
\coqdocindent{1.00em}
\coqdoctac{rewrite} \coqdocvar{H}.\coqdoceol
\coqdocindent{1.00em}
- \coqdoctac{reflexivity}.\coqdoceol
\coqdocindent{1.00em}
- \coqdoctac{assumption}.\coqdoceol
\coqdocnoindent
\coqdockw{Qed}.\coqdoceol
\coqdocemptyline
\coqdocnoindent
\coqdockw{Lemma} \coqdef{SECF.Noninterference.merge state pub equiv}{merge\_state\_pub\_equiv}{\coqdoclemma{merge\_state\_pub\_equiv}} : \coqdockw{\ensuremath{\forall}} \coqdef{SECF.Noninterference.pub:228}{pub}{\coqdocbinder{pub}} \coqdef{SECF.Noninterference.ss:229}{ss}{\coqdocbinder{ss}} \coqdef{SECF.Noninterference.ps:230}{ps}{\coqdocbinder{ps}},\coqdoceol
\coqdocindent{1.00em}
\coqref{SECF.Noninterference.pub equiv}{\coqdocdefinition{pub\_equiv}} \coqref{SECF.Noninterference.pub:228}{\coqdocvariable{pub}} \coqref{SECF.Noninterference.ss:229}{\coqdocvariable{ss}} \coqref{SECF.Noninterference.ps:230}{\coqdocvariable{ps}} \coqexternalref{::type scope:x '->' x}{http://rocq-prover.org/doc/V9.1.1/stdlib//Corelib.Init.Logic}{\coqdocnotation{\ensuremath{\rightarrow}}}\coqdoceol
\coqdocindent{1.00em}
\coqref{SECF.Noninterference.merge states}{\coqdocdefinition{merge\_states}} \coqref{SECF.Noninterference.ps:230}{\coqdocvariable{ps}} \coqref{SECF.Noninterference.ss:229}{\coqdocvariable{ss}} \coqref{SECF.Noninterference.pub:228}{\coqdocvariable{pub}} \coqexternalref{::type scope:x '=' x}{http://rocq-prover.org/doc/V9.1.1/stdlib//Corelib.Init.Logic}{\coqdocnotation{=}} \coqref{SECF.Noninterference.ss:229}{\coqdocvariable{ss}}.\coqdoceol
\coqdocnoindent
\coqdockw{Proof}.\coqdoceol
\coqdocindent{1.00em}
\coqdoctac{unfold} \coqref{SECF.Noninterference.pub equiv}{\coqdocdefinition{pub\_equiv}}, \coqref{SECF.Noninterference.merge states}{\coqdocdefinition{merge\_states}}.\coqdoceol
\coqdocindent{1.00em}
\coqdoctac{intros} \coqdocvar{pub} \coqdocvar{ss} \coqdocvar{ps} \coqdocvar{H}. \coqdoctac{apply} \coqdoclemma{functional\_extensionality}.\coqdoceol
\coqdocindent{1.00em}
\coqdoctac{intros} \coqdocvar{x}. \coqdoctac{destruct} (\coqdocvar{pub} \coqdocvar{x}) \coqdocvar{eqn}:\coqdocvar{Heq}.\coqdoceol
\coqdocindent{1.00em}
- \coqdoctac{rewrite} \coqdocvar{H}.\coqdoceol
\coqdocindent{2.00em}
+ \coqdoctac{reflexivity}.\coqdoceol
\coqdocindent{2.00em}
+ \coqdoctac{assumption}.\coqdoceol
\coqdocindent{1.00em}
- \coqdoctac{reflexivity}.\coqdoceol
\coqdocnoindent
\coqdockw{Qed}.\coqdoceol
\coqdocemptyline
\end{coqdoccode}
More specifically, we can only prove that an \coqref{SECF.Noninterference.sme while}{\coqdocdefinition{sme\_while}} execution
    implies a \coqref{SECF.Hoare.HoareAssertAssume.ceval}{\coqdocinductive{ceval}} execution: 
\begin{coqdoccode}
\coqdocemptyline
\coqdocnoindent
\coqdockw{Theorem} \coqdef{SECF.Noninterference.somewhat transparent while sme}{somewhat\_transparent\_while\_sme}{\coqdoclemma{somewhat\_transparent\_while\_sme}} : \coqdockw{\ensuremath{\forall}} \coqdef{SECF.Noninterference.pub:231}{pub}{\coqdocbinder{pub}} \coqdef{SECF.Noninterference.c:232}{c}{\coqdocbinder{c}},\coqdoceol
\coqdocindent{1.00em}
\coqref{SECF.Noninterference.noninterferent while}{\coqdocdefinition{noninterferent\_while}} \coqref{SECF.Noninterference.pub:231}{\coqdocvariable{pub}} \coqref{SECF.Noninterference.c:232}{\coqdocvariable{c}} \coqexternalref{::type scope:x '->' x}{http://rocq-prover.org/doc/V9.1.1/stdlib//Corelib.Init.Logic}{\coqdocnotation{\ensuremath{\rightarrow}}}\coqdoceol
\coqdocindent{1.00em}
\coqexternalref{::type scope:x '->' x}{http://rocq-prover.org/doc/V9.1.1/stdlib//Corelib.Init.Logic}{\coqdocnotation{(}}\coqdockw{\ensuremath{\forall}} \coqdef{SECF.Noninterference.s:233}{s}{\coqdocbinder{s}} \coqdef{SECF.Noninterference.s':234}{s'}{\coqdocbinder{s'}}, (\coqref{SECF.Noninterference.sme while}{\coqdocdefinition{sme\_while}} \coqref{SECF.Noninterference.pub:231}{\coqdocvariable{pub}}) \coqref{SECF.Noninterference.c:232}{\coqdocvariable{c}} \coqref{SECF.Noninterference.s:233}{\coqdocvariable{s}} \coqref{SECF.Noninterference.s':234}{\coqdocvariable{s'}} \coqexternalref{::type scope:x '->' x}{http://rocq-prover.org/doc/V9.1.1/stdlib//Corelib.Init.Logic}{\coqdocnotation{\ensuremath{\rightarrow}}} \coqref{SECF.Noninterference.s:233}{\coqdocvariable{s}} \coqref{SECF.Imp.:::x '=[' x ']=>' x}{\coqdocnotation{=[}} \coqref{SECF.Noninterference.c:232}{\coqdocvariable{c}} \coqref{SECF.Imp.:::x '=[' x ']=>' x}{\coqdocnotation{]=>}} \coqref{SECF.Noninterference.s':234}{\coqdocvariable{s'}}\coqexternalref{::type scope:x '->' x}{http://rocq-prover.org/doc/V9.1.1/stdlib//Corelib.Init.Logic}{\coqdocnotation{)}}.\coqdoceol
\coqdocnoindent
\coqdockw{Proof}.\coqdoceol
\coqdocindent{1.00em}
\coqdoctac{unfold} \coqref{SECF.Noninterference.noninterferent while}{\coqdocdefinition{noninterferent\_while}}, \coqref{SECF.Noninterference.sme while}{\coqdocdefinition{sme\_while}}.\coqdoceol
\coqdocindent{1.00em}
\coqdoctac{intros} \coqdocvar{pub} \coqdocvar{c} \coqdocvar{Hni} \coqdocvar{s} \coqdocvar{s'} [\coqdocvar{ps} [\coqdocvar{ss} [\coqdocvar{Hp} [\coqdocvar{Hs} \coqdocvar{Hm}]]]]. \coqdoctac{subst} \coqdocvar{s'}.\coqdoceol
\coqdocindent{2.00em}
\coqdoctac{assert}(\coqdocvar{H}:\coqref{SECF.Noninterference.pub equiv}{\coqdocdefinition{pub\_equiv}} \coqdocvar{pub} \coqdocvar{s} (\coqref{SECF.Noninterference.split state}{\coqdocdefinition{split\_state}} \coqdocvar{s} \coqdocvar{pub} \coqref{SECF.Imp.:::'true'}{\coqdocnotation{true}})).\coqdoceol
\coqdocindent{2.00em}
\{ \coqdoctac{apply} \coqref{SECF.Noninterference.pub equiv sym}{\coqdoclemma{pub\_equiv\_sym}}. \coqdoctac{apply} \coqref{SECF.Noninterference.pub equiv split state}{\coqdoclemma{pub\_equiv\_split\_state}}. \}\coqdoceol
\coqdocindent{2.00em}
\coqdoctac{specialize} (\coqdocvar{Hni} \coqdocvar{s} (\coqref{SECF.Noninterference.split state}{\coqdocdefinition{split\_state}} \coqdocvar{s} \coqdocvar{pub} \coqref{SECF.Imp.:::'true'}{\coqdocnotation{true}}) \coqdocvar{ss} \coqdocvar{ps} \coqdocvar{H} \coqdocvar{Hs} \coqdocvar{Hp}).\coqdoceol
\coqdocindent{2.00em}
\coqdoctac{apply} \coqref{SECF.Noninterference.merge state pub equiv}{\coqdoclemma{merge\_state\_pub\_equiv}} \coqdoctac{in} \coqdocvar{Hni}. \coqdoctac{rewrite} \coqdocvar{Hni}. \coqdoctac{apply} \coqdocvar{Hs}.\coqdoceol
\coqdocnoindent
\coqdockw{Qed}.\coqdoceol
\coqdocemptyline
\end{coqdoccode}
But we cannot prove the reverse implication, since a command
    terminating when starting in state \coqdocvar{s}, does not necessarily still
    terminate when starting in state \coqref{SECF.Noninterference.split state}{\coqdocdefinition{split\_state}} \coqdocvar{s} \coqdocvar{pub} \coqexternalref{true}{http://rocq-prover.org/doc/V9.1.1/stdlib//Corelib.Init.Datatypes}{\coqdocconstructor{true}}, as
    would be needed for proving \coqref{SECF.Noninterference.sme while}{\coqdocdefinition{sme\_while}}. 

 Yet it seems we can still do most of the things as in the setting
    without while loops, including SME (just not fully transparent).
    So is there anything special about loops and nontermination?


    Yes, there is! Let's look at our noninterference definition again:


Definition noninterferent\_while pub c := forall s1 s2 s1' s2',
  pub\_equiv pub s1 s2 ->
  s1 = \coqdocvar{c} => s1' ->
  s2 = \coqdocvar{c} => s2' ->
  pub\_equiv pub s1' s2'.


    It says that for any two \textit{terminating} executions, if the initial states
    agree on their public variables, then so do the final states. This is
    traditionally called \textit{termination-insensitive} noninterference (TINI),
    since it doesn't consider nontermination to be observable to an attacker. 

 In particular, the following program is \textit{secure} wrt TINI: 
\begin{coqdoccode}
\coqdocemptyline
\coqdocnoindent
\coqdockw{Definition} \coqdef{SECF.Noninterference.termination leak}{termination\_leak}{\coqdocdefinition{termination\_leak}} : \coqref{SECF.Imp.com}{\coqdocinductive{com}} :=\coqdoceol
\coqdocindent{1.00em}
\coqref{SECF.Imp.::com scope:'<x7B' x 'x7D>'}{\coqdocnotation{<\{}} \coqref{SECF.Imp.:com:com scope:'if' x 'then' x 'else' x 'end'}{\coqdocnotation{if}} \coqref{SECF.Imp.Y}{\coqdocdefinition{Y}} \coqref{SECF.Imp.:com::x '=' x}{\coqdocnotation{=}} 0 \coqref{SECF.Imp.:com:com scope:'if' x 'then' x 'else' x 'end'}{\coqdocnotation{then}}\coqdoceol
\coqdocindent{3.50em}
\coqref{SECF.Imp.Y}{\coqdocdefinition{Y}} \coqref{SECF.Imp.:com:com scope:x ':=' x}{\coqdocnotation{:=}} 42\coqdoceol
\coqdocindent{2.50em}
\coqref{SECF.Imp.:com:com scope:'if' x 'then' x 'else' x 'end'}{\coqdocnotation{else}}\coqdoceol
\coqdocindent{3.50em}
\coqref{SECF.Imp.:com:com scope:'while' x 'do' x 'end'}{\coqdocnotation{while}} \coqref{SECF.Imp.:com::'true'}{\coqdocnotation{true}} \coqref{SECF.Imp.:com:com scope:'while' x 'do' x 'end'}{\coqdocnotation{do}} \coqref{SECF.Imp.:com:com scope:'skip'}{\coqdocnotation{skip}} \coqref{SECF.Imp.:com:com scope:'while' x 'do' x 'end'}{\coqdocnotation{end}} \coqdoceol
\coqdocindent{2.50em}
\coqref{SECF.Imp.:com:com scope:'if' x 'then' x 'else' x 'end'}{\coqdocnotation{end}} \coqref{SECF.Imp.::com scope:'<x7B' x 'x7D>'}{\coqdocnotation{\}>}}.\coqdoceol
\coqdocemptyline
\end{coqdoccode}
And we can prove it ... 
\begin{coqdoccode}
\coqdocemptyline
\coqdocnoindent
\coqdockw{Lemma} \coqdef{SECF.Noninterference.Y neq X}{Y\_neq\_X}{\coqdoclemma{Y\_neq\_X}} : (\coqref{SECF.Imp.Y}{\coqdocdefinition{Y}} \coqexternalref{::type scope:x '<>' x}{http://rocq-prover.org/doc/V9.1.1/stdlib//Corelib.Init.Logic}{\coqdocnotation{\ensuremath{\not=}}} \coqref{SECF.Imp.X}{\coqdocdefinition{X}}).\coqdoceol
\coqdocnoindent
\coqdockw{Proof}. \coqdoctac{intro} \coqdocvar{contra}. \coqdoctac{discriminate}. \coqdockw{Qed}.\coqdoceol
\coqdocemptyline
\end{coqdoccode}
We use a lemma that is a homework exercise in Imp: 
\begin{coqdoccode}
\coqdocnoindent
\coqdockw{Check} \coqref{SECF.Imp.loop never stops}{\coqdocaxiom{loop\_never\_stops}} : \coqdockw{\ensuremath{\forall}} \coqdef{SECF.Noninterference.st:235}{st}{\coqdocbinder{st}} \coqdef{SECF.Noninterference.st':236}{st'}{\coqdocbinder{st'}},\coqdoceol
\coqdocindent{1.00em}
\coqexternalref{::type scope:'x7E' x}{http://rocq-prover.org/doc/V9.1.1/stdlib//Corelib.Init.Logic}{\coqdocnotation{\~{}(}}\coqref{SECF.Noninterference.st:235}{\coqdocvariable{st}} \coqref{SECF.Imp.:::x '=[' x ']=>' x}{\coqdocnotation{=[}} \coqref{SECF.Imp.loop}{\coqdocdefinition{loop}} \coqref{SECF.Imp.:::x '=[' x ']=>' x}{\coqdocnotation{]=>}} \coqref{SECF.Noninterference.st':236}{\coqdocvariable{st'}}\coqexternalref{::type scope:'x7E' x}{http://rocq-prover.org/doc/V9.1.1/stdlib//Corelib.Init.Logic}{\coqdocnotation{)}}.\coqdoceol
\coqdocemptyline
\coqdocnoindent
\coqdockw{Definition} \coqdef{SECF.Noninterference.tini secure termination leak}{tini\_secure\_termination\_leak}{\coqdocdefinition{tini\_secure\_termination\_leak}} :\coqdoceol
\coqdocindent{1.00em}
\coqref{SECF.Noninterference.noninterferent while}{\coqdocdefinition{noninterferent\_while}} \coqref{SECF.Noninterference.xpub}{\coqdocdefinition{xpub}} \coqref{SECF.Noninterference.termination leak}{\coqdocdefinition{termination\_leak}}.\coqdoceol
\coqdocnoindent
\coqdockw{Proof}.\coqdoceol
\coqdocindent{1.00em}
\coqdoctac{unfold} \coqref{SECF.Noninterference.noninterferent while}{\coqdocdefinition{noninterferent\_while}}, \coqref{SECF.Noninterference.termination leak}{\coqdocdefinition{termination\_leak}}, \coqref{SECF.Noninterference.pub equiv}{\coqdocdefinition{pub\_equiv}}.\coqdoceol
\coqdocindent{1.00em}
\coqdoctac{intros} \coqdocvar{s1} \coqdocvar{s2} \coqdocvar{s1'} \coqdocvar{s2'} \coqdocvar{H} \coqdocvar{H1} \coqdocvar{H2} \coqdocvar{x} \coqdocvar{Hx}. \coqdoctac{apply} \coqref{SECF.Noninterference.xpub true}{\coqdoclemma{xpub\_true}} \coqdoctac{in} \coqdocvar{Hx}.\coqdoceol
\coqdocindent{1.00em}
\coqdoctac{subst}. \coqdoctac{specialize} (\coqdocvar{H} \coqref{SECF.Imp.X}{\coqdocdefinition{X}} \coqref{SECF.Noninterference.xpubX}{\coqdoclemma{xpubX}}).\coqdoceol
\coqdocindent{1.00em}
\coqdocvar{invert} \coqdocvar{H1}.\coqdoceol
\coqdocindent{1.00em}
+ \coqdocvar{invert} \coqdocvar{H8}. \coqdoctac{simpl}.\coqdoceol
\coqdocindent{2.00em}
\coqdoctac{rewrite} (\coqref{SECF.Maps.t update neq}{\coqdocaxiom{t\_update\_neq}} \coqdocvar{\_} \coqdocvar{\_} \coqdocvar{\_} \coqdocvar{\_} \coqdocvar{\_} \coqref{SECF.Noninterference.Y neq X}{\coqdoclemma{Y\_neq\_X}}).\coqdoceol
\coqdocindent{2.00em}
\coqdocvar{invert} \coqdocvar{H2}.\coqdoceol
\coqdocindent{2.00em}
\ensuremath{\times} \coqdocvar{invert} \coqdocvar{H8}. \coqdoctac{simpl}.\coqdoceol
\coqdocindent{3.00em}
\coqdoctac{rewrite} (\coqref{SECF.Maps.t update neq}{\coqdocaxiom{t\_update\_neq}} \coqdocvar{\_} \coqdocvar{\_} \coqdocvar{\_} \coqdocvar{\_} \coqdocvar{\_} \coqref{SECF.Noninterference.Y neq X}{\coqdoclemma{Y\_neq\_X}}). \coqdoctac{assumption}.\coqdoceol
\coqdocindent{2.00em}
\ensuremath{\times} \coqdoctac{apply} \coqref{SECF.Imp.loop never stops}{\coqdocaxiom{loop\_never\_stops}} \coqdoctac{in} \coqdocvar{H8}. \coqdocvar{contradiction}.\coqdoceol
\coqdocindent{1.00em}
+ \coqdoctac{apply} \coqref{SECF.Imp.loop never stops}{\coqdocaxiom{loop\_never\_stops}} \coqdoctac{in} \coqdocvar{H8}. \coqdocvar{contradiction}.\coqdoceol
\coqdocnoindent
\coqdockw{Qed}.\coqdoceol
\coqdocemptyline
\end{coqdoccode}
\section{Termination-Sensitive Noninterference}



 We can give a stronger definition of security that disallows such
    nontermination leaks. It is traditionally called
    \textit{termination-sensitive noninterference} (TSNI) and it is defined
    as follows: 
\begin{coqdoccode}
\coqdocemptyline
\coqdocnoindent
\coqdockw{Definition} \coqdef{SECF.Noninterference.tsni while R}{tsni\_while\_R}{\coqdocdefinition{tsni\_while\_R}} (\coqdef{SECF.Noninterference.R:237}{R}{\coqdocbinder{R}}:\coqref{SECF.Imp.com}{\coqdocinductive{com}}\coqexternalref{::type scope:x '->' x}{http://rocq-prover.org/doc/V9.1.1/stdlib//Corelib.Init.Logic}{\coqdocnotation{\ensuremath{\rightarrow}}}\coqref{SECF.Imp.state}{\coqdocdefinition{state}}\coqexternalref{::type scope:x '->' x}{http://rocq-prover.org/doc/V9.1.1/stdlib//Corelib.Init.Logic}{\coqdocnotation{\ensuremath{\rightarrow}}}\coqref{SECF.Imp.state}{\coqdocdefinition{state}}\coqexternalref{::type scope:x '->' x}{http://rocq-prover.org/doc/V9.1.1/stdlib//Corelib.Init.Logic}{\coqdocnotation{\ensuremath{\rightarrow}}}\coqdockw{Prop}) \coqdef{SECF.Noninterference.pub:238}{pub}{\coqdocbinder{pub}} \coqdef{SECF.Noninterference.c:239}{c}{\coqdocbinder{c}} :=\coqdoceol
\coqdocindent{1.00em}
\coqdockw{\ensuremath{\forall}} \coqdef{SECF.Noninterference.s1:240}{s1}{\coqdocbinder{s1}} \coqdef{SECF.Noninterference.s2:241}{s2}{\coqdocbinder{s2}} \coqdef{SECF.Noninterference.s1':242}{s1'}{\coqdocbinder{s1'}},\coqdoceol
\coqdocindent{1.00em}
\coqref{SECF.Noninterference.R:237}{\coqdocvariable{R}} \coqref{SECF.Noninterference.c:239}{\coqdocvariable{c}} \coqref{SECF.Noninterference.s1:240}{\coqdocvariable{s1}} \coqref{SECF.Noninterference.s1':242}{\coqdocvariable{s1'}} \coqexternalref{::type scope:x '->' x}{http://rocq-prover.org/doc/V9.1.1/stdlib//Corelib.Init.Logic}{\coqdocnotation{\ensuremath{\rightarrow}}}\coqdoceol
\coqdocindent{1.00em}
\coqref{SECF.Noninterference.pub equiv}{\coqdocdefinition{pub\_equiv}} \coqref{SECF.Noninterference.pub:238}{\coqdocvariable{pub}} \coqref{SECF.Noninterference.s1:240}{\coqdocvariable{s1}} \coqref{SECF.Noninterference.s2:241}{\coqdocvariable{s2}} \coqexternalref{::type scope:x '->' x}{http://rocq-prover.org/doc/V9.1.1/stdlib//Corelib.Init.Logic}{\coqdocnotation{\ensuremath{\rightarrow}}}\coqdoceol
\coqdocindent{1.00em}
\coqexternalref{::type scope:x '->' x}{http://rocq-prover.org/doc/V9.1.1/stdlib//Corelib.Init.Logic}{\coqdocnotation{(}}\coqexternalref{::type scope:'exists' x '..' x ',' x}{http://rocq-prover.org/doc/V9.1.1/stdlib//Corelib.Init.Logic}{\coqdocnotation{\ensuremath{\exists}}} \coqdef{SECF.Noninterference.s2':243}{s2'}{\coqdocbinder{s2'}}\coqexternalref{::type scope:'exists' x '..' x ',' x}{http://rocq-prover.org/doc/V9.1.1/stdlib//Corelib.Init.Logic}{\coqdocnotation{,}} \coqref{SECF.Noninterference.R:237}{\coqdocvariable{R}} \coqref{SECF.Noninterference.c:239}{\coqdocvariable{c}} \coqref{SECF.Noninterference.s2:241}{\coqdocvariable{s2}} \coqref{SECF.Noninterference.s2':243}{\coqdocvariable{s2'}} \coqexternalref{::type scope:x '/x5C' x}{http://rocq-prover.org/doc/V9.1.1/stdlib//Corelib.Init.Logic}{\coqdocnotation{\ensuremath{\land}}} \coqref{SECF.Noninterference.pub equiv}{\coqdocdefinition{pub\_equiv}} \coqref{SECF.Noninterference.pub:238}{\coqdocvariable{pub}} \coqref{SECF.Noninterference.s1':242}{\coqdocvariable{s1'}} \coqref{SECF.Noninterference.s2':243}{\coqdocvariable{s2'}}\coqexternalref{::type scope:x '->' x}{http://rocq-prover.org/doc/V9.1.1/stdlib//Corelib.Init.Logic}{\coqdocnotation{)}}.\coqdoceol
\coqdocemptyline
\end{coqdoccode}
We can prove that \coqref{SECF.StaticIFC.termination leak}{\coqdocdefinition{termination\_leak}} doesn't satisfy TSNI: 
\begin{coqdoccode}
\coqdocemptyline
\coqdocnoindent
\coqdockw{Definition} \coqdef{SECF.Noninterference.tsni insecure termination leak}{tsni\_insecure\_termination\_leak}{\coqdocdefinition{tsni\_insecure\_termination\_leak}} :\coqdoceol
\coqdocindent{1.00em}
\coqexternalref{::type scope:'x7E' x}{http://rocq-prover.org/doc/V9.1.1/stdlib//Corelib.Init.Logic}{\coqdocnotation{\ensuremath{\lnot}}}\coqref{SECF.Noninterference.tsni while R}{\coqdocdefinition{tsni\_while\_R}} \coqref{SECF.Imp.ceval}{\coqdocinductive{ceval}} \coqref{SECF.Noninterference.xpub}{\coqdocdefinition{xpub}} \coqref{SECF.Noninterference.termination leak}{\coqdocdefinition{termination\_leak}}.\coqdoceol
\coqdocnoindent
\coqdockw{Proof}.\coqdoceol
\coqdocindent{1.00em}
\coqdoctac{unfold} \coqref{SECF.Noninterference.tsni while R}{\coqdocdefinition{tsni\_while\_R}}, \coqref{SECF.Noninterference.termination leak}{\coqdocdefinition{termination\_leak}}.\coqdoceol
\coqdocindent{1.00em}
\coqdoctac{intros} \coqdocvar{Hc}.\coqdoceol
\coqdocindent{1.00em}
\coqdoctac{specialize} (\coqdocvar{Hc} (\coqref{SECF.Imp.X}{\coqdocdefinition{X}} \coqref{SECF.Maps.:::x '!->' x ';' x}{\coqdocnotation{!->}} 0 \coqref{SECF.Maps.:::x '!->' x ';' x}{\coqdocnotation{;}} \coqref{SECF.Imp.Y}{\coqdocdefinition{Y}} \coqref{SECF.Imp.:::x '!->' x}{\coqdocnotation{!->}} 0) (\coqref{SECF.Imp.X}{\coqdocdefinition{X}} \coqref{SECF.Maps.:::x '!->' x ';' x}{\coqdocnotation{!->}} 0 \coqref{SECF.Maps.:::x '!->' x ';' x}{\coqdocnotation{;}} \coqref{SECF.Imp.Y}{\coqdocdefinition{Y}} \coqref{SECF.Imp.:::x '!->' x}{\coqdocnotation{!->}} 1)\coqdoceol
\coqdocindent{8.50em}
(\coqref{SECF.Imp.Y}{\coqdocdefinition{Y}} \coqref{SECF.Maps.:::x '!->' x ';' x}{\coqdocnotation{!->}} 42\coqref{SECF.Maps.:::x '!->' x ';' x}{\coqdocnotation{;}} \coqref{SECF.Imp.X}{\coqdocdefinition{X}} \coqref{SECF.Maps.:::x '!->' x ';' x}{\coqdocnotation{!->}} 0 \coqref{SECF.Maps.:::x '!->' x ';' x}{\coqdocnotation{;}} \coqref{SECF.Imp.Y}{\coqdocdefinition{Y}} \coqref{SECF.Imp.:::x '!->' x}{\coqdocnotation{!->}} 0)).\coqdoceol
\coqdocindent{1.00em}
\coqdoctac{assert} (\coqdocvar{HH} : \coqref{SECF.Imp.:::x '=[' x ']=>' x}{\coqdocnotation{(}}\coqref{SECF.Imp.X}{\coqdocdefinition{X}} \coqref{SECF.Maps.:::x '!->' x ';' x}{\coqdocnotation{!->}} 0\coqref{SECF.Maps.:::x '!->' x ';' x}{\coqdocnotation{;}} \coqref{SECF.Imp.Y}{\coqdocdefinition{Y}} \coqref{SECF.Imp.:::x '!->' x}{\coqdocnotation{!->}} 0\coqref{SECF.Imp.:::x '=[' x ']=>' x}{\coqdocnotation{)}} \coqref{SECF.Imp.:::x '=[' x ']=>' x}{\coqdocnotation{=[}} \coqref{SECF.Noninterference.termination leak}{\coqdocdefinition{termination\_leak}} \coqref{SECF.Imp.:::x '=[' x ']=>' x}{\coqdocnotation{]=>}}\coqdoceol
\coqdocindent{7.50em}
\coqref{SECF.Imp.:::x '=[' x ']=>' x}{\coqdocnotation{(}}\coqref{SECF.Imp.Y}{\coqdocdefinition{Y}} \coqref{SECF.Maps.:::x '!->' x ';' x}{\coqdocnotation{!->}} 42\coqref{SECF.Maps.:::x '!->' x ';' x}{\coqdocnotation{;}} \coqref{SECF.Imp.X}{\coqdocdefinition{X}} \coqref{SECF.Maps.:::x '!->' x ';' x}{\coqdocnotation{!->}} 0\coqref{SECF.Maps.:::x '!->' x ';' x}{\coqdocnotation{;}} \coqref{SECF.Imp.Y}{\coqdocdefinition{Y}} \coqref{SECF.Imp.:::x '!->' x}{\coqdocnotation{!->}} 0\coqref{SECF.Imp.:::x '=[' x ']=>' x}{\coqdocnotation{)}}).\coqdoceol
\coqdocindent{1.00em}
\{ \coqdoctac{clear}. \coqdoctac{unfold} \coqref{SECF.Noninterference.termination leak}{\coqdocdefinition{termination\_leak}}. \coqdoctac{constructor}.\coqdoceol
\coqdocindent{2.00em}
- \coqdoctac{reflexivity}.\coqdoceol
\coqdocindent{2.00em}
- \coqdoctac{constructor}. \coqdoctac{reflexivity}. \}\coqdoceol
\coqdocindent{1.00em}
\coqdoctac{specialize} (\coqdocvar{Hc} \coqdocvar{HH}). \coqdoctac{clear} \coqdocvar{HH}.\coqdoceol
\coqdocindent{1.00em}
\coqdoctac{assert} (\coqdocvar{H}: \coqdockw{\ensuremath{\forall}} \coqdef{SECF.Noninterference.x:244}{x}{\coqdocbinder{x}}, \coqref{SECF.Noninterference.xpub}{\coqdocdefinition{xpub}} \coqref{SECF.Noninterference.x:244}{\coqdocvariable{x}} \coqexternalref{::type scope:x '=' x}{http://rocq-prover.org/doc/V9.1.1/stdlib//Corelib.Init.Logic}{\coqdocnotation{=}} \coqref{SECF.Imp.:::'true'}{\coqdocnotation{true}} \coqexternalref{::type scope:x '->' x}{http://rocq-prover.org/doc/V9.1.1/stdlib//Corelib.Init.Logic}{\coqdocnotation{\ensuremath{\rightarrow}}}\coqdoceol
\coqdocindent{11.50em}
\coqref{SECF.Maps.:::x '!->' x ';' x}{\coqdocnotation{(}}\coqref{SECF.Imp.X}{\coqdocdefinition{X}} \coqref{SECF.Maps.:::x '!->' x ';' x}{\coqdocnotation{!->}} 0\coqref{SECF.Maps.:::x '!->' x ';' x}{\coqdocnotation{;}} \coqref{SECF.Imp.Y}{\coqdocdefinition{Y}} \coqref{SECF.Imp.:::x '!->' x}{\coqdocnotation{!->}} 0\coqref{SECF.Maps.:::x '!->' x ';' x}{\coqdocnotation{)}} \coqref{SECF.Maps.:::x '!->' x ';' x}{\coqdocnotation{x}} \coqexternalref{::type scope:x '=' x}{http://rocq-prover.org/doc/V9.1.1/stdlib//Corelib.Init.Logic}{\coqdocnotation{=}} \coqref{SECF.Maps.:::x '!->' x ';' x}{\coqdocnotation{(}}\coqref{SECF.Imp.X}{\coqdocdefinition{X}} \coqref{SECF.Maps.:::x '!->' x ';' x}{\coqdocnotation{!->}} 0\coqref{SECF.Maps.:::x '!->' x ';' x}{\coqdocnotation{;}} \coqref{SECF.Imp.Y}{\coqdocdefinition{Y}} \coqref{SECF.Imp.:::x '!->' x}{\coqdocnotation{!->}} 1\coqref{SECF.Maps.:::x '!->' x ';' x}{\coqdocnotation{)}} \coqref{SECF.Maps.:::x '!->' x ';' x}{\coqdocnotation{x}}).\coqdoceol
\coqdocindent{1.00em}
\{ \coqdoctac{clear} \coqdocvar{Hc}. \coqdoctac{intros} \coqdocvar{x} \coqdocvar{H}. \coqdoctac{apply} \coqref{SECF.Noninterference.xpub true}{\coqdoclemma{xpub\_true}} \coqdoctac{in} \coqdocvar{H}. \coqdoctac{subst}. \coqdoctac{reflexivity}. \}\coqdoceol
\coqdocindent{1.00em}
\coqdoctac{specialize} (\coqdocvar{Hc} \coqdocvar{H}). \coqdoctac{clear} \coqdocvar{H}.\coqdoceol
\coqdocindent{1.00em}
\coqdoctac{destruct} \coqdocvar{Hc} \coqdockw{as} [\coqdocvar{s2'} [\coqdocvar{Hc} \coqdocvar{\_}]].\coqdoceol
\coqdocindent{1.00em}
\coqdocvar{invert} \coqdocvar{Hc}.\coqdoceol
\coqdocindent{1.00em}
- \coqdoctac{simpl} \coqdoctac{in} \coqdocvar{H4}. \coqdoctac{discriminate}.\coqdoceol
\coqdocindent{1.00em}
- \coqdoctac{apply} \coqref{SECF.Imp.loop never stops}{\coqdocaxiom{loop\_never\_stops}} \coqdoctac{in} \coqdocvar{H5}. \coqdocvar{contradiction}.\coqdoceol
\coqdocnoindent
\coqdockw{Qed}.\coqdoceol
\coqdocemptyline
\end{coqdoccode}
More generally, we can prove that TSNI is strictly stronger than TINI
    (noninterferent\_while) 
\begin{coqdoccode}
\coqdocemptyline
\coqdocnoindent
\coqdockw{Lemma} \coqdef{SECF.Noninterference.tsni noninterferent}{tsni\_noninterferent}{\coqdoclemma{tsni\_noninterferent}} : \coqdockw{\ensuremath{\forall}} \coqdef{SECF.Noninterference.pub:246}{pub}{\coqdocbinder{pub}} \coqdef{SECF.Noninterference.c:247}{c}{\coqdocbinder{c}},\coqdoceol
\coqdocindent{1.00em}
\coqref{SECF.Noninterference.tsni while R}{\coqdocdefinition{tsni\_while\_R}} \coqref{SECF.Imp.ceval}{\coqdocinductive{ceval}} \coqref{SECF.Noninterference.pub:246}{\coqdocvariable{pub}} \coqref{SECF.Noninterference.c:247}{\coqdocvariable{c}} \coqexternalref{::type scope:x '->' x}{http://rocq-prover.org/doc/V9.1.1/stdlib//Corelib.Init.Logic}{\coqdocnotation{\ensuremath{\rightarrow}}}\coqdoceol
\coqdocindent{1.00em}
\coqref{SECF.Noninterference.noninterferent while R}{\coqdocdefinition{noninterferent\_while\_R}} \coqref{SECF.Imp.ceval}{\coqdocinductive{ceval}} \coqref{SECF.Noninterference.pub:246}{\coqdocvariable{pub}} \coqref{SECF.Noninterference.c:247}{\coqdocvariable{c}}.\coqdoceol
\coqdocnoindent
\coqdockw{Proof}.\coqdoceol
\coqdocindent{1.00em}
\coqdoctac{unfold} \coqref{SECF.Noninterference.noninterferent while R}{\coqdocdefinition{noninterferent\_while\_R}}, \coqref{SECF.Noninterference.tsni while R}{\coqdocdefinition{tsni\_while\_R}}.\coqdoceol
\coqdocindent{1.00em}
\coqdoctac{intros} \coqdocvar{pub} \coqdocvar{c} \coqdocvar{Htsni} \coqdocvar{s1} \coqdocvar{s2} \coqdocvar{s1'} \coqdocvar{s2'} \coqdocvar{Hequiv} \coqdocvar{H1} \coqdocvar{H2}.\coqdoceol
\coqdocindent{1.00em}
\coqdoctac{specialize} (\coqdocvar{Htsni} \coqdocvar{s1} \coqdocvar{s2} \coqdocvar{s1'} \coqdocvar{H1} \coqdocvar{Hequiv}).\coqdoceol
\coqdocindent{1.00em}
\coqdoctac{destruct} \coqdocvar{Htsni} \coqdockw{as} [\coqdocvar{s2'{}'} [\coqdocvar{H2'} \coqdocvar{Hequiv'}]].\coqdoceol
\coqdocindent{1.00em}
\coqdoctac{rewrite} (\coqref{SECF.Imp.ceval deterministic}{\coqdoclemma{ceval\_deterministic}} \coqdocvar{\_} \coqdocvar{\_} \coqdocvar{\_} \coqdocvar{\_} \coqdocvar{H2} \coqdocvar{H2'}).\coqdoceol
\coqdocindent{1.00em}
\coqdoctac{apply} \coqdocvar{Hequiv'}.\coqdoceol
\coqdocnoindent
\coqdockw{Qed}.\coqdoceol
\coqdocemptyline
\end{coqdoccode}
The reverse direction of the implication only works for programs that
    always terminate (such as most of our simple examples above). 
\begin{coqdoccode}
\coqdocemptyline
\coqdocnoindent
\coqdockw{Lemma} \coqdef{SECF.Noninterference.terminating noninterferent tsni}{terminating\_noninterferent\_tsni}{\coqdoclemma{terminating\_noninterferent\_tsni}}: \coqdockw{\ensuremath{\forall}} \coqdef{SECF.Noninterference.pub:248}{pub}{\coqdocbinder{pub}} \coqdef{SECF.Noninterference.c:249}{c}{\coqdocbinder{c}},\coqdoceol
\coqdocindent{1.00em}
\coqexternalref{::type scope:x '->' x}{http://rocq-prover.org/doc/V9.1.1/stdlib//Corelib.Init.Logic}{\coqdocnotation{(}}\coqdockw{\ensuremath{\forall}} \coqdef{SECF.Noninterference.s:250}{s}{\coqdocbinder{s}}, \coqexternalref{::type scope:'exists' x '..' x ',' x}{http://rocq-prover.org/doc/V9.1.1/stdlib//Corelib.Init.Logic}{\coqdocnotation{\ensuremath{\exists}}} \coqdef{SECF.Noninterference.s':251}{s'}{\coqdocbinder{s'}}\coqexternalref{::type scope:'exists' x '..' x ',' x}{http://rocq-prover.org/doc/V9.1.1/stdlib//Corelib.Init.Logic}{\coqdocnotation{,}} \coqref{SECF.Noninterference.s:250}{\coqdocvariable{s}} \coqref{SECF.Imp.:::x '=[' x ']=>' x}{\coqdocnotation{=[}} \coqref{SECF.Noninterference.c:249}{\coqdocvariable{c}} \coqref{SECF.Imp.:::x '=[' x ']=>' x}{\coqdocnotation{]=>}} \coqref{SECF.Noninterference.s':251}{\coqdocvariable{s'}}\coqexternalref{::type scope:x '->' x}{http://rocq-prover.org/doc/V9.1.1/stdlib//Corelib.Init.Logic}{\coqdocnotation{)}} \coqexternalref{::type scope:x '->' x}{http://rocq-prover.org/doc/V9.1.1/stdlib//Corelib.Init.Logic}{\coqdocnotation{\ensuremath{\rightarrow}}}\coqdoceol
\coqdocindent{1.00em}
\coqref{SECF.Noninterference.noninterferent while R}{\coqdocdefinition{noninterferent\_while\_R}} \coqref{SECF.Imp.ceval}{\coqdocinductive{ceval}} \coqref{SECF.Noninterference.pub:248}{\coqdocvariable{pub}} \coqref{SECF.Noninterference.c:249}{\coqdocvariable{c}} \coqexternalref{::type scope:x '->' x}{http://rocq-prover.org/doc/V9.1.1/stdlib//Corelib.Init.Logic}{\coqdocnotation{\ensuremath{\rightarrow}}}\coqdoceol
\coqdocindent{1.00em}
\coqref{SECF.Noninterference.tsni while R}{\coqdocdefinition{tsni\_while\_R}} \coqref{SECF.Imp.ceval}{\coqdocinductive{ceval}} \coqref{SECF.Noninterference.pub:248}{\coqdocvariable{pub}} \coqref{SECF.Noninterference.c:249}{\coqdocvariable{c}}.\coqdoceol
\coqdocnoindent
\coqdockw{Proof}.\coqdoceol
\coqdocindent{1.00em}
\coqdoctac{unfold} \coqref{SECF.Noninterference.noninterferent while R}{\coqdocdefinition{noninterferent\_while\_R}}, \coqref{SECF.Noninterference.tsni while R}{\coqdocdefinition{tsni\_while\_R}}.\coqdoceol
\coqdocindent{1.00em}
\coqdoctac{intros} \coqdocvar{pub} \coqdocvar{c} \coqdocvar{Hterminating} \coqdocvar{Hni} \coqdocvar{s1} \coqdocvar{s2} \coqdocvar{s1'} \coqdocvar{H} \coqdocvar{Eq}.\coqdoceol
\coqdocindent{1.00em}
\coqdoctac{destruct} (\coqdocvar{Hterminating} \coqdocvar{s2}) \coqdockw{as} [\coqdocvar{s2'} \coqdocvar{H'}].\coqdoceol
\coqdocindent{1.00em}
\coqdoctac{\ensuremath{\exists}} \coqdocvar{s2'}; \coqdoctac{split}.\coqdoceol
\coqdocindent{1.00em}
- \coqdoctac{assumption}.\coqdoceol
\coqdocindent{1.00em}
- \coqdoctac{apply} \coqdocvar{Hni} \coqdockw{with} (\coqdocvar{s1} := \coqdocvar{s1}) (\coqdocvar{s2} := \coqdocvar{s2}).\coqdoceol
\coqdocindent{2.00em}
+ \coqdoctac{assumption}.\coqdoceol
\coqdocindent{2.00em}
+ \coqdoctac{assumption}.\coqdoceol
\coqdocindent{2.00em}
+ \coqdoctac{assumption}.\coqdoceol
\coqdocnoindent
\coqdockw{Qed}.\coqdoceol
\coqdocemptyline
\end{coqdoccode}
Now for a more interesting use of TSNI: it turns out that
    \coqref{SECF.Noninterference.sme while}{\coqdocdefinition{sme\_while}} is transparent for programs satisfying TSNI. 
\begin{coqdoccode}
\coqdocemptyline
\coqdocnoindent
\coqdockw{Theorem} \coqdef{SECF.Noninterference.tsni transparent while sme}{tsni\_transparent\_while\_sme}{\coqdoclemma{tsni\_transparent\_while\_sme}} : \coqdockw{\ensuremath{\forall}} \coqdef{SECF.Noninterference.pub:252}{pub}{\coqdocbinder{pub}} \coqdef{SECF.Noninterference.c:253}{c}{\coqdocbinder{c}},\coqdoceol
\coqdocindent{1.00em}
\coqref{SECF.Noninterference.tsni while R}{\coqdocdefinition{tsni\_while\_R}} \coqref{SECF.Imp.ceval}{\coqdocinductive{ceval}} \coqref{SECF.Noninterference.pub:252}{\coqdocvariable{pub}} \coqref{SECF.Noninterference.c:253}{\coqdocvariable{c}} \coqexternalref{::type scope:x '->' x}{http://rocq-prover.org/doc/V9.1.1/stdlib//Corelib.Init.Logic}{\coqdocnotation{\ensuremath{\rightarrow}}}\coqdoceol
\coqdocindent{1.00em}
\coqexternalref{::type scope:x '->' x}{http://rocq-prover.org/doc/V9.1.1/stdlib//Corelib.Init.Logic}{\coqdocnotation{(}}\coqdockw{\ensuremath{\forall}} \coqdef{SECF.Noninterference.s:254}{s}{\coqdocbinder{s}} \coqdef{SECF.Noninterference.s':255}{s'}{\coqdocbinder{s'}}, \coqref{SECF.Noninterference.s:254}{\coqdocvariable{s}} \coqref{SECF.Imp.:::x '=[' x ']=>' x}{\coqdocnotation{=[}} \coqref{SECF.Noninterference.c:253}{\coqdocvariable{c}} \coqref{SECF.Imp.:::x '=[' x ']=>' x}{\coqdocnotation{]=>}} \coqref{SECF.Noninterference.s':255}{\coqdocvariable{s'}} \coqexternalref{::type scope:x '<->' x}{http://rocq-prover.org/doc/V9.1.1/stdlib//Corelib.Init.Logic}{\coqdocnotation{\ensuremath{\leftrightarrow}}} (\coqref{SECF.Noninterference.sme while}{\coqdocdefinition{sme\_while}} \coqref{SECF.Noninterference.pub:252}{\coqdocvariable{pub}}) \coqref{SECF.Noninterference.c:253}{\coqdocvariable{c}} \coqref{SECF.Noninterference.s:254}{\coqdocvariable{s}} \coqref{SECF.Noninterference.s':255}{\coqdocvariable{s'}}\coqexternalref{::type scope:x '->' x}{http://rocq-prover.org/doc/V9.1.1/stdlib//Corelib.Init.Logic}{\coqdocnotation{)}}.\coqdoceol
\coqdocnoindent
\coqdockw{Proof}.\coqdoceol
\coqdocindent{1.00em}
\coqdoctac{unfold} \coqref{SECF.Noninterference.tsni while R}{\coqdocdefinition{tsni\_while\_R}}, \coqref{SECF.Noninterference.sme while}{\coqdocdefinition{sme\_while}}.\coqdoceol
\coqdocindent{1.00em}
\coqdoctac{intros} \coqdocvar{pub} \coqdocvar{c} \coqdocvar{Hni} \coqdocvar{s} \coqdocvar{s'}.\coqdoceol
\coqdocindent{1.00em}
\coqdoctac{assert}(\coqdocvar{HH}:\coqref{SECF.Noninterference.pub equiv}{\coqdocdefinition{pub\_equiv}} \coqdocvar{pub} \coqdocvar{s} (\coqref{SECF.Noninterference.split state}{\coqdocdefinition{split\_state}} \coqdocvar{s} \coqdocvar{pub} \coqref{SECF.Imp.:::'true'}{\coqdocnotation{true}})).\coqdoceol
\coqdocindent{2.00em}
\{ \coqdoctac{apply} \coqref{SECF.Noninterference.pub equiv sym}{\coqdoclemma{pub\_equiv\_sym}}. \coqdoctac{apply} \coqref{SECF.Noninterference.pub equiv split state}{\coqdoclemma{pub\_equiv\_split\_state}}. \}\coqdoceol
\coqdocindent{1.00em}
\coqdoctac{split}.\coqdoceol
\coqdocindent{1.00em}
- \coqdoctac{intros} \coqdocvar{H}. \coqdoctac{specialize} (\coqdocvar{Hni} \coqdocvar{s} (\coqref{SECF.Noninterference.split state}{\coqdocdefinition{split\_state}} \coqdocvar{s} \coqdocvar{pub} \coqref{SECF.Imp.:::'true'}{\coqdocnotation{true}}) \coqdocvar{s'} \coqdocvar{H} \coqdocvar{HH}).\coqdoceol
\coqdocindent{2.00em}
\coqdoctac{destruct} \coqdocvar{Hni} \coqdockw{as} [\coqdocvar{s'{}'} [\coqdocvar{Heval} \coqdocvar{Hequiv}]].\coqdoceol
\coqdocindent{2.00em}
\coqdoctac{\ensuremath{\exists}} \coqdocvar{s'{}'}. \coqdoctac{\ensuremath{\exists}} \coqdocvar{s'}. \coqdoctac{split}.\coqdoceol
\coqdocindent{2.00em}
+ \coqdoctac{assumption}.\coqdoceol
\coqdocindent{2.00em}
+ \coqdoctac{split}.\coqdoceol
\coqdocindent{3.00em}
\ensuremath{\times} \coqdoctac{assumption}.\coqdoceol
\coqdocindent{3.00em}
\ensuremath{\times} \coqdoctac{apply} \coqref{SECF.Noninterference.merge state pub equiv}{\coqdoclemma{merge\_state\_pub\_equiv}}. \coqdoctac{assumption}.\coqdoceol
\coqdocindent{1.00em}
- \coqdoctac{intros} [\coqdocvar{ps} [\coqdocvar{ss} [\coqdocvar{Hp} [\coqdocvar{Hs} \coqdocvar{Hm}]]]]. \coqdoctac{subst} \coqdocvar{s'}.\coqdoceol
\coqdocindent{2.00em}
\coqdoctac{specialize} (\coqdocvar{Hni} \coqdocvar{s} (\coqref{SECF.Noninterference.split state}{\coqdocdefinition{split\_state}} \coqdocvar{s} \coqdocvar{pub} \coqref{SECF.Imp.:::'true'}{\coqdocnotation{true}}) \coqdocvar{ss} \coqdocvar{Hs} \coqdocvar{HH}).\coqdoceol
\coqdocindent{2.00em}
\coqdoctac{destruct} \coqdocvar{Hni} \coqdockw{as} [\coqdocvar{s'} [\coqdocvar{Hp'} \coqdocvar{Hni}]].\coqdoceol
\coqdocindent{2.00em}
\coqdoctac{rewrite} (\coqref{SECF.Imp.ceval deterministic}{\coqdoclemma{ceval\_deterministic}} \coqdocvar{\_} \coqdocvar{\_} \coqdocvar{\_} \coqdocvar{\_} \coqdocvar{Hp} \coqdocvar{Hp'}).\coqdoceol
\coqdocindent{2.00em}
\coqdoctac{apply} \coqref{SECF.Noninterference.merge state pub equiv}{\coqdoclemma{merge\_state\_pub\_equiv}} \coqdoctac{in} \coqdocvar{Hni}. \coqdoctac{rewrite} \coqdocvar{Hni}. \coqdoctac{apply} \coqdocvar{Hs}.\coqdoceol
\coqdocnoindent
\coqdockw{Qed}.\coqdoceol
\coqdocemptyline
\end{coqdoccode}
Unfortunately \coqref{SECF.Noninterference.sme while}{\coqdocdefinition{sme\_while}} does not \textit{enforce} TSNI and this is hard
    to fix in our current setting, where programs only return a result
    in the end, a final state, so we had to merge the public and
    secret inputs into a single final state. Instead, SME is commonly
    defined in a setting with interactive IO, in which public outputs
    and secret outputs can be performed independently, during the
    execution \coqdocvar{Devriese} \coqexternalref{and}{http://rocq-prover.org/doc/V9.1.1/stdlib//Corelib.Init.Logic}{\coqdocinductive{and}} \coqdocvar{Piessens} 2010 (in Bib.v). In that setting, it
    does transparently enforce a termination-insensitive version of
    noninterference later research has called Indirect TSNI
    \coqdocvar{Ngo} \coqdocvar{et} \coqdocvar{al} 2018 (in Bib.v). 
\begin{coqdoccode}
\coqdocemptyline
\end{coqdoccode}
\subsection{Optional: Counterexample showing that SME doesn't enforce TSNI}



 We build a counterexample command that does not satisfy TSNI and
    for which the same publicly equivalent initial states \coqref{SECF.Noninterference.s1}{\coqdocdefinition{s1}} and
    \coqref{SECF.Noninterference.s2}{\coqdocdefinition{s2}} can be used to show that it still does not satisfy TSNI when
    run with \coqref{SECF.Noninterference.sme while}{\coqdocdefinition{sme\_while}}.


    In particular, we choose \coqref{SECF.Noninterference.s1}{\coqdocdefinition{s1}} below so that the command terminates
    and so that zeroing out the secret variable Y has no effect on \coqref{SECF.Noninterference.s1}{\coqdocdefinition{s1}}.
    We choose \coqref{SECF.Noninterference.s2}{\coqdocdefinition{s2}} so that the command loops, which implies that it
    will still loop on \coqref{SECF.Noninterference.s2}{\coqdocdefinition{s2}} also when executed with \coqref{SECF.Noninterference.sme while}{\coqdocdefinition{sme\_while}}. 
\begin{coqdoccode}
\coqdocemptyline
\coqdocnoindent
\coqdockw{Section} \coqdef{SECF.Noninterference.TSNICOUNTER}{TSNICOUNTER}{\coqdocsection{TSNICOUNTER}}.\coqdoceol
\coqdocemptyline
\coqdocnoindent
\coqdockw{Definition} \coqdef{SECF.Noninterference.counter}{counter}{\coqdocdefinition{counter}} : \coqref{SECF.Imp.com}{\coqdocinductive{com}} := \coqref{SECF.Imp.::com scope:'<x7B' x 'x7D>'}{\coqdocnotation{<\{}} \coqref{SECF.Imp.:com:com scope:'while' x 'do' x 'end'}{\coqdocnotation{while}} \coqref{SECF.Imp.:com::'(' x ')'}{\coqdocnotation{(}}\coqref{SECF.Imp.Y}{\coqdocdefinition{Y}} \coqref{SECF.Imp.:com::x '=' x}{\coqdocnotation{=}} 1\coqref{SECF.Imp.:com::'(' x ')'}{\coqdocnotation{)}} \coqref{SECF.Imp.:com:com scope:'while' x 'do' x 'end'}{\coqdocnotation{do}} \coqref{SECF.Imp.:com:com scope:'skip'}{\coqdocnotation{skip}} \coqref{SECF.Imp.:com:com scope:'while' x 'do' x 'end'}{\coqdocnotation{end}}\coqref{SECF.Imp.:com:com scope:x ';' x}{\coqdocnotation{;}} \coqref{SECF.Imp.X}{\coqdocdefinition{X}} \coqref{SECF.Imp.:com:com scope:x ':=' x}{\coqdocnotation{:=}} 1 \coqref{SECF.Imp.::com scope:'<x7B' x 'x7D>'}{\coqdocnotation{\}>}}.\coqdoceol
\coqdocemptyline
\coqdocnoindent
\coqdockw{Definition} \coqdef{SECF.Noninterference.s1}{s1}{\coqdocdefinition{s1}}: \coqref{SECF.Imp.state}{\coqdocdefinition{state}} := \coqref{SECF.Imp.X}{\coqdocdefinition{X}} \coqref{SECF.Maps.:::x '!->' x ';' x}{\coqdocnotation{!->}} 0\coqref{SECF.Maps.:::x '!->' x ';' x}{\coqdocnotation{;}} \coqref{SECF.Imp.Y}{\coqdocdefinition{Y}} \coqref{SECF.Maps.:::x '!->' x ';' x}{\coqdocnotation{!->}} 0\coqref{SECF.Maps.:::x '!->' x ';' x}{\coqdocnotation{;}} \coqref{SECF.Imp.empty st}{\coqdocdefinition{empty\_st}}.\coqdoceol
\coqdocnoindent
\coqdockw{Definition} \coqdef{SECF.Noninterference.s2}{s2}{\coqdocdefinition{s2}}: \coqref{SECF.Imp.state}{\coqdocdefinition{state}} := \coqref{SECF.Imp.X}{\coqdocdefinition{X}} \coqref{SECF.Maps.:::x '!->' x ';' x}{\coqdocnotation{!->}} 0\coqref{SECF.Maps.:::x '!->' x ';' x}{\coqdocnotation{;}} \coqref{SECF.Imp.Y}{\coqdocdefinition{Y}} \coqref{SECF.Maps.:::x '!->' x ';' x}{\coqdocnotation{!->}} 1\coqref{SECF.Maps.:::x '!->' x ';' x}{\coqdocnotation{;}} \coqref{SECF.Imp.empty st}{\coqdocdefinition{empty\_st}}.\coqdoceol
\coqdocnoindent
\coqdockw{Definition} \coqdef{SECF.Noninterference.s1'}{s1'}{\coqdocdefinition{s1'}}: \coqref{SECF.Imp.state}{\coqdocdefinition{state}} := \coqref{SECF.Imp.X}{\coqdocdefinition{X}} \coqref{SECF.Maps.:::x '!->' x ';' x}{\coqdocnotation{!->}} 1\coqref{SECF.Maps.:::x '!->' x ';' x}{\coqdocnotation{;}} \coqref{SECF.Noninterference.s1}{\coqdocdefinition{s1}}.\coqdoceol
\coqdocemptyline
\coqdocnoindent
\coqdockw{Lemma} \coqdef{SECF.Noninterference.counter s1 terminates s1'}{counter\_s1\_terminates\_s1'}{\coqdoclemma{counter\_s1\_terminates\_s1'}}: \coqref{SECF.Noninterference.s1}{\coqdocdefinition{s1}} \coqref{SECF.Imp.:::x '=[' x ']=>' x}{\coqdocnotation{=[}} \coqref{SECF.Noninterference.counter}{\coqdocdefinition{counter}} \coqref{SECF.Imp.:::x '=[' x ']=>' x}{\coqdocnotation{]=>}} \coqref{SECF.Noninterference.s1'}{\coqdocdefinition{s1'}}.\coqdoceol
\coqdocnoindent
\coqdockw{Proof}.\coqdoceol
\coqdocindent{1.00em}
\coqdoctac{unfold} \coqref{SECF.Noninterference.counter}{\coqdocdefinition{counter}}, \coqref{SECF.Noninterference.s1}{\coqdocdefinition{s1}}. \coqdoctac{eapply} \coqref{SECF.Imp.E Seq}{\coqdocconstructor{E\_Seq}}.\coqdoceol
\coqdocindent{1.00em}
- \coqdoctac{eapply} \coqref{SECF.Imp.E WhileFalse}{\coqdocconstructor{E\_WhileFalse}}. \coqdoctac{simpl}. \coqdoctac{reflexivity}.\coqdoceol
\coqdocindent{1.00em}
- \coqdoctac{eapply} \coqref{SECF.Imp.E Asgn}{\coqdocconstructor{E\_Asgn}}. \coqdoctac{simpl}. \coqdoctac{reflexivity}.\coqdoceol
\coqdocnoindent
\coqdockw{Qed}.\coqdoceol
\coqdocemptyline
\coqdocnoindent
\coqdockw{Lemma} \coqdef{SECF.Noninterference.counter s2 loops}{counter\_s2\_loops}{\coqdoclemma{counter\_s2\_loops}} : \coqdockw{\ensuremath{\forall}} \coqdef{SECF.Noninterference.s2':256}{s2'}{\coqdocbinder{s2'}},\coqdoceol
\coqdocindent{1.00em}
\coqexternalref{::type scope:'x7E' x}{http://rocq-prover.org/doc/V9.1.1/stdlib//Corelib.Init.Logic}{\coqdocnotation{\ensuremath{\lnot}}} \coqexternalref{::type scope:'x7E' x}{http://rocq-prover.org/doc/V9.1.1/stdlib//Corelib.Init.Logic}{\coqdocnotation{(}}\coqref{SECF.Noninterference.s2}{\coqdocdefinition{s2}} \coqref{SECF.Imp.:::x '=[' x ']=>' x}{\coqdocnotation{=[}} \coqref{SECF.Noninterference.counter}{\coqdocdefinition{counter}} \coqref{SECF.Imp.:::x '=[' x ']=>' x}{\coqdocnotation{]=>}} \coqref{SECF.Noninterference.s2':256}{\coqdocvariable{s2'}}\coqexternalref{::type scope:'x7E' x}{http://rocq-prover.org/doc/V9.1.1/stdlib//Corelib.Init.Logic}{\coqdocnotation{)}}.\coqdoceol
\coqdocnoindent
\coqdockw{Proof}.\coqdoceol
\coqdocindent{1.00em}
\coqdoctac{unfold} \coqref{SECF.Noninterference.counter}{\coqdocdefinition{counter}}. \coqdoctac{intros} \coqdocvar{s2'} \coqdocvar{Hcontra}.\coqdoceol
\coqdocemptyline
\coqdocindent{1.00em}
\coqdoctac{assert} (\coqdocvar{NSTOP}: \coqdockw{\ensuremath{\forall}} \coqdef{SECF.Noninterference.s:257}{s}{\coqdocbinder{s}} \coqdef{SECF.Noninterference.s':258}{s'}{\coqdocbinder{s'}}, \coqref{SECF.Noninterference.s:257}{\coqdocvariable{s}} \coqref{SECF.Imp.Y}{\coqdocdefinition{Y}} \coqexternalref{::type scope:x '=' x}{http://rocq-prover.org/doc/V9.1.1/stdlib//Corelib.Init.Logic}{\coqdocnotation{=}} 1 \coqexternalref{::type scope:x '->' x}{http://rocq-prover.org/doc/V9.1.1/stdlib//Corelib.Init.Logic}{\coqdocnotation{\ensuremath{\rightarrow}}}\coqdoceol
\coqdocindent{12.50em}
\coqref{SECF.Noninterference.s:257}{\coqdocvariable{s}} \coqref{SECF.Imp.:::x '=[' x ']=>' x}{\coqdocnotation{=[}} \coqref{SECF.Imp.:com:com scope:'while' x 'do' x 'end'}{\coqdocnotation{while}} \coqref{SECF.Imp.Y}{\coqdocdefinition{Y}} \coqref{SECF.Imp.:com::x '=' x}{\coqdocnotation{=}} 1 \coqref{SECF.Imp.:com:com scope:'while' x 'do' x 'end'}{\coqdocnotation{do}} \coqref{SECF.Imp.:com:com scope:'skip'}{\coqdocnotation{skip}} \coqref{SECF.Imp.:com:com scope:'while' x 'do' x 'end'}{\coqdocnotation{end}} \coqref{SECF.Imp.:::x '=[' x ']=>' x}{\coqdocnotation{]=>}} \coqref{SECF.Noninterference.s':258}{\coqdocvariable{s'}} \coqexternalref{::type scope:x '->' x}{http://rocq-prover.org/doc/V9.1.1/stdlib//Corelib.Init.Logic}{\coqdocnotation{\ensuremath{\rightarrow}}}\coqdoceol
\coqdocindent{12.50em}
\coqexternalref{False}{http://rocq-prover.org/doc/V9.1.1/stdlib//Corelib.Init.Logic}{\coqdocinductive{False}}).\coqdoceol
\coqdocindent{1.00em}
\{ \coqdoctac{clear}. \coqdoctac{intros}.\coqdoceol
\coqdocindent{2.00em}
\coqdocvar{remember} \coqref{SECF.Imp.::com scope:'<x7B' x 'x7D>'}{\coqdocnotation{<\{}} \coqref{SECF.Imp.:com:com scope:'while' x 'do' x 'end'}{\coqdocnotation{while}} \coqref{SECF.Imp.Y}{\coqdocdefinition{Y}} \coqref{SECF.Imp.:com::x '=' x}{\coqdocnotation{=}} 1 \coqref{SECF.Imp.:com:com scope:'while' x 'do' x 'end'}{\coqdocnotation{do}} \coqref{SECF.Imp.:com:com scope:'skip'}{\coqdocnotation{skip}} \coqref{SECF.Imp.:com:com scope:'while' x 'do' x 'end'}{\coqdocnotation{end}} \coqref{SECF.Imp.::com scope:'<x7B' x 'x7D>'}{\coqdocnotation{\}>}} \coqdockw{as} \coqdocvar{loopdef}\coqdoceol
\coqdocindent{6.50em}
\coqdocvar{eqn}:\coqdocvar{Heqloopdef}.\coqdoceol
\coqdocindent{2.00em}
\coqdoctac{generalize} \coqdoctac{dependent} \coqdocvar{H}.\coqdoceol
\coqdocindent{2.00em}
\coqdoctac{induction} \coqdocvar{H0}; \coqdoctac{try} (\coqdoctac{discriminate} \coqdocvar{Heqloopdef}).\coqdoceol
\coqdocindent{2.00em}
- \coqdoctac{intros} \coqdocvar{HY}.\coqdoceol
\coqdocindent{3.00em}
\coqdoctac{injection} \coqdocvar{Heqloopdef} \coqdockw{as} \coqdocvar{H0} \coqdocvar{H1}. \coqdoctac{subst}.\coqdoceol
\coqdocindent{3.00em}
\coqdoctac{simpl} \coqdoctac{in} \coqdocvar{H}. \coqdoctac{rewrite} \coqdocvar{HY} \coqdoctac{in} \coqdocvar{H}. \coqdoctac{discriminate} \coqdocvar{H}.\coqdoceol
\coqdocindent{2.00em}
- \coqdoctac{intros} \coqdocvar{HY}.\coqdoceol
\coqdocindent{3.00em}
\coqdoctac{injection} \coqdocvar{Heqloopdef} \coqdockw{as} \coqdocvar{H0} \coqdocvar{H1}. \coqdoctac{subst}.\coqdoceol
\coqdocindent{3.00em}
\coqdoctac{inversion} \coqdocvar{H0\_}; \coqdoctac{subst}. \coqdoctac{eapply} \coqdocvar{IHceval2}; \coqdoctac{eauto}. \}\coqdoceol
\coqdocnoindent
\coqdoceol
\coqdocindent{1.00em}
\coqdoctac{inversion} \coqdocvar{Hcontra}; \coqdoctac{subst}. \coqdoctac{eapply} \coqdocvar{NSTOP} \coqdoctac{in} \coqdocvar{H1}; \coqdoctac{auto}.\coqdoceol
\coqdocnoindent
\coqdockw{Qed}.\coqdoceol
\coqdocemptyline
\coqdocnoindent
\coqdockw{Lemma} \coqdef{SECF.Noninterference.initial pub equiv}{initial\_pub\_equiv}{\coqdoclemma{initial\_pub\_equiv}}: \coqref{SECF.Noninterference.pub equiv}{\coqdocdefinition{pub\_equiv}} \coqref{SECF.Noninterference.xpub}{\coqdocdefinition{xpub}} \coqref{SECF.Noninterference.s1}{\coqdocdefinition{s1}} \coqref{SECF.Noninterference.s2}{\coqdocdefinition{s2}}.\coqdoceol
\coqdocnoindent
\coqdockw{Proof}.\coqdoceol
\coqdocindent{1.00em}
\coqdoctac{unfold} \coqref{SECF.Noninterference.s1}{\coqdocdefinition{s1}}, \coqref{SECF.Noninterference.s2}{\coqdocdefinition{s2}}, \coqref{SECF.Noninterference.pub equiv}{\coqdocdefinition{pub\_equiv}}. \coqdoctac{intros}.\coqdoceol
\coqdocindent{1.00em}
\coqdoctac{eapply} \coqref{SECF.Noninterference.xpub true}{\coqdoclemma{xpub\_true}} \coqdoctac{in} \coqdocvar{H}. \coqdoctac{subst}.\coqdoceol
\coqdocindent{1.00em}
\coqdoctac{repeat} \coqdoctac{rewrite} \coqref{SECF.Maps.t update eq}{\coqdocaxiom{t\_update\_eq}}. \coqdoctac{reflexivity}.\coqdoceol
\coqdocnoindent
\coqdockw{Qed}.\coqdoceol
\coqdocemptyline
\coqdocnoindent
\coqdockw{Lemma} \coqdef{SECF.Noninterference.not tsni counter}{not\_tsni\_counter}{\coqdoclemma{not\_tsni\_counter}} :\coqdoceol
\coqdocindent{1.00em}
\coqexternalref{::type scope:'x7E' x}{http://rocq-prover.org/doc/V9.1.1/stdlib//Corelib.Init.Logic}{\coqdocnotation{\ensuremath{\lnot}}} \coqexternalref{::type scope:'x7E' x}{http://rocq-prover.org/doc/V9.1.1/stdlib//Corelib.Init.Logic}{\coqdocnotation{(}}\coqref{SECF.Noninterference.tsni while R}{\coqdocdefinition{tsni\_while\_R}} \coqref{SECF.Imp.ceval}{\coqdocinductive{ceval}} \coqref{SECF.Noninterference.xpub}{\coqdocdefinition{xpub}} \coqref{SECF.Noninterference.counter}{\coqdocdefinition{counter}}\coqexternalref{::type scope:'x7E' x}{http://rocq-prover.org/doc/V9.1.1/stdlib//Corelib.Init.Logic}{\coqdocnotation{)}}.\coqdoceol
\coqdocnoindent
\coqdockw{Proof}.\coqdoceol
\coqdocindent{1.00em}
\coqdoctac{intros} \coqdocvar{Htsni}. \coqdoctac{unfold} \coqref{SECF.Noninterference.tsni while R}{\coqdocdefinition{tsni\_while\_R}} \coqdoctac{in} \coqdocvar{Htsni}.\coqdoceol
\coqdocindent{1.00em}
\coqdoctac{specialize} (\coqdocvar{Htsni} \coqdocvar{\_} \coqdocvar{\_} \coqdocvar{\_} \coqref{SECF.Noninterference.counter s1 terminates s1'}{\coqdoclemma{counter\_s1\_terminates\_s1'}} \coqref{SECF.Noninterference.initial pub equiv}{\coqdoclemma{initial\_pub\_equiv}}).\coqdoceol
\coqdocindent{1.00em}
\coqdoctac{destruct} \coqdocvar{Htsni} \coqdockw{as} [\coqdocvar{s2'} [\coqdocvar{D} \coqdocvar{\_}]].\coqdoceol
\coqdocindent{1.00em}
\coqdoctac{eapply} \coqref{SECF.Noninterference.counter s2 loops}{\coqdoclemma{counter\_s2\_loops}}. \coqdocvar{eassumption}.\coqdoceol
\coqdocnoindent
\coqdockw{Qed}.\coqdoceol
\coqdocemptyline
\coqdocnoindent
\coqdockw{Lemma} \coqdef{SECF.Noninterference.sme counter s1 terminates s1'}{sme\_counter\_s1\_terminates\_s1'}{\coqdoclemma{sme\_counter\_s1\_terminates\_s1'}} : \coqref{SECF.Noninterference.sme while}{\coqdocdefinition{sme\_while}} \coqref{SECF.Noninterference.xpub}{\coqdocdefinition{xpub}} \coqref{SECF.Noninterference.counter}{\coqdocdefinition{counter}} \coqref{SECF.Noninterference.s1}{\coqdocdefinition{s1}} \coqref{SECF.Noninterference.s1'}{\coqdocdefinition{s1'}}.\coqdoceol
\coqdocnoindent
\coqdockw{Proof}.\coqdoceol
\coqdocindent{1.00em}
\coqdoctac{unfold} \coqref{SECF.Noninterference.sme while}{\coqdocdefinition{sme\_while}}, \coqref{SECF.Noninterference.counter}{\coqdocdefinition{counter}}.\coqdoceol
\coqdocindent{1.00em}
\coqdoctac{\ensuremath{\exists}} \coqref{SECF.Noninterference.s1'}{\coqdocdefinition{s1'}}, \coqref{SECF.Noninterference.s1'}{\coqdocdefinition{s1'}}.\coqdoceol
\coqdocindent{1.00em}
\coqdoctac{split}; [|\coqdoctac{split}].\coqdoceol
\coqdocindent{1.00em}
- \coqdoctac{assert} (\coqdocvar{Hsplit}: \coqref{SECF.Noninterference.split state}{\coqdocdefinition{split\_state}} \coqref{SECF.Noninterference.s1}{\coqdocdefinition{s1}} \coqref{SECF.Noninterference.xpub}{\coqdocdefinition{xpub}} \coqref{SECF.Imp.:::'true'}{\coqdocnotation{true}} \coqexternalref{::type scope:x '=' x}{http://rocq-prover.org/doc/V9.1.1/stdlib//Corelib.Init.Logic}{\coqdocnotation{=}} \coqref{SECF.Noninterference.s1}{\coqdocdefinition{s1}}).\coqdoceol
\coqdocindent{2.00em}
\{ \coqdoctac{unfold} \coqref{SECF.Noninterference.split state}{\coqdocdefinition{split\_state}}, \coqref{SECF.Noninterference.s1}{\coqdocdefinition{s1}}, \coqref{SECF.Noninterference.xpub}{\coqdocdefinition{xpub}}.\coqdoceol
\coqdocindent{3.00em}
\coqdoctac{eapply} \coqdoclemma{functional\_extensionality}. \coqdoctac{intros} \coqdocvar{x}.\coqdoceol
\coqdocindent{3.00em}
\coqdoctac{destruct} (\coqdocdefinition{Bool.eqb} (\coqref{SECF.Maps.:::x '!->' x ';' x}{\coqdocnotation{(}}\coqref{SECF.Imp.X}{\coqdocdefinition{X}} \coqref{SECF.Maps.:::x '!->' x ';' x}{\coqdocnotation{!->}} \coqref{SECF.Imp.:::'true'}{\coqdocnotation{true}}\coqref{SECF.Maps.:::x '!->' x ';' x}{\coqdocnotation{;}} \coqref{SECF.Maps.:::'  ' '!->' x}{\coqdocnotation{\_\_}} \coqref{SECF.Maps.:::'  ' '!->' x}{\coqdocnotation{!->}} \coqref{SECF.Imp.:::'false'}{\coqdocnotation{false}}\coqref{SECF.Maps.:::x '!->' x ';' x}{\coqdocnotation{)}} \coqref{SECF.Maps.:::x '!->' x ';' x}{\coqdocnotation{x}}) \coqref{SECF.Imp.:::'true'}{\coqdocnotation{true}}) \coqdocvar{eqn}: \coqdocvar{B}.\coqdoceol
\coqdocindent{3.00em}
- \coqdoctac{reflexivity}.\coqdoceol
\coqdocindent{3.00em}
- \coqdoctac{destruct} (\coqdocdefinition{eqb} \coqdocvar{x} \coqref{SECF.Imp.Y}{\coqdocdefinition{Y}}) \coqdocvar{eqn}:\coqdocvar{HY}.\coqdoceol
\coqdocindent{4.00em}
+ \coqdoctac{rewrite} \coqdoclemma{eqb\_eq} \coqdoctac{in} \coqdocvar{HY}. \coqdoctac{subst}. \coqdoctac{rewrite} \coqref{SECF.Maps.t update neq}{\coqdocaxiom{t\_update\_neq}}.\coqdoceol
\coqdocindent{5.00em}
\ensuremath{\times} \coqdoctac{rewrite} \coqref{SECF.Maps.t update eq}{\coqdocaxiom{t\_update\_eq}}. \coqdoctac{reflexivity}.\coqdoceol
\coqdocindent{5.00em}
\ensuremath{\times} \coqdoctac{intros} \coqdocvar{Hcontra}. \coqdoctac{inversion} \coqdocvar{Hcontra}.\coqdoceol
\coqdocindent{4.00em}
+ \coqdoctac{rewrite} \coqdoclemma{eqb\_neq} \coqdoctac{in} \coqdocvar{HY}.\coqdoceol
\coqdocindent{5.00em}
\coqdoctac{destruct} (\coqdocdefinition{eqb} \coqdocvar{x} \coqref{SECF.Imp.X}{\coqdocdefinition{X}}) \coqdocvar{eqn}:\coqdocvar{HX}.\coqdoceol
\coqdocindent{5.00em}
\ensuremath{\times} \coqdoctac{rewrite} \coqdoclemma{eqb\_eq} \coqdoctac{in} \coqdocvar{HX}. \coqdoctac{subst}.\coqdoceol
\coqdocindent{6.00em}
\coqdoctac{rewrite} \coqref{SECF.Maps.t update eq}{\coqdocaxiom{t\_update\_eq}}. \coqdoctac{reflexivity}.\coqdoceol
\coqdocindent{5.00em}
\ensuremath{\times} \coqdoctac{rewrite} \coqdoclemma{eqb\_neq} \coqdoctac{in} \coqdocvar{HX}.\coqdoceol
\coqdocindent{6.00em}
\coqdoctac{rewrite} \coqref{SECF.Maps.t update neq}{\coqdocaxiom{t\_update\_neq}}; \coqdoctac{eauto}.\coqdoceol
\coqdocindent{6.00em}
\coqdoctac{rewrite} \coqref{SECF.Maps.t update neq}{\coqdocaxiom{t\_update\_neq}}; \coqdoctac{eauto}. \}\coqdoceol
\coqdocindent{2.00em}
\coqdoctac{rewrite} \coqdocvar{Hsplit}. \coqdoctac{eapply} \coqref{SECF.Noninterference.counter s1 terminates s1'}{\coqdoclemma{counter\_s1\_terminates\_s1'}}.\coqdoceol
\coqdocindent{1.00em}
- \coqdoctac{eapply} \coqref{SECF.Noninterference.counter s1 terminates s1'}{\coqdoclemma{counter\_s1\_terminates\_s1'}}.\coqdoceol
\coqdocindent{1.00em}
- \coqdoctac{eapply} \coqdoclemma{functional\_extensionality}. \coqdoctac{intros} \coqdocvar{x}.\coqdoceol
\coqdocindent{2.00em}
\coqdoctac{unfold} \coqref{SECF.Noninterference.merge states}{\coqdocdefinition{merge\_states}}, \coqref{SECF.Noninterference.xpub}{\coqdocdefinition{xpub}}.\coqdoceol
\coqdocindent{2.00em}
\coqdoctac{destruct} (\coqref{SECF.Maps.:::x '!->' x ';' x}{\coqdocnotation{(}}\coqref{SECF.Imp.X}{\coqdocdefinition{X}} \coqref{SECF.Maps.:::x '!->' x ';' x}{\coqdocnotation{!->}} \coqref{SECF.Imp.:::'true'}{\coqdocnotation{true}}\coqref{SECF.Maps.:::x '!->' x ';' x}{\coqdocnotation{;}} \coqref{SECF.Maps.:::'  ' '!->' x}{\coqdocnotation{\_\_}} \coqref{SECF.Maps.:::'  ' '!->' x}{\coqdocnotation{!->}} \coqref{SECF.Imp.:::'false'}{\coqdocnotation{false}}\coqref{SECF.Maps.:::x '!->' x ';' x}{\coqdocnotation{)}} \coqref{SECF.Maps.:::x '!->' x ';' x}{\coqdocnotation{x}}); \coqdoctac{reflexivity}.\coqdoceol
\coqdocnoindent
\coqdockw{Qed}.\coqdoceol
\coqdocemptyline
\coqdocnoindent
\coqdockw{Lemma} \coqdef{SECF.Noninterference.sme counter s2 loops}{sme\_counter\_s2\_loops}{\coqdoclemma{sme\_counter\_s2\_loops}}: \coqdockw{\ensuremath{\forall}} \coqdef{SECF.Noninterference.s2':261}{s2'}{\coqdocbinder{s2'}},\coqdoceol
\coqdocindent{1.00em}
\coqexternalref{::type scope:'x7E' x}{http://rocq-prover.org/doc/V9.1.1/stdlib//Corelib.Init.Logic}{\coqdocnotation{\ensuremath{\lnot}}} \coqexternalref{::type scope:'x7E' x}{http://rocq-prover.org/doc/V9.1.1/stdlib//Corelib.Init.Logic}{\coqdocnotation{(}}\coqref{SECF.Noninterference.sme while}{\coqdocdefinition{sme\_while}} \coqref{SECF.Noninterference.xpub}{\coqdocdefinition{xpub}} \coqref{SECF.Noninterference.counter}{\coqdocdefinition{counter}} \coqref{SECF.Noninterference.s2}{\coqdocdefinition{s2}} \coqref{SECF.Noninterference.s2':261}{\coqdocvariable{s2'}}\coqexternalref{::type scope:'x7E' x}{http://rocq-prover.org/doc/V9.1.1/stdlib//Corelib.Init.Logic}{\coqdocnotation{)}}.\coqdoceol
\coqdocnoindent
\coqdockw{Proof}.\coqdoceol
\coqdocindent{1.00em}
\coqdoctac{unfold} \coqexternalref{not}{http://rocq-prover.org/doc/V9.1.1/stdlib//Corelib.Init.Logic}{\coqdocdefinition{not}}, \coqref{SECF.Noninterference.sme while}{\coqdocdefinition{sme\_while}}. \coqdoctac{intros} \coqdocvar{s2'} \coqdocvar{H}.\coqdoceol
\coqdocindent{1.00em}
\coqdoctac{destruct} \coqdocvar{H} \coqdockw{as} [\coqdocvar{ps} [\coqdocvar{ss} [\coqdocvar{A} [\coqdocvar{B} \coqdocvar{C}]]]].\coqdoceol
\coqdocindent{1.00em}
\coqdoctac{eapply} \coqref{SECF.Noninterference.counter s2 loops}{\coqdoclemma{counter\_s2\_loops}}. \coqdocvar{eassumption}.\coqdoceol
\coqdocnoindent
\coqdockw{Qed}.\coqdoceol
\coqdocemptyline
\coqdocnoindent
\coqdockw{Lemma} \coqdef{SECF.Noninterference.not tsni while sme}{not\_tsni\_while\_sme}{\coqdoclemma{not\_tsni\_while\_sme}} :\coqdoceol
\coqdocindent{1.00em}
\coqexternalref{::type scope:'x7E' x}{http://rocq-prover.org/doc/V9.1.1/stdlib//Corelib.Init.Logic}{\coqdocnotation{\ensuremath{\lnot}}} \coqexternalref{::type scope:'x7E' x}{http://rocq-prover.org/doc/V9.1.1/stdlib//Corelib.Init.Logic}{\coqdocnotation{(}}\coqref{SECF.Noninterference.tsni while R}{\coqdocdefinition{tsni\_while\_R}} (\coqref{SECF.Noninterference.sme while}{\coqdocdefinition{sme\_while}} \coqref{SECF.Noninterference.xpub}{\coqdocdefinition{xpub}}) \coqref{SECF.Noninterference.xpub}{\coqdocdefinition{xpub}} \coqref{SECF.Noninterference.counter}{\coqdocdefinition{counter}}\coqexternalref{::type scope:'x7E' x}{http://rocq-prover.org/doc/V9.1.1/stdlib//Corelib.Init.Logic}{\coqdocnotation{)}}.\coqdoceol
\coqdocnoindent
\coqdockw{Proof}.\coqdoceol
\coqdocindent{1.00em}
\coqdoctac{intros} \coqdocvar{Htsni}. \coqdoctac{unfold} \coqref{SECF.Noninterference.tsni while R}{\coqdocdefinition{tsni\_while\_R}} \coqdoctac{in} \coqdocvar{Htsni}.\coqdoceol
\coqdocindent{1.00em}
\coqdoctac{specialize} (\coqdocvar{Htsni} \coqdocvar{\_} \coqdocvar{\_} \coqdocvar{\_} \coqref{SECF.Noninterference.sme counter s1 terminates s1'}{\coqdoclemma{sme\_counter\_s1\_terminates\_s1'}} \coqref{SECF.Noninterference.initial pub equiv}{\coqdoclemma{initial\_pub\_equiv}}).\coqdoceol
\coqdocindent{1.00em}
\coqdoctac{destruct} \coqdocvar{Htsni} \coqdockw{as} [\coqdocvar{s2'} [\coqdocvar{D} \coqdocvar{\_}]].\coqdoceol
\coqdocindent{1.00em}
\coqdoctac{eapply} \coqref{SECF.Noninterference.sme counter s2 loops}{\coqdoclemma{sme\_counter\_s2\_loops}}. \coqdocvar{eassumption}.\coqdoceol
\coqdocnoindent
\coqdockw{Qed}.\coqdoceol
\coqdocemptyline
\coqdocnoindent
\coqdockw{End} \coqref{SECF.Noninterference.TSNICOUNTER}{\coqdocsection{TSNICOUNTER}}.\coqdoceol
\coqdocemptyline
\end{coqdoccode}
\end{document}
